% ============================================================
%  Harald Høffding, MINDRE ARBEJDER [I].  København: Det Nordiske Forlag
%  (Bogforlaget Ernst Bojesen), 1899.  VIII, 243 pp.
%  "Shorter Works" -- TRANSLATION (English).
%
%  GENERATED by ebook-method/mindre-arbejder/tr_build.py from the parts in
%  texts/hoeffding/mindre-arbejder/.parts/tr/ (one per essay, in the print's order);
%  edit the parts, or this file only after the last build.
%  Translated from transcription.tex (Danish) in this directory, the text of record
%  (transcribed and collated 2026-10-05; see its header), as it stood on 2026-10-05.
%  The English of three essays made earlier from Wikisource's text (Realisme,
%  Vitalisme, Forholdet mellem Tro og Viden) was revised against that text.
%
%  DRAFT COMPLETE (2026-10-05): the whole book -- Preface (V-VIII), Contents, the twelve
%  essays, pp. 1-243 -- with every \opage of the Danish once and in order and the
%  footnotes page by page as in the Danish (tr_build.py checks both).  Made by eight
%  translator agents working from TRANSLATOR-BRIEF.md (in ebook-method/mindre-arbejder/),
%  one or two essays each; Realisme, Vitalisme and Forholdet revised from the earlier
%  English (Realisme's „Idealisme“ (p. 12) now „Realisme“ as printed; Vitalisme's
%  „Forundring“ (p. 48) now the print's reading, no longer a conjecture).  The
%  translators' odd spots in the Danish were read at the image and the transcription
%  corrected first (see its header).  Not yet read by an independent reviewer; terms
%  were harmonized only through the shared brief (Sibbern's coinages, Tönnies' Gemeinschaft/
%  Gesellschaft = Community/Society, are as each translator settled them).
%
%  CONVENTIONS (beyond TRANSLATION-PLAYBOOK.md §2):
%   * \opage{N} at the English word where printed page N begins (V-VIII: Forord).
%   * Italics and letterspacing in the Danish (\emph) -> \emph.
%   * Footnotes at the same anchor, numbered per page as printed by the class.
%     "Smlgn." / "Jfr." -> "Cf."; "Se" -> "See"; "Overs." -> "trans."; page and
%     volume numbers as printed; titles cited as cited.
%   * Quotations Høffding gives in Danish are translated from his Danish; what he
%     quotes in another language stays as printed; Danish verse is translated in
%     the text with the Danish in a comment below.
%   * »...«, „...“ -> ``...''.  "d. v. s." / "ɔ:" -> "i.e."; "f. Eks." -> "e.g."
%   * Misprints (PRINTED AS IS in the Danish): translate the sense, comment.
%  TERMS: as in the repo's other Høffding translations (humane-ethik, Etik,
%  Psykologi): Erkendelse = knowledge (cognition where set against feeling or
%  will); Viden = knowledge; Tro = faith (belief where not religious);
%  Livsanskuelse = view of life; Aand / aandelig = spirit / spiritual (mental
%  where plainly psychological); Sindsbevægelse = emotion; Følelse = feeling;
%  Forestilling = idea; Udvikling = development (evolution where Darwinian);
%  Samfund = society; Fremskridt = progress; Personlighed = personality;
%  Livsførelse = conduct of life; Tvivl = doubt; human = humane.
% ============================================================
\documentclass[12pt,a4paper,openany]{book}
\usepackage[utf8]{inputenc}
\usepackage[T1]{fontenc}
\usepackage{libertinus}
\usepackage{libertinust1math}
\usepackage{textalpha}
\usepackage[english]{babel}
\usepackage[protrusion=true,expansion=true]{microtype}
\usepackage{enumitem}
\usepackage{setspace}
\usepackage[margin=1.2in]{geometry}
\usepackage{fancyhdr}
\setlength{\headheight}{28pt}
\usepackage{hyperref}
\hypersetup{pdftitle={Shorter Works (1899)}, pdfauthor={Harald Høffding}, hidelinks}
\pagestyle{fancy}
\fancyhf{}
\fancyhead[LE,RO]{\thepage}
\fancyhead[RE]{\textit{Harald Høffding: Shorter Works}}
\fancyhead[LO]{\textit{1899}}
\fancypagestyle{plain}{\fancyhf{}\fancyhead[LE,RO]{\thepage}}
\setstretch{1.3}
\usepackage{marginnote}
\newcommand{\opage}[1]{\marginnote{\footnotesize\textit{[#1]}}}
\newcommand{\essay}[2][]{\chapter*{#2}\addcontentsline{toc}{chapter}{#1}}
\newcommand{\streg}{\begin{center}\rule{3em}{0.4pt}\end{center}}
\begin{document}
\frontmatter
\begin{titlepage}\centering\vspace*{2cm}
{\large HARALD HØFFDING}\\[1ex]\rule{3em}{0.4pt}\\[3ex]{\Huge SHORTER WORKS}\par\vfill
{\small Translated from the Danish}\\[0.6em]
{\small \textit{Mindre Arbejder.} København: Det Nordiske Forlag (Bogforlaget Ernst Bojesen), 1899}
\end{titlepage}
% --- p. V ---
\opage{V}
\chapter*{PREFACE.}

This book contains a series of essays that appeared in their
time partly as lectures, partly as critical reviews. The outward occasion for
my having collected them was that I --- in preparing for
a work I have in view --- felt a need to see how far back
I could trace the conception of the fundamental questions of science, society, and
religion that has now for a good many years formed
the foundation of my philosophical work. There then marked itself off
for me a period from 1882, when I gave a lecture at the first
meeting of the Students' Union on ``Realism in
Science and Faith,'' to 1898, when I gave a lecture on vitalism
at the Biological Society. These utterances, and all the
smaller works between them, bear witness to one and the same
fundamental conception, which it is now my wish --- if time and ability
make it possible --- to work through and carry out. It was then
a natural wish to have these various works collected
in print, partly because they may serve as an introduction to the
work I have in view, partly because they otherwise also contain
a good deal that might still be of interest, and I am grateful
to the highly esteemed publishers for having entered into my
idea.

% --- p. VI ---
\opage{VI}These works now go out to a different generation from the one
to which the first of them were addressed. When
the series of these works began, I still felt myself as belonging to
the younger men. Now I shall soon belong to the veterans; but I
have no reason to believe that I should not, for what I have here
to say, find receptiveness in one or another of the
generation that is now the young one.

It has been a special pleasure to me to reprint
the essay on F. C. Sibbern. I should like his
work and his personality to be preserved in honored memory,
and I should be glad if my own activity could in any
degree deserve to be regarded as a continuation
of his. ---

In and of themselves, several older smaller works might very well
have been included. Thus the lecture ``Darwinism and
Philosophy'' (Nær og Fjærn 1874) and my detailed critique of
Martensen's Ethics (Nær og Fjærn 1878). But
a limit had to be set, and the latter work was altogether
predominantly of a polemical nature. For this last reason
a series of articles in ``Tilskueren'' from November
1889 to May 1890 has likewise not been included; it polemicized against the
in my eyes exaggerated significance that had been attached to
Nietzsche's ideas and writings. ---

On one of the reprinted works I wish to
make a few remarks, namely on the one I have called
``Social Pessimism'', % print: Pessimisme» (misprint)
and which was occasioned by Ferdinand
Tönnies' book ``Gemeinschaft und Gesellschaft.'' The author
of this work, whom I now have the pleasure of being able to call
my personal friend, made, immediately after the essay had
been printed (1890), an objection to the title. ``If social
% --- p. VII ---
\opage{VII}pessimism,'' he wrote to me (July 26, 1890), ``is to mean
a practical tendency, then I do not hold with it; on the contrary, I
am inclined to regard all kinds of reform with
hopes.'' And this misgiving about the title he has
expressed again recently, when he heard that the essay was to be
reprinted. His intention with the book was a purely sociological
investigation and, besides, a drawing of parallels between the successive
social forms and psychological conditions. Although
this investigation in reality leads to a sharp criticism
of the forms of society now existing, he nevertheless looks with
hope and confidence toward a new culture and holds that work
on it is in full swing. In any case his view of
conditions has grown brighter since the publication of the book. ``I
readily admit to you,'' he writes (May 14, 1899), ``that when the
book was written I saw the future in a darker light, as regards
a possible rebirth, and that traces of this are to be found in
the book; but practically I stood even then very close to
the standpoint of Social Democracy.'' --- I must therefore ask the
reader to take the word ``pessimism'' as a characterization of
Tönnies' view with every possible reservation. It was now
really not so much my meaning to express that
Tönnies saw the future in a dark light, as that he saw
the past development --- as it has run its course since
the end of the Middle Ages --- in a dark light. Naturally
the two things are fairly closely connected. It is
difficult to understand how it can be bright for the toe when
it is wholly dark for the heel. There is here a
difficulty which Tönnies' conception shares with Marxism, a
view to which on the whole he stands close. I must add that
it is only against the title of the essay that my esteemed friend
% --- p. VIII ---
\opage{VIII}has made objection. ``I have,'' he wrote in the last
cited letter, ``again read your essay, and I again acknowledge
the excellent thoroughness and care that you have
bestowed on my book, to which I could not wish a better
philosophical interpreter, wherefore the new edition too is
for me very gratifying.'' I therefore dare to rely on it that my essay
will be a good preparation for anyone who might wish to
study Tönnies' book, which slowly but surely has
won an outstanding place in sociological literature.

\emph{May 22, 1899.}

\textbf{Harald Høffding.}


\chapter*{CONTENTS.}
\streg{}
\begin{flushleft}
On Realism in Science and Faith (1884) \dotfill\ 1\\
The History of Doubt in Modern Times (1891) \dotfill\ 15\\
On Vitalism (1898) \dotfill\ 40\\
The Physiology of the Emotions (1886) \dotfill\ 51\\
Frederik Christian Sibbern (1885) \dotfill\ 61\\
Alexander Hamilton and the North American Federal Constitution (1889) \dotfill\ 113\\
Social Pessimism (1890) \dotfill\ 142\\
Pagan Seekers of Truth (1892) \dotfill\ 158\\
Philosophy as Art (1893) \dotfill\ 180\\
The Relation between Faith and Knowledge in Its Historical Development (1885) \dotfill\ 195\\
The Religious Problem (1897) \dotfill\ 216\\
The Struggle between Old and New Ideas. A Retrospect (1895) \dotfill\ 225\\
\end{flushleft}
\begin{center}\rule{10em}{0.4pt}\end{center}

\mainmatter

% --- p. 1 ---
\opage{1}
\essay[On Realism in Science and Faith]{ON REALISM IN SCIENCE AND FAITH\footnote{Essentially a lecture delivered at the ``Studentersamfundet'' on 2 September 1882.}.}

\streg{}

The word realism is nowadays most frequently used to designate a tendency in poetic literature. It is sometimes employed in a reproachful sense, sometimes in a laudatory one, but rarely with a clear awareness of what it actually expresses, or of the limits of the validity of the concept it names. A clear account of this may perhaps only emerge once the modern aesthetic movement has more fully and distinctly unfolded its character and its consequences. It is not aesthetic realism that I intend to discuss here. Rather, I shall attempt to bring out what is distinctive about the realist tendency that, in our time, manifests itself no less clearly in the domains of science and of the view of life. It follows of itself that there must be an inner connection among the tendencies at work in these different domains, for just as little in the spiritual world as in the physical is there anything accidental or isolated. The human mind cannot be partitioned in such a way that, in its imagination, it follows paths that have absolutely nothing analogous or parallel to those it takes in its thinking and in its faith. Only the most doctrinaire and rigid among the opponents of modern aesthetic realism assign art and poetry such a secluded place in the world of the spirit that currents from other domains cannot --- or ought not --- reach them. By examining the nature and justification of realism in one domain, we indirectly contribute to the understanding of its appearance in others.

% --- p. 2 ---
\opage{2}
\section*{I.}

In the most comprehensive sense, we may define realism as \emph{the principle of natural causes}. If realism has opponents, and if it has only slowly worked its way forward to clear consciousness, this must therefore be because the explanation of phenomena and events has been sought --- and is still being sought --- by routes other than those of natural causes.

The human drive to find the causes of things has expressed itself long before the founding of a realistic science. No period can be pointed to in the developmental history of the human mind in which it stood entirely passive before sense-impressions, without the slightest attempt to explain them in some way or other. The most primitive savages have their explanation of nature, which, seen from their standpoint, may have a perfectly rational character. And if we could sufficiently enter into the consciousness of an animal, we should surely find that it too interprets the impressions it receives according to certain points of view. This already follows from the fact that the laws governing the connection of ideas (association of ideas) prove to apply equally to animal and to human consciousness. The history of knowledge does not, therefore, begin with a period in which single impressions are merely accumulated, to be combined and worked over in a subsequent period. At every stage of the development of knowledge, a more or less pronounced striving manifests itself to unite what is manifold and to fuse what is kindred.

The contrast between different periods in the history of knowledge is instead characterized by the different \emph{explanatory or interpretive principles} applied to experience. And here, above all, two great principles stand in sharp opposition to one another.

The one interprets natural phenomena according to rules drawn from \emph{outside} nature; the other interprets nature through \emph{nature itself}.

Under the first kind of natural explanation belongs all
% --- p. 3 ---
\opage{3}mythology. The mythological conception of nature explains phenomena as expressions of the will of personal beings. For it, nature is populated by more or less human-like spirits, whose intervention is traced above all in remarkable and significant events. From the land of fantasy and dreams, the real world is peopled. At the primitive stage of mind, where the inhabitants of dreamland are every bit as real as the phenomena of the actual world, it is in itself entirely consistent and justified to lay them at the foundation of natural explanation. In the theology of the higher national religions, the ultimate explanatory principle for everything is the will of a supernatural being --- likewise entirely consistent, once one has established the reality of such a supernatural being. But the natural explanation thereby attained does not regard nature in its own light. Just as Archimedes wished for a point outside the earth in order to move it, theology has sought a point \emph{outside} nature in order to explain it.

Something similar holds as well of many forms of philosophical idealism. If one finds the true explanation of a natural phenomenon in the purpose it is assumed to serve, or in the idea it is supposed to express, one is not explaining nature through nature itself. For whence do we know the purpose or the idea? Experience itself does not teach us these; knowledge of them must therefore be drawn from a source other than that through which natural phenomena become accessible to us. We must thus possess a special faculty for knowing ``the Ideas,'' and our faculty of knowledge splits into two branches whose connection cannot be demonstrated. We then have on the one side a faculty for apprehending what is given in experience, and on the other a faculty for entering into relation with a higher world and drawing from it the principles for understanding the world of experience. Insofar as the principle of natural causes is acknowledged at all, it is regarded as belonging only to the lower knowledge. True knowledge we obtain only when we recognize the Ideas in the world of experience. This conception, founded by \emph{Plato}, the
% --- p. 4 ---
\opage{4}progenitor of speculative idealism, recurs in various forms right down to the most recent times.

Realism, with the principle of natural causes, by contrast always derives natural phenomena from other natural phenomena, the more complex phenomena from the simpler. When it is said that nature must be explained \emph{through itself}, it must be remembered that nature is an unsurveyable world of phenomena unfolding for us in time and in space. Every point within this world we determine by its relation to other points. We do not go outside nature, since we have more than enough to do in connecting the infinite number of elements within it. Or more precisely: the expressions ``within'' and ``without'' now lose their meaning. However far we advance from cause to effect, or however far we go back from effect to cause --- no limit comes into view. Our inquiry may come to a standstill because it lacks the means to proceed further, as for instance the hypothesis of Kant and Laplace takes us no further back than the rotating primordial nebula. But wherever we begin or end, a new question always stands before us. The principle that nature is to be explained through itself therefore sets an infinite task. Nature as a whole, as the totality encompassing all natural phenomena in their mutual connection, is an ideal for our knowledge, which it approaches step by step but will never be able to realize completely.

\section*{II.}

When we think through the concept of realism in this way, the absolutely hostile opposition to idealism falls away. It is one of idealism's fundamental thoughts that we cannot rest at the level of individual, isolated facts, but must seek a bond that can bring coherence and unity to bear. Idealism went wrong in assuming that we could reach this unity and coherence by any route other than the continued working-through of the facts of experience. In its impatience it anticipated an ideal that can
% --- p. 5 ---
\opage{5}only be reached slowly and approximately; and since it did not find this ideal fully grounded in experience, it derived it from a higher sphere.

Realism and idealism become complete opposites only when realism entirely abandons the maxim that true knowledge does not consist in accumulated experiences, but is insight into the inner connection among experiences. But this maxim realism \emph{cannot} abandon, since it is the first and most important principle of all science. When realism turns polemically and skeptically against idealist speculation, it can only be because it holds that idealism has misunderstood and misused its own maxim. Its doubt is the expression of a deeper faith: alte dubitat qvi altius credit! Doubt and negation are not science's last word. There are of course many popular prejudices that dissolve into nothing under scientific criticism, many problems and questions that on closer examination prove to owe their origin to illusory presuppositions. The human mind is originally sanguine. It does not wait patiently for experience's instruction, but hastens ahead with its dogmas and hypotheses. With advancing experience there therefore also follows advancing criticism. The development of empirical science brings with it a damming up of those connections of ideas that hitherto spread securely in all directions like a river that has overflowed its banks. But this negative, inhibiting activity is only one side of the matter.

The conflict between the different conceptions of nature is decided in the same way as the conflict between dream and waking life. I attribute reality to waking life and its contents because it is more comprehensive and more coherent than the dream, which is always more or less fragmentary. The greater and firmer coherence makes the lesser and looser coherence intelligible, not the reverse. Only by giving us a more coherent conception of the world than speculative idealism and theology were able to provide will realism be able to prevail in the conflict. About particulars there will always be room for dispute. Because a series of experiences must be explained in accordance with realism's principle, other
% --- p. 6 ---
\opage{6}experiences might in themselves well be explained idealistically or theologically. Realism therefore becomes a distinctive conception of the world only when it ventures beyond what is merely given, or given hitherto, and opens the prospect of a total and thoroughgoing application of the principle of natural causes. This is what has occurred in the great scientific hypotheses, whose series is opened by the theory of Kant and Laplace and closes with those of Darwin and Spencer. These hypotheses have grown naturally as consequences of the course of development of science. In them, realism comes forward clearly as a universal scientific principle. It is not always the case that researchers in particular fields become conscious of this principle. The division of labor in science has meant that the individual researcher can with ease, and apparently also with good conscience, dig down into a single point without troubling himself over the extent to which the principles and methods he follows have a more general significance and contain consequences of far-reaching importance for the whole conception of the world. Many natural scientists' attitude toward Darwinism finds its explanation here. But one who fixes his gaze on the general course of development of the sciences and sees how every advance in understanding has been won by applying the principle of natural causes, regards this principle as the one called to govern our entire conception of nature. It must unfold all its consequences; we must appropriate it; there is no way around it.

Is realism, then, proven? Not at all. And it can be proven in the strict sense just as little as any principle can be. We arrive at a principle when, by following a developmental course, we discover a dominant rule or tendency that becomes more clearly visible the further the development proceeds. Every principle is, to that extent, a hypothesis that we try to apply to reality. We try one principle after another, until we find the one best able to render the given intelligible. The gods of mythology and the Ideas of Platonism would still reign in the sciences, if the advancing experience were explained just as well by them as by the natural
% --- p. 7 ---
\opage{7}causes accessible to experience itself. The validity and justification of scientific realism will depend on whether it truly indicates the only principle capable of bringing unity and harmony to our thinking, and on whether it continues to receive confirmation as science advances. Provisionally it indicates the only way in which our general world-view can come into harmony with the spirit in which particular investigations are conducted --- which follows from the fact that it is itself nothing other than the very spirit of modern science.

\section*{III.}

The expression ``natural causes'' requires closer explanation. If by the natural we understand that which can become an object of secure human experience, then spiritual phenomena belong to nature just as much as material ones. When a lively feeling sets the imagination in motion, or when by an exertion of will we recall something in our memory, we have here a natural causal relation just as much as when bodies expand upon the supply of heat or are compressed under strong pressure. Nature accordingly encompasses two groups or series of phenomena, which we can bring into relation with one another only by way of hypothesis or speculation: the material phenomena unfolding in space, and the phenomena of conscious life, which become immediately accessible to us only through self-observation, and which can only be represented in spatial form in a symbolic way. For both groups the principle of natural causes holds. Scientific psychology seeks the mutual connection among mental phenomena, just as physics and chemistry seek the mutual connection among material phenomena. In both cases we explain nature through nature itself.

The difficulty lies, as already noted, in finding the connection between the two domains of experience. Materialism cuts through the knot by overlooking their difference and simply making the spiritual a form or product of the material. But it must firmly be maintained that
% --- p. 8 ---
\opage{8}the question of the relation between the spiritual and the material is an open question within realism. Realism can very well acknowledge that there is such a difference between the two groups of experience that the one cannot be reduced to the other, even though they prove to stand in intimate connection with one another. We know of no spiritual phenomena that are not bound up with material phenomena, but we have no common measure that could make clear to us how a material motion can produce a thought, or conversely.

In earlier times this difficulty did not always appear so clearly. One let the brain produce thoughts and feelings, and conversely the will set the limbs in motion, without concern, so long as one had neither a clear idea of the laws of material causation nor of the fundamental distinctive character of mental phenomena. But just as the principle of natural causes has made its way through science as a whole, so in natural science in particular the principle has been increasingly applied that for every material phenomenon a material cause must also be demonstrated, and that a material cause can only have a material effect. Physiology too --- the science of organic life --- works according to this principle and advances the demand that all organic phenomena --- including, therefore, everything that takes place in our brain when we think, feel, and will --- be explained in accordance with the principle of material causes. Physiology concedes that there is still a very long way to go before this demand can be met, but it rejects any explanation that does not rest on this principle. Physiology properly concerns itself with mental phenomena only insofar as they are bound up with certain material functions. It is these material functions that physiology is to explain; that they are accompanied by thoughts and feelings is, strictly speaking, a matter of indifference to it.

But if we can no longer allow thoughts and feelings to intervene in the series of material causes, what place and significance are we to assign to them? Their reality no one can deny; it is immediately certain to us, more certain than
% --- p. 9 ---
\opage{9}anything else, for it is through thinking that we attain all well-grounded certainty. We might rather be driven to the assumption that the world of thought and feeling is the only real one, and that what we call the material world is merely a product of thought that forms itself for us unconsciously to ourselves and by virtue of an unavoidable illusion. But even then the question would recur: how should we conceive the relation between our consciousness-phenomena and the material phenomena we now, once and for all, picture as closely bound up with them?

There is probably no other way out than to conceive of consciousness-phenomena as inner, mental expressions of the same thing that, for the outer senses, appears as a material phenomenon. Real existence is thus only one, but we apprehend it (at certain points at least) in a double form. Which of these forms is the original, the constitutive one --- and whether there might perhaps be deeper-lying principles from which both are derived --- these are questions we shall not enter into here. Only this much should be emphasized: that realism does not here foreclose more far-reaching speculations, provided they have firm starting-points and supports in real experiences.

Realism therefore need not contest what is the genuine motive and deepest interest of the idealistic world-view. It is quite compatible with an effort to uphold the value and significance of spiritual life in existence. It cannot, to be sure, accept that the spiritual is something that descends from a world of Ideas into the material world as into a prison; but it does not deny or diminish its reality when it conceives of it as the inner and, for us, most significant manifestation of the same power that is at work in outer nature. It cannot --- but neither has idealism hitherto been able to --- teach us \emph{why} and \emph{how} spiritual states and activities are bound up with material ones; it restricts itself to the modest assumption that, since causal connection and lawfulness appear to prevail in the world, it is presumably no accident that consciousness and brain activity are bound up with one another, as experience shows us.
% --- p. 10 ---
\opage{10}It does not know whether there is a goal toward which all existence steers, but it finds no impossibility in itself in the assumption that this whole great natural interconnection, with its unbreakable laws and its firmly linked chain of causes and effects, is like a great household, a world order in which development and dissolution alternate, but through all the surging there nonetheless occurs an increase of what is valuable through the arising of higher forms of life.

\section*{IV.}

Already in the last reflections we have been led from the domain of science over into that of faith. For the most common distinguishing mark between \emph{science} and \emph{faith} is that the former investigates the laws and inner coherence of existence, while the latter asks about the ethical significance and value of existence.

When we define the concept of faith in this way, no necessary conflict exists between science and faith. Science grounds and explains; faith evaluates. The common foundation of both is the reality accessible to our experience. Its laws are sought by science, supported by the principle of natural causes, and faith expresses the value that this law-governed reality must have for us in virtue of our ideals. Explanation and evaluation need not exclude one another. What stands before us as high and glorious, as beautiful and good, may nonetheless have had its natural developmental history. Only false values disappear when the explanation is won. When we believe that world-development is more than a play of atoms, that --- even if it occurs through much struggle --- something good and valuable works its way forward through it, this faith does not exclude an equally firm conviction that this result is attained only through a firmly coherent causal sequence.

The reason we nonetheless distinguish between faith and knowledge is partly that ideal evaluation in and of itself says
% --- p. 11 ---
\opage{11}nothing about the real conditions for the arising and subsistence of what is valuable, and partly that the conviction of the subsistence and validity of what is valuable has its root more in feeling and will than in cognition. Only when I feel ideal tasks and ends as the true center of gravity of my own conduct of life do I have occasion to conceive of natural development as progressive realization of an ideal content. The ethical ideal I acknowledge casts its light out over nature. As I am --- thinking, feeling, and willing, with my ideals and my capacities --- I have emerged from the great natural development. What stands fixed for me as the unconditionally true and good is certainly conditioned by my finite nature and my vanishing position in the great whole; but thereby it is not rendered meaningless and accidental. It is by its fruits that one shall know the tree, and from the fruits that nature has allowed to ripen in and for us, we can learn that forces are at work within it that not only operate in the same direction as our highest strivings, but have themselves from the outset led us into these strivings. On such a foundation, however, no theoretical system of doctrine can be built, as speculative philosophy assumed. In its further elaboration, poetic symbols would soon take the place of scientific concepts.

By \emph{realism in faith} I mean a view of life that holds fast to the belief in the ideal value of world-development, while at the same time being equally deeply convinced that this ideal value is realized through the operation of natural causes. Once this faith has been won, it is independent of theoretical vicissitudes, for it is from the outset in harmony with the spirit of science itself.

The customary meaning of the word faith, by contrast, excludes such a harmony. It demands a more or less supernatural apparatus for the realization of its ends, and sets up fixed forms and inviolable symbols. On this rests the difference between humane or philosophical faith and theological faith. Theological faith must, in accordance with its principle, put forward theoretical doctrines that come into conflict with science. It has as its object something that is supposed to be
% --- p. 12 ---
\opage{12}elevated above theoretical explanation and to defy the principle of natural causes. Under these conditions, the conflict between faith and knowledge can never come to an end. So long as these presuppositions are maintained, we shall not escape repetitions of the disturbing spectacle that the history of science shows us of faith's anxiety and flight before science. Each time that science, in experiment, proof, or hypothesis, advances to a new domain, an outcry is raised on the part of theological faith. Copernicus, Galilei, and Darwin, Giordano Bruno, Spinoza, and J.\ G.\ Fichte are examples of how unprejudiced observation and logical consistency lead into conflict with a faith that has not yet learned that evaluation and explanation do not cancel each other out. Only a faith that can show the resignation to bow before the principle of natural causes leads beyond this conflict. Here too it holds: \textit{alte dubitat qui altius credit}. For it is surely the deepest faith that does not require supernatural interventions and revelations in order to maintain the conviction of the significance and validity of the ideal in the world. ---

In an entirely different sense of the word, one might also speak of a realist tendency within theological faith in our day. Realism is everywhere characterized by mistrust of speculation; it seeks to have firm ground underfoot and therefore holds to what is positively given, to the factual --- or to what each variety of realism regards as factually given. Now when theological faith is in its own way gripped by the spirit of realism and looks about for its facts to hold fast to, it encounters above all its inherited formulas and confessional writings, in which it has given to it the content sanctioned by absolute authority. It will therefore be clear that when a realist tendency dominates an age, it will within the domain of theological faith naturally tend to appear as intensified faith in authority. The same tendency that produces materialism outside the domain of theological faith produces literalism within it. Theology has therefore also in our day abandoned every speculative grounding and retreated to the ecclesiastically sanctioned dogmas.
% --- p. 13 ---
\opage{13}This creates a great distance between a theological thinker like Bishop Martensen and the younger theological generation. At the end of his book on Jakob Bøhme one can read his warning words to ``the young theological realists.'' He is surely right, for this narrow-minded clinging to inherited dogmas is not sustainable in the long run.

Nothing stands still in the world. Even where on the surface there seems to be rest and unchangeability, the most significant processes may be taking place in the depths, and layers being deposited that will in time come to light. The universally human experiences exert their quiet influence. Through the theological-religious development since the days of the Middle Ages there can be traced a constant tendency to reach the goal by as short and direct a path and with as simple means as possible --- the same tendency that on scientific ground expresses itself in the rule that we must not multiply our principles and hypotheses without compelling necessity. According to theological faith, there are a very great many large things one must know about in order to live as a human being and not as one who eats and drinks because he knows he will die tomorrow. The advancing development leads here slowly but steadily toward greater economy and modesty. Protestantism already marks a certain reduction relative to Catholicism, and relative to old Lutheran orthodoxy, Pietism, Rationalism, and Grundtvigianism are all definite advances in concentration and simplification. In this respect there is a parallel between theology and medicine. The older medicine distinguished itself by an enormous number of medications and a large number of heroic cures. But one gradually manages both with fewer medications for one's bodily health and with fewer dogmas for one's spiritual health. \emph{Where} the simplification will end is difficult to say. It will surely not end at the same point for all individuals. But in principle the goal will surely be reached as soon as the harmony between evaluation and explanation has properly dawned on the general consciousness. Then the significant kernel in the higher forms of theological faith will also be preserved,
% --- p. 14 ---
\opage{14}while the husks crumble away. For what animates the higher national religions and has given them their mighty power over thousands of generations is ultimately the ideal impulse that wants to see existence as something other and more than a sum of facts, and this impulse is far nobler and more important than any of the forms and symbols in which it has expressed itself. Humane faith can and will preserve this impulse; indeed, its very striving is directed toward purifying it and freeing it from all inhibiting constraints.

Only jadedness throws away the kernel along with the husks. Such jadedness is surely the most dangerous enemy of spiritual development in our age, more dangerous than all intolerance and fanaticism, for in the latter there is at least life, while jadedness is itself petrified and petrifies everything around it. But jadedness, for which everything is equally good so long as it is factual, naturally ends by no longer troubling itself very much to distinguish between real and illusory facts; it thus finally stifles the very sense for reality from which it derives its own origin --- petrifies its own source. It arises easily in times of strong upheaval; but since it rests on a confusion of the accidental with the essential, it will disappear, like the foam that a strong breaking of waves produces.

\streg{}

% --- p. 15 ---
\opage{15}
\essay[The History of Doubt in Modern Times]{THE HISTORY OF DOUBT IN MODERN TIMES.}

\begin{center}\emph{C. N. Starcke:} Skepticismen som Led i de aandelige Bevægelser siden
Reformationen. Copenhagen 1890.\end{center}

\streg{}

It has often been pointed out what a large place the
wavering and doubting natures occupy in modern literature.
The Hamlet figures recur constantly in countless forms
and have perhaps scarcely ever flourished to the degree they do
in our own day, although their dimensions grow smaller and their
content more insignificant, the more modern the portrayal is. Such
a standing type points to an essential
peculiarity of modern development, the same
peculiarity which makes it the case that, when the modern age is to be characterized,
it is as a rule negative expressions such as liberation, emancipation,
and the like that present themselves first. What is most striking is
the break from the old, just as what takes place at
birth is first and foremost the child's break from the
maternal organism; only afterward comes the question of the
newborn's viability and content of life. Whether that which hitherto
was nourished in its mother's life can now stand on its own legs and
lead its own life must naturally be decided by practical trial.
That one can live one proves best by living. And yet
man, who, as Homer says, is a being that ``looks
forward and back,'' as soon as consciousness is developed,
will naturally come to think about the means that are at his
disposal in the struggle for life. So it goes also with the
modern age. While its life unfolds in the
various directions, reflection asserts itself, and
conflicting thoughts arise, when the question is raised
what it is that bears our life, and whether that which bears it
has strength enough. There are those who hold that only in the old,
% --- p. 16 ---
\opage{16}whose undisputed dominion ceased with the Renaissance and
the Revolution, is there strength and peace. In opposition to this
there are those who hold that the fault is precisely that the remains of the
old weigh upon us, and that ``we sail with a corpse in the hold.''
Others again hold that the break was necessary, but that it
has inflicted wounds which can be healed only slowly, and has inhibited
sides of the spiritual life which only gradually can
develop fully anew. Finally there are those --- and to
their number belongs the author of the above-mentioned work ---
who hold that the fault lies in a lack of self-confidence and
consistency within the new view of life, in that the doubt which
quite naturally had to arise at the break with the
medieval view of life and ordering of life has been for far too long
regarded as justified and insuperable, and that the principal task
must therefore be to drive doubt out of every hiding place in which
it can conceal itself.

Dr. Starcke is well known in philosophical and sociological
literature. Besides his doctoral dissertation on Ludwig
Feuerbach he has published an investigation of ``the primitive
family,'' which now exists in German, English, and
French, and which I reviewed in its time in ``Tilskueren,''
and a treatise on ``The Theoretical Foundation of Ethics,'' which is
included in the 6th series of the Transactions of the Academy of Sciences. In
the present work the author has made great progress
in powers of presentation. One receives from the book a beneficial
impression of independent thinking and of a sober and yet
warm-hearted conception of human life and its
conditions. --- The author has shown me the honor of putting my
name on the first leaf of the book together with a Finnish
man of science. I see in this no hindrance to my being able here to
state both what I have to acknowledge and what I have
to dispute in the work before me.

It is a good idea that Dr. Starcke has had, to regard
the skeptical tendency in philosophy as a kind of thermometer
of spiritual development. He investigates what
was the cause of a series of skeptical authors arising soon after the
Reformation, and what fundamental thoughts and
points of view came forward in them. By then showing how
% --- p. 17 ---
\opage{17}in the later authors, who may be said to fall under
the concept of skepticism, motives and fundamental thoughts change, he gets
a measure of what strength has been won, and
what prospect an independent view of life has of asserting
itself. Dr. Starcke comes by this route into
a series of epistemological, psychological, and ethical
problems, the treatment of which he holds has been inhibited by
the skeptical tendency, which has persisted longer than
necessary, and he himself gives able contributions to their
treatment.

The work falls, as will be seen from this, into two main parts,
a predominantly historical part and a predominantly theoretical one.
The author has, it seems to me, not kept them sufficiently
apart. In particular he theorizes too
much already in the first part, although this was after all to
constitute an essential part of his body of evidence. His
presentation would certainly have gained in clarity and interest
if the historical treatment of skepticism had been kept
more to itself, and if there had then also been given a
more thoroughgoing psychological and historical characterization of
the individual authors. Thereby it would certainly also have
proved necessary to include more authors; one misses at
least some who are of importance for the questions
that are discussed in the last part of the book. The author gives so
much that is good that we naturally wish for more, especially as it would
accord with the well-conceived plan.

To characterize the book I shall dwell on both
main divisions, that is, first consider skepticism as a historical
phenomenon, and next discuss Dr. Starcke's view that
every trace of skepticism can and shall disappear from
thought and life.

\section*{I.}

Skepticism consists, according to Dr. Starcke, in the
theoretical domain in a despair of man's ability
to form a knowledge of the innermost real nature of the world and of man, and in the practical domain
% --- p. 18 ---
\opage{18}in a doubt arising about man's ability to learn on
his own to know the good and to love it.
Such a doubt can arise only when neither the theoretical
view of the world nor the practical view of life and
conduct of life any longer is determined throughout by faith in holy
authorities. In the century of the Renaissance and the Reformation
there was therefore a fruitful soil for it. It takes on a
very different character and direction, according as it is
motivated chiefly by love of the old and is used as a weapon
against the new, as one from indolence or fear wishes to keep
oneself away from life, or springs from love of the
new, combined with fear that this new cannot be
carried through and be sufficient to itself. In the foreground, however,
there always comes to stand the question whether there is an independent
theoretical science and an independent ethics.

According to Dr. Starcke's conception, skepticism in the
form in which it appeared in the 16th century and in
the time that followed has an essentially different character from
the skeptical tendency within Greek philosophy. The skepticism of antiquity
had, he holds, a more theoretical character
and can in no way be put alongside the modern one,
in which a strong sensation shows itself of having
come off worst in its struggle for existence. He finds the
difference expressed in the different use that is made
of the arguments for the relativity of knowledge: with the Greek
skeptics these arguments are used, according to his conception,
to demonstrate the impossibility of certain knowledge, whereas
in modern times they are used to show that the practical insecurity
in the face of nature cannot be remedied. --- I have doubts
whether the author is right in this. One has become too much
accustomed to regard the Greek skeptics (whose own
writings we do not possess) from the purely negative side. Much
indicates that their dismissive attitude toward dogmas
and speculations was connected with their interest in
the method of experience. There is a historical connection between the
skeptical philosophers and the empirical school of medicine. Galen uses
of this the expression that the empirical physicians were in medicine
what the skeptics were in philosophy, and he presents Pyrrho,
% --- p. 19 ---
\opage{19}the founder of the skeptical school, as a model of
the true empiricist\footnote{I cannot here go more closely into this, but must refer to
the chapter ``Die Erfahrungslehre der Skeptiker'' in \emph{P. Natorp:}
``Forschungen zur Geschichte des Erkenntnissproblemes im Alterthum''.
(Berlin 1884).}. That the science of experience in antiquity
was early inhibited in its growth is true; but it made its
attempts, and a comparison between the Greek and the modern
skeptics would therefore not lack justification
and interest. ---

In the first place Dr. Starcke discusses \emph{Michel de Montaigne}
(1533---1584), the chief representative of modern
skepticism, in so far as this denotes ``the revolt of the individual against
authority.'' This expression seems to me somewhat too strong
where Montaigne is concerned. With him there is no question
of revolt. His manifold knowledge, his ingenious
juxtapositions, and often really profound thoughts he uses
in an easygoing way to make room for the unfolding
of individualities according to their different nature. This very
striving for individual liberation, of which there is
incidentally an interesting description in the book, appears
--- what I believe one as a rule overlooks when Montaigne
is spoken of --- not in a purely negative way, but is connected
with Montaigne's appeal to Nature, this ``our great
and mighty mother,'' whom men have forsaken in order to
submit to dogmas and laws that arose by chance and were sanctioned by
custom. Montaigne glorifies ignorance when
it is combined with a good head, and praises the natural
instinct, which helps better than all clever
inventions. In each individual Nature expresses itself in a different way: Everyone
has his ruling faculty (forme maistresse). What matters then is
to bring Nature to its right, which does not happen by
folding one's hands in one's lap, but demands reflection and
self-control. --- There is a whole great positive background in
Montaigne, which even at a certain point recalls his
contemporary Giordano Bruno's enthusiastic nature-pantheism.
The name skeptic is not exhaustive here. He stands in
% --- p. 20 ---
% print p. 20: Religionen. da (a stop for a comma) -- misprint; the sense is translated
\opage{20}French literature as a forerunner of Rousseau and Diderot by
his appeal to Nature. A skeptic he is nevertheless, in so far as
his passive clinging to the old sharpens his criticism
of the new, and because he feels a shyness, motivated by indolence,
of the unrest which the new tendencies aroused both
in the domain of religion (Luther) and in the domain of
science (Copernicus, Paracelsus). It was not in his
nature suddenly to turn all his ideas
upside down.

Montaigne's successor \emph{Pierre Charron} too must, I believe,
be conceived from a more positive side than happens in Dr.
Starcke's presentation. Charron is remarkable in the history of ethics
in that we find in him, the Catholic clergyman, who uses
the skeptical method to show that man cannot
measure everything by his natural reason but must hold to the
old doctrine of the Church, --- that we find in him nevertheless the first
express assertion of the independence of ethics from religion.
Religion and morality, he teaches, rest each on its own motive
(or mainspring, \emph{ressort}) and must therefore not be confounded.
Above all goodness must not be made dependent on religion,
for it thereby becomes accidental and does not spring from \emph{le bon
ressort de la nature}. ``I want one to be a good
man even if there were no Paradise or Hell.
It is for me a dreadful saying: `If I were not
a Christian, if I did not fear being condemned, I should
do or leave undone this or that.' I want you to be
a good man, because Nature and reason will it so,
because it is demanded by the universal order of things, of which
you are only a part, --- because you cannot act against yourself,
your being, your end. Let what may come of it.'' (Charron: De la sagesse, Bordeaux. 1600. p.
342). --- There is in Charron, as in Montaigne,
decidedly a positive background behind the skeptical form and
method.

The same holds for two thinkers whose
skeptical method finds application especially in the theoretical
domain. --- \emph{Sanchez} did indeed set up the proposition that
we know nothing, even in the title of his chief philosophical work
% --- p. 21 ---
\opage{21}(Qvod nihil scitur) But one must not lay too much
weight on a paradoxical manner of expression. The main thought in
Sanchez is the contrast between the ideal of knowledge in the
sense of perfect insight into the nature of things (perfecta
cognitio rei) and our actual knowledge. He does not have the great
hopes of what improved methods would be able to bring
that his contemporary Bacon of Verulam had; but dejection was
not the only mood in him. He urges energetically
``to go to the things themselves'' and not to entangle oneself in scholastic
chains of inference. And the very work in which he develops
his skepticism is an introduction to his medical
works. He has moreover himself presented a philosophy of nature
which recalls the one earlier set up by the Italian Telesio,
and which admittedly has no lasting significance, but yet
bears witness that the courage of knowledge was not broken in him.

This holds still more of \emph{Joseph Glanvil}. He has,
like the preceding authors, a clear eye for everything
that inhibits the development of our knowledge, and gives a broad
presentation of it. But he declares expressly that
skepticism is for him a method, not a system. He wants
to be rid of dogmatism, both in its scholastic and in its
materialistic form, and uses the skeptical method to
undermine it and to widen and liberate the mind; for all
dogmatism is for him narrowness of spirit. Toward the end
of his presentation he says (Scepsis scientifica. Chap. XXVI):
``I have fought against dogmatism under the banner of
skepticism, but I do not wish that the purpose
of my whole investigation be regarded from this point of view.'' It is,
he says in another place (Chap. X), only dogmatic
ignorance that counts the cautious reserve from
every untested assertion as skepticism. What Glanvil especially
insists on is the necessity --- and at the same time the difficulty
--- of an empirical confirmation of all our hypotheses. In his
dedication to the ``Royal Society'' he says that without
facts of experience ``our hypotheses are but dreams and
romances, and our science mere conjectures and opinions. For
when we form systems without consulting the phenomena,
we build only in the air and describe a self-made and imagined
% --- p. 22 ---
\opage{22}world, which is but little akin to the real one \ldots.
And it is possible that all the hypotheses hitherto devised were built
on too narrow a view of things.'' Such a
warning every age needs, not only the age that
lay before Newton. --- And the positive side appears
very definitely even in Glanvil, in that he adheres to the
Cartesian philosophy and physics, to which his whole work
can be regarded as an introduction. He has really only
given a broad development of the skeptical method which
Descartes himself had used in his ``Discours sur la méthode.'' ---

By these examples I have wished to show that one easily comes
to read too much into the word skeptic, and that several
of the older authors who give themselves this name
or figure under it in the history of philosophy nevertheless stand in
a positive relation to knowledge and its tasks. I
believe that Dr. Starcke's conception here is somewhat one-sided, and that
this has influenced his train of thought throughout the whole book.
For just as he conceives the older skeptics,
e.g. men such as those just named, purely from the negative
side, so his conception of ``skepticism'' on the whole
is that it is something which arose under conditions of ferment of
a passing nature, and that it now can and shall disappear
completely. His book is a hunt for skepticism in order
to drive it out of every nook in which it can hide
itself. In opposition to this I see the significance of those ``skeptics''
in their asserting the scientific method in all
domains, which could not happen without strong weight being laid
on the distance from the ideal of knowledge. And I
cannot see that all skepticism is evil, even in
our own day, when we have yet seen so much that would have
stood as a wild dream for the authors of the 17th century
established as certain knowledge. But more
of this later. ---

Excellent and, as far as I know, new is the light in which
\emph{Pascal}'s skepticism is placed. Dr. Starcke distinguishes
between thinkers like most of those mentioned above, whose skepticism
rests chiefly on a lack of confidence that man
by his own help can build his world-view, and a
% --- p. 23 ---
\opage{23}Pascal, with whom what is decisive is ``the horror which seizes
man when he sees the possibility that on his own
he can find his way, but at the same time sees this
possibility tied to a condition, the mechanics of the world, which
seems to attack the value of life at its innermost root.'' It
is now ``already the very method of research and the scheme
of the universe which this presupposes that awaken concern.
The thought of a world which, like a clockwork, runs by
eternal mechanical laws, these men seem to find must
fill us with dread.'' In this way Dr. Starcke explains
the passion of Pascal's polemic against natural
reason, in that he sees in him ``a vivification of the great
tragedy which was played out in the collision between the old
and the new.'' --- I add here only that it is a tragedy
which has often repeated itself later in the rhythmic
development of thought. In our own literature we have an original
``vivification'' of this kind in S. Kierkegaard. And he will
perhaps not be the last.

In a following section the English \emph{Deism} is treated in
an interesting way. --- Deism was a compromise: one
threw a number of dogmas overboard which one thought had
a purely external, historical origin, and kept back a small
chosen circle of religious ideas (especially of a personal
God and immortality), which one thought sprang from
``reason'' and was compatible with a purely scientific conception of the world. That was the standpoint on which the 18th
century's ``Enlightenment'' stood. Dr. Starcke seeks to show that
Deism dissolves, since the consequence is drawn that man in
his conduct of life must be able to dispense with these remaining dogmas
just as well as with the rejected ones, and that ethics is just as
independent of ``natural'' religion as of positive
religion. It might in this connection have been pointed
out that Deism was also attacked from an entirely different side.
In his remarkable work ``Analogy of Religion, natural and
revealed, to the constitution and course of nature'' (London
1736) \emph{Joseph Butler} showed that the most important objections
which can be raised against revealed religion can also be
raised against faith in a providential government in nature, since one
% --- p. 24 ---
\opage{24}by the way of experience cannot prove that the author of Nature
possesses wisdom, justice, and goodness. When one takes offense
at revealed religion, then there is, Butler teaches,
quite as much to take offense at in Nature. In Nature the
innocent, e.g., often suffers for the guilty, just as in
the ecclesiastical doctrine of atonement. Butler's work, which --- in
contrast to the prevailing optimism --- emphasizes the
dark, disharmonious side of Nature, especially such features
as are now brought under the point of view of the blind struggle
for existence, played a decisive role in the religious
discussion of the following century and has, among others, strongly
influenced James Mill's and Stuart Mill's religious views.
--- Another want, too, one feels here in the historical
presentation. It would have been natural here to mention
\emph{Mandeville}'s Fable of the Bees, with its fundamental principle that society
can thrive and civilization flourish only at the expense
of morals, a contribution of not inconsiderable importance in the
history of modern skepticism, and which would in particular have
been of interest to Dr. Starcke, who in what follows so
strongly emphasizes the problem how the individual's
pursuit of happiness can be reconciled with the demands of society. --- Finally it would
have been interesting if the author, who dwells on several
less significant men from the literary circle that surrounded
\emph{Frederick the Great}, had instead illuminated
the philosophizing king's own relation to the questions treated in
the book. The king's passage from Wolff's dogmatism
through Bayle's skepticism to Locke's empiricism would be
of interest for illuminating ``Skepticism as a Link in the
spiritual movements since the Reformation,'' just as his
definitive views denote a peculiar modification
of the ``natural religion'' prevailing in the Age of Enlightenment. ---

But all these remarks do not shake the value
of the historical part of the book. The author has --- by
carrying his presentation through to Hume, Rousseau, Kant, and
Stuart Mill --- clearly proved his main thesis, that
skepticism ends by having an entirely different character and taking up
a different position than at its beginning. It began
% --- p. 25 ---
\opage{25}as an expression of distrust, modesty, or caution.
It ends with the conviction that a whole series of
speculative and theological questions can safely be put aside,
since both positive science and practical
human life go their way and reach their goal, whether those
questions are solved or not. The conviction reached through many spiritual
movements and struggles, which the author especially
finds represented by Rousseau and Kant, he expresses
aptly in the following words: ``In life itself the means must be sought
to solve the tasks of life; if they are not found there, they are found
nowhere \ldots{} Only on human ground can the way
be laid \ldots{} After Kant and Rousseau skepticism will not
any more appear in such clear forms as hitherto, because
no choice is any longer posed at all as to whether
this way on human ground shall be followed or not \ldots{}
We no longer meet with pure, consistent skeptical
standpoints; there is in most standpoints a skeptical
tinge, and it therefore becomes our task to gather all these
isolated features into distinct pictures.'' These last words
denote Dr. Starcke's conception of skepticism in our day.
We are, he holds, ``broken out beyond the boundary marks
of skepticism,'' and what now matters is only ``to draw
the last sting out.''

This result appears not merely in
the sciences of experience and the purely theoretical results of critical philosophy.
It shows itself especially in the changed position which religion
takes up in practical respects. At the beginning of the development
both the principles of the world-view and of the conduct of life
were derived from religion. There has now formed a
world-picture formed in a purely scientific way. The significance of religion
is therefore restricted more and more to the practical,
ethical domain. But here too a change takes place.
In the Middle Ages morality became dear to man through
religion. In modern times the relation first changed to the effect
that religion was asserted as the source of morality. Later, when
the self-sufficiency and independence of ethics was carried through,
religion was no longer its source but merely its guarantee
and sanction. Religion is now defended as necessary, not
because through it one learns to know morality, but because
% --- p. 26 ---
\opage{26}it is thought that without faith in divine retribution
morality gains no power in the human world. --- This
essential change is one of the most important results
that the history of skepticism can show.

\section*{II.}

The last part of Dr. Starcke's book seeks to show
the necessity and the possibility of ``the skeptical sting
being drawn out.'' He treats on this occasion a series of
epistemological and ethical problems in an acute
manner, but in such detail that I cannot at this place
follow him, but must keep to a few main points. And I choose here, as in what has gone before,
especially those on which I depart from him.

First and foremost there arises the question whether he
is right that the ``skeptical sting'' which is left can and
shall be drawn out. The decision of this will naturally depend in
part on how much one puts into the word
skepticism. I have in what precedes sought to show that Dr. Starcke
takes the ``skeptical'' authors of the older age somewhat too heavily,
in that he makes the skeptical method one with absolute
skepticism and underestimates the positive fundamental thoughts for
which one by doubt wished to make room. If one puts into the
concept of skepticism unrest and dejection as
necessary features, then there can hardly be dispute that
skepticism is a subordinate and passing phenomenon.
It is then due essentially to a kind of contrast effect.
Dogmatism promised gold and green forests, answers to all
questions, a conclusion of the work of thought. The modern
science of experience and philosophy cannot give such promises.
It promises us progress in detail, new facts
and greater clarity as regards the principles, and it points
toward a great inner connection in everything; but it raises
at the same time ever new problems, calls for ever new work,
because ever new barriers present themselves to be moved. As
there is no absolute rest to be found in external nature, but everything is
in motion and activity, so this holds especially
% --- p. 27 ---
\opage{27}for thought. It is not permitted to rest in contemplative
beholding; the provisionally concluding ideas at which it halts
soon show themselves to be symbols, which hint and signal,
but do not immediately reproduce the reality of existence. He
who is suddenly brought face to face with this conception of our
knowledge, after having previously rested securely in the
bosom of dogmatism, must indeed feel uneasy and dejected. But
only so long as the break is fresh, and the contrast to the
dogmatic faith is felt in its whole sharpness, will this
disquiet last. It will vanish, the more he lives himself into
the new conditions. Resignation will always be necessary.
Such resignation no world-view whatever can dispense with,
since everywhere what matters is submission to presupposed
fixed conditions. There lies absolutely nothing skeptical in
the acknowledgment of the necessity of resignation.

Life must always begin with a surplus of strength and
courage, for the application of which experience successively sets its
limits. Every life is a stream, which partly breaks through
what stands against it, partly is dammed in by what cannot be
broken. The final bed of the stream is determined by the relation
between these two sides. To get to know life is to experience
how far one with the forces and drives one commands
can extend the given barriers, and in what degree one
on the other hand is compelled to submit to them. The necessity of
resignation lies in this, that the final bed of the stream is not \emph{merely}
determined by the force of the source, but also by the
opposing, damming forces. This holds both for
our theoretical and our practical striving.

The scientific method does not consist in a passive
accumulation of observations from which results are later
drawn. The first condition for fruitful scientific
work is a guiding idea, a provisional hypothesis, which becomes
the motive for seeking out new experiences. These new experiences
often bring with them the necessity of the idea being altered or even
exchanged for other ideas, the possibility of coming to apply which one
perhaps had not at the outset thought of.
Theoretical resignation is conditioned by the necessity of
procuring confirmation by experience for our ideas and yet of
% --- p. 28 ---
\opage{28}beginning with ideas. Here is a fundamental relation which cannot
fall away, but which perhaps even in the future will
appear in a sharper form than hitherto.

Also in practical respects one had to begin with
a surplus, with anticipations, which are rectified in many ways.
Without vigorous hopes and longings and without
great models nothing great can be accomplished. But
realization is reached only through friction and interaction with
the real conditions. Much that was at the outset regarded as
valuable must be given up, and much must be taken up and accepted which
at the outset was an object of aversion.

In theory the opposition between hypothesis and
confirmation by experience, in practice the opposition between ideal and
realization: these oppositions will always assert themselves
with more or less strength, according to the nature of the tasks.
They lie in the very constant conditions of life. And they must not
least assert themselves in philosophy, which is after all only the
clear reflection upon the principles of our theoretical and
practical striving. ``Skeptical tinges'' or a ``skeptical sting''
would therefore not come to be lacking. Hopefully then
not: for they are the forces that drive forward, without which
development would come to a standstill. ---

In the \emph{epistemological} domain it holds that the
last presuppositions of all our knowing must be of
a hypothetical nature. Such a presupposition is e.g.
the causal principle, the assumption that every phenomenon has a
cause. It lies at the basis of all inductive research, but
cannot itself be directly proved. We assume it because it makes
real science possible, but a proof can be given only in so
far as it can successively be shown that the phenomena
in fact arrange themselves in such regular sequences as
the causal principle demands. A proper proof it is not,
since the causal relation itself cannot be observed; what we
observe is only relations of succession. It is and remains then a
peculiarity of our knowledge that it from first to
last works with hypotheses. --- This would be only
skepticism if by hypothesis one understood an arbitrary
assertion or fancy. But in the hypotheses which lie
% --- p. 29 ---
\opage{29}at the basis of our knowledge, the very nature of our spirit expresses itself,
its innermost need, the need of unity and
connection, its striving not merely to be filled with a rich and
manifold content, but to be able to order and master it.
The question is only to what degree this need can be
satisfied.

Dr. Starcke seems to hold that the very fact that one still
in philosophy speaks of a problem of causation and of
the difficulty of grounding the causal principle bears witness that there
still sits too much skepticism in our bodies. ``Let us
suppose,'' he says, ``that this question of the validity of the law of causation
had not come down to us from a time when one still
did not believe in it; how would man then take up
his position toward this whole problem? Would it then not rather
be he who directed a doubt against the causal principle who
would have to prove that he had a right to doubt, than he who believed
in it, who would have to ground his faith?'' To this it must be answered
that when it is the task of philosophy to bring forth the hidden
presuppositions which condition the validity of our knowledge, then
it is indifferent on whom the burden of proof rests. The problem
is there and must be treated, whether the individual thinker
approaches it more from the negative or more from the
positive side. And the grounding of doubt is indeed a criticism
of faith, just as conversely, so it does not much matter
how one proceeds. Several of the older ``skeptics''
stand as forerunners of modern epistemology by
the way in which they have attacked faith in the
principles of knowledge.

When one speaks of a time that did not yet believe in
the causal principle, this is an inexact manner of expression.
There has hardly been any such time. The struggle for
existence drives one to seek causes in order to gain power
over things, and a need in that direction lies moreover in
the nature of human consciousness. Even in the wildest
superstition this need expresses itself. On the other hand there have been
times when one did not assume \emph{natural}
causes for all phenomena. The various stages of development in the intellectual
respect are characterized by what causes are regarded as valid.
% --- p. 30 ---
\opage{30}The miracle does not conflict with the causal principle, since it is
expressly conceived as the effect of a supernatural cause.
What the struggle was about in the beginning period of modern
science was whether there are \emph{natural} causes of everything that
happens, that is to say causes that themselves can be shown
as real phenomena. And the burden of proof fell here so much
the more on those who asserted natural causes, as
the demonstration of these causes is precisely the positive
task of science. It goes with the principle of natural causes as
it goes with the causal principle in general: only progressive
verification is possible.

Dr. Starcke maintains that because our knowledge is unfinished,
it is not therefore uncertain, and he holds that ``the skeptical
attitude'' in modern epistemology is connected with
a confusion of these two concepts. And yet he himself comes
in the same breath (p. 236) to state that the
finality of knowledge would only have the value of making its \emph{certainty}
final. Now it is after all clear that a certainty which is not
``final'' is no full certainty. It is one thing that we
are sure that there is only one way to go to win
knowledge, another how far this way can lead
us. This we do not know as long as we have not gone the way
to the end. One must always venture and hope; the distance between
ideal and reality in the domain of knowledge will not
vanish. In that the old skeptics are right. ---

Besides the problem of the causal principle the author
also discusses the problem of the relation of knowledge to things
and of the relation between soul and body, in order here too to
hit ``the skeptical trait that runs through the
epistemology of our time.'' I cannot here go more closely into
these points, although it would have its interest. On the
whole I can moreover agree with what the author develops
concerning the questions mentioned. Only I believe that it is a
``tinge'' of spiritualism that makes him lay down that
every state of consciousness must be the state of an atom (although,
as he seeks to show, all our states of consciousness
need not therefore be states of the same atom). The atom
is a concept that has been constructed in order to make visualizable to oneself
% --- p. 31 ---
\opage{31}what goes on in the chemical and physical processes.
There are after all eminent physicists who would prefer to be
free of it and to keep to describing the movements
or activities themselves. In any case
natural science has nothing to do with and knows nothing of what
goes on in the interior of the atom. It is the relation between the atoms,
their rearrangements and mutual attraction and repulsion,
that plays a role in the explanation of nature. And it is these
same processes to which experience shows us the
phenomena of consciousness linked. The phenomena of consciousness are then
probably connected with what goes on between the atoms,
with their rearrangements and oscillations, not with mystical
states in the atoms ``themselves.'' --- Yet this is something of which only
a longer development could give full clarity. ---

Not only in the domain of epistemology, but also in that of \emph{ethics}
skepticism has, according to Dr. Starcke, still too much
scope.

Rousseau and Kant had, as already mentioned, asserted
the possibility of an ethics on a purely human ground. But
they did not, in Dr. Starcke's conception, overcome
skepticism. This expresses itself in them --- and later in the whole
rationalistic tendency --- in a disharmony between the individual's
nature and the ethical demand, a disharmony that becomes so
much the stronger as, according to Kant, the ethical demand
springs from the very nature of man. The faith that the
ethical will endure and triumph in the world can then be held fast only
when one takes theology to help and in the form of religious faith
lays down what one cannot ground by way of knowledge.
Kant set up, as is well known, in the form of religious
postulates on ethical grounds assumptions which he in his Critique of Reason
had deprived of every theoretical foundation. ``One believes,'' as
Dr. Starcke expresses it, ``in a God governing the world,
because only thus can one safeguard one's trust in the power
of morality.'' It is the lack of a theoretical foundation that gives
this view a skeptical character.

But also in other modern ethical tendencies the author finds
a similar disharmony and skepticism. Both
the morality of the doctrine of evolution and the morality of welfare have according to his
% --- p. 32 ---
\opage{32}conception skepticism in them still. They ground morality from
the standpoint of the race. But why shall the individual sacrifice himself
for the race? What motive shall keep him from, as far as
possible, rooting out or counteracting the inherited dispositions in
the direction of sympathy and conscience which he may possess?
With what right is it demanded of him that he give up what
he regards as his own happiness, in order that at some time in the future
other individuals may enjoy their happiness? --- It is of no use
here to refer to sympathy. ``It is indeed precisely the
sympathetic impulses that one conceives as the great deception
which Nature practices on man, and by which it gets
him fooled into serving its ends, in that man under the
influence of these impulses believes he lives for his happiness
and yet lives only for the race with sacrifice of himself.''
There is then required a proof, which has not yet been given, that ``on
the whole the moral life, which is lived in society
and the family, is that which makes it possible for the individual to bring
his sufferings and renunciations down to a minimum and at the
cheapest price to get the greatest possible security for
personal gain.'' And --- what is the main point --- this
proof must be carried out from the standpoint of the single, \emph{isolated} individual:
``The demands of social life must be formulated
precisely and exhaustively, and these demands must be capable of being appropriated
by the isolated individual as his own demand for happiness.''
The grounding must thus be egoistic in the sense that it
shall be shown that when the individual will follow his well-understood
interest (his egoism in the wider sense of the word) it
does not pay him to be an egoist in the narrower
sense of the word, that is to say by satisfying himself at the expense of
others.

No one is further than Dr. Starcke from playing with words
or hunting for ``witty'' turns of phrase. But I wonder
that he in his development of thought uses the ambiguous word
egoism, as he does, notwithstanding that he expressly points
out its double meaning. That this should make
full clarity possible I am unable to see. But about
words we would in any case not quarrel. --- Yet there is,
% --- p. 33 ---
\opage{33}apart
from usage of language, a good deal which makes me unable to
appropriate the author's train of thought in one and all. In the
first place I must remark that it would have been desirable if
he had given his critique of the morality of evolution and
the morality of welfare a definite address and backed it with definite
citations. In a book of especially historical character it would
be natural. In what precedes the author has carefully
named even several very insignificant literary phenomena
which seemed to him to fall under the concept of skepticism;
but here we get not a name, not a reference. Therefore
there becomes something vague about the whole development. It gets
the character of an estimate, and if I am to set against this estimate
my own, then I must say that his complaint does not seem
to me to be historically justified. One cannot say that the
problem of the relation between individual and society has been
overlooked by the ethical tendencies the author here criticizes.
It has on the contrary right up to the most recent time been discussed
with great zeal, among others in \emph{Henry Sidgwick}'s ``Methods of Ethics''
and the discussions which this excellent work has given
occasion to.

But the main point is how matters stand with
this very problem, as the author sets it up.

For my part I avow myself to be a
``skeptic'' in ethics, by Dr. Starcke's measure, in that I
do not believe that the account can be settled between individual and
society in the manner he demands. There is something
irrational about the first principles of ethics, in so far as the basis
from which these are laid down and acknowledged is not that
on which all individuals to whom these principles are to be
applied already in reality stand. The ethical judgments,
the judgments about good and evil, of which the individual becomes
the object, cannot always be passed from his own
present standpoint. The goods which ethics (the
practical ethics) protects are not always goods in the single
individual's eyes, and he can often only through a long
education (perhaps discipline), through a long development of character
learn to regard them as goods. A proof cannot always be
given to him just as he stands and walks. In the case of children, the
% --- p. 34 ---
\opage{34}insane and criminals this is especially evident; them we subject
to a practical treatment before we trust in the power of
reasonings. But also in individuals other than those who
must be reckoned among the classes mentioned, the ground must often
be prepared before rational ethics can begin its work. What
ethics lays down is then here the means and ways, the pedagogical
system (in the widest sense), which is to be applied in order that
a properly ethical treatment can take effect.

The pure, consistent egoist could at the most
be brought to regard the welfare of other men as a \emph{means}
for his own. But he would then continue to stand in a purely
external and mechanical relation to social life, and to all
life outside his own. And is it to be supposed that a proof could be given
that social life, as it is at the given moment,
presents harmony between the egoistic (``isolated'')
interests of all individuals? Even if there could, whom would such a
proof really interest? Not the real egoist; for him it is
enough if society is merely a means for \emph{his} interests,
--- no matter how things then go with the other egoists;
indeed, perhaps he does best by others not being able to
settle the account in the same way.

Can there at all be \emph{isolated} individuals?
On that everything depends in this connection. In every
man there stir various feelings, which drive toward
a valuation of life and its relations. How many of them
shall the individual divest himself of in order to become isolated? Is it
in and of itself settled that sympathy and the feelings related to it
shall be laid aside when the account is to be settled? --- Man
exists only as a historically given being, of a definite stock
and race, with definite dispositions and interests, and cannot
arbitrarily withdraw from this, nor avoid
having his valuation determined thereby. By
isolation, if it were possible, he would precisely deprive himself of some of the most
necessary items in the account. Let me develop this
a little more closely.

There is harmony of interests between two men where they
can regard each other as means, and each for himself stands
as end. Thus e.g. between the dealer and
% --- p. 35 ---
\opage{35}his customer. The former needs the latter's money, the latter the former's
wares. There was harmony of interests between the tailor Mustafa
in Ispahan and his friend of the same trade, for as Morgiane
says:

\begin{verse}
``The other often praised your father's cloaks; \\
In return his breeches were praised again. \\
Ah, that was the best time of the tailors' guild!''
\end{verse}
% „Den Anden roste tit din Faders Kapper;
% Til Gengæld blev hans Bukser rost igen.
% Ak, det var Skræderlavets bedste Tid!“

To a great extent society subsists by means of
such a harmony of interests. That it is not present in
full measure and in the desirable degree is seen from the constant
complaints about society and from there being a ``social
question.'' One cannot therefore appeal to it
in the given, present moment. A work is required
to bring it about, and a special
proof must be given that this work can be carried out merely by virtue of
the motives that can be found in the ``isolated'' individual,
and by virtue of what such a one conceives as his ``well-understood
interest.''

Dr. Starcke himself assumes that the account cannot \emph{yet}
be made to balance: ``It is an advantage for society,''
he says, ``that there are men who set their virtue
above their advantage; but it is the fault of society that there can
be found such a disproportion between virtue and happiness, and its
task must be to see to rectifying that fault.'' The whole
question must then stand open until the fault is rectified. But what
motives is one then to refer to, in order that individuals may
be brought to work at getting the fault rectified? It might
well be that the full harmony of interests (the harmony between
the interests of ``isolated'' individuals) were established only late in the future,
and then Dr. Starcke's theory would be subject
to the same criticism he himself directs against the morality of evolution,
in that it would demand sacrifice and work for a happiness that
fell to a future race, but not to ourselves. The ``isolated''
individuals would then perhaps answer that since their
descendants in remoter generations have done nothing for them, they will
likewise do nothing for these descendants. Après nous
le déluge! ---

% --- p. 36 ---
\opage{36}But even if such a harmony of interests should some day
be brought about completely, so that everyone could
regard all others as means for himself, then thereby
neither an ideal society, nor the most inward and strongest
kind of social life would have been attained. A vigorous society
exists only when in living together and cooperation there develops
fellow-feeling and real sympathy, so that the one does not
separate his cause from the other's, does not merely regard the
other as means, but also as end. Only under this
presupposition will the account referred to be capable of being made to
balance. It is then no longer settled in an external,
mechanical way. For in fellow-feeling and sympathy itself there is had
one of the greatest goods, about which one has not merely to ask
what it is useful for, but whose usefulness is wholly immediate
(immanent), in that it fills the mind with one of the most inward
feelings of pleasure. The individual's own being is strengthened and widened
through devotion; the horizon is widened, and new possibilities
for acting, enjoying, and suffering show themselves; the whole of life is felt in
a richer and deeper way. In every relation of love this
holds true, most strongly in the narrowest, most
immediate and personal, in decreasing degrees in larger circles, but
always with a reflected glow from those. There opens a new
point of view for all other feelings and faculties: all the feelings and
faculties that lead in the same direction as the feeling of fellowship and
sympathy receive a sanction they before lacked, an accent they
could not have as long as they were valued merely from the
``isolated'' individual's lonely turning about its own axis. --- This
rotation on an axis, incidentally, does not cease because the whole globe
now shows itself to stand under the influence of more deeply lying
world forces.

All this will fall away if the individual merely
regards himself as ultimate end. Joseph Butler used
in his time the expression that ``when we settle down in a cool
hour'' (when we sit down in a cool hour), we could not
justify our striving after the right and good, if
we are not convinced that it will serve our happiness
or at least not hinder it. He uses this
assertion to demonstrate the necessity of ethics
% --- p. 37 ---
\opage{37}in the end leaning on religion and its promises of
the beyond (as later Kant and in our day
Sidgwick). But his standpoint is in reality the same
as that which comes forward in our author's demand for
a grounding from the standpoint of the isolated individual.
Toward Butler one must raise the question why
the cool hour should here be right and the hour of
enthusiasm, where the warmth of devotion and self-forgetfulness rules,
necessarily wrong? It must be remembered that here it is not
a matter of logic and mathematics, not of philosophy about life, but
of experience and ascertainment of the real elements of life itself.
If ``coolness'' leads to some of the most essential
elements not being able to assert themselves, then the account
cannot possibly come out right.

But I concede, and it is even an essential part
of my ethical theory, that here the ways may part, without
agreement being attainable for the time being by theoretical means\footnote{In the 3rd chapter of my \emph{Etik} (1887) I had motivated this standpoint more closely. In the first chapter of my \emph{Etiske Undersøgelser}
(1891) I took the question up for fresh treatment. In the second
edition of \emph{Etik} (1897) an attempt has again been made at several points % print: Punker (misprint)
 to attain
the requisite clarity with regard to this point. (Later note.)}. It
is possible to concentrate in oneself, to regard this
self-concentration as one and all, and from this
point of view to carry through a consistent ethical system (which I
should, however, rather call individualism than egoism). In
favor of the tendency that makes sympathy the foundation of
the ethical system it may well be adduced that thereby one is led by a
more direct way to the harmony of interests than when one
starts from the separated, isolated individuals. But a compelling
proof this is not.

Nor does it avail here to emphasize what was
shown above, that the individual's life becomes richer and deeper
through fellowship and sympathy, thus precisely by his
not making himself the sole end. Of that the individual can
convince himself only by the way of practical experience.
And experience here means nothing less than a
% --- p. 38 ---
\opage{38}change of character. He who surrenders himself to the sympathetic
impulses and lets his egoism be damped thereby gets a wholly new
point of view for his valuation, so that he comes to draw
the consequences of his experiment from a wholly different
foundation than the one he began from.

There will therefore always be room for skepticism with
regard to the possibility of a rational ethics. Since the
psychological and historical foundation is different, the
systems that are built upon it must also become different, even if they
all take their starting point in real human nature.
This is in process of becoming and development and comprises manifold
and often conflicting elements, of which now one, now
another thrusts itself forward with special force. Whether the historical
development will some day bring complete unity in the
foundation of ethics, I shall not say anything about; but before that
``the skeptical sting'' in any case cannot be wholly drawn out. Ethical
thinking therefore does not lose its significance. It shows
us the light in which actions and forms of life appear,
seen from a definite stage of development, a definite fundamental motive given in
real human nature. Thereby it illuminates
the possibilities that lie in human nature, and reacts
upon the very foundation that is built upon; through the
consequences we first learn to know our nature aright.
Also a presentation from the individualistic
standpoint asserted by Dr. Starcke will have its great interest. Such a one
has basically not been seriously attempted yet, whereas we
have analogous attempts in theological ethics,
political economy, and jurisprudence. It would be extraordinarily interesting
to see in what points it meets with presentations from
other starting points. A serious thinking always brings
its fruit. From the rich treasure of psychological and historical experience
we all draw; but precisely because it is so rich, one can
gain access to it from different sides.

The individual investigator must regard himself as a miner
who makes his shaft into the mountain. When he feels
lonely during this work, he rejoices at hearing
his comrade's hammer in the neighboring shaft. He knows that the
other is working toward the same goal, although he has chosen
% --- p. 39 ---
\opage{39}another starting point. During my occupation with Dr.
Starcke's book this feeling of comradeship has not left
me, however much criticism I have had to let come to
utterance. It is the best of the much good that I take with me from
his book\footnote{In his book ``Samvittighedslivet'' (1894---1897) Dr. Starcke has later
given a detailed presentation of his ethics. (Later note.)}.

\streg{}

% --- p. 40 ---
\opage{40}
\essay[\emph{On Vitalism}]{\emph{ON VITALISM}\footnote{A lecture delivered in the Biological Society on 27 January 1898.}.}

\streg{}

It is more than a figurative expression when one speaks
of thoughts living and dying. Like all organic beings,
thoughts have their determinate conditions, on which their
arising, their development or their disappearance depends. Thoughts
must, like organisms, struggle in order to live. Only those
thoughts live in the long run which can subsist with, and be the expression of,
large circles of well-tested and coherent experience. And just as a natural
selection takes place among the organic germs, tentative beginnings and
variations which are present at every point of development, so too a choice is
constantly made among the thoughts and points of view that are available at any
given time. There is a reciprocal relation between thinking and experience. For
the involuntarily formed thoughts and points of view among which experience
chooses have of course not themselves arisen without the influence of
experience. But just as in nature far more germs and variations arise than are
selected for continued development, so thinking too begins with energetic and
luxuriant interpolations and anticipations, which are gradually sifted, limited,
developed or discarded.

This holds especially as regards the understanding of organic life. The struggle
between speculation and verification moves in rhythmic fashion through the
history of biology, and in the most recent period a swell seems to be astir
which is leading to the taking up again of points of view whose time seemed
decidedly to be past. Life has been thought to be found in old thoughts
% --- p. 41 ---
\opage{41}whose decease had already been announced. In what sense that opinion is
justified, we shall examine more closely.

There are two fundamental conceptions which here stand over against one another.
The one rests on this, that the clearest understanding we can have of a
phenomenon is the one we have where we can resolve it into its elements and see
it again as the product or result of those elements' cooperation --- that is,
when we go from the parts to the whole. The so-called mechanical conception of nature
maintains that this mode of explanation holds not only for the physical and
chemical phenomena but for all natural phenomena. Against this it is urged from
the other side that what is peculiar to the organic phenomena is that the
constitution, the form and the mode of operation of the elements, of the parts,
cannot be understood except by taking as one's ground the thought of the whole
which they form, or are to form. The organism cannot be explained as a mere
product of organs or cells, nor the cell as a mere product of molecules. We must
then here assume a peculiar force, operating according to other laws than the
mechanical ones. This is the vitalistic conception. --- And since the only way in
which we can conceive a whole as determining the constitution and the
cooperation of the parts is to conceive it as a purpose that is striven after,
a vitalistic conception will always remain more or less teleological.

\streg{}

When history shows us a rhythmic alternation of these opposed ways of seeing,
the advance of knowledge is by no means thereby excluded. It is the new
experiences that let the phenomena appear now more from the one side, now more
from the other, and that also gradually transform the earlier forms of the two
fundamental conceptions.

A mechanical conception of nature first comes forward at a relatively advanced
standpoint. Primitive man's conception is vitalistic or animistic, explaining
all phenomena that awaken attention by the intervention of personal beings. All
later vitalism has a certain kinship with this primitive conception. --- In
Greek antiquity
% --- p. 42 ---
\opage{42}the standpoint of mechanism is asserted by the \emph{Atomists}
(whose doctrine was later set forth in Lucretius' famous poem de rerum
natura).
\emph{Aristotle}, who was himself a brilliant investigator of nature in the
domain of organic life, maintained, in opposition to atomism, a conception which
regarded nature as a great scale of stages, where every member is matter and
means in relation to a higher one, while it is form and purpose in relation to a
lower. Nowhere in nature could he conceive the whole as arising as the result of
the elements, the form as the effect of the matter. He therefore extends
vitalism to the whole of nature.

It is thus characteristic of the position of the problem of life in antiquity
that no sharp boundary was set between
inorganic and organic phenomena. For the Aristotelian conception as for the
atomistic there are in nature only differences of degree.

For the time being the Aristotelian conception was victorious, and it became the
prevailing one in the Middle Ages. A turning point comes with \emph{the founding
of modern natural science} by Leonardo da Vinci, Kepler and Galilei. The
proposition was now set up that all change is motion, and it was demanded that
this be carried through in every domain, so that all laws of nature might be
derived from the principles of the mechanical theory of motion. \emph{Cartesian
philosophy} generalized this conception and carried it out into wide circles. The discovery
of the circulation of the blood and of reflex movement was an empirical support
for the conviction that the domain of organic life could form no exception.

As a reaction against this whole tendency there arises the conception one now
chiefly thinks of when one speaks of vitalism. \emph{G. E. Stahl} combated
mechanical physiology and maintained (especially in Theoria medica vera
1708) that without a soul's intervention organic growth and the organism's
self-preservation were inexplicable. He thus goes further back than to
Aristotle, in that he renews the primitive animism. He held that since the
conscious soul can be the cause of voluntary movements, it must also be able to
be the cause of the involuntary processes in the organism, even if it here
operates in a more instinctive manner, and even if it does not
% --- p. 43 ---
\opage{43}always behave rationally --- which, after all, it does not always do in its
voluntary actions either. --- The well-known medical school at Montpellier
modified the theory by distinguishing between soul and life, so that between
soul and body a peculiar life-principle was inserted. In this form vitalism is
presented by \emph{Barthez} (Nouveaux Elements de la science de
l'homme. Montpellier 1778). He combats two sects, the Mechanists and the
Animists, and represents vitalism in the narrowest sense of the word. He defines
the life-principle as ``the cause which brings forth all the phenomena of life
in the human body.'' It follows laws that lie above the laws of physics
(d'un ordre transscendant par rapport aux lois
de la Physique). Its
proper nature, and especially its relation to the soul proper, is unknown. But
Barthez plainly enough conceives it on the analogy of the life of the soul proper;
he ascribes to it instinct, temperament and sensory capacity. It intervenes in the
life-process and directs it in a determinate direction. It lays in the muscles
the force needed to overcome a determinate resistance, and shows its power
especially in convulsive attacks. It dilates the muscle fibers in the ducts of
milk and of semen, and it acts upon the humors, so that the substances are
formed which are needed for repair or for reproduction. It regulates respiration
and the development of heat, and conditions a ``sympathetic'' connection between
the organism's different parts. --- Different both from Barthez's vitalism and
from Stahl's is the theory which the famous anatomist \emph{Bichat} set up (in
Recherches physiologiques sur la vie et la mort.
Paris 1800). He
conceives the vital force as a property of the organic tissues. It was an
incipient decentralization of vitalism.

\emph{Classical vitalism} --- so I designate the conception represented, in
different shadings, by Stahl, Barthez and Bichat --- has this in common with the
Aristotelian, that by assuming a vital force it takes itself to have given a
sufficient explanation of the peculiarity of the phenomena of life. It differs
from Aristotelian vitalism in sharply emphasizing the opposition between the
inorganic and the organic, and even assuming a hostile relation
% --- p. 44 ---
\opage{44}between them. The soul, or the vital force, must struggle with the mechanical
and chemical conditions in order to get them to operate in a determinate
direction; indeed Bichat even defines life as the sum total of the forces that
struggle against death.

A new and magnificent rise of the mechanical conception of nature in its application
to biology came about on the basis of the work of Italian, French and Swiss
investigators of nature at the close of the eighteenth century --- work which,
however, only came to be properly recognized towards the middle of our own
century. About the year 1840 \emph{Dumas} and \emph{Liebig} reformed plant
physiology and demonstrated the circulation of matter in nature, and with it the
connection between organic and inorganic nature. Soon after, \emph{Robert Mayer}
discovered the law of the conservation of energy and applied it at once to the
organic domain (in his work Die organische Bewegung 1845). A
proposition from that work may stand as the motto of the whole new biological
conception: ``During the process of life there takes place only a transformation of
matter and of force, but never a creation of the one or the other.'' A whole
series of eminent physiologists, German ones above all (notably the disciples of
Johannes Müller --- though the master himself was still a vitalist), showed the
fruitfulness of the new points of view. Physiology now seemed able to be
conceived as applied physics and chemistry. The assumption of a vital force was
stamped as mysticism and superstition. \emph{Helmholtz} uses (in his work
Das Denken in der Medicin 1878) vitalism as the opposite of natural
inquiry, and regards the question of the vital force as one with the question of
the perpetuum mobile. And in his memorial address on Helmholtz (1897)
\emph{Dubois Reymond} declared that the doctrine of the conservation of energy
gave the phantom of a vital force its final blow. The period in biology which
began about 1840 is one of the most fruitful in its whole history, especially
through its many applications of physical and chemical laws to the organic
domain. A great and coherent conception of nature might hereby seem to have been won
once for all.

In the latest years one has begun to speak again of vitalism as a necessary
standpoint. Connected with this in a certain way is the fact that one speaks of
natural science
% --- p. 45 ---
\opage{45}having disappointed the expectations that had been placed in it --- indeed, that one
has even spoken of ``the bankruptcy of science.'' The facts on which so-called
\emph{neo-vitalism} rests were in the main already known earlier, and were
emphasized by critical investigators. Men like Lotze and Virchow, Littré and
Claude Bernard, who all maintained that only the methods of physics and
chemistry could lead to insight into the laws of vital phenomena, nevertheless
by no means held that life's innermost nature could be fully understood from
those laws alone. They regarded life as a fact, and saw it as biology's task to
investigate the single members of the life-process as it actually runs. The
occasion of neo-vitalism's appearance can be sought only in this, that many
parts of the life-process have shown themselves to be \emph{still} more
intricate than had been thought beforehand. There are three points especially
that come into the foreground for neo-vitalism.

1) Even if at every single transition in the life-process there occurred only
such transformations of matter and force as could be explained by chemical and
physical laws, there would still remain the fact that the chemical and physical
processes in growth and nutrition run in determinate \emph{directions}. This is
all the more striking in that growth takes place sporadically, from different
points, so that the joining into a form of the whole comes about only gradually.
Should it not then be justified to speak of a formative drive (nisus
formativus)?

2) In \emph{the taking up of nutriment} into the cells a choice is exercised.
Neither the quantity nor the quality of what is taken up can be explained in a
purely chemical or physical way. The cells seem here to act as little beings
whose need and wish determine their relation to the surrounding world. Several
neo-vitalists therefore personify the cells, as the classical vitalists
personified the whole organism. Since we know only from ourselves what need and
wish are, one of the most zealous neo-vitalists, \emph{Bunge}, has even declared
that psychology ought properly, here, to come to physiology's aid.

3) Organic tissue differs from the crystal in that its molecules \emph{never}
come into \emph{complete equilibrium}, as
% --- p. 46 ---
\opage{46}long as life lasts. To be able to die is therefore a mark of the living. ---

It is beyond doubt that attention to these points has been sharpened in recent
years. It is doubtful, on the other hand, whether a new stage has really thereby
emerged in the old rhythmic process, such that the name neo-vitalism could be
justified. How far one will extend the name may in any case be doubtful.
Investigators like \emph{Giulio Fano} (La fisiologia come scienza
autonoma. Genova 1884) and \emph{Hans Driesch} (Die Biologie als
selbständige Grundwissenschaft. Leipzig 1893), who wish only to maintain
biology's independence over against chemistry and physics, can be called
neo-vitalists only improperly. The name does fit \emph{G. Bunge}, on the other
hand (the chapter ``Vitalismus und Mechanismus'' in ``Lehrbuch der
physiologischen und pathologischen Chemie''. Leipzig 1887), who almost renews
Stahl's animism. From standpoints very different from one another,
\emph{Haeckel}, \emph{Verworn} and \emph{Rindfleisch} have expressed themselves
in a similar direction\footnote{The significant difference between the
conceptions that are often designated together as neo-vitalism was already
pointed out by \emph{Felice Tocco} in his essay ``Le disfatte della scienza''
(Nuova Anthologia. 1 March 1896).}.

\streg{}

It will have its interest to compare neo-vitalism with classical vitalism, in
order to see whether it is, as some hold (for example \emph{Mosso} in his essay
``Materialismo e Misticismo'' in Nuova Anthologia 1 December 1895), a sign of
retrogression, or whether, in relation to earlier and more or less kindred
tendencies, differences can be shown which point to advance. Such a comparison
can be led back to three principal points.

1) Classical vitalism holds that with the assumption of a ``vital force'' a
conclusive, positive explanation has been reached. But for the most critical
investigators, whom neo-vitalism claims
% --- p. 47 ---
\opage{47}for itself, it is merely a matter of maintaining that biology has its independent
foundation in experience, and is something other and more than a sum of physical and chemical
deductions. In so far as a vital force is spoken of, the word means an $x$, an
unresolved remainder, and stands as the expression of the composite conditions
under which the general laws of nature find application in the organic domain.

2) Classical vitalism was called by Claude Bernard ``a doctrine that makes
lazy'' (une doctrine paresseuse). In its confidence in the vital force
it disdained exact investigation of the particular relations. Barthez held, for
example, that it was unnecessary and unreasonable to seek an anatomical basis
for the connection between sensory and motor functions; the unity of the vital
force was sufficient explanation of that connection. He thus rejected in advance
investigations such as those which later led Charles Bell to his discovery of
the difference between sensory and motor nerves. (See Barthez: Nouveaux
Éléments. p. 142---153; 207 f.)\footnote{This point was already brought out by
\emph{Louis Peisse} in his critique of the Montpellier school. (La
Médicine et les Medecins. Paris 1857. I. p. 291---293).}. Neo-vitalism,
by contrast, holds fast --- despite its inclination to psychological analogies ---
that the only method which can actually be in question for use in biology is the
one the mechanical conception has proclaimed. It regards the problem of life as
more intricate than classical vitalism did.

3) The most striking difference between neo-vitalism and classical vitalism
consists in this, that the vital force is not, or not only, ascribed to the
organism as a whole, but to the cells, the organism's elements. In place of a
genius directing the organism's life as a whole, one is now inclined to assume a
kind of cell-geniuses, which decide what and how much is to be taken up from the
inner medium. \emph{Bunge} says outright: ``Wir übertragen das aus dem
eigenen Bewusstsein geschöpfte auf die Objecte unserer Sinneswahrnehmung, auf \emph{die
Organe}, auf \emph{die Gewebselemente}, auf \emph{jede kleine Zelle}.'' It
undoubtedly marks an advance that the
% --- p. 48 ---
\opage{48}concept of force is used of the single cooperating elements instead of
en bloc of the whole organism. The problem has become more localized; the
riddles have been pushed further back. To be sure, the great problem then
remains, how the life of the many cells is combined into the organism's
collective life --- and it was this above all that awakened the wonder whose
theoretical expression the old vitalism was. Neo-vitalism too is the expression
of a wonder, but especially of the one awakened by what is mysterious in the
processes that go on in the organism's smallest parts. One had not expected that
the problem of life would come again at the very elements into which one thought
one had resolved the whole of life. ---

The difference between neo-vitalism and classical vitalism is so great that it
is really unwarranted to say that it is one and the same conception that has
come again. With respect to what is scientifically asserted, and to the method
actually employed, no change has really taken place in relation to the
immediately preceding period in the history of biology. All talk of retrogression
(Mosso) or of bankruptcy (Brunetière) is therefore unwarranted, and the name
neo-vitalism itself had best be avoided.

\streg{}

In conclusion it will be of interest to bring out the analogy between the
problem which makes biology something other and more than applied chemistry and
physics, and a whole series of other problems which all rest on the opposition
between the rational and the empirical --- or, as one may also put it, between
the formal and the real in our knowledge. The more rational and formal our
knowledge is, the further it can move forward through definitions and
deductions, so that every transition takes place with intuitable necessity.
Empirical or real knowledge, by contrast, must often stop at facts which can
indeed be described and analyzed, but not defined and derived from other facts.
Such a fact is life; and therefore biology belongs, and will perhaps always
belong,
% --- p. 49 ---
\opage{49}to empirical knowledge. It was not only Stahl who in his day missed a
definition of what life is (a lack which he did, however, believe himself able
to supply); it has been said in our own day as well (by Lotze and by Littré, by
Bernard and by Fano) that no such definition has been found. But life does not
stand alone in this respect. At a whole series of points
the rational continuity in our knowledge is replaced by empirical data which for
the time being stand as unresolvable. They are kinks in the thread, or perhaps
even knots.

1) Let us suppose that the proposition that all change is motion could be
completely carried through, and that the ideal of the mechanical conception of nature
could so far be realized. The empirical kinks or knots in the thread would not
on that account be lacking. At every point --- however far back we go --- the
motion would still appear with a certain given direction and a certain given
velocity, and the outer causes which alter that direction and that velocity
would operate from certain determinate places and in certain determinate degrees
and directions. --- The sum of energy expended in the motion, and in the causes
that alter it, is moreover of a determinate given magnitude, and its distribution
among the single operative elements can likewise be known only empirically.

2) Energy appears, further, in qualitatively different forms, and the operative
elements likewise have their qualitative peculiarities. Even if, by the law of
the conservation of energy and of matter, equivalents can always be shown at
every transformation from one form of energy into another and at every change of
matter, there arise at the same time in these transformations new properties
which cannot be derived from the antecedent conditions, but are attached to them
purely empirically.

3) In the organic domain the sum of empirical elements becomes greater still,
and more pronounced. It is attention to these elements that often brings with it
a renewed inclination towards vitalistic views. As such elements we have in the
foregoing mentioned direction of growth and structure, qualitative and
quantitative choice in the cells' change of matter, and the maintenance of the
unstable equilibrium.

4) If we take the psychological domain along with the rest,
% --- p. 50 ---
\opage{50}the phenomena of consciousness ---
from the most elementary sensation to the highest thoughts, feelings and acts of
will --- stand as an empirical given which cannot be derived from the material
presuppositions. Psychology is to a still higher degree than the other parts of
biology a purely descriptive science. And within psychology there then come
forward again new qualitative differences and transformations, where the same
problem presents itself anew.

It is clear that if, at all the points where the empirical, the qualitative and
the individual put to the test the rational continuity our knowledge strives
after, one were to assume new independent and original ``forces,'' one would be
deceiving oneself. The limits we thus constantly come up against are always at
the same time to be regarded as new material, setting new tasks. It is
existence's richness and depth that produces these kinks and knots in the
thread, and thereby makes necessary a constant interaction % print: Vekselvirking (misprint)
 between experience
and thinking, induction and deduction, analysis and synthesis --- that doubleness
in method which Goethe compared with the relation between breathing in and
breathing out. The struggle between mechanism and ``vitalism'' is, rightly
understood, a struggle between two tendencies between which there can at bottom
be no opposition. We should not be more perfect beings if we could breathe out
without breathing in. That doubleness makes both life and thought the richer.

\streg{}

% --- p. 51 ---
\opage{51}
\essay[The Physiology of the Emotions]{THE PHYSIOLOGY OF THE EMOTIONS.}

\begin{center}\emph{C. Lange:} Om Sindsbevægelser. En psykofysiologisk Studie. 1885.\end{center}

\streg{}

It is an old prejudice that whatever stands as beautiful,
glorious, or sublime to our feeling must lose its luster
and its value by having its coming-into-being, development, and
continuance explained according to simple natural laws. It was
a scandal to the Greeks of antiquity when philosophers and
students of nature began to say that the sun and the moon were
material bodies like the earth; hitherto they had regarded
them as gods, who move across the heavens majestically and by
supernatural powers, to the admiration and benefit of the
inhabitants of the earth. Attempts at scientific explanation have very
often met with the same prejudice and the same resistance again.
But how strong must the resistance become when it is
the life of feeling itself that one wishes to explain in a
scientific manner? To many it will seem as though that were to kill
feeling at a stroke, --- as though to explain it were
one with explaining it away. The present work, which
aims at a purely physiological investigation and explanation of
the emotions, will hardly escape being met with such
prejudices. And yet it is certain that it only applies
more closely and carries through a conception that has long been
asserted in scientific psychology and physiology.

By emotions the author understands \emph{the simplest, least
composite phenomena of feeling}, which form the constituents
of which the proper feelings and passions consist.
Such emotions are grief, joy, fear, anger,
tension, etc. They recur as elements in such more
composite feelings as love, hatred, admiration,
% --- p. 52 ---
\opage{52}contempt, etc. This division agrees on the whole with that
which modern psychology takes as its basis, as it is indeed also
the most natural procedure to proceed from the
elementary to the more composite. One might, however, perhaps raise
objections to the way in which the author in particular
draws the boundary between the two groups of phenomena of feeling.
Wonder ought also to have been counted among the simple feelings that are
called emotions. There is an element of
wonder every time a sense-impression comes to consciousness,
since it is a condition for the arising of any sensation
that the state which the impression calls forth stands in a certain
contrast to the preceding state. The more an impression
differs from the preceding one, and the more suddenly it
comes, the stronger becomes the element of wonder that
asserts itself, and which can form the introduction to very
different states of mind (fear, disappointment, contempt, and
love, etc.). There is so much the more reason to
emphasize this, as the boundary which the author wishes to maintain between the two
groups of phenomena of feeling can hardly be upheld
unless one adds to the distinguishing marks he has given
yet another, namely that the emotions have more the character
of sudden outbursts, whereas the proper feelings or
passions are firmer, more lasting, more ingrained states.

These emotions the author now wishes to study by means
of \emph{the physiological method}. He had originally posed his
problem by asking about the effects of the emotions
on the bodily expressions of life; but he came to the
conclusion that this was an entirely wrong way of posing the
question. The emotion is not one thing, its bodily
effect another. We have, physiologically viewed, not here a
principal phenomenon, ``the'' emotion itself, and a secondary phenomenon,
its effects on the organism. The physiologist cannot
concede that these effects are something subordinate; for
him it is precisely they that make up the whole emotion. One
proceeds from a preconceived and unfounded distinction between
soul and body when one assigns the emotion to
the soul and only after it has arisen there lets it act
on the body. Physiologically viewed we can only study the
% --- p. 53 ---
\opage{53}emotion through what goes on in the body; an
assertion that this is caused by something that goes on in
the soul is unscientific and unprovable. We must in any
case see how far we can get by means of mere
physiology.

The author shows this by a physiological description of some
of the most important emotions. This part of his book
is written with true mastery. One notices the practiced
clinical observer, and one comes to think that the
penetrating psychological sense --- which makes itself felt in the midst of
the physiological characterization --- is a heritage in the author's
family, now already in the third generation. Here I can only
communicate some basic features of the full exposition.

In \emph{grief} the capacity for voluntary movement is weakened.
Hence the mourner's tired appearance, bowed posture, and
faltering gait. The involuntary muscles, on the other hand, contract,
especially those found in the walls of the blood vessels and
set in activity by nerves that issue from the
``vessel-moving'' (vasomotor) center in the brain. The blood
is thereby pressed out of the fine vessels, so that the tissues and
organs of the body become poor in blood; hence follow pallor,
sensations of cold, and slackness of the features. The secretion of tears
seems explicable as a consequence of a relaxation of the
vessel muscles after the first contraction; therefore weeping is
a relief and fails to occur in the most violent grief. In a similar
way the sigh seems to be the natural reaction against
the shortness of breath produced by the contraction of the fine vessels of the lungs.
--- In \emph{joy}, by contrast, the capacity for
voluntary movement is heightened and the fine blood vessels dilate. A sensation
of lightness, warmth, and strength arises. And whereas the
bloodlessness of the brain in grief shows itself in mental weariness and
dullness, so the abundant blood supply of the brain in
joy shows itself in the flight of thought and imagination. --- \emph{Terror}
stands physiologically close to grief, only that the symptoms are stronger,
and that not merely the vessel muscles but also all other
involuntary muscles seem to contract. --- \emph{Anger} stands
in about the same relation to joy as terror to grief.

Of these various symptoms the author now regards
% --- p. 54 ---
\opage{54}\emph{the constriction of the blood vessels} as the most important, the primary, of which
the physiological consequences are all the others. To be sure, this
assumption cannot yet be completely proved, and the author puts
it forward as a hypothesis that sets future research a
significant task. But there are nevertheless important facts
that for the time being speak for it. In particular one knows that the brain,
spinal cord, and the whole nervous system are affected to a high degree
in their capacity to function by the change in their abundance of blood.
The disturbance in the vascular system will in that way be able to gain
influence on all the more important functions of the body. Individuals
whose vasomotor system is easily set in motion show
themselves also especially inclined to emotions.

When it is now to be shown that in the phenomena
here described we have the whole nature of the emotion before us, the
author lays particular weight on the fact that \emph{the same states as are
as a rule designated as purely mental feelings can be called forth
by purely bodily influences}, when these merely in a certain
definite way gain influence on the vasomotor apparatus.
Wine makes one glad, just as much as good news; anger can
be brought about by the consumption of fly agaric; ipecacuanha produces
a depressed mood resembling grief or fear; anxiety
can arise in nervous diseases, etc. To this belong
also such cases where outbursts of feeling, e.g. laughter,
are called forth by causes that do not seem to fit this
effect\footnote{The author has not seen this phenomenon mentioned by other authors, but
mentions two cases from his own practice. In my Psychology
(2nd ed. p. 337, 4th ed. p. 324), I have collected a number of examples of it.
--- I may perhaps be permitted to remark at once
that I have taken the kinship emphasized here by the author between the
feelings called forth by bodily causes and those called forth by ``mental'' causes
as the basis of my presentation of the psychology of feelings, in that I proceed
from the simplest forms to the more composite. I follow e.g.
fear (anxiety) through its forms, as a mere expression of an
organic state (through impeded respiration or cardialgia), as
motivated by sense-impressions (instinctive starting at sudden
impressions) and as motivated by ideas % print: Forestilinger (misprint)
 (that is, by central
processes in the brain).}. On the other hand ``mental''
emotions can be relieved or removed by bodily means.

% --- p. 55 ---
\opage{55}Here the author is certainly right that it is unjustified
to call the emotions called forth by bodily means merely
apparent, those called forth by mental means true
emotions. Why should the former not be just as true
as the latter? True is only what really \emph{is}, that is to say here
what is \emph{felt}. A joy or a grief can in and of itself be neither
true nor false. A feeling is present or is not
present; the ideas the individual attaches to it can
be true and false, but to the feeling itself this
distinction has no application. The feeling can have a more or
less solid foundation. The joy and the sympathetic
mood awakened by wine or opium do not have the
sound and firm ground of that which is due to a happy
temperament or self-acquired clarity and ability or
happy conditions of life independent of the favor of chance and of the moment.
The feeling of happiness produced by opium is afterward replaced
(when it has been cultivated for a longer time) by the most
dreadful anguish, whereas the feeling of happiness produced in the natural way
can last throughout life. Only in this
respect can one find the standard for an appraisal of
the emotions.

But I believe all the same that the author overrates the
physiological way of viewing things. First and foremost we know
the emotions from \emph{direct, subjective observation}, and without
this the author would not have been able to give his presentation.
His meaning must really be this: what we in ourselves
feel when we say that we grieve would, to an external observer
who could see all that went on in our nerves and muscles,
appear as tension of some muscles and slackening of
others; anything else the external observer cannot, by the nature of the case,
catch sight of. But such an observer always infers
something he does not \emph{see}, but \emph{imagines} in \emph{analogy} with certain
experiences he has had in himself. Only if he knows grief from \emph{his own
subjective experience} can he \emph{as a clinical
observer} see the grief in the physiological phenomena he has before him.
Otherwise they would be for him merely indifferent twitchings in
an indifferent stuff. This inference from analogy is as a rule made easily
and quickly, so that we most often do not become clearly conscious of it; but
% --- p. 56 ---
\opage{56}it is of decisive importance in theory to call
attention to it.

One can easily err in this interpretation of the
physiological phenomena. Should the author not, for instance, be mistaken
when from the agreement between some important
symptoms he concludes that anger stands nearer to joy than to
grief? Most people would surely call anger a feeling of displeasure
and for that reason place it nearer to grief than to joy.
It is only in so far as anger is connected with preliminary
steps toward revenge that a feeling of pleasure comes to play
an essential role in anger. Anger is grief, but revenge is
pleasure, says \emph{Aristotle} already; revenge makes anger
cease, in that it calls forth pleasure in place of grief. It is
indeed also clear that no one becomes angry unless a
displeasure or a grief has been inflicted on him. Anger differs from grief in that
it is of an active, not a passive character, and thereby it approaches
joy, to be sure. --- However uncertain subjective
observation may be, it is nevertheless a necessary
presupposition.

As to \emph{the relation between that which presents itself
to subjective observation and that which the external observer can
see}, I am fully prepared to go along with the author's
conception, which does not permit one to assign the emotion
to the soul and its effect to the body. I regard
every appeal to the intervention of the soul in the physiological
connection as unscientific. It is general theoretical
considerations that have led me to assert \emph{the physiological
continuity}. For if the process that goes on in
the nervous system and brain were interrupted at some points in order to
propagate itself over into the soul and later resumed again by
the soul's reaction, then we should here have a breach of the law
of the conservation of physical energy, in that first
physical energy would perish and later physical energy be created.
Therefore I have adopted the hypothesis according to which
the phenomena of consciousness, as subjective observation
teaches us to know them, are inner mental forms of the
same activity that for external observation appears
as a material phenomenon. In this way both
% --- p. 57 ---
\opage{57}psychology and physiology come into their full right. It is
really a very old assumption that we here embrace.
``The student of nature and the philosopher,'' says Aristotle already, ``will give
different definitions of every [mental] phenomenon, e.g.
of anger. The one will namely say that it is a desire
to retaliate for an unpleasant feeling or something similar;
the other will say: it is a boiling of the blood or of the
heat around the heart.'' Aristotle here expresses by
means of the physiology of his time a theory that seeks to let both
the physiological and the psychological side of the phenomenon
come into their full right. In more recent times Spinoza above all others
has asserted this view, which in our days has so
many adherents. ---

The assertion of \emph{the physiological continuity} is a doctrine
that can stand independently of the special assumption that
the author maintains: that \emph{the essential thing about the emotion is precisely
the organic effects,} in which the usual conception
sees only secondary phenomena, --- that \emph{emotion is thus really only
sensation of the organic effects}. An American
psychologist, \emph{William James}, has in the journal ``Mind'' (April
1884) put forward a similar theory. For it to prove
its truth, it would have to be shown that the emotion does not
arise \emph{before} the organic effects are sensed. This would
agree well with the fact that feelings of pleasure and displeasure on the whole seem
to need longer time for their appearance than sensations
and ideas. But the condition for those organic
effects being able to occur is that there takes place a radiation
in the brain from the centers (for sense-perception or idea)
that awaken the emotion. Only when this radiation
reaches certain other centers (according to the author the vasomotor ones),
does emotion arise. Might there not now in consciousness
be something that makes itself felt corresponding to this
radiation? And would it not be justified to distinguish
between two stages within the course of the emotion, in such a way
that the one corresponds to what goes on in the brain during that
process of radiation, the other to the way in which this
brain process is strengthened and modified by the impression of the organic
effects on the brain? Here there would then take place
% --- p. 58 ---
\opage{58}a circuit: from the brain out into the body and back to the brain.
In its full peculiarity we have the emotion only
when this reaction back is completed. But it could
nevertheless be that a feeling of pleasure or displeasure arose already before
the effect propagates itself beyond the region of the brain, even
if the time difference is only very slight. Conversely, after a strong emotion
the organic effects can sometimes still continue, although the mind is otherwise at rest. Here
the second stage thus continues while the first is
concluded. So for instance when we ``cannot recover''
from a fright or from an outburst of laughter.

To the idea of \emph{physiological continuity} and to the theory
of \emph{the organic reflexes as the principal physiological phenomenon
in emotion} there comes, further, in the author his special
theory of \emph{the vasomotor reflex as the primary one}, from
which all other physiological effects are derived and
conditioned. The same theory has recently emerged from another
quarter, namely in a work on the physiology of fear by the Italian
\emph{Mosso} (La peur. Paris 1886). Such an assumption lies close to
the physiology of our time, partly because so many facts
establish the influence of the vascular innervation on the functions of the other nerves,
whereas a converse influence cannot be demonstrated,
partly because in recent years central organs have been demonstrated
in the brain for the nerves of the vessel muscles, not only in the medulla oblongata,
but according to some in the cerebral cortex itself,
thus among the highest nerve centers. If this
assumption continues to find confirmation, the physiology of the emotions
will round itself off in a beautiful and perspicuous way.
Of the prospects for this only physiologists by
profession can have a well-grounded opinion, and I shall therefore not
go further into it. I will only remark that even if it
should have ever so great importance for the physiological symptoms of emotion
that the movement within the brain reaches
the vasomotor center, it is not thereby
excluded that there takes place an equally immediate influence on
the centers for the nervus vagus situated in the medulla oblongata,
whereby the strong movement in the brain at once gains
influence on the functions of the heart and thereby again on the
% --- p. 59 ---
\opage{59}circulation of the blood, first and foremost on the blood supply to the
brain itself. \emph{Claude Bernard} laid the main weight on this
circumstance in his treatise on the physiology of the heart (reprinted
in the collection of his treatises published under the title
``La Science Expérimentale''). Also by this route
it would be the change in the blood supply to the different
organs that became the most important physiological phenomenon in
emotion. In my Psychology I have laid the main weight
on the point emphasized by Bernard, but have nevertheless also
emphasized the significance of the vasomotor centers, and I shall,
without needing to change anything in my psychological theory,
be able to make use of the author's beautiful presentation, if it
continues to find confirmation. ---

\emph{The future of the emotions} looks gloomy according to the author's view.
Both the education and the cultivation of the race and of individuals
aim at pressing the emotions back, at inhibiting the organic reflexes
to which they are attached. Will this not then lead, through inheritance, to the
vascular innervation becoming so sluggish that in the end there can no longer
be any question of emotions? To this is added that the life of the understanding
plays an ever greater role in the development of culture.
The life of the understanding and the life of feeling stand in opposition to each other.
They both draw on the brain's blood supply, and the more
the one kind of brain activity consumes, the less
is left over for the other. Will not in the end
the life of the understanding lay claim to the whole, so that the
time of feelings is past? --- With this pessimistic result the
author concludes his book.

I cannot believe otherwise than that he means this only
as a warning, --- a warning that we too greatly
need. It is quite certainly a great one-sidedness in our whole
modern culture that it has predominantly developed the
faculties of the understanding. It is its strength, but also its defect.
The consequence is, I believe, not so much a weakening
of the life of feeling as a whole, as a disharmony between our
knowledge and our feeling. Our thoughts embrace far more
than our feelings can as yet manage to fill and to warm. Our
knowledge goes all too often in one direction, our faith in another.
% --- p. 60 ---
\opage{60}Hence the wavering and divided characters. The whole of life is no
longer, as formerly, directed with all its energy toward one goal
that was able to satisfy all mental needs.
In order that the life of feeling may be energetic enough, sympathy must
be quite otherwise lively and full than it has
yet shown itself. The tendency toward individualism and egoism, which
the modern development of the understanding cannot be said to be free from
having favored, is what narrows feeling. The right
education will aim at repressing the egoistic
emotions, but on the other hand seeking to favor the sympathetic ones
and to lead them along the most beneficial paths. It can only be a
great misunderstanding if the race should think that it had now
come so far that it could dispense with the emotions.
These will always show themselves as the living forces that drive
the work and overcome resistance.

The mutual interaction between the development of the understanding
and the development of feeling will certainly change
the character of the emotions more and more. Their elementary vehemence and
violence will be replaced by milder forms. This has to some extent
happened, and is it always a loss? Is it a loss that
human beings now know erotic love in place of the animal's brutal instinct,
or that hunger and thirst for justice can replace the
blind, ruthless instinct of revenge? It is quite certainly not
always that the gain is so easy to point out, and a warning
therefore always comes with a certain justification. We must in any
case remember that the more numerous and the more comprehensive the objects
that lay claim to our feeling, the greater the energy that must
be able to be applied, if life is not to be weakened. The
more various expenses we have, the greater the total capital
we must command, if we are not to fall short. ---

The progress of psychological research takes place through contributions
from very many different sides. The author has from his
standpoint made a significant contribution, significant not
only by the result he has arrived at in it, but also
by the thoughts it sets in motion in various directions.
And it must further be added that his book also by its
attractive form and its beautiful language deserves to win
many readers.

\streg{}

% --- p. 61 ---
\opage{61}
\essay[Frederik Christian Sibbern]{FREDERIK CHRISTIAN SIBBERN.}

\begin{center}1785---1872.\end{center}

\streg{}

There are eminent men who seem to have succeeded completely
in giving themselves in their works. Sibbern does not belong
to their number. What he himself often complained of, the
readers of his scientific writings have also found to be true:
that he had difficulty in shaping and expressing what stood
before his inner eye. This was probably connected in part with
an imperfection in his scientific standpoint, which, despite his
constantly renewed labor, he never got beyond; but it was also
connected with the fact that his personality was so rich, warm,
and peculiar, also peculiarly angular, that it could not find
vent through a single form. The thinker, the poet, the man gripped
by religion, the cordial friend of humanity, and the enthusiastic
social prophet were different selves in him and each demanded its
own kind of satisfaction. He worked throughout his life with these
different powers, of which now one, now another pushed itself
forward into the foreground; the unity, the connection between
these different endeavors he has left to posterity to grasp, if
it is able to. And an attempt must be made at it, while the
memory of him still lives among us.

He does not belong to those who have strongly and thoroughly
influenced the course of the spiritual and public life of our
fatherland. But he is one of the soundest and at the same time
deepest human beings who have lived here at home. And he
lived through his time fully and independently, not merely taking
into himself what it offered, great and small, but taking it all
in his own peculiar way, putting in his word,
% --- p. 62 ---
\opage{62}privately or publicly, and by his warm sympathy and
penetrating understanding he has been extraordinarily much
to so many of the best men here at home, as he himself
must be counted among the best men of our country. In him the whole
spiritual and public life of Denmark from 1820 to 1870
is mirrored in an instructive way. Instructive --- not only because he
lived it with warm participation, but also because
he himself was always one who strove, who right up to his advanced old age
went forward in the search for truth. His youth and first
development fell in the period of reaction against the eighteenth
century. He felt keenly the quickening effect of the
fullness of feeling and imagination that replaced the
``enlightenment'' of the understanding and the dominion of the dry morality of reason.
But he took only what he could use, and he could only
use what had a true human kernel in it; the shells
he left lying. It was therefore not a relapse, but
the result of a natural and sound course of development, when he
later put forward religious and political views which, by
their freedom and radicalism, astonished and offended many
of his earlier friends, and would have astonished still
more, had they become better known in wider circles.
It was his living sense for what is living that made him
break through every form that the spirit of the age set up as
alone saving, whether it was in the domain of religion, of
politics, or of philosophy. There is an inner
connection between his polemic against the Hegelian philosophy,
his liberation from all dogmatic faith, and his combating
of the superstition about paper constitutions and legal
formalities.

By his concluding religious and social ideas he meets
the freer endeavors in these domains in the most recent
times. By his living sense for real life, its
conditions and its demands, he meets the best in the modern
realistic spiritual tendency. But with him everything has grown out of
the deep mind's need to seek out and immerse itself in life
under all its forms, on all its stages. One can hardly
point to external influences that should have called forth and
determined his development in the later part of his life. He
% --- p. 63 ---
\opage{63}was prepared for the Revolution precisely by fully and wholly
living through the Reaction. And he \emph{grew} from the Reaction
into the Revolution without any noticeable break and without sudden
catastrophe. He, who himself always warned against plucking
the unripe fruit of the tree of life, also let his own fruit
have the necessary time. Unfortunately, in the meantime evening
came for himself; he was now unable to win attention and recognition for what was
the soundly developed result of his life and striving. As long as
he himself lived, it came mostly to stand as an old man's strange
ideas, and his old friends and admirers preferred to behave as
if they knew nothing of them. Perhaps he did not himself wish it otherwise.
He was the man of quiet development and an enemy of all haste. All the more encouraging
and instructive is his image for the younger generation, which has
the task of carrying the thread further, where his old hand
had to let go of it. It calls upon them both to go
forward and to lay a deep and firm foundation for
the work of progress. When in our days we see so many
unsavory examples of old champions of freedom --- often without themselves
wanting to be aware of it --- denying the ideals of their youth,
this must come from the fact that the establishment of these ideals has
been hasty work without a solid foundation. It is
work, the quiet, assiduous, and patient work, first and
foremost the work upon oneself, that has been lacking. But there
has also been lacking the living sense for, and with it the
living faith in, life as that which never stands still. The
form under which one first conceived of life, one demands
that it shall put on in all ages. The older generation
thereby comes to weigh like a nightmare upon the younger
generation, one social stratum upon another, instead of
paving the way for its successor and helping him forward. But
life is always right; that was precisely one of Sibbern's
principal propositions.

Now that I here wish to pay off a debt of gratitude
by attempting a characterization of F. C. Sibbern, I
must begin by depicting his personality, as far
as it may succeed in gathering the scattered traits about certain
% --- p. 64 ---
\opage{64}main points. We shall thereby see how his nature
predestined him for his work. What he has accomplished
as psychologist, as philosopher, as religious and social thinker,
will only be rightly understood when we first --- to use one of his
own favorite expressions --- seek out ``the inner source'' from which
it all has welled up.

\section*{I.\\  Personal Characterization.}

We come to know Sibbern through his letters, both the
real ones and those he published under the name of Gabrielis's
Letters. The picture of him that can thus be gained is confirmed
by his other writings, but contains at the same time several
traits which serve for a better understanding of these.

A turning point in his youth was the long
stay in Germany from 1811 to 1814, from which he returned
home to take over a professorship in philosophy at the
University. We follow him on the journey through his letters, and we
come to know his mood and condition in the time immediately after
his return through the first part of Gabrielis's
Letters, which was written 1813---1814, although it was only
published in 1826.

He had gone abroad to study philosophy and to train himself
as an academic teacher. He made the acquaintance of Fichte, Schleiermacher,
Steffens, and Schelling and was especially strongly and
lastingly influenced by the two last named. He sought out Goethe, who
stood for him as an exalted genius and was the object of
his cult, indeed almost of his infatuation. He undertook
walking tours through the most beautiful parts of Germany, lived
with mountains, forests, and lakes and was refreshed by the beauty of nature.
But amid all the new and glorious things he saw and
came to know, his heart clung to home. He clung
to the Danish nature, to the city of his fathers, whose great
misfortunes he had lived through, and where he hoped that a new
and significant period for science and art had
been inaugurated. But there was One thing that bound him more firmly than
anything else to home, namely his love for Sofie
% --- p. 65 ---
\opage{65}Ørsted, to whom the most numerous and most interesting of his
travel letters are addressed. ``No one should ever travel,'' he writes
to her, ``before the fatherland has bound him with
the firmest, strongest bonds.'' --- In the gallery in Berlin he met
a young girl, with whom he often talked, and whose image
he preserved and often dwelt upon in his thoughts. ``But that
I can do this \emph{with calm}, for that I thank my fatherland,
thank it heartily; it let me early behold and grasp
high beauty, and it thereby bestowed on me the happy
gift: everywhere to be able to rejoice in every beauty I meet
on my way, to rejoice undisturbed by desire and craving,''

From his youth he was much in A. S. Ørsted's house. For
the great jurist he felt an extraordinary admiration
and a warm devotion, which he has often given beautiful
and cordial expression, most beautifully perhaps in his ``political
Intelligence Leaves,'' on the occasion of Ørsted's speech at the opening
of the assembly of the consultative Estates in Roskilde. He
was himself Ørsted's disciple, since his first studies were
juridical. Ørsted's wife, Oehlenschläger's beautiful and gifted
sister, gathered no less than her husband a large circle
about her. Sibbern loved her with all the warmth and
enthusiasm of his nature. Even in his old age he was moved
in a conversation with Sofie Ørsted's biographer, Nikolaj Bøgh, by
the memory of the love of his youth. He walked with
her and read to her in the Odyssey, Sophocles, Plato, and
Goethe. He was all joy and admiration at her
beauty and at her lively interest and clear gaze.
He had gone abroad in jubilant joy over what he had here
found. It did not belong to him, but the joy over it
belonged to him.

Already from the travel letters, however, one sees that it was not
quite so easy for him merely to take the relationship from this side.
Sofie Ørsted did not love him; she treated him as a younger
brother; and besides she was attached to a man she esteemed
and loved as highly as anyone else. Despair at
the hopelessness of his love had to break forth, and it
was only when Sofie Ørsted was dead that he could say that
``it was sweet and joyful merely to know that she existed, and
% --- p. 66 ---
\opage{66}to be known by her.'' There were times when his heart
demanded more than that. In one of his last travel letters to
her he mentions how in Breslau, where he
visited Steffens, he had felt overwhelmed by his pain and
had found relief in weeping with his whole heart. ``And
truly it is,'' he adds, ``as if there, with my
dear, precious countryman, I first found the peace and
rest that Orestes once found in the goddess's grove, in the pure
sister's nearness.'' --- There seems here to be an allusion to a
struggle between the flaring-up desire and the glad feeling, unaffected by all
craving, for what he had in her whom he
loved. He has fought his way to the resolution which
he lets his Gabrielis express: ``If there is anything I definitely
will, it is this, not to see her again until
I know that I am calm and will be able to rejoice without craving
in the contemplation.''

``The hour I return home,'' he had written, ``it
will bring me the flower of all joys.'' But how it
turned out for him he lets his other self, Gabrielis, describe to the
friend to whom he writes. One will find in it some
turns of phrase which have already been quoted from one of the real
letters to Sofie Ørsted.

``You know how long I was away, but you do not know
how faithfully I clung to her. All my longings
belonged to her and my home, and a thousand times I said to
myself: Happy one! yes, you can travel; for the fatherland holds you fast by
the most beautiful bonds, and whatever meets you there, the most glorious that
the earth has to give, home gave you. ---
Alas! and also the most bitter. I came back; my feelings
still blazed in me with unconsumed intensity and strength.
I found also in her the accustomed kindness toward me.
Inexpressibly I thanked her for receiving me so
kindly, so gladly. I lived again some joyful days. But
it soon grew fainter. No, I lost what I had only now
hoped to win. --- With each day I grew poorer.
In an endless passion's, I might say jealousy's never
resting anxieties my feelings were consumed, and what I
still possessed, from that my bewildered conduct soon parted me.''

% --- p. 67 ---
\opage{67}It was the first great struggle he had to take up with
life. He lived through its whole bitterness, although it
often seemed to him as if his life were wasted. It was a
heavy winter for him, the first after his return home, the one
in which ``Gabrielis'' was written. He writes himself to his
sister on the 5th of March 1814: ``I must confess to you,
I have not yet spent any winter with such dejected,
often bitter feelings as this one.'' Besides the
cause on which we have here dwelt, and which he does not mention
to his sister, there were two other causes. --- He had
taken over the philosophical teaching post, but felt that things
did not go for him in his new position as he wished. He
struggled, as we already see from the travel letters, with doubt about
himself. He feared that nothing much would come of
him, felt weary, forsaken of spirit; everything within him seemed to
come to a standstill. In the period of forsakenness and dejection that had
come over him, this doubt and feeling of impotence had to
grow, and had at the same time to be felt all the more strongly, the more
he needed to immerse himself in activity in order to forget
his pain, and the greater the demands his new position made
upon him. --- Finally it was also this winter in which
the unhappy Seven Years' War was concluded with the equally
unhappy Peace of Kiel. He had experienced at close quarters
the outbreak of the War of Liberation in Germany and the joy over the victories;
the wretchedness and misfortune here at home stood in contrast
to it in all the more glaring a light. ``While the whole
rest of Europe except France can break out in loud jubilation,
we alone must suffer in so shameful a way!''

Such were the experiences through which Sibbern's
character first appears to us. When there now appears in his
letters, in his writings, and in so many oral utterances
of his, as well as in his whole bearing, such
a cheerful joy in life, an unshakable optimism, which sucks
bliss out of everything, it is already clear from what we have now
mentioned that this joy in life and optimism was not
built upon sand, but was won through struggle with much
and deeply felt opposition. He lets his Gabrielis console himself
with the thought that the sea of tears that pressed for a
% --- p. 68 ---
\opage{68}breakthrough in him would one day be followed by a sea of light. Thus
for Sibbern there is always a deep and strong shadow that gives
the light its full effect. He rejoiced over the good time
because he knew the bad, and especially when he had
obeyed his teacher Goethe's admonition: not to work himself up
in the bad time.

In the crisis that has been mentioned we see him in particular struggling
to preserve his joy and admiration even there where
possession is impossible. Admiration and sympathy were on the whole
the feelings that had the deepest ground in Sibbern's nature.
Life was for him, if only he could subdue the
bitterness that could rise up in his own interior, essentially a
source of joy. In his 24th year he writes to his sister:
``Make for yourself out of what surrounds you a heaven full of joy!''
and even in his advanced old age he felt (as he writes in
a letter from the year before his death, kindly left to me)
``summeriness'' in himself. External and internal pressures and sufferings
could darken and embitter the mind; and pressure was in his
experience worse than suffering; --- but even in such times one
can, he maintained, procure for oneself glimpses of something brighter or a
``foretaste of the Muse's fragrance.'' The worst enemy of
the poetry of life, of the preservation of youth, he found in desire
that is attached to something definite and particular, and which begets
bitterness and pain when this cannot be attained. ``No, turn
to what is everywhere to be had. There is wonderful music of life
enough everywhere; one need only become quiet within oneself and listen.''
Thus he has Gabrielis (in the second part of the Letters)
answer the question how youth and the poetry of life
are preserved. And now follows the whole wonderfully beautiful
development, which everyone ought to know at first hand, of how
one can find joy in the contemplation of the gleam of sun on
the roof, of the sparrows hopping about, of the scenes in
the street, of human life in its various forms. One
shall live with ``the whole great life-movement of nature,'' ``make a
familiar acquaintance with winter as with summer, so
that one can feel within oneself the melodies they strike up for us
when they come to visit us.'' We understand here Sibbern's great
joy in walking tours. As he had tramped about in Germany
% --- p. 69 ---
\opage{69}on his great journey, he tramped in later years about in
Zealand and on Funen, and in Gabrielis's Letters there are found several
traits of what one might call the Epicureanism of walking tours,
the rocking in moods and thoughts that solitude in
field and forest favors, and the cozy rest at
evening. It was not merely nature he reveled in, and
whose small and great scenes he watched and often describes with
great freshness. It was, as already touched upon, also
human life. He found it beautiful when Lavater had published his
great work on human physiognomies ``for the promotion of
love of humanity.'' One shall love human beings as
the botanist loves plants: for him there is no weed.
One shall sit down by the craftsman, watch his doings
and talk with him about them. ``Are there not in the craftsman's
work, in the shopkeeper's bustle, and in all human doings
in markets and in streets great powers of life present?'' In
his later years he mentioned in particular cottagers and
country folk with predilection. He, who lived in our people's
finest and most nobly educated circles and himself had immediate
access to the treasures of science and poetry, ends his (posthumously
published) Moral Philosophy with the words: ``I could
adduce more than one example from the cottager class
in the country, likewise from servants, with whom living together can have
a most beneficial effect on the attentive.'' Also in other writings
he often returns to the cottager class in a manner
that shows that he speaks from personal acquaintance.

One sees: the poetry of life that he preaches and wants to have
eyes opened for is borne by a cordial and warm sympathy.
It is no fantasy-dallying as in so many forms of Romanticism,
and (I may add) in so many forms of ``Realism,'' which
in this respect are only disguised Romanticism. He sought
nourishment, not entertainment. Life in all its conflicts
has value for him, and that for its own sake. It was
his faith that the same life, whose smallest stirrings outside
him could have so great a worth in his eyes, also stirred
in himself, --- and that it was therefore that he could get an eye
for it. The gospel of the poetry of life is here, as will later
appear, closely connected with his philosophical view and
% --- p. 70 ---
\opage{70}his religious faith, as he indeed on the whole does not make the
great distinction between these three things that we poor modern
human beings are forced to.

It was not merely nature and human beings from which he drew nourishment
and refreshment. He also lived in the poets, in
the Bible, in old, pithy edifying writings, which he
still valued and recommended after dogmatic faith
had long since become wholly foreign to him. Favorite passages from
such authors he often had with him on excerpted sheets
or written down on slips for use on walking tours, in order really to
make them living and present to himself. Even the year
before his death he wrote down a number of such favorite passages,
especially from Schiller and Goethe, together with some of his own
favorite thoughts, for a young girl who was to go out as
a governess.

As said several times, the joy in life that was so
peculiar to Sibbern was not superficial and frivolous; it was
not merely often dearly bought, but it was to him a serious matter,
a matter of heart and will. He had, as we have seen,
times when the pain itself was the only thing that held the mind
together and prevented its bursting. ``Nothing joyful,''
says Gabrielis, ``had I that could give my existence
continuity; so pain and despair became that which
held my existence together and preserved me from giving
my soul up to disintegration.'' There lies a sound and
manly thought in these words. Spiritual pain has, like
all pain, a tendency to dissolve our soul. But it
is itself, where it announces itself, a necessary link in
the series of our inner experiences, and to push it away, to
want to repress it at any price, is to give oneself up
to a spiritual morphinism and is the sure road to
ruin. Pain must be lived through, if one cannot
change the conditions that call it forth. Only thereby does the mind grow
and harden. It was this that Sibbern again and again, right
to the last, impressed, that ``what hinders, that
furthers, when it is taken rightly.'' He knew from his own experience
what Goethe calls the ``secretly formative power'' of pains.
One must fight one's way forward with endurance and
% --- p. 71 ---
\opage{71}patience from the bad time to the good; only so does the
contrast between light and darkness work in the right way. Sibbern's
view of life was essentially a faith in life, and the last
thought with him was probably this, that one must let life have
its course and also in oneself wait for the good times
to come. The bad times, especially the times of spiritual forsakenness, of which
he so often complained, he regarded as times in which
the mind lay fallow and could gather strength for new activity.
But he by no means paid homage to a quietism that led to
folding one's hands in one's lap. How energetic his
conception was on the contrary, can be seen from the fact that he often opposed
the scriptural passage which says that one cannot add a
cubit to one's stature. ``But an inch,'' he writes even
the year before his death, ``one can surely add to it, when one
keeps oneself brisk and erect.'' He did not want the limits
sharply staked out for how far one can get. And although the
highest for him was devotion to the life in us and outside
us, he did not conceive of this devotion as something merely
passive, but as the condition for the fullness of life being able to become
a power in us. It must go as one with our own willing.

To this it must be added that there was perhaps indeed something sentimental in
his nature, but that he did not belong to the gentle and meek.
He was a decidedly choleric temperament. A psychological
characterization of him would be very defective, if it did not
bring this out, both because it stamps the whole physiognomy
of the man, and because we have here another inhibiting and
dark power to contend with. It is known from several
anecdotes how he could flare up in violence. He is said once,
while Martensen was professor, to have berated him
at the University in the students' hearing, when they were in dispute
about a lecture hall; and yet Martensen was his good friend.
We have various self-confessions from him in this direction.
--- One of the most beautiful letters in the published collection is that
to Lecturer C. F. Molbech in Sorø, who had been seized by
a mental illness, which had altered his formerly loving and
gentle character and made him bitter and unmanageable. When
the worst time seemed to be over, but the sufferer still
had to contend with the after-swell, Sibbern wrote to
% --- p. 72 ---
\opage{72}him to help and strengthen him. ``Precisely I will write to
you, because I too know and from my own interior know what
all dwells in the abyss of our soul and can rise up out of it and
darken our mind. One I know especially, which is especially evil,
because it is most destructive for those who are nearest
to us .... This one is the bitter irascibility, the dark
discontent and dejection. Do not ask me what can
provoke it and nourish it in me. The least thing is
sometimes enough for it, when the mood is accordingly. But the cause
is to be sought early in my life. Already from my mother
I inherited it; and did not our whole childhood's education and
in our youth everything aim at inflaming us in an ambition and
a self-esteem, which at that time perhaps drove us forward, but
now halts and hampers?'' There was thus a constant source of
bitterness and anger, even with vanishing external causes.
But what Sibbern saw around him in public life often made
the inner fire flame up violently in him, in later
years especially the conduct of the National Liberal politicians and journalists.
He then reveled in fantasies of what he would
do to them, if he had Aladdin's lamp or could
make himself invisible. One could not wish to be in the place of
Lehmann, Ploug, and Bille, if these fantasies had
become reality. --- In ``Communications of the Contents of a
Writing from the Year 2135'' there is told of a book that had
appeared in the last quarter of the nineteenth century,
and which contained an old man's autobiography
of himself. In it the author had told how he,
without others having noticed it, had suffered from a
terrible disease of the soul, which he called the vengeful-irascibility sickness.
It arose only from what he read and pictured to himself in
his ideas. Often he could for that reason not bear
to read the daily papers. The sickness was nourished by everything that
he read or himself experienced of theological and political
fanaticism both in past and present; he had often thereby got
fits of irascibility, of which he had even felt bodily
effects. ``The book,'' it says, ``appeared only after his
death, and all declared that they had never thought that he
% --- p. 73 ---
\opage{73}could have borne such a vengeful-passion irascibility in himself.''
This man is plainly enough Sibbern himself.

This propensity, he knew well, was grounded in his nature;
but he emphasizes sharply and decidedly that it is nourished and cultivated
more or less in all of us by the way in which public affairs
are handled in State and in Church. In his
old age, as a spectator of political life, especially after
1864, it was often difficult for him to hold fast to what
he had let his Gabrielis pronounce, that there is really no
weed, not even in the human world. He made
the experience that ``in order to preserve one's goodwill toward human beings,
love of human beings, esteem for human beings, one must at times let it
take refuge under contempt for human beings and seek here a
place of shelter.'' There are thus, so it seems it must be
understood, certain excrescences one must reject, in order to be able
to preserve one's faith and love toward humanity as a whole.
But this kind of contempt, which can serve as a
preservative for love of humanity, suits only the old,
those who have gone through the school of life. In the younger
there must be found ``a differently strong contempt, bordering on hatred,
than that under which love of humanity is to seek
shelter.'' (Samfundsbetragtninger. Second Number, 1870.) It was
especially indignation over unfair judgments springing from party fanaticism and
ignorance that called forth such utterances
from him.

This man did not take life lightly. It was
also one of his principles, ``that one must take life
\emph{duly} to heart.'' He had enough to contend with, and that he
could bear witness to the value of life, as he has done, is a
significant fact. That his faith in life was won through a
struggle makes his result all the more weighty. And
finally he built upon an insight into the laws of the spiritual life
as his theoretical foundation. He was led early to the
conviction that the spiritual life, like all life, develops
sporadically, that is to say from different, often
apparently opposed points of departure. Very different dispositions
and impulses may express themselves at one and the same time, and as they draw
the mind each in its own direction, there arises the ferment, unrest, and
% --- p. 74 ---
\opage{74}pain that come so easily especially in the significant
years of youth in deeper natures. No one has described this with
such clear understanding as Sibbern. It was one of his
specialties both in theoretical and in practical
psychology. That he could be so much to so many younger people
in their time of ferment was owed not least to this insight.
He understood what was breaking within them and could instill in them
patience to let the different powers unfold in
calm and in confidence that out of the strife and unrest there would develop
a so much the richer harmony.

He himself did not feel happiest in his youth. He
looked forward to the time when the calm of evening should replace the day's
unrest and fussiness, and he wished fervently to become
an old man in order to be able to look upon life without ``ardor and craving.''
His wish was fulfilled. It rested, as is sufficiently
well known, on an illusion, so common in old
people, when in a letter from the year 1863 he wrote that he gave
his lectures with just as much intellectual power as before.
It was hard to get him to withdraw from the
University; he felt so hale that he could not
think otherwise than that it was his political opponents
who wanted to remove him. Yet for his own sake one ought
to have intervened more energetically. It was a scandal that
the excellent old man, whose senses were
weakened, and who could not assert his authority,
was exposed to the young listeners' saucy boyish pranks. But if he
thus had illusions with regard to his capacity for outward-going
activity, there was another respect in which
no illusion was possible, namely the life and freshness of feelings,
the ``summeriness,'' as he expressed it himself, which stayed
with him to the last.

This freshness of feelings was connected with his whole
childlike nature. A great, precious naïveté marked everything
he thought, felt, and said. It is characteristic of naïveté
to go straight at the kernel of the matter without any
formalities, with open eye for the real connection of life.
This naïveté naturally often had to have a comic effect. When
he, e.g., at a doctoral disputation, which in its time made
% --- p. 75 ---
% print p. 75: at »Gabrielis« (at for af) -- misprint; the sense is translated
\opage{75}a great stir, concluded by very cordially wishing the doctoral candidate
joy that he had now finally for once attained
that which he had so many times tried in vain,
the listeners' jubilation knew no bounds. It was not meant ironically.
For irony was not in Sibbern's line at all. He wrote in
the twenties an ``ironic'' defense of the English expedition to
Copenhagen in 1807; but half the readers took it in earnest,
and he had to explain publicly that it was meant ironically. ---
If this faculty was denied him, he had full compensation
for it in the immediate warmth and enthusiasm with which he
took life and strove to get others to take it.

Unfortunately Sibbern did not succeed in any single work
in giving himself wholly, nor in giving any single work
full completion as regards the content of thought. There is
no book by him that one can point to as that in which
one has him wholly and entirely. Perhaps some selected pieces
of ``Gabrielis,'' of ``Psychological Pathology'' (``The Doctrine of the
Human Feelings and Passions''), and not least of
``Communications of the Contents of a Writing from the Year 2135'' would
be best suited to characterize him from that side from
which he certainly has his greatest significance. For whatever
opinion one may have of him as a scientific philosopher, he
was in the finest sense of the word a philosopher of life, a man
who both lived richly and deeply and was able to give us
valuable contributions toward understanding our life and toward practicing the
difficult art: to live.

\section*{II.\\  Sibbern as Psychologist.}

The most numerous and greatest of Sibbern's scientific writings
treat psychology, and he brought to this subject in several
respects excellent qualifications. A fresh and lively
sense of nature and power of observation, especially for phenomena belonging
to the domain of feeling and impulse, and no small
knowledge of natural science and physiology. But his
psychological works (the first of which appeared in 1819) bear
strongly the stamp of the fact that at that time one was to a far higher degree than
% --- p. 76 ---
\opage{76}now unclear about the position of psychology and its relation to other
sciences. Sibbern's endeavor is indeed to
conceive psychology as an empirical science; but he
does not sufficiently separate it either from speculative
philosophy or from moral philosophy. Some of the most important
psychological concepts still stand with him encumbered with and
determined by elements that derive from speculative or
metaphysical assumptions, with which psychology in and of itself
has nothing to do. The merely observing and describing
method was not enough for Sibbern; he found in the most important
psychological observations such significant data that they
in his opinion could find their explanation only in the light
of a higher idea. He distinguishes indeed between ``a psychology
from the standpoint of observation'' and ``a speculative psychology,''
but after having acknowledged this distinction he adds that
``while in the domain of natural science one ought constantly to
keep close to experience, one cannot even in observing
psychology rise quickly enough to the higher
idea of reason in order to consider all of man's inner mental
or spiritual utterances in its light.'' Partly we are only
by starting from such an idea to understand how the
different powers and impulses can develop in accordance
with the conditions under which individuals live, partly
man's nature, his dispositions and faculties, are supposed to be so closely connected
with his moral and religious vocation that the
one cannot be understood without the other. Many parts of
his psychology might just as well be found in an ethics. Since
ethics ought always to be built on psychology, there is in and of
itself no fault in this; but the fault comes when one without
more ado makes ethical distinctions into psychological distinctions.
What in ethical respect stands miles apart may
yet psychologically belong very closely together. Therefore it is
of great importance that psychology preserve its
independence as an empirical science, independent of ethics and
speculation.

If there is a speculative philosophy, it can only
be in the sense that there are certain fundamental assumptions and
fundamental presuppositions that form the ultimate consequences of
% --- p. 77 ---
\opage{77}all our knowledge. It may have its significance to try to
bring them forth and shape them; thereby we first get to see clearly
what lies in our knowledge, and thus first thereby come
to complete self-consciousness. But whatever significance one
may wish to attribute to such speculations, which can never
become anything but hypotheses, the empirical science ---
and that the empirical science of spiritual phenomena
just as well as the empirical science of material
phenomena --- must as far as possible be kept independent of them. They must be
built on the consequences of the empirical sciences, not
the reverse. This was not how one conceived the matter at the
time in which Sibbern's philosophical development took place. One
wanted first to deduce and demonstrate the highest idea, which expressed
the innermost essence of existence, and then from it again to deduce, or
at any rate to view in its light, the special phenomena.
And in this highest idea one found at the same time the ultimate
goal toward which life must strive, if it is to reach
its proper perfection.

As a good disciple of Schelling and Steffens Sibbern stood
also essentially on this standpoint, and it has
hampered him in the clear, natural working out of
psychology. A single but very important example will
show this.

There is hardly any psychological concept in whose
treatment the different psychological views are better
characterized than in the concept of impulse. In the wider sense
of the word, impulse is any obscure urging, especially when it
goes in a definite direction. How are such
urgings now explained? They arise indeed often prior to all experience of
and enjoyment of that at which they aim! --- Sibbern's explanation is
this: ``Life wants to come forth in us and wants to express itself in us in the fullness
of effects in which it belongs to life to come forth,
and it cannot then do otherwise than carry with it and call forth
allusions to all that which is a condition for all its
effects.'' Here a speculative intimation is put in the place
of a real explanation. That life wants to come forth in us, we know
only from the involuntary manifestations of life themselves, which we call
natural urgings. But whence comes it that life wants to come forth in
% --- p. 78 ---
\opage{78}a certain definite direction? It is no answer to say that this
definite direction belongs to life. It does so quite
certainly, since every natural urging in a certain direction is a
manifestation of life. One concludes then: all manifestations of life, thus also all
natural urgings, must have their ground; now there are
natural urgings that aim at nourishment, others that aim
at the satisfaction of higher needs, such as living together with
other beings, or lead to moral actions, etc.;
therefore the general impulse, in accordance with which life
works its way forward, must appear partly as an impulse of nourishment, partly
as an impulse of communication, partly as a moral impulse, partly as
a religious impulse, etc. One thus endows human
nature with as many original impulses as there are
different directions in which life unfolds. Thus
Sibbern says: ``There is nothing that belongs to the content of life
or to the sum total of effects and forms
in which the perfected life expresses itself, without its being comprised
in the natural urging, and there is therefore nothing at which
man's life and striving aim, up to the very highest,
without a corresponding impulse being to be named.''

By this way of presenting the matter one turns
in a circle, in that one creates as many original
impulses as there are special directions of life. One stands here in
psychology at the same standpoint as one in physiology
stood at, so long as one assumed a peculiar vital force for
the explanation of that which constitutes the difference between the phenomena of organic
and inorganic nature. Only then would one
have got further than to such a circular course, if one
from the general idea of life had been able to deduce the special
forces and impulses in their mutual relations. But this
has always proved impossible, and Sibbern does not try
it either; he could not, as we shall later see, given his whole
conception of philosophy, try it either. When such a
deduction is now not possible, there remains only to
give a cataloging enumeration of the different impulses
with a constant reminder that they all have their place and
significance in life. Sibbern in his later writings
and editions especially goes strongly in for such a cataloging and
% --- p. 79 ---
\opage{79}rubricating, and it is for a large part this that gives his writings
something forbidding for many readers.

Nevertheless Sibbern's psychology has a naturalistic character.
That lies already in the way in which it lets life
work its way forward from its highest to its lowest stage under
% --- p. 80 ---
\opage{80}the form of impulse. It is seen also from the zeal with which he
made use of what the physiology of that time offered for the elucidation
of the life of consciousness. But one of the reasons for
the imperfection of the psychology of that day was precisely the unclear
standpoint at which physiology then still stood. It had not yet
at that time constituted itself as a mechanical natural science,
following the same method and recognizing the same general
principles as physics and chemistry. The doctrine of the ``vital force''
still furnished the general point of view for the organic
phenomena. Because of this indeterminate, half-mystical character of the leading
physiological principles the theory of the relation between
soul and body also could not take shape clearly and definitely.
The problem came to lack its real sting. Only when
physiology sets up the demand that \emph{all} phenomena in
the organism, in the brain and nervous system as well as in the circulatory and
nutritive systems, are to be explained according to the general
mechanical principles, does the difficulty of understanding
the connection between the life of consciousness and organic life stare one in the face.
Do the mechanical principles hold for this connection? If
not, have we not here a breach of the general laws of material nature?
It is precisely the doctrine of the persistence of physical energy
that in our time has anew given the old
question of the relation between soul and body a great
interest. Sibbern's psychological works were not yet
influenced by this point of view. Even his book ``On the Relation
between Soul and Body'' (1849) was composed at a
time when the works of Mayer, Joule, and Helmholz had not yet
gained influence on physiology. Nevertheless it is
interesting to see how closely his conception approaches
theories which in the most recent times are advanced by several
investigators on this point.

He declares himself against every dualistic conception and
lays at the foundation the conception that life is one as
life of consciousness and as material life. Only this unity of life makes
it possible to understand how certain
phenomena of consciousness can correspond to certain phenomena of the brain and vice versa:
life in the form of consciousness is according to its nature in
agreement with life in the form of matter, if it is one
% --- p. 81 ---
\opage{81}and the same life that appears under the two forms.
``Bodily activity and mental activity are essentially
connected \emph{coordinate} effects of one common cause acting simultaneously in
both.'' He lays special weight on the fact that bodily life exists only under an incessant
activity of reproduction, in that parts of the organic
tissue are constantly excreted and others are taken up and formed in their place. It is
really only an appearance when the body seems to have a more solid
basis for its existence than the soul. In both places we have
only an activity before us: the soul too reveals itself only in a constant production and working over of
ideas, feelings, and strivings.

Not all of Sibbern's utterances on this point are equally
clear and consistent, but the general conception he
points toward seems to be clear enough. To be sure, in this
theory he is more governed by the general interest in
maintaining the unity of life under all forms than in finding out
which hypothesis best satisfies the
most important results and leading thoughts of empirical science; but he has for his
time understood with great certainty how to grasp some of the
most essential points of view. He is among other things a decided opponent of
the spiritualistic doctrine of the soul as a substance or
a being by itself: ``Soul is only mental living.'' The
dogmatizing by which one fixes \emph{that} as being and resting
which is always given only in becoming and development, has contributed
much to making the psychological problems more difficult than
they are in reality. Such a fixing and
substantializing was moreover contrary to Sibbern's whole world-view.
``Our whole existence consists in being a function in
the all-organization.'' ``The individual activity and living is only
a function of the all-effecting.'' The individual issues from the
whole great activity of nature and cannot tear itself so
loose from this that it should become an independent whole
by itself. It stands in the same relation to the whole activity of nature
as the single idea to the whole activity of consciousness.
Thereby it also becomes, according to Sibbern's conception,
understandable how in the single individual something can stir
which is of universal nature: instincts and tendencies that point
% --- p. 82 ---
\opage{82}far beyond the individual's domain. But before I go more closely
into Sibbern's general philosophical world-view,
I will bring forward some points of his psychology which certainly
have lasting interest and significance.

His psychology clearly marks the reaction against the
rationalistic period by the weight that is laid on
``the heart'' in contrast to ``the head.'' He depicts under
various forms and turns of phrase the difference between the purely
intellectual, understanding-based assent to what one
knows and acknowledges, and the living devotion by which
we become One with the matter itself. ``That which goes beyond
reason,'' namely the powerful and warm feeling for that which
reason has grasped, belongs to what he depicts best. Many
sides of the psychology of feelings, especially the great role
that relations of contrast and mixture play here, he has in broad outline
clearly brought out\footnote{The significance of Sibbern's conception on this point is
now recognized also in foreign works (such as \emph{Ribot:} La Psychologie des
Sentiments. Paris 1896. p. 265), after it had been referred to in my Psykologi
(1882) and in A. Lehmann's ``Hovedlovene for det menneskelige
Følelsesliv'' (1892).}.

Connected with this is also the significance which he
attaches to the unconscious mental life, or, as he expresses himself:
``that which goes on in the deep ground of the soul.'' Only a slight part
of our nature steps clearly forth before our consciousness. Just as
we come into being at the start as individuals by means of a
natural activity of which we cannot become conscious, so
this natural activity unconscious to us must always be assumed
to continue, as long as our spiritual life lasts. There is
indeed also in the physical domain no absolute boundary
between generation and growth. In instinct and in genius we have
striking examples of something that works in us but only in
its effects, often remote effects, comes up to
consciousness. Our spiritual growth goes on predominantly
unconsciously, so that it is with us as with the man in the parable, who
``casts seed into the ground and sleeps and rises night and day, and
the seed grows and becomes tall, he knows not himself how.'' There
is constantly something that forms the foundation of consciousness without
% --- p. 83 ---
\opage{83}coming forward before consciousness; and there is a constant
reciprocal relation between the unconscious and the clearly conscious. It is
certainly a perfectly correct remark that Sibbern makes about
remorse (which he on the whole depicts well and without the
indeterministic dogmatism to which the older psychology was
inclined), --- that its psychological and ethical significance
rests predominantly on its influence upon the whole unconscious
mental life; man's character can through the feeling of remorse
and by its purifying power become another. One is reminded
here of Gabrielis's utterance that pain was the only thing that
held his consciousness together. Thus remorse can be
an unavoidable intermediate link, without which the development of character
cannot go forward.

What from this unconscious foundation now influences and
determines the life of consciousness announces itself often scattered from
several isolated points and in very different directions
at once. And likewise consciousness receives from the outer world
a multitude of different, mutually unconnected
observations and ideas. Finally by experiences
a multitude of different impulses, wishes, and resolutions are awakened. All
this can now fight in motley confusion in consciousness, and
that precisely the more, the richer the individual's nature is. There will then
come to be ``a certain stir of life, a ferment of life, a debate of life,
a struggle of life,'' before a definite center
and a harmonious order in the mind can form. ``Although strongest in the
youthful period of life (the age of young manhood), which for the rest in
different individuals can fall at different times,
it nevertheless persists, more or less, through the whole of life. There
are also individuals who in this respect remain as
constantly young, indeed some there are who even through the whole of life
remain so entangled in the ferment naturally connected in youth
with this struggle of life, that it is as if
the whole of life were for them a spiritual labor of birth.''

This is the doctrine of sporadic development, which in Sibbern's
world-view plays no less a role than his faith
in the unity and wholeness of life, and which he finds not least
clearly confirmed in the psychological domain. He rightly
points out that if development did not proceed
% --- p. 84 ---
\opage{84}sporadically, it would not proceed through struggle. His
observation and his train of thought lead him here straight
to the modern realistic doctrine of the struggle for existence.
This point of view lies so close to him that he emphatically impresses
(already in a treatise from the year 1829) that the proposition:
life is a struggle, must not be misunderstood, as if it were only
the external conditions that forced life to fight; no,
life exists only in the struggle, cannot subsist in and of itself
without it. By the energy with which he, not only in the
psychological domain, but, as we shall soon see, in his
world-view as a whole, holds fast to this point of view,
he has given interesting contributions to the general doctrine of
development.

\section*{III.\\  Sibbern's Philosophical World-View.}

Sibbern was strongly influenced by German philosophy, and
his doctrine is one of the soundest forms in which this
philosophy has gained influence on the course of thought in the Danish
world of spirit; but his sense for the concrete, his
interest in observation and his warm sympathy for individual
life made it impossible for him to become an adherent of the
abstract and a priori deducing and constructing in which the
older school ended. He took the field decidedly against Hegel,
and already Fichte and Schelling he took in his own way.
He decidedly dissociates himself from the Schellingian philosophy of nature.
In a letter from 1824 it says: ``It is advisable for every
student of nature to keep as near to experience as possible
and to be wary of every attempt to want, placing
oneself over on the higher standpoint, to elucidate and interpret
the things of nature by a consideration \emph{from above}: for to that
the considerations of nature are still far from being matured.'' As one recalls,
he did not carry this principle through in the case of psychology.
And yet it is characteristic of his whole
conception of philosophy.

He saw, to be sure, the highest ideal in a knowledge
by which we from a highest idea ground and explain
our whole knowledge of existence. A construction a priori
% --- p. 85 ---
\opage{85}was also for him the highest goal of philosophical striving.
But philosophy must, according to his conviction, always proceed
from a given content, and only through the working over of this
can the last, leading idea gradually develop.
Philosophy must, as he expresses it, be explicative before it
can become speculative, that is, it must proceed from
considerations of the given in order to find in this intimations
of an idea in whose light it might find its explanation.
And the speculative knowledge must itself again constantly seek
its confirmation through continual comparison (``Konferens'')
with the really given.

If we ask how it is possible that such an idea
can be found, and such an explanation to a certain degree attained,
since we indeed know and survey only so slight a part of
existence, Sibbern answers that we ourselves are indeed a member of
existence, and that the same that constitutes the innermost
essence of existence as a whole also asserts itself in us, when the ideas
in the consideration and working over of the given experiences
work their way forward in us. But that the goal will not
be completely attainable within a foreseeable time lies in the sporadic character of
existence. The world steps forth before us from the many
points of its fullness at once; the center around which all
these points can be seen in their connection can only
be found through ``a debate of all against all,'' which can hardly
be thought concluded. It is then no wonder that there are so many
disputes among the philosophers; we are indeed still engaged in
trying our way forward, and therefore now one, now another
main point must come to step forth with one-sided strength.

Those who know the more recent history of philosophy will find that
Sibbern's whole conception of the essence and method of philosophy
reminds not a little of Lotze. Also otherwise a certain
kinship appears between the views of the two thinkers.

Sibbern paves his way from his theory of knowledge to his
metaphysical theory by showing that the possibility of knowledge rests
on the fact that existence presents a connection of the one
with the other, so that everything in the end forms an all-embracing
whole. If there is not in existence an all-embracing
order of things, according to which each single thing is determined
% --- p. 86 ---
\opage{86}by the others, then no single thing can find its full
explanation. There must be given a total consequence that bears the
single things and phenomena. This is the fundamental
philosophical thought, as yet in indeterminate and contentless form. It is already
the idea of an all-being and all-effecting, which through the
all-embracing order of things bears and determines the single
phenomena.

But this presupposition is not sufficient. For the
all-being and all-effecting, the general order of things
must at each single point work differently according to the given
which is determined by it. All arising and becoming comes about
through two factors: the one is the order of things working in all,
the other is something definitely given, which brings it about that each
time something definite and not something else comes into being. There are
constantly in nature different points of departure, from which
the processes take their course. And these points of departure are
each for itself finite and relative; they stand in mutual
interaction, so that they are not only points of action, but also
points of reaction. Therefore nature appears to us as
a great process of interaction, a great debate and struggle of
all against all, and only through this great world-debate
does the principle of existence unfold its whole fullness.

This second presupposition one might call the
realistic presupposition, as one might call the first
the idealistic, although Sibbern himself uses other terms.
By this presupposition he steps from a new side into a
decided opposition to Hegel and the German speculative school.
He reproaches Hegel that, on account of his one-sidedly
idealistic point of view, he could not obtain any really \emph{historical}
conception of existence. Historical development becomes first
possible when there are separate points of departure in mutual
interaction. Development then consists in bringing
harmony among the sporadically emerging processes, in
organizing the material given in scattered form. Experience
shows us, as we have already emphasized with regard to psychology,
that everything in nature at first forms itself sporadically.
In the process of crystallization, such as that which takes place
when water freezes, the needles arise from different points,
% continues on p. 87: "men gaa efterhaanden sammen og danne en kompakt Enhed." (but gradually merge and form a compact unity)

% continued from p. 86: Ved Krystallisationsprocessen ... opstaa Naalene ud fra forskellige Punkter,
% --- p. 87 ---
\opage{87}but gradually merge and form a compact unity.
The formation of the embryo takes place in several respects from
different points of organization, from which the processes meet
in order finally to end in a result that bears a striking stamp of
unity. The social world shows us something similar: within the
individual state the endeavors of different individuals and from
different quarters work toward forming a common social organism,
and the different states are in turn points of departure for effects
which, through many collisions, lead us forward toward a
general human society. If we go back to the origin of the globe
and of the whole solar system, then every particle in the
primeval mist must have played its part, must have been a point of departure
for attraction and repulsion, and through the swarm of
interactions thereby brought about our globe and
our solar system received its present form. In every domain
a great, advancing process of organization can be traced in
unbroken connection. ``The very first stirring in the
process of the globe was the beginning of the constitution of the whole
world of plants, world of animals and spirituality now subsisting on the earth''.
Existence is one great, continuous process of development,
within which the coming-to-be of man marks a single
stage. ``How, then, did the first human being come to be?
Without doubt thus: that a certain animal formation, from a first
slight beginning through a series of generations or
offspring, as in the progeny the process of formation
went on continually further and further, reached the point at which its latest
offspring became a being which, so to speak, opened
the eye of spirituality, in that it said I to itself and felt itself
in its selfhood''\footnote{A similar view of the coming-to-be of man had already been
expressed by Kant (in his Anthropology), and among us by Treschow (in his ``Historiens
Filosofi'').}.

Thus, despite the sporadic points of departure, there is a
progress toward ever greater harmony and toward higher forms of life.
But at the intermediate stages forms
and phenomena can often appear which seem to point in quite another direction
than toward a progress to higher perfection, just as
% --- p. 88 ---
\opage{88}the fetus at certain stages of its development appears to the
untrained observer to be a malformation. Nature seems
often to go the ways of absurdity. All such things are nevertheless only something
provisional and transitory, a little bend in a road
which the development of the world must traverse. Existence
is as yet on the way, indeed even at the beginning stage of the way
to the goal. There still shows itself in our experience an
opposition between the central principle, the principle of unity in the world,
and the manifold peripheral points of departure. But it will
more and more be replaced by a relation of harmony. Yet only for
a limited part of the world can there be any question of a
complete realization of the world's ends; since in
the world as a whole an inexhaustible fullness of content is to be realized,
an infinite time will be needed for it.

Such were the philosophical ideas. Sibbern expressed them in his
remarkable little work ``Spekulativ Kosmologi'', which appeared in 1846,
the only one in which he has presented in connected form his
general world-view. It still repays those
who are interested in philosophical speculations to study
this work, which, as will have been seen, does not lack
points of contact with the theories and outlooks of the most recent time.
For a closer characterization I shall here add only a few
remarks.

One of the most important aspects from which a philosophical
world-view can be characterized is the relation it posits
between unity and manifoldness in existence. Sibbern seeks
to give both what is due to them, each for itself, and in
his presentation the motives which give them
interest and significance in philosophical respects are seen very clearly. We are led to
presuppose unity by the fundamental fact that existence,
wherever we are able to penetrate into it, appears to us as
a connection ordered by law. There must then, it seems, be a
single principle which works in all. But on the other hand
experience shows us just as well a manifoldness of working
elements, of separate points of departure. Although now this
latter standpoint is stressed very strongly by Sibbern, especially in the
doctrine, so peculiar to him, of sporadic development,
yet he decidedly subordinates manifoldness to
% --- p. 89 ---
\opage{89}unity. The real points of departure are indeed designated as something
``Given'', upon which ``the All-Working'' exerts its influence; but
this Given is itself supposed to be only the All-Working's own
primitive activities, and the influence of the All-Working is supposed to
coincide entirely with the activity of the separate elements. A
conception such as Leibniz's doctrine of monads, or such as the absolute
atomism, which ends in a swarm of independent single beings as the
last thought, Sibbern rejects decisively. Every single being
is only a moment in the infinite being, every molecule
only the seat of one of the functions of the All-Working. It is
supposed to be only our imagination that cannot help ascribing to
the elements or molecules a reality other than that of
being mere points of departure of activity for the All-Working,
just as our individual ideas are so many forms
of the activity of spirit and have no existence apart
from it.

But with this it cannot be reconciled that the sum total of
the elements is called \emph{something standing over against}, an \emph{Other}, an \emph{Objective},
which is to be worked upon and shaped by the All-Working. Here
a duality appears, an opposition which is not
overcome. When an idea is to be realized, says Sibbern in
this connection, everything cannot proceed from its inner
wellspring, but material and means must be taken from without. Here, then,
a duality is present from the start; the two presuppositions
on which the construction rests cannot be traced back to a unity,
and there lies here a threatening danger for Sibbern's optimistic
world-view.

How can he be so certain that there will come
ever greater harmony between the principle of unity and the sporadic
processes of development? Only so short and slight a part of the
development can we survey, and we cannot on that basis cast
the horoscope of the whole development. We may console ourselves with the thought
that the absurdities and disharmonies are due only to our lack
of acquaintance with the real laws and conditions of development:
but then it is only a faith or a hope to which we here take
refuge, no scientific insight. It is, if one will,
a philosophical faith that we can have, that discord and
disharmony, through the struggle, through the great debate of all against
% --- p. 90 ---
\opage{90}all, lead over to an organization so much the richer;
but more than that it can never become. It is an energetic
conviction that stirs in the enthusiastic thinker who has so
strongly stressed the sporadic and fragmentary, when
he constantly holds fast to the idea of a harmony working itself out
through strife; but it cannot be scientifically
grounded and carried through.

Perhaps it is connected with this that Sibbern does not
hold fast to his doctrine of development all the way. His faith in the
All-Working leads him not merely to ascribe to it an
infinite superiority over the individual elements that
are to be organized; but he even exempts it from being
subject to the process of development. The real development
consists for him in the organizing of the elements; the organizing
power, the All-Working, is not itself becoming but eternal.
Or, as he expresses it with a theological turn: God
does not become, but only God's Kingdom becomes. Such a
distinction cannot be maintained, and Sibbern feels it himself.
How can God be eternal and without development, when his
Kingdom is not so? God and his Kingdom belong so closely
together that they come to be only with one another; there cannot
be a relation of indifference between them. If God's
Kingdom becomes, then God too is caught up in becoming. Sibbern has here
taken over more of theology than the mere words. To be sure,
the Christian doctrine of a suffering God points with magnificent
profundity toward what a thoroughgoing philosophy of development
demands; but theology's usual inclination to halve
thoughts has made it shrink back from this bold
consequence.

Although Sibbern likes to use theological expressions,
his doctrine is utterly heretical. God is for him not a being
separated from and different from ``the world''. He distinguishes between
existence's ``world side'' and ``divine side'': the former we have
in the manifoldness of the elements, the latter in the working of the principle of unity.
The Godhead works through the activities proceeding from the elements.
Yet pantheism (which he, by the way,
speaks of in warm words as a serious and legitimate view of life)
does not satisfy him, because it teaches ``that there is no other God than God's
% --- p. 91 ---
\opage{91}Kingdom, or that behind the Kingdom I here call God's Kingdom
there is no personal God to be sought''. It is nevertheless a
question how many of those who believe in a personal
God would find themselves satisfied by the proposition in
which Sibbern finally expresses his view in the philosophy of religion: ``Does not
everything lead to this, that one must in the end
acknowledge that the innermost ground of the world's existence must be
centralized in itself, that the world is, so to speak, one great individuality?''

By what has just been discussed we are led naturally over to Sibbern's
religious views.

\section*{IV.\\  Sibbern's Religious Development.}

Sibbern's youth fell in the time of the great religious
reaction against the eighteenth century. His philosophical
teachers, Fichte, Schelling and Steffens, were, when he came under
their influence, already on the way to theologizing and
found themselves in any case in a certain reaction against their own
youthful views. Here at home he stood in close personal
relation to Grundtvig and Mynster. He even stood for a long
time as the only teacher at the University with whom those
who felt dissatisfied with Rationalism could find
refuge. His position in this respect was determined by
his whole personality and by his fundamental psychological
conception. He felt deeply the need of a richer content of feeling and
imagination than Rationalism had been able to offer, and he
saw clearly that in these the human spirit stirred far
more forces and impulses than the flat, intellectualist and
unhistorical ``religion of reason'' had acknowledged.

Yet he always made his reservations. Although he
protested against Heiberg's proposition that salvation was to be sought in
philosophy, and although philosophical inquiry was for him only
a means, a preparatory trial, whereas the real
view of life rests on personal being-gripped and heartfelt
adherence, he nevertheless made the validity of every dogmatic
assumption dependent on philosophical discussion. It is one
thing that a doctrine is grounded in the tradition of the Church, and that
it sets our feeling in strong motion, it is another whether
% --- p. 92 ---
\opage{92}it contains anything that has validity and reality.
Philosophy is to examine this independently of all tradition and all % print: uathængigt (misprint)
influences of feeling. ``On the whole,'' it says in an
essay on the relation of philosophy to theology, to be found
in ``Filosofisk Arkiv'' (1829), ``philosophy stands in the same relation to
what is positively or de facto given in religious respects as to
what is given in reality in other ways; and that whether
the former is given as an external, so-called positive
revelation, or as the inward revelation so called by some.'' In both cases
it is the real significance and genuineness that are to be demonstrated. And it is,
in Sibbern's view, not a presumption on the part of an outsider when philosophy
wishes to test and examine the dogmatic doctrines. On the
contrary: in these doctrines an old philosophy is contained.
``I believe I can assert,'' he says, ``that if one were to extract from
dogmatics, not merely as it now is, but also as
it has formed itself from the earliest times on, all
that has sprung from philosophizing, following it in its
finest ramifications into the premises of the dogmatic
doctrines, then it would collapse.'' Historical
criticism vindicates Sibbern in this assertion, in that it
shows Platonism's great influence on the development and formulation of
the Christian central dogmas. But he was not, when
he wrote down the sentences cited, conscious of the reach of
their content. He still thought himself essentially in
harmony with the prevailing theology and with the teaching of the Church.

Only on a few points had he come to have
misgivings, and these points contained the germ of the decided
change in his religious standpoint. It was especially
Lindberg's conduct in the Grundtvigian controversy that drew his
attention to them. There were two things that offended him.
First the theological fanaticism to which that controversy
seemed to him to bear witness; next the weight Lindberg laid on
belief in the Devil and belief in eternal damnation. Of this
belief Sibbern had to concede that it was found in Scripture, but
he rejected it nevertheless in the most decided manner. This
had to remove him decisively from the champions of Orthodoxy.
The violent controversy in the theological field here at home in
% --- p. 93 ---
\opage{93}the twenties, thirties and forties turned in dogmatic
respects essentially on how far a man could be
priest or teacher in the Danish Church if he did not believe
in a personal Devil and in eternal damnation. The
orthodox tendency has, as is well known to all, essentially won out;
the dogmas mentioned are the most characteristic of the theological view
now prevailing in our fatherland; Sibbern's
adherence to the religious movement ceased at the point
when he noticed what weight was laid on them.

After he had for a longer time not expressed himself on
religious questions, he set forth in a series of
university programs (1846, 1847 and 1849) the new standpoint
to which he had now attained. Of the change that
had taken place in him he has himself spoken thus in a letter
from the year 1864: ``Everything in Christianity that appeals to the spirit, begets souls, draws all its
strength from the spirit's inner accord
is for me entirely what it was. Only that which was
to draw its validity from positivity did I \emph{respect}
at that time as a Given ...... To this dependence % print: Athængighed (misprint)
I will now no longer concede any validity at all.'' --- What he
had stated in the programs mentioned in a general way
he carried further partly in ``Meddelelser af et Skrift fra
Aaret 2135'', partly in the moral philosophy published after his death.

Several different motives can be traced that in Sibbern
had worked together to bring about and shape his new
standpoint. --- First, there stirs in him a
philosophical-historical interest in bringing the rise of Christianity
and the personality of Christ into natural connection with
the development of the world. Already in his ``Kosmologi'' he had
stated that, just as the coming-to-be of man had from the beginning to be
explained as the result of a further and higher development from
lower forms, so the rise of Christianity had in turn to
be conceived within human development as a peculiar turning point,
yet one that arose in continuity with the preceding development. If one
wishes to speak of an incarnation, then the whole
development of the world, in particular the development of the earth, must be conceived
as an advancing incarnation. At the point in the
development of the primeval mist at which the globe began to separate itself
% --- p. 94 ---
\opage{94}from the rest of the world-mass, we must seek the first
germ of spiritual life on the earth as a whole and therefore
also of the highest forms in which spiritual life
has historically appeared. The historical Christ has emerged
from the midst of mankind as a man in whom the wellspring of life
stirred in a peculiar manner, so that he could become
a spiritual leader for millennia. Only when he is conceived in a
purely human way have his struggle and suffering, his
striving and working true and natural significance. One ought
not to have distorted this by grounding it on anything
other than the inner might and power of his activity itself. ---
To this is joined the conviction that religious life
consists essentially in an inner relation to the
fundamental source of existence, a feeling of heartfelt unity with and adherence to
this ``inner Giver''. It is often thought that this is an unsafe
ground to build on, and that one must have an objective
foundation in a dogmatic confession and in an outwardly appearing
and organized Church, if religious life is not to lose
itself in merely subjective and arbitrary impulses and illusions.
But to this it must be answered that everything spiritual that is to gain
dominion in the world precisely goes the way of subjectivity, stirs
with a power coming from within, gripping and animating,
in the individual personalities. The real Objective is
here precisely the inner animating power which stirs in the
individual persons. Historical development and the historical
forms do not thereby lose their significance, for external
influence and help are needed in order that the inner may be set in
motion. On the other hand it is an illusion in which the Church and
theology have been ensnared, that anything should become objective,
that is, valid in and for itself, because it is acknowledged by large
circles and made the foundation of institutions in the external
world. The inner power that carries spiritual life forward
stirred quite as vigorously in Rationalism and Pietism as
in the orthodox Church. ``Did not Rationalism break
mightily forth, and dare anyone say that humanity, care for human beings,
the care of feeding the hungry and quenching the thirsty
did not take place just as much and just as warmly on this
side as on that?'' --- One sees here how Sibbern has
% --- p. 95 ---
\opage{95}come beyond the reaction against Rationalism and is able
to place it in its rightful position among the
contending tendencies. He finds also in this
connection application for his favorite standpoint: the debate
of all against all. Through the mutual conflict between all
the different tendencies and personalities religious
life advances. One is not to take up an attitude of \emph{indifference} toward these
different views, but one is to ascribe to them \emph{equal
validity}, in so far as each for itself bears witness to the inner
wellspring of life and has a right to assert itself and make its
contribution to the process of development of human life. In the second part of
Gabrielis' Breve, which otherwise has so ecclesiastical a stamp, and
where the hero ends by becoming a priest, the constant
undercurrent is nevertheless that every strong and living faith is in and for itself
justified as a vigorous expression of life's inner power. Even
of those who after full conviction reject faith in
God and immortality one must, it is said, say that
they have a right to do so, when it is ``a living sense for the
ideal, or at least for that which ought and should be'' that stirs
in them; they too think of the welfare and salvation of souls.
The form of religiosity which Sibbern himself depicts with warm
adherence is a faith in a personal Godhead, ``a life in
God as the inner Fundamental Worker and Fundamental Giver, by whose
inner workings and givings every good, enlivening and
fruitful thought and feeling arises,'' in such a way that the individual's
own striving and the living of the inner Giver in it constantly
go together as One. It is a kind of mysticism which reduces the
dogmatic assumptions to something relatively small. In the Bible
he finds indeed constantly many profound and appealing thoughts, but
far from full clarity and truth. With all its depth, it is said,
Holy Scripture often gives only half thoughts. It
gives us, for example, no clear and adequate information
about the existence beyond. It cannot then help to
base the significance of so-called revelation on the
information it is supposed to give us in this respect, especially since that
which Scripture plainly enough teaches about it, namely the eternal
damnation of so many, is enough to make our moral
feeling reject its authority. Instead of holding
% --- p. 96 ---
\opage{96}to one such single, transmitted book, let one hold to
the Book of Nature, ``the great book of the great earthly fullness of life
appearing both in the bodily
and in the spiritual''. And if it is now asked whether by our natural experiences we can
find information about everything we could wish to know, particularly
about another life, Sibbern answers that whether there is
an immortality or not, we shall live life as we live
it; our moral conviction is not built on any
doctrine of the beyond.

The tendency that runs through the history of the Church, to
dogmatize, that is, to fix the truth one believes
one has won in definite, prescribed formulas, is in several
respects pernicious. One thereby tears truth loose
from its living root, its ground of life, and makes a mummy of
it, so that it may in this mummy form be
accepted by all. This tendency expresses itself perhaps still more strongly
in orthodox Protestantism than in Catholicism. One has
in Protestantism set a holiness of dogma in the place of
Catholicism's holiness of works. Instead of hedging in in this way
the ways by which the Spirit works, let it
have free play! Let no one force his way into the individual's
sanctuary with polemics or with mockery! Let each keep
his fetish, that in which he has found his spiritual refuge!
But Christianity has not conceived the matter thus. Although
it was supposed to be the religion of love, it has nevertheless in all
its time far more fostered and nourished the spirit of presumption,
of domineering and of hatred than the spirit of love of humanity.
It has abhorred and combated the gods of paganism and
its deities, without thinking whether among the heathen too
there might not be found such little ones as one above all
must not offend, and without examining whether
polytheism too might not have its good justification. What matters
is that love of humanity unfolds, and one
must therefore advise everyone to stay with his faith, when it
is a consolation to him and promotes what duty and
care for humanity require.

Within Christianity, in Sibbern's opinion, men have in this respect gone by such pernicious and wrong
% --- p. 97 ---
% print p. 97: en stort (for stor) -- misprint; the sense is translated
\opage{97}ways that he has the future Europeans he depicts in
the book from the year 2135 refuse to call themselves Christians.
``Could we will to bear a name in common with those who
in the name of Christianity and under its banner have filled
centuries with presumptions and atrocities?'' --- He
distinguishes between Christliness and Christendom. The Christly
is the true and sound kernel in Christianity, but to
Christendom belongs all that is dogmatic and ecclesiastical.
He attaches himself essentially to Neo-Rationalism, which he
characterizes in his moral philosophy in the following way: ``It
wants to have the ecclesiastical dominion of Christianity thrust aside
in order to put in its place the disposition of free devotion to God,
which I, in contrast to Christendom, have called
Christliness, letting the word Christianity stand as indeterminate''. ---

What seems to me to be the most
significant thing in Sibbern's religious standpoint is the magnificent
faith in the right of personality. He has in far fuller
measure and in a far more natural manner than S. Kierkegaard
carried through the proposition that subjectivity is the truth.
S. Kierkegaard gave this proposition an all too narrow
application and execution. We meet here again Sibbern's firm
faith in life in its sporadic working in connection with
his warm sympathy for every personal expression of life. It
is not a proposal of tolerance grounded in embarrassment and skepticism
that we have before us here; but a demand set up on the basis of
a great and deeply conceived world-view.
Here there is likewise no anxious search to see whether there might not
be found a hiding place where we could conceal our dogmas
so well that criticism could not find them. There is
on the contrary an endeavor to bring about as great
agreement as possible between faith and knowledge, between
view of life and world-view. And we have seen that
essential features of his world-view find confirmation in
the newer theories and outlooks; therefore his
religious standpoint too will be of interest to the present, apart from the
great interest it has as the result of the inquiry and search of a soundly thinking
and warmly feeling human being. The religious
questions have here at home in recent years, on the part of those
% --- p. 98 ---
\opage{98}who cannot join Orthodoxy, certainly been far too much
treated critically and polemically. In Sibbern
criticism and polemics are always induced and borne by the
positive fullness for which he works. In this he can be a model
also for those who arrive at results that differ from his.

\section*{V.\\  Sibbern's Political and Social Views.}

Sibbern's political and social views bear in a high degree
the stamp of his whole personality. This by no means deprives them
of their significance and lasting interest, for as one of
his chief peculiarities we have come to know the
living sense for and love of life in its reality
and individuality. He seeks everywhere the content and the kernel,
and if he has any fault, it is that of overlooking that in
this world there are once for all needed certain definite forms
and shells, often hard and unpalatable shells, even for its
most valuable content.

What so often made him unpractical when he
took part in the political strife was not least that he
appeared as a psychological pastor of souls. He had a keen
eye for what is irritating and spirit-consuming in political
strife, and not least for how in the course of the strife
much may pile up which hides what it is really
a question of. The form of political activity that most
appealed to him was that which corresponded to Heracles's way of
cleansing the Augean stable. Just as Heracles accomplished his
work by leading the stream in and then leaving it to
flush out into all corners, so Sibbern considered it his
calling, when in the course of the strife polemics on the one hand and
loss in details on the other had darkened the view of most people, then
to work as directly as possible to gain entrance for the
main point itself. Thereby he thought to work not only
best for the cause, but also to procure for minds the refreshment
and rest which they needed. It would also be good, he thought,
if in politics one could let, for example, every
fifty-first week be a week of jubilee, when one let the disputes
% --- p. 99 ---
\opage{99}rest under pleasant personal intercourse. There was, in
Sibbern's opinion, absolutely no haste in politics. He
rejoiced in his ``Politiske Intelligensblade'' (1835) that
the consultative Estates first met four years
after the ordinance on their introduction had come. In
this great interval there was room for many-sided and
thorough deliberation. Therefore he also regarded the
separate assemblies of Estates for the Islands and Jutland as a
fortunate arrangement, since every matter was thereby considered twice.
Later he rejoiced that the free
constitutional form of government did not come before 1848, when conditions were suited
to it. And both during the strife over the
change of the unitary state (1854) and during the strife over the revision of the
Constitution (1865) he appeared with the counsel to put the matter in abeyance
or at rest for some years, so that it might be considered and
worked over in quiet, and so that the more important ends of the state, in particular
the care for the welfare and civil liberty of ``the commonalty'',
might be promoted.

That in the heat of the strife men were not inclined to listen
to such counsels cannot surprise. And likewise it follows
of itself that the great constitutional decisions which are
determining for the whole manner in which, and the conditions
under which, all the most important affairs of state are in detail
dealt with and decided, cannot be postponed too long, and that such
a postponement will in any case hardly free minds from
tension and unrest.

It is on the whole characteristic of Sibbern's political
views that he is so careful that deliberation
may be thorough and many-sided. His politics here hangs together
with his philosophy. In the whole of existence development advances, after
his doctrine, through a great debate of all against all.
The same holds of political development.
The goodness of a constitution therefore depends on whether it opens
access for all forces and tendencies in the people to come
to speak and to work as co-determinants in development. In that
he sees the characteristic of the \emph{free} constitution. On the other hand
it is in his opinion of more subordinate significance where
the final decision, the determination of the will after deliberation,
% --- p. 100 ---
\opage{100}lies. He considers it most fortunate that it lies with the King.
The freest, most many-sided deliberation could then be united
with the most energetic decision, thoroughgoing
decentralization with thoroughgoing centralization. Sibbern's
ideal of a constitution remained always ``an
absolute king with a merely consultative assembly of the people's spokesmen,
which was to seek all its weight and strength in the goodness and
well-consideredness of its grounds''. Precisely when the people's representatives do not
take part in the decision, the greatest freedom can prevail in
public discussion, and the people's judgment will be able to gain
a far greater influence on the final determination than
when this rests on a vote in the assembly of
representatives. The power of government is to stand under the control of the Estates.
A control of all over against all is in every constitution
the last resort, and the so-called constitutional guarantees are in the
last analysis all of a moral nature, in a
``free-state constitution'', where the representation has a vote in the decision,
as well as in a purely monarchical constitution. So far
is it from being the case that such complete freedom of expression of opinion
and public discussion should harm the monarchy, that nothing
at all supports so-called absolute monarchy more than
a republican spirit both in the people and in the government. The people feel
thereby that they stand near their King.

By a constitution one understands as a rule an outline of
a state order written down on paper. But it is more natural
to take the word in its medical sense, thus here
of the real condition and order of life of the people and the state. Only
a small part of this can at any time be formulated
and written down. And such a written constitution will always
be a fetter on life. ``The more provisions a
fundamental law contains, the greater guardianship has been assumed
over the future and its generations, the less is the state itself
truly free, the more, on the other hand, are life in the state
and the whole national development, the national endeavor and the nation's
fresh stirring burdened with fetters which easily become
troublesome drawbacks and entanglements'' (Dikaiosyne, 1843, p. 82).
Such a written constitution was, in Sibbern's opinion,
as pernicious for political life as a written and
% --- p. 101 ---
\opage{101}prescribed confession of faith for religious life. Of all
forms of absolutism he considered that which is committed to
a piece of paper the worst. ---

It is a strange opposition that Sibbern posits
between deliberation and decision, and it is certainly just as
little tenable in politics as in psychology. Every motive
that comes forward during deliberation also gets its
determining influence at the decision, and if we could wholly
see through what goes on in our inner life during the act of will,
we should often find that the decision stands as the conclusion of
a process in which the earlier determines the later throughout.
It is always a sign of inner discord and disharmony
when the decision is determined by something that has not been able
to prevail during deliberation, and it will as a rule be for
the worse when it happens. Sibbern's political ideal will in
a corresponding way lead to a disharmony, or else to
the royal power's drawing the conclusion from what has come forward
in the public debate. --- In the first case --- that is,
when the royal power and public opinion are at odds ---
the question will arise whether as a rule the King or the people
has the truest view of things. If one assumes that the King
as a rule is right, what then becomes of ``the republican
spirit''? Then the King can at the most be obliged
to take precautionary considerations, or a pedagogical consideration
for public opinion, but the debate and the control of
all by all becomes illusory, and it will then surely be
most expedient to stop the idle discussion. Frederik
the Sixth's famous saying: ``We alone know what serves
the true welfare of the subjects'', in his negative reply to
a petition for extended freedom of the press, is the logical
consequence of this conception. If, on the other hand, one assumes that
public opinion as a rule is right, one will conversely get in
the royal power a constant hindrance to progress and a
danger to the state. --- One will then be driven over to the other
possibility, that the royal power is to exercise and execute the
result of the public debate. But in that case it will be
most expedient that this result be expressly formulated and
fixed, and that can only happen by a vote, that is,
% --- p. 102 ---
\opage{102}by the representatives of the people (or the people through its
representatives) having a share in the power of the state.

\emph{One} proposition one must grant Sibbern is right, and that is this,
that all constitutional guarantees are moral. It gives no
security that the people's elected representatives get a share in the
power of the state, when it depends only on the government's own will
whether it is to take any notice of their vote.
Had Sibbern lived to see the centenary of his birth,
he would have received an experiential confirmation of his proposition.

But in order rightly to assess his political conception, one must
lay the chief weight on the republican spirit which he
demands shall prevail in the monarchy. What he wishes to exclude
and prevent is the sharpening of distinctions of estate and the formation of
an aristocracy. He wanted neither an aristocracy of nobility nor one of officials, and he
feared that in a ``free-state constitution'' it would be ``the well-to-do'' who came to
rule over ``the commonalty''. It had, in his view,
been the endeavor of the Danish absolute kings to work
for the emancipation of the people. The change of government in 1660
he regarded as a liberation of the burgher estate. From
this year on he therefore dated freedom in Denmark. In particular
Frederik the Third and Christian the Fifth, and later
again Crown Prince Frederik's government from 1784 worked immensely
much for the benefit of the commonalty and for general equality
among the citizens of the country. Nothing angered Sibbern more than
when he heard the liberal politicians speak of the time before 1848
as a time of unfreedom. He had himself undertaken a thorough
study of Danish history of the seventeenth and eighteenth centuries.
Still in his ``Samfundsbetragtninger'' (First Issue 1865)
he says: ``If in my great age I dared think of
such a thing, I should surely attempt a historical-philosophical
survey of the history of Denmark from 1536 on ..... It is especially for
those who take part in the politicizing that I should wish for depictions
of Denmark's inner \emph{history of the development of state and people}, written
so that the most important sources, the royal ordinances,
and what among the rescripts must be placed together with
them, are rightly examined through, and what is significant in
them brought forth.'' He rightly emphasizes that therein lies
% --- p. 103 ---
\opage{103}the continuity in the history of the Danish people. He would highly
have rejoiced at Edv. Holm's various works on
Denmark's inner development under absolutism and would have
found a confirmation of his views in them. He has
here a more thorough historical view than the National Liberal
politicians, who were so occupied with their own political wisdom that
they saw nothing but darkness before their own time, just as they
also have thought that nothing but darkness would follow
after them. On the other hand he did not see that absolutism developed
an aristocracy of officials, which in its way was bound to become a new drag on
development. But in this the National Liberal politicians
saw no more clearly than he. ---

Now fate would have it that Denmark should after all have
a free-state constitution, and although Sibbern did not regard it
as the ideal constitution, he nevertheless saw that it could not
be otherwise. He protested only against its being
freedom that was bestowed on the people in 1848. Freedom cannot be
bestowed. It was power that was divided between King and
people. That it came to this was due, in Sibbern's view,
to several different circumstances. The prevailing political
movement in Europe propagated itself to us too, and the zealous
liberal agitators (``the stir-up men'', as Sibbern in his language
called them) were not content with the old constitution.
Besides, Frederik the Seventh was not fitted to maintain
the absolute power. ``One must say that for him the new
organization of the state was a great piece of good fortune; for partly it freed him from
duties of government for which he was not at all made --- nor
had he sought to train himself --- to be able to fulfill
even tolerably, as such things may be required of men, partly
it put him in a very good and at the same time humanly free
relation to the people.'' It is only a pity that this ``little
king-wise'' man did not understand how to use the part of the power
which he retained! (Meddelelser af et Skrift fra Aaret 2135.
Første Del, p. 645 and several other places).

The people was \emph{suited} to the new constitution. Under
absolutism equality had been furthered, and emancipation
carried through. It was a great piece of luck that the new
constitution % print: Foratning (misprint)
 did not come until it could be founded on a
% --- p. 104 ---
\opage{104}very broad basis, with so extensive a franchise. Denmark was
thereby spared the opposition which puts
other countries in such great and constant unrest: between the
enfranchised and the unenfranchised. ``So-called
democracy stands everywhere behind the Estates halls and wants in;
happy the country that has opened the doors wholly to it, and
at the same time opened an arena of tumult and combat not merely
for it, but also for its opponents and adversaries.''
``\emph{Suited} the people was, not ripe; for a people becomes ripe only
by means of the constitution itself. All might have gone
well, if only the National Liberal leaders had not
lacked the maturity which must be required of those who wish to
steer the ship of state; maturity to draft the sea chart they had
perhaps indeed'' (Samfundsbetragtninger. Første Hefte, p. 18, 42 f.).

How sincerely Sibbern had attached himself to the spirit of
the new constitution, when historical development had brought
it with it, he had two occasions to show, under the
Ørsted Ministry and during the Constitutional struggle of 1865---1866.

The work he published during the Ørsted Ministry's
struggle with the Rigsdag (``Om og i Anledning af Kampen
mellem de tvende højeste Statsmyndigheder i Danmark'') is
of no small interest precisely in our own days. One finds there
a series of complaints which have come forward anew here
at home in the present political strife. --- Sibbern
complains that the Ministry has drawn the King into the strife.
Free access to the King no longer stands open,
as it for a long time had been in Denmark. ``It
was,'' says Sibbern, ``scarcely either right for the state or
wise for the state that free access closed when men wished to come
with addresses from many men very devoted both to the King's person and to the King's
royal dignity. The government's ends
are scarcely furthered thus; far better without doubt in the opposite
way .... One has thereby touched the people at a very tender
spot, perhaps especially the people in Jutland.'' It was indeed precisely what
Sibbern had always emphasized as the task of kingship:
to stand in close and living contact with the whole people. He
cites in another work with enthusiasm the well-known song
from ``Kong Salomon og Jørgen Hattemager'':

% --- p. 105 ---
\opage{105}
\begin{verse}
Constantly there was a way in Hertha's grove \\
From the people's cottages up to the prince's halls. \\
The humblest himself drew near the people's lord \\
And brought thither his heart's petitions \\
And felt himself, in the King's castle, \\
That he was one of the Land's sons.
\end{verse}
% Bestandig var der Vej i Herthas Lund \\
% Fra Folkets Hytter op til Fyrstens Sale. \\
% Den Mindste selv sig nærmed Folkets Drot \\
% Og bragte did sit Hjertes Bønner \\
% Og følte selv i Kongens Slot, \\
% At han var en af Landets Sønner.

That was now over, and under a free-state constitution! Under
absolutism, in Sibbern's opinion, it could not have gone
thus!

Through the many press cases the fear had awakened that the government
meant to attack freedom of the press. The freedom and standing of the civil service
were shaken by various dismissals. The standing of the Rigsdag was
shaken, and yet it belongs just as much
as the government to the authorities to which Danish citizens are
subject, and which Christians, according to the 13th
chapter of Romans, are to obey! --- Sibbern finally points out
that not all lawful means of putting an end to the strife
have been tried, so long as it has not been tested what influence
a change of ministry would have\footnote{This whole strife tormented Sibbern especially because he saw A. S. Ørsted,
whom he so highly admired, on the other side. ``I cannot
tell you,'' he writes to Wegener, ``how strongly this rescript
[which contained a ministerial reprimand to Wegener for his
conduct in the succession case, although the Ministry had in vain
accused him before the Supreme Court] has affected me, most of all for Ørsted's
sake. I have now for 48 years stood in the best relation to him,
enjoyed much good in his house, held to him, believed in him, until
last year; and now this again!'' --- After the Ministry's departure Sibbern writes
(to Ørsted's sister): ``He [A. S. Ørsted] has by the Ministry's
dismissal come out of a not good position. It is
generally known that much of what happened, and to which he had to give his name,
indeed at times put his name under, by no means had his approval,
but was in conflict with his view. But he was in the Ministry
in the minority and was outvoted.''}.

Some ten years later Sibbern raised his voice against a
``assault on the Constitution'' of another kind. There appeared in the year
1865 three different small political writings by him, all concerning
the Constitutional struggle. It revolted him to see that immediately
after the sorrowful war men made use of juridical
subtleties ``by a detour that looks like a crooked path, to
% --- p. 106 ---
\opage{106}impair Denmark's good free, popular Constitution''. It was for
him only ``trickery'' when men would not admit that
the part of its power which the Rigsdag had in its time ceded
in order to make it possible to carry through a common constitution
that also embraced the Duchies, --- that this power after
the cession of the Duchies must immediately fall back to
the Rigsdag. Men cried for legality, but went the way of
illegality, because they had their definite ends which they wished to
attain: ``Why did men, when the conditions for
the restriction of the Constitution of 1849 were entirely gone, … let a
piece of play-acting enter which had to be offensive to the sound sense of
the whole Danish people? Without doubt it was because
men saw therein a means of bringing the Landsting reform forward,
and then the rest of the comedy had to be performed along with it''
(Samfundsbetragtninger. Første Hefte, p. 70).

Here the one condition had been shaken under
which Sibbern could sympathize with a ``free-state constitution'',
namely that it be built on a democratic foundation.
The Landsting reform was a reaction against the development of
the feeling of equality which especially in the last century had taken
place. Men divide all the sons of the Land into two main classes,
the well-to-do and the not well-to-do, and do not see that
they thereby sharpen the opposition between Folketing and
Landsting, so that the new constitution contains in itself the germ of
long disputes. There is accumulated in it a considerable
material for continued agitation, and those are mistaken who think that
the adoption of the constitutional change has settled the strife. --- Sibbern
stood in this matter wholly on the same standpoint as
Tscherning and Grundtvig and sympathized with the steps they took
to maintain their view. On the whole there is no slight
likeness between the political ideas of these three men. A
kingship on a very broad democratic basis, with a leveling of
the distinctions of estate, was the ideal of all three. Under the changing
conditions Sibbern pursues this ideal with great assurance and
sound judgment of personalities and situations. His
political pamphlets are of no slight psychological and historical interest.
That they had no effect on their contemporaries lay partly in the author's
idealistic conception, partly in the oddities with which they swarm,
% --- p. 107 ---
\opage{107}and which it was easy to score points off. But for him who looks
back on the events it has great interest to see
how they are conceived and judged by a sound, honest spirit endowed with
warm feeling, which at the same time brings with it a not
common historical and philosophical sense. I shall make no
attempt to state how Sibbern would have taken his stand
if he had lived to see the present political strife. He
was a declared opponent of parliamentarism; but he
had also in his time been a declared opponent of
the ``free-state constitution'' and yet learned to regard it as a
natural and necessary consequence. His genuinely democratic
disposition would have placed him far otherwise in this strife
than the many who are led by pride of estate and bureaucratic
prejudices. He would have felt painfully how the two
elements in his ideal of the state, kingship and democracy,
came into sharp opposition to one another; but he would not
have found any way out in the
rule of an aristocracy of landowners and officials\footnote{For those who think that on account of party strife everything
good and noble in Denmark is going downhill, it will be of interest
to know that it is in any case an old trouble. On the 5th of January
1825 Sibbern writes to his old Rector, Brorson, at Herlufsholm:
``Truly, when I think of you, there meets me the image of a
time when there was in many respects a finer life in
Denmark than now. Urbanity and an agreeable manner were still
counted a virtue; the great men of the country respected everyone; if anything
really excellent and good appeared, then every heart celebrated a national feast . . .
Now, on the other hand, everything is so immensely critical that for this critical
sickness almost nothing can come up; to set one's egoism thus
against the whole world, that men did not understand when you were young.'' --- In the year
1828 he speaks (in his ``Psykologiske Patologi'') of ``the propensity to
nourish party spirit and to be moved by party spirit, which is so
prevalent in our time''. --- In a letter from December 1835 Paludan-Müller
(the later Bishop) expresses his indignation at ``the tone
of the daily papers'' as well as at the fact that men use the new Assembly of Estates
to make opposition, ``as men in general take pains to
bring the people to regard all existing conditions in an opposing
spirit''. --- Of the National Liberals it was Sibbern's opinion that they
had much in common with the Schleswig-Holsteiners. ``They were trained
in the same school. Respect for the people's mentality, care for the
preservation of the plain good spirit of the people they did not have. They went so far as to tear down the country's most learned, insightful,
highly esteemed men, and thereby themselves undermined respect for learning
in such a way that, when afterwards the scholars,
the so-called doctrinaires, came into hatred, one could say to them: You have
yourselves worked at it with all your might'' (Meddelelser af Indholdet af et Skrift fra Aaret
2135. Første Del, p. 695 f.).}. ---

% --- p. 108 ---
\opage{108}All the various political and constitutional
questions concerned now, properly speaking, very little what lay nearest
Sibbern's heart. Already in a university program from
1849 he maintains that the social question was far
more important than the constitutional one, and that it was this which statesmen
were first and foremost to make themselves familiar with. The constitutional
business has, he thought, glitter or halfness before and behind,
whereas the idea of Communism indeed has darkness before and darkness behind,
but yet has reality for its foundation and reality for its end.
The question on whose decision everything here depended was for
him whether it belongs to the essence of the community of the state to work by
external power, or whether this must not be regarded as something
provisional, something that holds for mankind's present
stage of life, but will be replaced by an ideal life of the state, where there
indeed is found a firmly organized community for the attainment of common
human ends, but where external, coercive power
has fallen away, because it is not needed. The state grounded on power
and working with power does not lead us beyond
the state of nature. ``Our states are so far from abolishing that war
of all against all which men have supposed to prevail in the state of nature,
that the constitutions of states and the national constitutions connected with them
on the contrary in many ways have made such war of
all against all stronger and sharper, keener, more comprehensive
and more intrusive, not only outward, but also in the interior of every state and
nation in the midst of the so-called condition of peace. One
sees % print: er (misprint)
 in this, to be sure, no bloodshed, but all the more
one sees, when one has the whole of Europe in view, the struggle for
livelihood, indeed the struggle even for the bare support of life,
lead to countless wastings and ruinings of manifold
streams of nerve force, indeed of whole moral existences.''
Should this struggle for power or wealth or even for
the bare existence be eternal? Will not so great a
% --- p. 109 ---
\opage{109}communal feeling develop that external power can fall away and
blind, hard-skinned competition can be replaced by a rational
ordering? Is it not after all the task of the state to look after
the care of all? Could it not be determined that what the
individual acquired beyond a certain measure should be handed over to
the state's common working capital, so that the state could in earnest
appear as the organ of distributive justice, as a
true organization of mankind, which does not let things take their
course under the wild struggle of forces and drives, but orders
and helps in all domains? ---

It is characteristic of Sibbern that he ascribes to the
social question greater significance than to the political. We
have also seen how he repeatedly urges
that the political strife be stopped in order to care for ``the commonalty''.
Yet he sees clearly the connection between both questions,
and it was not his opinion that one could leap straight
into a communistic order. On the contrary, so long as
the coercive state exists, Communism cannot be realized.
Sibbern looked forward to a time when European
mankind had come beyond its present conditions of life and
could arrange itself in an entirely new way. This
fantasy of the future it is that he carries out in ``Meddelelser af Indholdet af et
Skrift fra Aaret 2135''. Sibbern himself intimates that this book
must be regarded as a companion piece to Plato's ``Republic''. He
has thus wished to give vent to his ideal fantasy by depicting
how a human society would look under conditions, inward and outward, as favorable
as possible. The book contains besides
several ``Retrospects'' on the nineteenth century. In the
foregoing we have made use of it in the presentation of Sibbern's
religious and political views. Now we shall consider the
social Utopia which he depicts. ---

Toward the end of the nineteenth century the
evils and dark sides of European civilization rose to
a height which made life a burden. All joy in life was
gone; not even music was able any longer to enliven anyone.
Pessimism reigned. The only refreshment that was found
was to express one's feeling of wretchedness and one's despair
in words as strong as possible. A prince left in the year
% --- p. 110 ---
\opage{110}1896 unhappy Europe with the words: ``I cannot
bear it any longer, this game with worm-eaten souls for
worm-eaten goods; I have long enough toiled in this treadmill
and threshed this kernelless straw.'' And it was proposed that
over all Rigsdag and government buildings these
words should stand: ``Come hither, all ye that are fresh of life and glad, here
shall ye become weary and heavy-laden''. But as distress,
despair and weariness of life were at their highest, salvation came.
A torpor spread over European
mankind, after death had first raged in a mysterious way
through the countries. There came a long time of rest, in
which all the forms of civilization were as good as dissolved
and the few survivors returned to the simple
forms of natural life. At the same time the feelings were softened; there came calm and
love of humanity into men's minds. Weariness was not the only
cause why ``the great salvation of sleep'' came. Great changes
in nature, the breaking forth of subterranean vapors, etc., worked
together. The time seemed to have come for a breakthrough both
in the physical and in the spiritual life on the globe. When
the time of reawakening and revival came, the
youthfully fresh, reborn mankind could begin its existence
under new conditions. Constitution of the state, legislation,
constitution of the Church, school system, monetary system, system of trade, property,
penal system, differences and inequalities among men --- of all
this men had been freed, and did not wish it back,
but felt horror and dread at reading about it or seeing it
among other peoples. Men lived out of the inner
communal feeling. All took part in bodily and spiritual labor. First
and foremost it was a matter of not cowing or hampering the
mental life of anyone. --- Each brought to the public stores
what he produced beyond what he himself needed; and from this
all took what they required. When a house was built,
he moved in who had most desire. Government was not needed.
The riots of boys were the only thing to be guarded
against. Just as little as men wished to have any
holy book or any Church between their feeling and the inner
source of life in which they believed, just as little would they have
any legal constitution as an artificial fence for life. They
% --- p. 111 ---
\opage{111}held that all the acuteness and reflection which
had been applied to jurisprudence, legislation and
application of the law had been of very little real furtherance to
the intellectual life. Crime men had long since regarded
as a sickness which, as soon as it was established by judgment,
was to be cured by humane care and improvement.
The penal institutions became medical institutions or infirmaries. But
with the cessation of private property and of all cowing and provoking
institutions crimes ceased too. ---

Behind the fantastic in these ideas and their closer
working out one perceives a warm and humane mind and a
strong indignation at all oppression and cowing. It was
sorrow and disgust at the actual that drove the thinker, so concerned for human
happiness, to seek in an imagined world
what he did not find in the real one. In the story of
the period of death, of the great sleep and the reawakening he
has admitted that he is unable to show how
the communistic state of bliss could develop
in a natural way out of the present conditions of society.
There must take place an inner transformation of human nature,
so that the feeling of community fully gains
the mastery over egoism. And even then there would have to be
a happy harmony between the inclinations of individual men,
since, for example, a newly built house was to be inhabited by him
who had most desire for it, and everyone was to be able to take
from the public stores what he required. How
is one to decide who has most desire for anything,
or how will it be possible to make the different
needs fit together?

It has hardly been Sibbern's meaning, however, to wish to
relegate everything he sets forth in this book to the land of fantasy, eternally remote from
reality. That is seen already from the fact that in
his university program of 1849 he proposes that the state shall
gather a working capital and appear as a distributive
authority. He is here close to what in our days is called
state socialism. The German national economist Wagner has
carried further what Sibbern intimates in his program. Wagner
regards it, namely, as an essential point of view and aim
% --- p. 112 ---
\opage{112}in taxation, that the state through it acts regulatively
on the distribution of fortunes in the country. In state socialism
there appears the distributive justice which Sibbern demands of the state.
But the coercive state he would not entrust with its exercise,
and how are we to get a state which is not a coercive state? And
how are we to get in the leaders of the state the insight, the
knowledge of human beings and the love of humanity which must be
present in order that we may permit them to intervene
distributively? Even if state socialism should have the future
on its side, development in this direction will be able to go only very
slowly and with renewed trial at each single point.
Sibbern has himself taught us that all development takes place sporadically,
from different individual points of departure. Thus it will
surely continue to go, and it would already be a great
good fortune and a great progress if we could get so far
that the different individuals did not collide too hard
in their different and often opposed endeavors. More
than to act helpingly and supportively in this direction
the state can scarcely do, so long as it has not become something
wholly Other than it now is, --- and when it has become that, it
is not said that there is anything to remedy. Sibbern's
firm and unshakable faith in development comes here again.
He was, as we have seen, certain that the sporadic
and disharmonious would vanish, and that the absurdity and
painfulness which experience presents come only from the fact that
the world's development is as yet at one of its first stages. This
thought sustained him and gave his humane mind
peace. It has probably also made him not strictly
distinguish between Utopias and real ideals in his depiction
of the Europe of the future.

\streg{}

% --- p. 113 ---
\opage{113}
\essay[Alexander Hamilton and the North American Federal Constitution]{ALEXANDER HAMILTON AND THE NORTH AMERICAN FEDERAL CONSTITUTION.}

\streg{}

\section*{I.}

It has often been remarked that the historical conception
which prevails in our century stands in sharp opposition to
the one that ruled in the previous century. In the eighteenth
century one was inclined to explain historical events, above all the
founding of institutions that had proved to be significant, by the conscious
deliberation and intention of individual personalities. And in natural connection
with this one cherished a great, almost unbounded confidence in
what the rightly ``enlightened'' human reason could do
to order social conditions in the best way. The noblest
and the best that the eighteenth century produced is due to
this confidence. The now prevailing historical mode of consideration
proceeds, on the contrary, from the thought that everything great happens, or at
any rate is initiated, in silence, in unconsciousness. More thorough
knowledge of history has shown how social and political
institutions develop slowly, and without there needing from the first
to be any consciousness of their significance
and their remoter effects. History stands before us as
a great unconscious course of development, in which the moving
forces that prevail from the first may be wholly other than
those that later lead to the preservation and appreciation of what
has developed. The significance an institution acquires at
the height of its development cannot be assumed to have
made itself felt in the consciousness of those who laid the first
foundation of it. And we are accustomed to add that it is healthy when
development thus proceeds predominantly unconsciously, just
% --- p. 114 ---
\opage{114}as the child's development proceeds most fortunately when one takes as
little notice of it as possible, and above all when it
takes as little notice of itself as possible. When the strength
and the forms are developed, then it is time for self-feeling and
self-consciousness to awaken, so that what has developed may be grasped and
continued with full understanding of its worth. But this
more thorough and healthier historical conception, by which our
century stands so high above the previous one, has on the other hand
robbed us of the latter's great confidence in the power of thought in the social
and political domain. We have seen so many examples that
the most cleverly devised arrangements do not endure where the
natural soil is lacking in the people's experiences and in its character, grounded
through long ages. It is no
wonder that one has become skeptical of the human
spirit's capacity to intervene consciously and in a beneficial
way in the ordering of social and political conditions. --- Such
a doubt, a doubt that seems so well grounded in the whole
modern conception of history, has something weakening and
disconsolate about it. Are we then referred to letting ourselves drift
with the current? Might we just as well leave off thinking
and willing?

Even a simple consideration (which has already been hinted at
in what precedes) shows that our conception may easily become just
as one-sided as that of the eighteenth century was in its way.
Conscious thought and will do indeed not help when
no natural conditions and capacities are present. They
can build or create nothing out of themselves. But when
thought is used precisely first and foremost to understand what
conditions are present and to what they can be put,
--- what dangers threaten, and how they can be avoided
and reduced, then it is itself an important link in development,
a link that cannot be dispensed with when the course of life is to be
great and healthy.

An interesting confirmation of this is offered by
the founding of the American Federal Constitution. Here the consequence was drawn with clear
consciousness from a long development, running its course in smaller circles. The keystone was set on a structure
that had grown up in concealment. With forward-looking vision
% --- p. 115 ---
\opage{115}means were provided to forestall dangers, of whose gravity the history of the following century has
amply borne witness. And there was a definite presentiment that the
work founded would come to furnish the frame for
a grand development, which would be of significance not only for
the single people, but for all mankind\footnote{At the Convention in Philadelphia, where the Federal Constitution was drawn up,
\emph{James Madison} was especially active for its adoption. He ranked
below Hamilton as a thinker, but had more eye for the
compromises which the political conditions made necessary.
The proceedings of the Convention were held secret, but Madison
wrote at night a detailed account of them, in order --- as he
said --- ``to preserve materials for the history of a Constitution, on
which \emph{the happiness of a people --- great even in its infancy ---} rests, nay
on which perhaps \emph{the cause of freedom in the whole world} rests.'' (Curtis:
History of the Constitution of the United States. I. p. 427). The
first number of ``The Federalist'' (27 Oct. 1787) opens with
the following statement of \emph{Hamilton:} ``After sufficient experience of the inadequacy of
the existing federal constitution, you are now called upon,
fellow citizens, to deliberate on a new constitution for the
united Free States in America. The importance of the subject is clear. It
includes in itself no less than \emph{the existence of the Union}, the security
and welfare of its individual parts, \emph{the fate of the empire which in many
respects is the most remarkable in the world}. It has often been remarked
that it seems to have been reserved for the population of this country,
by their conduct and example, to decide \emph{the important question,
whether human societies are really capable of establishing a good
government by deliberation and choice, or whether they are forever destined
to depend on accident and force}, as regards their political
constitutions .... This idea must, by uniting the impulses of love of humanity
with those of love of fatherland, increase the
concern which all thoughtful and good people must feel
as to what the outcome will be.''}. This
foreseeing consciousness, too, history has confirmed. It
has shown us the spectacle, unique in modern times, of a democratic
great power. A hundred years have now passed since
the Federal Constitution came into force; and it will surely have its
interest to dwell upon the thoughts, and the external conditions, through
whose interaction it came into being.

The credit for this great work belongs first and foremost to
a whole circle of distinguished men, who were inspired by
% --- p. 116 ---
\opage{116}the new thoughts that the English and French literature of the eighteenth century
had set in motion, and who were no less
familiar with the conditions of the North American states and with
the national spirit that had developed there, --- men who
in war and peace had fought for the cause of liberty and self-government.
But among them there is again one man whose name is
in a special way linked to the Federal Constitution, because he
was the first to conceive the fundamental thought on which
it rests, and with the greatest clarity and energy grounded and
developed this thought, just as he also worked with great force
for its execution. This man was \emph{Alexander
Hamilton}, one of the greatest statesmen and at the same time one of the
greatest political philosophers that ever lived. It
is high time that his name be brought forward before the general
consciousness. He deserves a place beside Washington
and Franklin in the history of North America, and in the history of political ideas
he deserves a place in the first rank.

Also by his personality he is worthy of a more
general attention than has hitherto fallen to
his lot. The significant thoughts he fought for found in
him not only one of the most brilliant and energetic, but
at the same time one of the noblest champions that a great
and good cause has ever had. He staked his life on the
purity of his cause. Amid all the strife of the day, amid the muddying of all
conditions, all the falsehood and pettiness which the political
life of the present offers in so many places, it is invigorating to turn
the gaze to a sight which history by no means offers
on many of its pages: a great and practical idea, grasped with
clear consciousness and carried through by a pure and energetic will,
which succeeds at the decisive moment in getting the people to
listen to its voice.

\section*{II.}

The North American Free States owe their origin to
the drive for independence and the feeling of liberty. They were founded by
men who wished to flee the religious and political tyranny
in Europe. New England (the northernmost of the Free States)
came into being, e.g., through the emigration of the so-called ``Pilgrim Fathers''
% --- p. 117 ---
\opage{117}(1620). A band of Puritans had emigrated in 1612 from England
to Holland, but later resolved to cross over to the new
world in order to be able to live according to their own convictions.
Under great hardships and dangers they settled there.

The constitution of the new colony was in the first period wholly
independent of England's. All the members of the congregation assembled
from time to time to hold counsel, and they chose for themselves
a governor, who with a few advisers conducted the current
business. When the population later increased, they established
(1639) a representation, which was to make laws. What
appeared most distinctly in New England developed in
kindred forms in the other colonies. The English government
sought in vain to subdue this so naturally arisen
self-government. Through participation in the internal administration of
public affairs at the various levels there developed
a public sense and a civic spirit, which at that time was
unique in the world. Even the most advanced spirits in
Old England did not rightly understand what went on over
there. One sees this, e.g., from the constitution which John
Locke drafted for Carolina. True, several of the draft's
more liberal provisions were altered by the Lords to whom
the colony had been granted by Charles II; but even apart from this
the draft ran counter to the whole direction that development had
taken over there, and it therefore never came into actual application.

Just as the colonies' citizens thus had occasion
to practice themselves in internal self-government, so the frequent
struggles with Indians and Frenchmen offered occasion to learn
the significance of union and unity among the individual
states. Already in 1643 five colonies thus joined
together as ``the united Colonies of New-England'' in an
offensive and defensive league. While each colony
retained its self-government, all questions of war and
peace and in general all matters of common interest were to be decided by
an assembly consisting of two delegates from each colony.
A century later (1754) Benjamin Franklin drafted
a plan for a union of all the colonies under the leadership
of a governor-general appointed by the English government.
But this plan was rejected in the colonies, because it was thought
% --- p. 118 ---
\opage{118}that it contained too great prerogatives for the Crown, and by
the English ministry, because it was too democratic.
Franklin thought (as appears from his Autobiography) that if
his plan had been followed, the colonies would have been
strong enough to defend themselves; the English government
would then have lacked any pretext to tax them, and
the bloody struggle between the colonies and England would have been
avoided. ``But,'' adds the old sage, ``such
mistakes are nothing new; history is full of the errors of states and princes.
Look about you in the world! How few understand
their own welfare or follow it, when they understand it! Those who
stand at the helm of the state have for the most part much to do
and do not give themselves the trouble of considering and carrying out new
plans. The best public measures are therefore
seldom adopted on account of prior wisdom, but are forced upon us
by circumstances.''

Thus the American people had both occasion
to gain experience and to exercise its powers by independent
participation in the ordering of internal and external affairs. It had
received the best education a people can receive, that which
consists in acting itself, face to face with the real
conditions, without the tutelage of a hereditary government or
a self-wise bureaucracy stepping in between in every matter.

The American Revolution therefore also took on a different
character from the French. In America one had exercised
freedom before one officially received it or took it. Through
long historical development a solid substructure had been brought
into being, which could hold even if the scepter were taken away. With
the English government in America not every government fell;
one had learned to govern oneself. Continuity was not
broken. The people were not accustomed to receive their password from London,
and so it made no difference that the connection was severed.
The difficulty in the coming into being of the Federal Constitution rested, on the contrary,
on the fact that one was so accustomed to almost unlimited
self-government in the smaller circles that one felt a reluctance to
establish a central authority. But life did not dissolve because
the supreme authority was thrust aside. Otherwise in France.
There centralization had penetrated to such a degree that all
% --- p. 119 ---
\opage{119}in the end hung by one thread, so that it had to fall to the ground
when this one thread was cut. The continuity in the people's
life was at one stroke broken and has never since been fully
restored. The development of absolute monarchy
had led to pressing down all classes and depriving all circles
of every initiative. No fellowship was felt. The King and his
officials had to intervene everywhere, if all were not to
come to a standstill. In a report to Louis XVI Turgot wrote:
``Your Majesty is obliged to decide everything yourself or through your
officials. One awaits your express orders before
contributing to the public good, before respecting the rights of others,
sometimes even before acting in one's own affairs''\footnote{Cf. \emph{Tocqueville:} L'ancien régime et la revolution. 7. ed. p. 158.}.
At the Revolution one retained the centralized
mode of government, and the Jacobin holders of power as well as later the
military usurper found their way thereby paved.

What the Americans fought for in their war of independence
was old English liberty, the principle that a people cannot
be taxed unless it has itself through its representatives
granted the tax. This was also publicly recognized by
many Englishmen. The elder Pitt thus declared in
the House of Lords (20 January 1775): ``The spirit which now offers resistance to
taxation in America is the same that formerly
opposed loans, privileges, and ship-money in England,
the same spirit that made all England rise up and
through the `Bill of Right' maintained the English Constitution, the
same principle that established the great, fundamental, and
essential maxim on which our liberties rest: that no
English subject can be taxed except by his own consent.
The honorable spirit of Whiggism animates three millions
of people in America, who prefer poverty with liberty and
will die in defending their rights as men,
as free men.'' --- Even in their Revolution the
Americans had historical ground under their feet. True, in their famous
Declaration of Independence they appealed to universal
and original human rights, and the author of the Declaration, Jefferson, was accused of having written it out
% --- p. 120 ---
\opage{120}of Locke's work ``On Government''. But the American
leaders had a different idea of their
``human rights'' than the French revolutionaries. They did not think that
everything must first be razed, so that one might then on bare ground
erect a new building according to certain abstract principles.
The ``original'' principles one regarded as the innermost
motives and tendencies in the historical development itself, and
what one demanded was merely that this development should not
be halted by any fiat of power.

The history of the war of independence we shall not dwell on here.
Through alternating victories and misfortunes it led to a point where
England's strength was exhausted. Despite its immense expenses,
despite the cruelties and the in part treacherous means which
were employed, England had to give up. --- The Americans stood
at the goal. But to become a people it is not enough to
cast off the external yoke. Great internal difficulties were still
to be overcome, and with right an American author
has recently called the time from the conclusion of peace to the adoption of the Federal Constitution
(1783---1789) the most critical period in America's
history.

\section*{III.}

While English rule lasted, the
colonies had had no steady connection with one another.
Each of them had had to do with the English
government and had fought to assert its self-government against it;
but apart from occasional attempts at union,
the thought of union was not very lively. Just as little as there was
any steady political connection was there any considerable
commercial connection among the colonies. North America
was essentially an agricultural country. Journeys from one
colony to another were, moreover, difficult and rare;
whereas thousands now travel daily between Boston
and New York, two post coaches were then sufficient to
hold all the travelers. Since the building of the Pacific railroad one
can travel more quickly from Boston to the northwesternmost states
than one could a hundred years ago travel from Boston to
Philadelphia.

% --- p. 121 ---
\opage{121}During the war of independence (as formerly during the
struggles with Indians and Frenchmen) the want of
a firm central government had been felt. A Congress, consisting of
delegates with full powers from the various states, had
assembled. This Congress had issued the
Declaration of Independence and had directed the course of affairs during the war and at
the conclusion of peace. But the individual states had constantly felt a certain
jealousy toward the Congress, and one was afraid of giving it
too great authority. The Congress could indeed pass resolutions,
but not carry them out. It could only \emph{recommend} to the
governments of the individual states whether they would carry them out! Actual
governmental power was known only to the individual states, not to the whole
Confederation. In each individual state there was a
free form of government, which dated from the founding of the colonies and which had developed in constant struggle with
the royal governors sent from England. The legislative assemblies
elected in each individual state were in the people's
eyes the only authorities that had the right to impose
taxes. It therefore depended on them whether the contributions desired by the Congress
for the conduct of the war were paid. Every
time a grant was to be made, they all had to be asked, and
often no answer at all was given. When the Congress in the year 1781
proposed, for the payment of the debt, an import duty of 5 per cent,
it aroused great indignation in the individual states. One asked
why one had resisted the Stamp Act and the tea tax,
if taxes were now nevertheless to be imposed by a power
outside the state (the individual state)! Only by virtue of his great
patience and unselfishness was Washington able to
endure the painful situations into which he was often brought by want
of money. He declared later that the war would have ended
sooner and have cost less blood and money if the government
of the Union had had the right to impose taxes. On the 1st of January
1777, quite shortly after Washington's famous crossing of the
Delaware and the attack on Trenton (Christmas Day 1776),
Robert Morris, who during the war had the disagreeable
post of providing for the common finances, had to go from door to door in
New York to collect money for the continuation of the campaign.

It became worse still after the war, when the enthusiasm
% --- p. 122 ---
\opage{122}which again and again flared up and made it possible to end
the war in so fortunate a manner was over. One turned
now each to his own work, and the sense for the great and common
concerns was lost. In English times one had indeed been
accustomed to having nothing to do with what lay outside
the colony's borders. A shortsighted self-love so often made itself
felt, which refused to give up anything
of the individual state's sovereignty for the benefit of the Union. Yet
there stirred at the same time in many, in part among the most
distinguished men, a conviction of principle, which feared
that the establishment of a firm central power would crush life
in the smaller circles and in the end even the free,
independent life of the individual personalities. Thus
General George Clinton, whom New York for nine years in succession
chose as its governor, regarded the sovereignty of the individual states as
a necessary condition of the people's liberty, and considered those who
worked for a strong Union government as ``the
spokesmen of despotism''. Another esteemed American statesman, Thomas
Jefferson, who nevertheless approved the adoption of the later
Federal Constitution, writes in his Autobiography: ``It is not through
consolidation or concentration of power, but through its
distribution, that good government is brought about. Were this great country
not already divided into states, such a division would have to be made,
in order that each state by itself; might look after what directly
concerns it, and what it can do far better than a distant
authority. Every state is again divided into counties, of which
each shall look after what lies within its borders;
every county again into circles or districts for
the care of smaller particulars, and every district into farms,
each governed by its individual owner. If we were to
receive orders from Washington [the city W.] as to when we should
sow and reap, we should soon lack bread. Only through this
division of labor, descending by degrees from the general to the special,
can human affairs be
managed to the weal and happiness of all''\footnote{\emph{Th. Jefferson's} Works. I. p. 82.}. Jefferson was so ardent a
democrat that by his own account he was not a friend of any
% --- p. 123 ---
\opage{123}energetic government anywhere in the world. After the internal
disturbances that took place at the beginning of the eighties in
Massachusetts, he wrote to a friend that such disorders were
necessary if public liberty was to prevail. Even
if liberty should degenerate into licentiousness, narrow limits
had to be drawn for the government's power. Every individual state
ought to have exclusive authority in its own affairs,
in particular the right of taxation\footnote{Cf. \emph{Curtis}: History of the Constitution of the United States.
II. p. 507.}. --- When such views
were held by distinguished men and had roots in the historical
conditions, it was no wonder if the experiences one had
had of the necessity of union and of the Union did not
receive sufficient attention. There seemed to be an
\emph{Either}---\emph{Or}: either the sovereignty of the individual states or the sovereignty
of the Union, and of these two seemingly irreconcilable
opposites the first had for the time being popular feeling on its side.

Toward foreign countries one had at the conclusion of peace appeared
as one state. One had incurred obligations toward
England; but the individual states took measures that
were wholly contrary to these obligations. Thus those who
had sided with England during the war, and who as a rule were
called Tories, were persecuted and lost their property, although
it had been guaranteed to them by the peace treaty. The Congress
lacked means to carry out its own obligations.
--- It was important to set foreign trade on its feet.
But the foreign powers refused to conclude
commercial treaties with a people where there was no authority
that could win respect. Foreign countries even speculated
on the individual states' conflicting commercial interests. There
arose a veritable tariff and trade war between neighboring states
among themselves. --- Furthermore, the army was to be disbanded, and officers
and soldiers were to receive what was owed them. But the Congress
could not provide the necessary means. Mutinous
soldiers then gathered outside the Congress's hall in
Philadelphia and uttered threatening and scornful cries against the
members of Congress. The governor of Pennsylvania would not
% --- p. 124 ---
% print p. 124: han mødrene (han for hans) -- misprint; the sense is translated
\opage{124}call out the militia, because it was only the Congress that
was insulted! --- Also against the governments of the individual states there were in several
places risings, and at the same time an Indian war threatened.

It was then no wonder that both abroad and
at home one began to harbor doubts about North America's future.
The foreign powers, not only England but also France
and Spain, which in the war had fought on North America's
side, would have preferred to see the country divided. Or it was thought that the
only salvation lay in the introduction of the monarchical
constitution, whether it was to be a French or an English prince
who should come over to hold the country together.

And it was likewise no wonder that among the
distinguished men in North America there prevailed despondency and
misgiving in the first years after the conclusion of peace. One felt
what was lacking: a strong central power; but the task was
how it was to be established, and how it could be
reconciled with the people's democratic habit of thought and with
self-government. It was a question whether the motto adopted by the Congress:
E pluribus unum! was to become reality. And
it was a question whether in North America a new great
power was to take shape, whose popular institutions could form a home
for liberty and self-government, or whether its free states were to
perish in mutual strife, each absorbed in its limited
interests, an easy prey to the intrigues of foreign powers.

\section*{IV.}

One of the men who contributed most to the
happy decision of this question was \emph{Alexander Hamilton}.

He was born on the West Indian island of Nevis on the 11th of January
1757. His paternal family was Scottish; his maternal kin
was a French Huguenot family. His mother, who is said to have
been a gifted and high-minded woman, died early, and as
his father's affairs declined, some of his
maternal relatives took charge of him and had him brought up
on St. Croix. Here he went to school until he was twelve years old,
and has presumably received only the most necessary school knowledge.
Thereafter he was placed in a countinghouse on the same
island, and showed himself here so capable that the head of the business
% --- p. 125 ---
\opage{125}left the management of all affairs to the thirteen-year-old lad, while
he himself for reasons of health went on a longer journey. Business letters
from him have been preserved from this time, which
bear the stamp of early-developed shrewdness and energy. His
leisure he used for reading, above all of economic,
historical, and philosophical writings, and thereby laid the foundation
of the many-sided knowledge he came to need in his later
life. His aspiration already now went beyond the narrow
circumstances in which he lived. In a letter of the year 1769 he writes
to a friend: ``I must confess that ambition is the
predominant thing in me, so that I despise the humble position of a clerk,
to which my fate condemns me, and would gladly
stake my life, though not my character, to win a higher position.
I know well that my youth
excludes me from every hope of immediate promotion,
nor do I wish for such; but I will prepare the way for
the future. I am, as you see, no philosopher, and people
will say that I build castles in the air; my folly makes
me ashamed, and I beg you not to mention
it. Yet one has seen such plans lead to the goal when
the originator was persevering. I will conclude by saying
that I wish there would come a war.'' Seldom is
a youth's ambition expressed in so clear and definite a manner
and with at the same time so distinct a consciousness of what it costs
to reach the goal. No definite goal stands before him, but he
feels his strength. He beats his wings, though he does not
know whither the flight is to go.

The governor of St. Croix took notice of him on
the occasion of a description of an earthquake he had given
in a newspaper. Provision was now made that he might continue
his education, so energetically begun. He came to
New York and studied at King's College. When the war of independence
broke out, the young students were naturally strongly seized by
the movement. At a meeting held outside the city on the
6th of July 1774 to exert pressure on the leading men and
to get them to give up their hesitation, Hamilton appeared for the first
time as a speaker. The fiery and clear eloquence of the seventeen-year-old
captivated the assembly. It was resolved that New York
% --- p. 126 ---
\opage{126}should follow Boston's example and break off all commercial connection
with England. With this Hamilton had entered upon the political
career. In the immediate future he wrote several
pamphlets against a clergyman who had defended the English
ministry's procedure. He maintained here civil
liberty as agreeing with the nature of man and as
necessary for the welfare of society. The English Parliament's
right derives from its electors and is therefore limited to
Great Britain. An English elector has nothing to say in America.
An arbitrary regime is carried on everywhere where a people
is governed by laws in whose making it has itself had
no share. Great Britain is only a free country because its
inhabitants have a share in legislation. True, in every
civil society there must be a supreme power; but for that reason
it is not necessary that one part of the realm should rule over the other parts\footnote{These pamphlets are to be found in the 2nd part of Hamilton's collected Works.}.
--- This pen-feud had a peculiar conclusion. Some
zealous friends of liberty carried off Hamilton's opponent, the clergyman
Seabury, as a prisoner and attacked the printing house from which
his writings had issued. This aroused great indignation, not
only in the English party, but also in the Opposition.
Hamilton himself joined the militia when it pursued those
who wished to end the debate in such a manner. But for the rest
the time of peaceful negotiation was now past. The first blood had
already flowed at Lexington and Bunker's Hill (April and June 1775),
and the struggle was soon in full swing. Hamilton and several other
students from King's College took part in capturing some
cannon from a fort outside the city. They took part in the exercises of the
volunteers, and zealous military studies replaced for Hamilton the
earlier peaceful studies. After having undergone an
examination he was appointed captain of a battery which
the State of New York equipped. The young captain (not twenty
years old!) soon drew by his zeal and ability the
attention of the commanding generals, in particular of Washington.
He became lieutenant colonel, and Washington made
him his adjutant and secretary. In the army he went by the
name of ``the little lion''; Washington liked to call him
% --- p. 127 ---
\opage{127}``my boy''. The close relationship that was now formed between
the two distinguished men lasted throughout life, although
a slight collision caused Hamilton in the year 1781 to leave
Washington's staff in order to take up an active command again.

Hamilton's position in Washington's vicinity gave him
occasion to follow the development of the military and the political conditions
at close range and from an elevated vantage point. He received a
deep impression of the necessity of a firm ordering of the states'
common affairs. Often the possibility must --- amid the external and the
internal difficulties --- have presented itself to him
that the outcome would not be good. In a letter to his fiancée
he wrote at this time that if things should go wrong, they would
settle in Switzerland; only in a free country could they think of
living.

In the midst of the occupations and exertions of the war
his thoughts turned on \emph{how} the Union between the states could be
ordered and made firm. In a letter to Duane, member of
Congress for New York, of the 3rd of September 1780, he sets forth
that the Congress lacks power on account of the individual states'
jealousy, sprung from a misunderstood spirit of liberty. The Congress
has become dependent on the individual states instead of standing above
them. It is a self-contradiction that it has exercised some of
the most important rights of sovereignty: the issuing of
the Declaration of Independence, the raising of army and navy,
the coining of money, the alliance with France, etc. --- and yet
has not possessed authority to procure what was necessary
to bring the war to an end! If the Congress will not itself
assume the authority which manifestly in the ground of things already
is conferred upon it by the nature of circumstances, an
assembly of delegates from all the states must be convened, which
with deciding authority can adopt a new ordering,
in particular the establishment of a vigorous executive power, and the people
must be prepared to agree to its decision by rousing
popular writings. The government of the Union should then have
sovereignty in respect of military affairs, commercial policy, customs
and coinage.

Here for the first time the thought was uttered on which
the Federal Constitution came to rest. The genius of this
% --- p. 128 ---
\opage{128}letter, which Hamilton by his own account wrote in his tent in the midst
of an army that might rather be called a rabble, an army that
lacked clothes, provisions, money, morale, and discipline, ---
the genius lies in the thought that the Union is founded on
the assent of \emph{the whole people}, not merely of \emph{the states}. The elements
of the North American federal state were to be the same as
the elements of the individual states. ``The stream of national power
shall'', as Hamilton expresses it in another place, ``flow
immediately from the source of all legitimate authority''. In each citizen's mind
the Union is to take root. It was this thought that led
to founding the sovereignty of the Union, just as that of the individual states,
on universal suffrage, and that led to the formation of
a \emph{federal state} instead of the earlier \emph{confederation of states}.
The individual state was not to be an intermediary between the individual and
the Union, but the individual was as a citizen to live immediately in
both, just as he lives immediately in family,
commune, and state alike. --- This thought marks a great political
discovery. Neither antiquity, the Middle Ages, nor later
times knew anything wholly corresponding. Since then the Swiss
federal constitution (of 1848) and the constitution of the German Empire
(of 1867) have been founded on the same principle\footnote{\emph{Tocqueville} declared in his book on ``Democracy in America'' that
language lacked a word to designate the peculiar mixture
of national and federative constitution on which the government of the Union rests.
--- \emph{Bluntschli} has (in his ``Politik als Wissenschaft'' 1876)
placed it alongside the Swiss and German constitution, after
a Swiss author \emph{(Rüttimann)} had already earlier made a
comparison between North American and Swiss state constitutions.}.

But there was still a strenuous labor to be done before
the thought stepped out into reality. The year after the
writing of the letter mentioned Hamilton took part in the storming of Yorktown,
one of the finest exploits of the war. He led the right
wing and was the first man in the enemy redoubt. Shortly after
he left the army and established himself as a lawyer. For
a short time he was a member of the Congress, but he withdrew
from it because he was unable to maintain his dignity.
Now came (after the peace) the sorrowful years when lassitude,
selfishness, and shortsightedness replaced the energy and enthusiasm of wartime.
% --- p. 129 ---
\opage{129}``The more I experience'', said Hamilton at this time, ``the
more reason I find for all who love this country to
weep over its blindness''. In vain had he already
before the peace (1781---82) in a series of pamphlets, which appeared
under the title ``The Continentalist'', sought to influence
public opinion. And in vain did Washington, at
Hamilton's request, express himself in the same direction in a circular
letter to the governors of the states, when he laid down the command. As
Hamilton wrote in a letter to Washington: the centrifugal
forces were stronger than the centripetal.

Yet it was not without significance that the thought had come
forward. It had to wait for the conditions for its
realization to come into being; but when attention to a thought
has been awakened, possible points of departure for its
realization are more easily discovered. A point of departure was now found in the material
interests, and in a connection where one would not have
suspected it. After the conclusion of peace Washington occupied
himself with plans for a connection between the Ohio valley and
the river valleys of the East Coast, since he clearly saw its importance for
the connection between East and West, especially as the mouth of the Mississippi
was still in Spain's possession. Delegates from the states
interested in such a canal connection assembled at
Mount Vernon, Washington's manor, at the beginning of the
year 1785. He used the occasion to initiate a
deliberation concerning questions of customs and coinage. These
negotiations led to all the states agreeing to send
delegates to Annapolis in 1786 to discuss a new, common
ordering of the commercial legislation. The meeting was not
sufficiently well attended; but those present resolved to address
to all the states a call, drawn up by Hamilton, to
send delegates to an assembly in Philadelphia in May 1787,
which was to work out proposals for such a change in
the Federal Constitution as would enable it to
satisfy the demands that conditions made upon it. Thus
the matter got under way.

The Convention in Philadelphia assembled on the 14th of May 1787.
North America's most famous men appeared as delegates from
their states. Washington was elected president. Besides
% --- p. 130 ---
\opage{130}him there were present, among others, Benjamin Franklin, James Madison, and Hamilton.
The oppositions here confronted one another sharply.
While Hamilton put forward a proposal that pressed the
individual states strongly down and established a vigorous executive
power, from the opposite side it was demanded that the authority of the
government of the Union be restricted to a minimum. Of New York's
delegates Hamilton was the only one who was in favor of a
real Federal Constitution; his two colleagues opposed it
and left the Convention when they saw that it was becoming serious
with it. James Madison showed himself here as the practical
statesman, who accommodates himself to the real conditions. Only through
various compromises could it succeed in getting a draft
adopted. Among these compromises may be mentioned the one by which
it was determined that while the Congress's House of Representatives was to
be elected by the people, its Senate was to consist of two delegates
from each state. The Senate thus expresses the equal rights
of all the individual states, while the House of Representatives emerges from the people
as a whole, without regard to the individual states. Furthermore
it was agreed that the slave trade might be prohibited by the Union
from 1808 on, while slavery continued to exist and even
procured for the slave states greater political power, in that a slave
was reckoned as \textsuperscript{3}/\textsubscript{5} of a human being in the determination of how many
representatives each individual state was to send. By this
last provision one got for the time being past the
difficult slavery question, which already now divided North and
South into two parties, which had an entirely different view.
The controversy was postponed\footnote{A chief reason for the postponement was no doubt also the expectation
that one cherished, that slavery would soon fall away of itself.
One did not suspect that an immense upswing of
cotton cultivation was impending, occasioned by the invention of machines in the English
cotton industry. At that time no cotton at all was yet sent
from the Southern states to England, and on the whole not much
cotton was grown. With the upswing of cotton cultivation the
politicians of the Southern states soon became fanatical opponents of the emancipation of the Negroes.
Cf. \emph{John Fiske:} The critical Period of Amercian History. % print: Amercian (misprint)
London 1888. p. 266 f. --- An interesting example of how
economic conditions come to influence ethical views!}, until it could find its decision
through a clear and energetic public opinion. And even then
% --- p. 131 ---
\opage{131}the decision came only after a hard struggle, which threatened
to burst the Union asunder --- and perhaps would have burst it, if
the foundation had not been laid in so wise and solid a manner.

Many had been inclined to compromise still more, because they did not believe that popular feeling was for so
great a centralization of power as the finally adopted
draft laid down. When remarks fell in this direction,
Washington rose from his chair and said: ``It
is only too probable that no plan we propose will be
accepted. Perhaps yet another hard struggle must be endured. But
how can we afterwards defend our work, if we,
to please the people, offer what we ourselves disapprove?''

After long debates and sharp collisions the draft of the constitution
was adopted. The Convention had then sat for four months.
It was a solemn moment when the draft was to be signed.
As Washington took the pen, he said: ``If this
excellent Constitution is rejected by the states, there will hardly
be occasion to put another in its place; --- the
next will be written in blood!'' But old Benjamin
Franklin had brighter hopes. Pointing to
a sun that was depicted above the chairman's seat, he said:
``While we sat and deliberated, I was often in doubt whether it
was a rising or a setting sun there; now I am certain
that it is a rising sun!'' ---

Yet there was still a great labor to be done. The draft
was to be adopted by popularly elected assemblies in the individual
states. Here in several places a very strong opposition
was to be overcome. Hamilton now again stepped into the lists for
his favorite thought. Together with Madison and Jay (an
American jurist who had distinguished himself during the
peace negotiations at Versailles) he published from October 1787 to
August 1788 a series of pamphlets under the common title of
``The Federalist'', in which all the provisions of the Constitution were
explained % print: orklaredes (misprint)
and grounded, partly by general considerations proceeding from
human nature, partly by historical examples, partly by regard to the actual conditions. Hamilton
himself wrote by far the greatest part of this work, which
one has called the Bible of Republicanism, and which deserves
% --- p. 132 ---
\opage{132}a place beside the writings of Aristotle, Locke, and Montesquieu.
It even distinguishes itself above these in that
it stands nearer to practical reality, --- proceeds from
it and in turn would lead over into it. It relates to
the Federal Constitution as the motives relate to a draft of a law.
The draft did indeed not coincide with Hamilton's original thought;
but his patriotism and his sense of reality made him
enter so vividly into the adopted draft that he
could now appear as its most important spokesman. In his own
state of New York there was a violent struggle. At the convention chosen by the
population, which assembled in June 1788, two thirds of the
delegates were at the beginning against the draft.
By his enthusiasm and his eloquence Hamilton nevertheless got a
majority won for the cause, and on the 26th of July the draft was
adopted. With New York's accession the Constitution was secured,
although North Carolina acceded only in 1789 and Rhode Island
only in 1790. In the spring of 1789 the Constitution came into force,
and Washington was elected the first President of the Union.

The great procession which after the adoption of the draft passed
through the streets of New York paid its homage especially to Hamilton.
His picture was to be found on the banners, partly alone, partly together
with Washington's. It was a well-deserved honor for the
man whose thought and whose word had accomplished such great things.

\section*{V.}

Before we follow Hamilton further on his course, we shall
dwell a little on the remarkable work of which he was the
chief author, and bring out its main thought.

The political discovery Hamilton already in 1780 enunciated
was the founding of the Union on a popular basis, instead
of a statal basis. In somewhat more abstract form
it can be expressed thus: there was found a more thorough form of
\emph{the connection of unity and liberty} than had hitherto been found.
In the republics of antiquity there was civil and political liberty;
but it seemed as if history had demonstrated that this
liberty could exist only with a limited territory.
As soon as a great realm was to be formed, liberty had to
vanish before unity. Yes, even in the ancient republics
% --- p. 133 ---
\opage{133}individual liberty was in many ways severely curtailed, in order
that the unity of the state should not suffer. \emph{Montesquieu} taught
expressly that the republican constitution was possible only with a
limited territory. Only there could the feeling of unity,
the immediate absorption in the common interests, arise, which
the republican form of government presupposes. All citizens
could take part in the decision of affairs\footnote{Esprit des Lois. III, 3, cf. IX, 1.}. \emph{Frederick the Great}
remarked in the year 1782 to the English envoy in Berlin that he
was convinced that the American Union could not
last long in its republican form; the very great
extent of the country would be an obstacle, for a republican
government has never existed long where the territory
was not limited and concentrated\footnote{\emph{Bancroft:} History of the Constitution of the United States. I p. 71.}. Even in our century
it has often been said by European authors that only
monarchy could be the form for a comprehensive state organism\footnote{See e.g. \emph{Hegel:} Philosophie der Geschichte. 3. Aufl. p. 379. ---
\emph{J. N. Madvig:} Blik paa Oldtidens Statsforfatninger.
(Universitetsprogram). København 1840. p. 12 and 21.}.

It was an objection which the Union's opponents
did not fail to make use of. Liberty's cause was for these bound up
with the individual states' absolute sovereignty. Hamilton answers
that Montesquieu knew only the small republics of antiquity, and
that history had not yet shown any application of the
representative system in a great country with a republican
constitution. The republics of antiquity presupposed all citizens'
\emph{personal} participation in the decision of affairs and knew only
representation to a very limited extent. If one will therefore
appeal to Montesquieu, one must be so
consistent as to demand that several of the greatest North American
individual states be divided; for they are too large to be republics in the
sense in which Montesquieu took this word. Moreover they overlook
that he expressly recommends a league between free
states as the only means of extending the republican
constitution over a larger area. Such a league
has now been tried in America and has proved insufficient. It has
shown that if it is in earnest about giving the Union strength
% --- p. 134 ---
\opage{134}and permanence, the federal power must not be based on the good
will of the individual states' governments, but it must be able to
exercise its sovereignty immediately over all citizens, just
as well as the individual states do. It is merely a matter of
dividing the rights between the Union and the individual states, so
that no collision becomes possible. But the Union's
authority must in the end spring from the same source as the
individual state governments' authority: from the people's very will. Unity
must spring from liberty itself\footnote{In the German Empire universal suffrage has been introduced as a necessary
means to strengthen the Empire's power over against the individual states. Cf.
\emph{Bismarck:} Reden. III. p. 194. But whereas in America one has
let unity exist so far as it was compatible with liberty,
in Germany the opposite way has been gone, in that liberty only
gets leave to exist so far as it is compatible with the
strict unity which the existence and power of the new Empire seem to demand.}; then it will not show itself
to be hostile to the latter, and then it will be strong enough
to let the smaller circles exist without absorbing
them. A great society will far more easily let self-government
flourish than a small one. A single party cannot so easily
seize an exclusive and lasting power in a large as in a
small country, if only the institutions are sufficiently free.
Only justice and the general welfare will be able to gather
very great majorities in a country of great extent.

Another main thought which appears in ``The Federalist''
is that a conservative tendency in the constitution is very well
compatible with free and republican institutions. Just as
the American Federal Constitution is characterized by
the connection of \emph{unity and liberty}, so it is peculiar
in the connection of \emph{stability and liberty}, and has precisely in the
most recent time awakened attention among conservative authors.

The Federal Constitution rests on the sovereignty of the people.
Hamilton laid great weight on the preamble to
the Constitution coming to read: ``We, the people of the united Free States,
ordain and establish this Constitution''. The Constitution was to
that extent founded not only on the people's power, but also on
confidence in the manner in which the people would use this power.
``Just as there is'', says Hamilton, ``a depravity in
% --- p. 135 ---
\opage{135}human nature which makes caution and distrust necessary,
so there are other qualities in human nature
which justify a certain degree of esteem and confidence. The
republican constitution presupposes the presence of these
latter qualities in a higher degree than any other
constitution. If the descriptions which some authors have
given were true renderings of human character,
one would have to draw the conclusion that there is not sufficient
virtue in men for self-government, and that only the
fetters of despotism can prevent them from destroying and devouring
one another''\footnote{That Hamilton himself had not exactly reaped good experiences of
men is seen from a characteristic statement of his in a letter to his
friend Meade (in the year 1782): ``Nothing shall interrupt our friendship. My
friendship for you is built on the solid foundation of a full
conviction that you deserve it, and that it is mutual; and it
is so much the more firmly grounded as you have few rivals.
Experience is a constant commentary on the unworthiness of the human race, and the few exceptions we find have so much the greater right
to be prized, the more seldom they are. I know few who are
respectable, fewer still who deserve to be loved; and when I meet
any of the latter kind, it is not in my power to hold back my
feelings.'' (Works I. p. 298).}. The republican principle thus rests on
a confidence in the people. ``But it does not demand unconditional yielding
to every flaring-up passion, or every passing
impulse, which might stir among the people through
the artifices of such men as flatter its prejudices
in order to betray its interests. It is a correct observation
that the people in general \emph{intend} the public good. In their very
errors this often shows itself. But its sound sense
would despise the flatterer who would maintain that it always
\emph{thinks rightly about the means} of promoting the public good.
It knows from experience that it sometimes goes astray; and it is a wonder
that it goes astray as seldom as it does, although it
is constantly exposed to the arts of parasites and false informers,
to the snares of the ambitious, the covetous, and the unscrupulous,
to the cunning of those who possess its confidence without deserving it.
When cases arise in which the people's interests conflict
% --- p. 136 ---
\opage{136}with its inclinations, it is a duty for those whom the people
has chosen to guard its interests, to oppose
the momentary error in order to give it time and
opportunity for calm and considered deliberation''\footnote{The Federalist. Philadelphia 1882. p. 427, 533.}.

From this consideration is derived the necessity of
conservative elements in a republican constitution. The most important
of these elements are the following.

1) An executive authority which can act with vigor
and firmness both outwardly and inwardly is the most important
condition. Vigor and firmness in the executive power rest
in turn on the power's being placed in one man's hand, on its having
a certain duration, and on the power it commands being
sufficiently great. When the President of the Republic is elected by
the people (indirectly, through electors who can deliberate on
the most worthy) for four years and can be called to account
for his conduct of office, then the two conditions will be
united which seem to many wholly irreconcilable, namely extensive
authority and dependence % print: Athængighed (misprint)
on the people. (Yet Hamilton had
originally preferred to have a president elected for life,
who could be removed only by sentence of a court of impeachment). When the
free constitutions in Europe restrict the executive power
and cherish a constant suspicion toward it, this comes of
their being agreements between a hereditary dynasty and
the people. They have a certain negative character and consist
essentially in an enumeration of what the ruler may \emph{not} do, of the
rights he does \emph{not} have. Thus Magna Charta, Petition
of Rights, Bill of Rights. The American constitution is,
by contrast, expressly founded on the people's power and is to be maintained
by its immediate representatives and servants. Therefore
there is no need for so great anxiety. Hamilton would even
find it questionable if the Federal Constitution enumerated
particular ``rights of man'' which the government might not
violate; for one could then easily draw the conclusion that what
was not found in this enumeration was permitted. This is
Hamilton's answer to those who missed in the draft an express
``bill of rights''\footnote{The Federalist, p. 360. --- This remark is apt, especially
when one thinks of the assumption prevailing in many European countries, that what is not in plain words forbidden to a government
is permitted to it.} He had here a clear eye for the
% --- p. 137 ---
\opage{137}peculiar conditions under which the American
constitution came into being: as a positive enactment of the people's collective
power, not as an accord or a contract between contending
powers in the people. The ground for jealousy of the
executive power which was found, and is constantly found, in Europe, and which
has made so many free constitutions unfruitful, was not found in
America. Locke's and Montesquieu's theory of the division of powers
could therefore be carried through there in purer and more consistent form.
The American Republic is no ``parliamentary'' republic.

2) The second conservative element lies in
the bicameral system. By the provision that the Senate's members
should be elected by the individual states, two from each state, it succeeded
in letting the independence of the states come forward in the
Constitution. But at the same time the bicameral system gives occasion to
more thorough deliberation than the unicameral system. A single,
numerous assembly easily yields to a sudden impulse.
At the same time expert knowledge can be specially represented in the Senate.

3) Thirdly, constitutional changes are
difficult to carry through. They must be proposed either by the Congress,
when two thirds of both Houses are agreed thereon, or by
a Convention, which is summoned by the Congress at the request
of the legislative assemblies of two thirds of the individual states.
For the adoption of the changes it is required that three fourths
of the legislative assemblies of the individual states, or three
fourths of conventions summoned in the individual states, vote
for them. The provision that the Senate shall consist of an equal
number of members from the different states cannot be
changed. --- Whereas the English House of Commons by simple law
can alter the English constitution, so that this finds itself
in an uninterrupted becoming, greater obstacles have here been set
to changes.

4) The courts are wholly independent of the executive and
the legislative power. The supreme court of the Union can declare
laws which conflict with the Constitution invalid. It has
thus the right to interpret the Constitution. The purpose of
% --- p. 138 ---
\opage{138}this provision is to prevent the people's will expressed in particular, special
laws from coming into conflict with the permanent and general will
expressed in the Constitution. The
authority which the legislative power has must be
subordinate to the authority which has put the Constitution itself into
force. --- Like the distinction between the executive and the
legislative power, so also the distinction between the
judicial power and the two other authorities is more sharply
carried through in North America than in European countries. It is an
element in the Federal Constitution which we have here reckoned among
its conservative elements, but which could also be reckoned
among the elements of liberty. Montesquieu (Esprit des lois XI, 6)
says rightly that there is no liberty where the judicial
power is not separated from the legislative and executive. ---

The most important of the fundamental thoughts here mentioned
Hamilton himself conceived from the first. The others (in particular that concerning the
composition of the Senate) he adopted and sought to ground
more closely. Only a few components of the rich content of ``The
Federalist'' have here been brought forward. On a whole series of more
subordinate points one finds detailed and clear grounding in
the remarkable work. One must wish the European
peoples that the time may come when their conservative parties
will support themselves on so clear, sound, and popular a way of thinking
as that which runs through ``The Federalist''. The book is a
combination, unique in its kind, of a work of political philosophy
and a tract of agitation. But seldom has a tract of agitation
appeared in so pure and noble a form. To give an
idea of the elevated standpoint from which its
chief author viewed the so strongly strained situation, in which
there was no lack of coarse and hateful words from
both sides, I shall quote a statement from the first
number.

``Among the most formidable obstacles which the new
Constitution will have to overcome, one will easily note the
manifest interest which a certain class of men in every
state has in opposing all changes that could
entail a diminution of the power, revenue, and influence
which they possess by their offices under the state institutions, as
% --- p. 139 ---
\opage{139}well as the corrupted ambition of another class
of men, who either hope to rise higher in the confused
conditions of their country or flatter themselves with obtaining greater
power by the division of the realm into various smaller confederacies.

``It is not, however, my intention to dwell on
considerations of this nature. I perceive that it would be ignoble
to explain without distinction the opposition of any class
of men by selfish or ambitious motives,
merely because their positions might make them the object of
such suspicion. Sincerity compels us to admit
that such men too may be animated by honest
intentions, and it cannot be doubted that much of the
opposition which has already appeared, or will hereafter appear,
springs from sources which cannot be blamed, if they
cannot be esteemed: honorable errors, which are called forth
by preconceived suspicion and fear. So numerous and powerful
indeed are the causes that can lead the judgment astray that we on
many occasions see wise and good men as well on the
wrong as on the right side in questions of the very greatest
importance to society. This circumstance will, when
it is rightly considered, always afford a salutary lesson to those
who are involved in some dispute, however convinced
they may be of being in the right. And a ground for caution
lies also in this, that we are not sure that those who
defend the truth are animated by purer principles than their
opponents''.

One cannot better make clear what speaks
for holding a debate in purely objective form. And better has this
form hardly been observed in any political debate.

\section*{VI.}

In Washington's first ministry Hamilton entered
as Secretary of the Treasury. It thereby fell to his lot to make firm
the work whose fundamental thought he had first conceived, and whose
nature he had most thoroughly developed. He consolidated the
public debt, so that the Union assumed the individual states'
debts. The common public debt was to serve as a new means
to hold the Union together by creating common interests.
% --- p. 140 ---
\opage{140}Moreover he sought to procure revenue for the Union by a law
on excise, and he had a National Bank established. It is
acknowledged from all sides that it was due to his wise
measures that the country's standing and credit rose, and that
trade and industry came to flourish.

When Talleyrand in 1794---1795 sojourned in America (because
he was permitted to stay neither in England nor in France),
he associated very intimately with Hamilton, for whom he
cherished the greatest admiration. He regarded in his old age
Napoleon, Fox, and Hamilton as the greatest men he had
met on his remarkable and winding path, and were he to place
one of them highest, he was inclined to give Hamilton that
place. He said of him that ``he had divined Europe'',
whereby he alluded to Hamilton's keen eye for the
political situation and the possibilities it contained.

But Hamilton's personality was not suited to
special politics with its party struggles and its
personal collisions. He was unlucky as a party leader, and he
understood better how to fight on his own than to get a close-knit
host to follow him through the many
vacillations of the political struggle. But as unlucky as he in this respect
was, so calmly and resolutely did he go the way along which
his conviction led him, whether he was solitary
or was followed by many.

Party struggle began soon again after the adoption of the
Constitution. In Washington's Cabinet there sat opposite Hamilton
Jefferson as his zealous opponent, who maintained the
diametrically opposite political conception. He looked with
misgiving on Hamilton's endeavors to make the Union firm.
These endeavors must, Jefferson feared, lead in the direction
of a monarchy. He maintained the individual states' right to
refuse obedience to the government of the Union, when they thought that
it exceeded its competence. This was the so-called
principle of nullification, which in the following ages' struggle between North
and South came to play so great a role.

In the year 1795 Hamilton left the ministry and
lived now the rest of his life as a lawyer. Only once did
he again enter public service, when during the troop
% --- p. 141 ---
\opage{141}levy which with Washington as commander-in-chief was undertaken in
the year 1798 on account of the strained relation to France,
he was at Washington's wish appointed second in command.
When Washington died, he became commander-in-chief in his stead.

As Hamilton's party, the so-called Federalists, did not
have fortune on its side against the ever mightier
``Republicans'' (thus Jefferson's party, which corresponds to the
later ``Democrats'', called itself as a protest against the Federalists'
alleged monarchical disposition), many of its adherents showed
an unpatriotic fanaticism which Hamilton was far from sharing.
He thus supported Jefferson's election as President, when one
had only the choice between him and the criminal
adventurer Colonel Burr, for the latter of whom the Federalists voted
out of hatred of Jefferson. And when Jefferson succeeded in
acquiring Louisiana from France, Hamilton approved this
step, while his fanatical party comrades were against it,
because they feared that the South, where the ``Republicans'' constantly
won more adherents, would thereby become mightier. Some
Federalists, among others Burr, even conceived the thought of
dissolving the Union and forming a separate Union for the Northern states.
Nothing could conflict more with Hamilton's conviction, and
he therefore combated with all his might and with strong words Burr's
attempt to be elected governor in New York, since he
feared that Burr would use this position to work
for his treasonable plans. This gave Burr occasion
to challenge Hamilton, and although the latter considered duels
improper, he yet thought he could not decline the challenge,
Only 47 years old he was severely wounded by Burr's shot
on the 11th of July 1804 and died the day after. He sacrificed his life
for the purity of his cause.

\streg{}

% --- p. 142 ---
\opage{142}
\essay[Social Pessimism]{SOCIAL PESSIMISM.}

\begin{center}\emph{Ferdinand Tönnies:} Gemeinschaft und Gesellschaft. Leipzig 1887.\end{center}

\streg{}

\section*{I.}

When an optimistic and a pessimistic view stand
opposed to each other, it is as a rule probably moods that are
decisive, moods that have their
origin either in the temperament of the individuals or in the limited
and accidental experiences they have had or think they have
had. Yet there may sometimes be question of more
than mere mood. It is thus clear that the processes in
an organic body can take a course that makes it
possible for the expert to predict a dissolution of life.
And physicists have sought to show that the time will come when
all life and all motion in our solar system will cease, in that
an equal distribution of heat will set in, and no
conversion of heat into motion will any longer take place.
They derive this from the fact that, while all motion can be converted
into heat, not all heat can again be converted into motion,
but a part of it will spread to the surroundings. If
one imagines this process completed, we get standstill. Should it
not, in a similar way, be possible to decide whether
in the domain of human social life a
dissolution and a decline or a progress is taking place?

One might think that the matter must be clear in the
eyes of those who embrace the hypothesis of development. But one would then
confuse development and progress. Development is a
process of adaptation, by which individuals and species assume
such forms and become able to exercise such
activities as the conditions of life, the external circumstances, demand.
% --- p. 143 ---
\opage{143}But such an adaptation can under certain conditions lead
to forms and faculties being diminished and impaired instead
of unfolding more richly and more variously. In parasites
there are thus found traces of a definite decline, in that organs
that in related animal species appear with a rich structure
are shrunken to feeble vestiges, and the corresponding
faculties have likewise vanished. Human beings too can
adapt themselves to such slight and simple conditions that they lose
the higher faculties that were innate in them. It is not
merely in passing from more various and greater conditions
to narrower and slighter ones, however, that such a decline can occur.
In passing from simple and even, but easily surveyable
conditions to more composite and comprehensive, but at the same time
also more artificial conditions, faculties of great value
may likewise be lost. There can be advance in variety, and
yet an essential retrogression take place in concentration,
in inner connection and strength. This happens easily, for example, when
the limited, but sure and secure instinct is replaced by the
more comprehensive, but often doubting and erring reason. In
man, therefore, excesses and unnaturalnesses become possible
which the animal led by instinct does not know. According to Darwin's
conception a retrogression took place shortly after ``the much
earlier time when man had just become man'',
and when reason on so many points was beginning to replace
instinct\footnote{See ``Menneskets Afstamning og Parringsvalget''. Danish trans. II
p. 364, 379, 400. --- In this connection one may recall Kant's
and Schiller's interpretation of the myth of the Fall as an expression of
the transition from instinct to reason.}. Human development, both in the
single individual and in whole peoples, often shows the
decline that in any case on certain points --- can
be connected with a transition from simple and limited
to composite and various conditions of life. The
objections to the value of culture that from time to time come forward
with great emphasis have their relative justification in
this circumstance.

The investigations of recent years into the development of conditions of society
and of ideas of right among the Indo-European
% --- p. 144 ---
\opage{144}peoples have now even led to the setting up of a whole social
pathology. It has been sought to establish the claim that the whole
modern development of culture is in reality one great
process of dissolution, which advances quietly and imperceptibly, revealing
itself only to the comparative and consistent
investigator. The above-mentioned work by \emph{Ferdinand Tönnies} (from
Husum, lecturer at Kiel) is an interesting and, as it
seems to me, very instructive contribution in this direction.
The author, who has earlier made himself known by his studies
on Hobbes and Spinoza, works that belong to the
best that has been achieved in the domain of the history of philosophy, has
the same objective calm and undisturbedness in his consideration of
the phenomena of human life that distinguishes those older
thinkers with whom he has occupied himself so much. One does not notice
in him the passion of a Rousseau or a Schopenhauer.
He investigates and ascertains in a purely theoretical way, and demands to be
met in the same way. He
combines sociology and psychology in a peculiar manner,
in that he shows how social development necessarily
is connected with and has its counterpart in a corresponding
development of the human mental faculties. The rich material
on which he builds in his investigation he has so made
his own that he has converted it into a circle of thoughts
that are often developed in very abstract form and with the use
of newly formed concepts and expressions, but always in such a way that
one feels firm ground under one's feet. Although the book demands
penetrating study in order to be understood, it rewards the labor
spent on it. Much that was hitherto known comes
to stand in a new light, and several new points of view
are advanced. Whoever in earnest wishes to understand his age and its
tendencies ought not to shun the labor of trying a bout with
this author. I shall here engage in such a bout,
in that I shall give an account of his standpoint and then
examine what can be urged against it.

\section*{II.}

According to the more recent investigations, the present moral and social
conditions of the European peoples have developed out of
% --- p. 145 ---
\opage{145}an earlier order of things in which the sharply marked
legal order, with a definite determination of every person's claims
and duties, did not exist, any more than did the great
struggle of interests in which each individual's gain
or loss is determined by the relation between supply and
demand. Neither the juridical nor the
economic distinction between individuals existed which has developed
among the Indo-European peoples, who from their home in
the distant East gradually settled in the western regions.
Remnants of the old order of society are found nowadays
in the country more than in the towns, among the common people more than among
the so-called educated, in the eastern countries more than in the
western. In India, for instance, there still exists the old village order
founded on the relation of family and neighborhood. It was
the bonds of blood and of spatial proximity that united
human beings. Immediate sympathy and understanding, personal
acquaintance and a cooperation strengthened by living together were the
forces that bore society. It was a natural,
self-evident, involuntary living together, not entered into with clear consciousness
and in the hope of definite individual advantages. The individual
persons felt themselves immediately as parts of a whole and
ascribed to themselves no existence outside it. The rules for their
conduct were given in ancient customs, ordinances, and
traditions, which were not discussed, but followed as an
order of nature and regarded with reverence. Also the larger
course of the world, on which the fate of the family and the tribe depended,
was apprehended and appraised not with testing inquiry and
contemplation, but with lively imagination, feeling, and faith.
Here too one saw an order of things that met the
individual persons, at one and the same time as something near and
homely, and yet as something venerable and ruling.

The course of culture has now --- as Tönnies shows --- more
and more led beyond this condition. A progressive
emancipation has taken place. The individual has obtained an
independent right, apart from the narrow, natural society to which he
belongs. Definite boundaries have been set between the
individuals' spheres of dominion. They are torn out of their immediate
connection, have become centers each for himself instead of having
% --- p. 146 ---
\opage{146}a common center. Intercourse and connection do indeed take place
between them, but intercourse and connection are determined
now more and more by what each individual thinks he can
obtain from the others. \emph{Society} (Gesellschaft) takes more and more
the place of \emph{Community} (Gemeinschaft). In place of
instinct and the feeling immediately determined by tradition, authority, and brotherhood
there now comes prudent,
calculating deliberation. External interests connect human beings,
without their therefore standing in any essential inner connection.
The one wishes, with the least possible sacrifice on his own part, to obtain
the most possible of what is desirable to him that the other
possesses. \emph{The commercial relation} is the most prominent example
of Society as opposed to Community, and trade
is at the same time the most important means by which Society replaces
Community\footnote{Also according to \emph{Leist} (Alt-arisches Jus Gentium. Leipzig 1889 p.
523---525) the ancient Aryan household order perished through
the rise of trade and the money economy, which had as a consequence the sharp
development of the purely private concept of property. It is ``the
demonic power of capital'' that transforms conditions. In place of
the household community and the occupational activity bound to it,
self-interest now presses forward, which is hostile to every lasting
community. Leist points out aptly how intelligible
and historically significant in this connection Aristotle's
polemic (in his doctrine of the state) against every pursuit of gain that goes
beyond what is necessary for the household comes to appear
to us. It was nevertheless a reactionary line of thought that was not able to
halt development. --- In a similar way trade has worked
in the transition to modern times. It was the changed
conditions of sale, the formation of a world market in place of the earlier
limited conditions of sale, that led to the break with the old
guild order, to the replacement of handicraft by large-scale industry, to
the invention of machines and to all modern competition. Cf.
\emph{L. Brentano:} Über die Ursachen der heutigen socialen Noth. Leipzig
1889 p. 16. Brentano reproaches Marx (whom Tönnies on this
point follows rather too much) with turning the relation around, in that he
lets large-scale industry develop out of itself and then
produce the world market, just as Marx in general in his theories
overlooks the significance of demand.}. In ``Society'' one abstracts from all original
or natural relations among human beings. A
Society presupposes only a plurality of naked persons who
% --- p. 147 ---
\opage{147}are in a position to render services, and consequently also to promise. In return for
getting means for \emph{his} purposes himself, the one renders the other
means for \emph{his} purposes. In the transition from Community
to Society we approach a condition in which everyone is
a merchant.

In \emph{Community}, in the old society, the connection
between individuals was given before they could reciprocally render
each other services, and before they could reckon on such
reciprocal services. In \emph{Society}, by contrast, individuals for
the most part have to do with one another only when their interests
meet in a harmonious or disharmonious way. An artificial,
mechanical connection takes the place of the inner one.
The joint-stock company is the characteristic type of Society, as
the family is of Community. In place of agriculture,
handicraft, and art come industry and trade. In place of
the personal relation between the craftsman and his customer
comes the impersonal relation between the mass producer and
his accidental, unknown buyers. The metropolis with its
motley, international life takes the place of the house,
the family, and the village. Society culminates finally in
the relation between capitalist and worker as a pure relation of purchase and
contract. To use Sir Henry Maine's expression:
the natural, inherited condition \emph{(status)} is more and
more replaced by that arbitrarily, with conscious deliberation brought about
\emph{(contractus)}.

There takes place here, according to Tönnies, a pathological
process of development, which can be altered, but not prevented, by
philanthropists and legislators. Whether that ``dualistic
construction of Society'', which shows itself in the opposition between
capitalists and workers, is the only possible one, --- what
will happen when the line has run out, and all relations have
become relations of contract, --- these are questions that
the author declares to be no concern of his.

That it is a \emph{pathological} development that is here described
is seen especially when we consider the psychological side of
the development, not merely the social. As Society is related
to Community, so \emph{arbitrary willing} is related
to \emph{essential will}. By essential will Tönnies understands instinct,
% --- p. 148 ---
\opage{148}immediate impulse, which goes with inner necessity
toward its goal. It works predominantly unconsciously, according to
natural disposition and habit or determined by imitation and
immediate sympathy. In arbitrary willing, by contrast, a
definite distinction is made between end and means, and it becomes a matter
of deliberation how much it is worth to sacrifice on the means
in order to attain the end. Every arbitrary willing is a purchase,
sacrifices something in order to gain something. The peculiarity of
arbitrary willing comes clearly to the fore when the means is willed solely
and alone for the sake of the end. There is therefore in
arbitrary willing something unnatural, a certain self-constraint, a
falsity. When one makes clear to oneself what one sacrifices,
and what one thereby wishes to attain, egoism will at the same time lie
close at hand and will displace immediate devotion. In
arbitrariness it is not goodness, but only prudence and
consistency that are to be admired. Life no longer stands
as a calling in which one is to work, but as a
business that is to be done.

The development of Society favors the tendency toward
arbitrary and egoistic willing at the expense of essential will and
devotion. Under the constant rivalry the ruthless use
of arbitrariness becomes a condition of the individual's
subsistence. Clear understanding takes the place, in the whole spiritual
life, of immediate feeling, fantasy, and faith.
A progressive intellectualism takes place; scientific
reflection takes more and more the place of
childlike faith. The religious and the social problem are
thus branches of one and the same great problem, which Tönnies
treats partly from the sociological, partly from the
psychological side.

\section*{III}

No one will be able to deny that such a
sociological and psychological development as Tönnies describes is taking place.
Culture penetrates where before there was only nature;
the forms of life become ever more composite and to a higher degree
the fruit of arbitrary ordering; ever greater
% --- p. 149 ---
\opage{149}demands are made on shrewdness and calculation. We seem in this respect ever
more to move away from the primitive, the natural, and the involuntary.
Will the line not then run out, so that human life
petrifies or dissolves, as according to a physical theory it will
go with the life and motion of the terrestrial globe and our solar system
as a whole? Or might the development described by Tönnies
be merely one side of the matter, in that to the attentive
eye other tendencies of development will appear, which can
hold their own against it, so that it does not come to run
the line out?

Let us see briefly what experiences and considerations
can lead us to win from the matter a brighter side. ---

Pessimism is here, as so often, akin to
\emph{Romanticism}. It sees the past in a radiant light, while
the present seems cold, empty, and dark. The old conditions of society
prior to the great process of emancipation, which fills so
large a part of history, were by no means ideal. They
appear to us with the character of unity and closedness
and at the same time with a certain idyllic stamp. Human life and the
human feelings appear in simple, intense, and
elementary forms, which are beneficial to an eye wearied by the motley,
divided, and contending life of the present. In Tönnies
it seems to me in particular that one notices the pious
attachment to life as it shapes itself in the country and in the
little hamlets; he portrays it and its forms with a
warmth which this clear thinker is not able to
maintain where it is a matter of understanding the conditions of the metropolis and of the
composite life of culture, and which is quite certainly
also harder to preserve there. But the idyll is often only
a semblance --- in the human world as in the rest of
nature. Under the old patriarchal system, for example,
much cruelty and oppression could be in vogue.
The father of the house ruled without control, subject only to the gods'
punishment; but they did not always make themselves heard. Of the worst
history perhaps keeps silent, for the wholly rightless get no
history. It was in reality no fall when the
ancient Aryan family system, built in part on priestly tradition and
priestly rule, was replaced by the clearly developed
% --- p. 150 ---
\opage{150}Greek-Roman civic system, within which private law
underwent its fundamental process of development. It was not
merely egoism and selfish interest in gain that burst the
ancient Aryan household order and in a later time the
medieval guild order. It was not least an --- awakened and nourished
by the changed economic conditions --- more explicit
recognition of the individual personality's worth and capacity.

Tönnies calls attention, no doubt with full right, to the
great difference between the thralls within the ancient Aryan
families and the slaves of later times. The thrall's (der Knecht)
position often approached that of the son or the friend; he was
a member of the household as well as wife and children. ``Since
he shares the family's joy and sorrow, renders his master
the reverence of a grown son and enjoys the trusted position of a helper or even of an
adviser, he is by his whole
nature a free human being.'' Only by way of exception is he pressed down by
maltreatment to a truly slavish position. Later
Leist has emphasized that the institution of thralldom in antiquity, still among the
Greeks, must be regarded and judged neither from a purely
private-law standpoint nor from the standpoint of public
law. ``In the household order we have before us an Aryan primal legal institution
which is older than the distinction between private and public
law. It is the juridical microcosm from which
private law and public law have only gradually developed
themselves.'' As with the pursuit of gain, so we understand
also here Aristotle's standpoint (in particular his well-known
defense of the institution of thralldom) as an attempt to
hold fast the old household order against the advancing
individualism.

This is an important consideration, by which much in
history is judged in a more just way than an abstract
philanthropic and natural-law view is able to do. But
precisely when the thrall's position has often and perhaps for long periods
been comparatively good and secured by patriarchal
care, it would be wholly incomprehensible how wishes could
arise to leave this position, if there were not
abuse and oppression. It has presumably
gone here as with the development of the free class of workers
% --- p. 151 ---
\opage{151}at the end of the Middle Ages, which was called forth by the
egoism with which the guild masters sought to keep their
advantages and privileges restricted to as small a circle as
possible, while at the same time (as already touched on in another
connection) the market expanded and ever more
labor forces had to be seized upon\footnote{\emph{L. Brentano} describes (in his work: Die Arbeitergilden I p. 59) this
matter thus: ``The flourishing of industry and the greater employment
of capital in it, which reciprocally called forth and conditioned each other,
necessarily drew the serfs in multitudes to the town. Thereby
provision was indeed made on the one hand for the greater number of laborers
that the progress of industry required, but on the other hand there arose
in every new worker a possible future competitor. The yield
of capital investment in industry was thus threatened; anxiety
over this awakened the monopoly spirit of the craftsmen, and there arose a multitude of
regulations that restricted the competition of those newly emerging families.
The payment for admission to the guild was raised; on the
Continent there arose the custom that of everyone who wished to practice
a handicraft independently a for the most part costly
masterpiece was required, whereas sons and sons-in-law as well as those who
married a widow belonging to the guild enjoyed
easements.''}.

The old system was not merely a \emph{patriarchal
household system;} it was at the same time a \emph{system of compulsion}, or it had in
any case at every moment the possibility of passing over into
a system of compulsion. But whether it worked by
solicitous mildness or by compulsion, one thing it was not suited
for: to develop \emph{independence} in faculty and drive and the \emph{self-regard}
connected with it in the individual persons. These stood ---
with the exception of those who were authorities (fathers, priests,
prophets, etc.) --- as more or less selfless parts of the
whole. But they were not, each for himself, small wholes.
The transition to the more individualistic order (from Community
to Society), whether it was most called forth by the
egoism of the authoritarian holders of power or by the economic conditions
changed through external circumstances, has in any
case brought about the good that there came to be a possibility of, and
a drive was awakened to, the unfolding of faculties and powers that hitherto lay
slumbering in the individuals. Thereby
there was certainly also made possible isolation and externality in the relation between the
% --- p. 152 ---
\opage{152}individuals. But there was nevertheless set the great task for each
individual of becoming an independent and distinctive
personality; a multitude of independent points of departure
was laid down for the life of the race. A greater inwardness in human
life, wholly new fields of experience of great value were thereby
opened up. Even if the activity of individuals does not go so
easily and instinctively as in the old society, a
great good has nevertheless been attained. In the course of development
something is unavoidably lost every time something is gained.
That is how it is. But it is unjustified merely to look
at what is lost. The highest thing would certainly be to
unite the valuable from the earlier stages of development with
the valuable from the following stages. But that is an ideal
that cannot fully be attained, although it has its great significance
to keep it clearly before one's eyes. The mission of the pessimistic view
is to prevent us from forgetting it.

Not only a greater independence and inwardness becomes
possible, but also a \emph{greater society}. The state and
humanity, those are concepts that could develop only when the
patriarchal order releases the individual persons and
permits them to feel themselves in human relations, even if they
are separated from the original home. This side of
development too presents its dangers; the narrower circles, with
their immediate and therefore more heartfelt contacts, are
easily pushed aside for the larger, but more abstract relations, and
life can thereby suffer in strength and warmth. But in itself
it was nevertheless a great and significant expansion of
the feeling of humanity that took place when state and
humanity became leading ideas in human life. Barriers
were torn down that in many ways brought hindrance and
strife\footnote{In ``\emph{Dansk Bondeliv}, saaledes som det i Mands Minde førtes, navnlig
i Vestjylland'' (København 1889) Pastor \emph{H. F. Feilberg} has given
an interesting portrayal of life, including the social life of
the village, as it was led before the modern reforms, the
extended means of communication, enlightenment, and the political movement brought
new conditions and a new spirit. The greater the love and
understanding with which the author dwells on the old order of things, the more interesting it is to hear his judgment of
the change that is in the course of taking place and to some extent has already
taken place. He emphasizes what narrow outlook and what
class prejudices had to arise under the isolation and limitation
under which life was led. One saw only ``to the boundary of one's own field,
or at most to the parish boundary; strange it would have been and almost inconceivable
if one had met people with a view of the great society, land and people''. There has now,
as a consequence of the causes mentioned, not least of the political
movement, even in the poorest homes come ``\emph{something of a
feeling of community, such as surely has never been there before}''. The worker and
the cottager have learned that ``they belong to a people, that they have common
interests with their people, that they have a right and a responsibility'' (p. 130,
137). --- If the author is right in this, we see how the dissolution
of the old town communities with their communal life and their humane and social connections resting on
neighborly relations can be
a transition to a larger and more comprehensive common feeling and to
a consciousness of greater tasks and points of view than the old
order's limited horizons made possible.}. ---

% --- p. 153 ---
\opage{153}The difficulty of uniting the goods of the earlier stages of development
with those of the later is now to a large extent relieved by the fact that
\emph{the race in a certain sense ever begins afresh}. It goes
here as in nature, where we have more and more learned that the
forces that have worked in the earlier periods of the earth are not
essentially different from those that to this day are in
operation around us. Also in the human and social
domain the principle of actuality holds. All possible emancipation
and individualism cannot abolish the simple fundamental relations
that ever anew put Community in the place of Society.
The family must ever remain a community.
The relation between parents and children and between husband and wife
lays the foundation of the strongest and most lasting feelings that
human life knows. Within the society that is determined
by these feelings every human being begins his life. Here
takes place a living together that is the constant home of sympathy and of
the feeling of humanity. From it there ever anew proceed
the forces that hold human life together. Here egoism is bent and
tempered; here brutality undergoes a quiet
change. No legal mechanics, however shrewdly devised, and
no enlightenment, however thorough, can replace the
% --- p. 154 ---
\opage{154}wellspring that here has its source. And this inwardizing,
mollifying, and cohesive power is not weakened because
the external order of life has become other than it was when one
knew no other forms of society than those that were founded
on blood and neighborhood. The family first comes into its
own precisely when it is no longer the only, or as good
as the only society. The relation of power then recedes;
those elements are separated out which in the old family too much
recalled a relation between ruler and ruled\footnote{The Russian dramatist \emph{Ostrovsky} has portrayed family life
as it was forty years ago in the towns of inner Russia.
One sees here to what degree even the family under the old
order of things was a system of compulsion. His ``Oeuvres dramatiques''
therefore have not only an exceedingly great poetic, but also a
very great cultural-historical and ethical interest.}.
The general process of liberation has also benefited the family;
its true nature has been able to unfold only
after the \emph{old}, patriarchally led community had
ceased. Thus not only do the same forces still work in human
life as millennia ago; but there is
even reason to believe that they work under more fortunate
conditions! ---

And even if individuals outside the family have come
to stand more in relations of contract than in a natural
relation of belonging together to one another, --- even if \emph{contractus} spreads at the
expense of \emph{status}, --- yet these altered
conditions again lead to new forms of community through \emph{free
association}. For not all associations have the loose
and selfish character of the joint-stock company. Where there is not merely a common
purpose, but where people work and strive, suffer and endure in
common, there will gradually develop a
feeling of community, a comradeship, which can gain a greater swing and
possess a greater strength than the purely egoistic
expectation of the attainment of definite external purposes can explain.
Even if the association is from the first entered into on the part of
the individuals for predominantly egoistic reasons, there will in the course of
living together and working together take place a shift of
motives, so that the original purpose
% --- p. 155 ---
\opage{155}recedes before the feeling of community. In this way the
individual becomes a sharer in the common life of the race; when he puts
his little weight in the scale, it is not in the thought of what he
himself will attain thereby, but it is the common fate
that he thereby will do his part to bend in the direction he
considers most beneficial for the whole to which he belongs.
The best and most familiar example of this is offered by
the history of the workers' trade unions. What we can learn from
this I have expressed in another place (Etik p. 244)\footnote{2nd ed. p. 336.} thus:
``The comradely feeling, sympathy with others, grows. The
individual sees that his conduct during and outside work, his
ability, diligence, self-control, and thrift become
of significance for far more than himself. There is formed, so to
speak, an ethical sphere around him, a great family, of which
he feels himself a member. He comes to know brotherhood,
which before was almost forgotten for freedom and equality. He learns to
subordinate his own interests to the common ones. And he gets
through the common experiences and the survey of trade
and factory conditions that determines the policy of the trade unions
a clearer conception of the workers' position relative to the other
classes of society, learns better to know both his rights and duties as
a member of the race. It is in short a whole
education from egoism to sympathy, from blind rawness to clearly
seeing strength, from struggle to peaceful negotiation, that here
takes place. And all this happens by the way of freedom. There is
no better answer to those who regard our time as a mere
time of dissolution than to refer them to the formation of society
and the ethical development that has here taken place.'' ---

Only with a few words can I here mention the
\emph{psychological} development which according to Tönnies runs parallel with and
at the same time stands in real connection with the social. That too
was indeed of a pathological nature: understanding and reflection
replace imagination and faith; in place of the immediate
warmth of feeling comes the cool clarity of thinking; the
spiritual life no longer moves according to natural impulses,
but according to shrewd calculations. But here too it holds that
% --- p. 156 ---
\opage{156}new natural and immediate states and motives can
form, after one has broken loose from the instinctual
and the traditional. What we at first only embrace with cognition
can at last again set feeling and fantasy in motion. The
more we live ourselves into the cognition won, the more
will it too be able to become a framework or a foundation for a new
immediate spiritual life. The modern world-view
has its grandeur and sublimity, which --- when the preparatory
inquiry and criticism, which for the time being lays claim to the best
powers, have done their work --- will call fantasy and feeling
into activity. When we consider how short a time scientific
cognition has worked, and to how narrow circles it is still
confined, it is with full right that we assume that we live
in a time of transition. Strife between knowledge and faith is characteristic
of such a time of transition. But there is no reason to
assume that this strife --- in its present form --- will last
forever. The harmony that the individual now laboriously
strives for in his view of life, when the break with tradition
has occurred, will in due time step out into common symbols, into thoughts
around which large circles can gather. We live in a ``critical''
period; but it will --- this is not an unfounded faith ---
be replaced by an ``organic'' period. ---

Will all these considerations now make any impression on
the social pessimist?

As mentioned at the beginning, behind the strife between
optimism and pessimism there lie two different fundamental moods. This
holds even here, where we stand before a clear and energetic
thinker, who commands a rich material of facts
and the analysis with great insight. Behind the grounds that
can be formulated there always hide grounds of mood
that cannot be formulated. As regards a single point
(a certain romantic predilection for the primitive, natural,
limited forms of life), it has been pointed to in what has gone before.

Moreover the conditions that are here in question are so
complicated that no one can survey or calculate the power of the various
cooperating and counteracting tendencies. Neither
sociology nor psychology will be able here to set up any
wholly certain deduction that could decide the strife.

% --- p. 157 ---
\opage{157}Our opponent will answer us: ``Even if all that
has here been put forward is correct, even if there are for example constantly
new formations in the direction of immediate social and spiritual
life, --- the opposite tendency, which goes in the direction of
mechanization, isolation, and intellectualization, works on an ever greater
scale; must it not in the end run the line out, so that life
dissolves?''

Well, yes, it is a matter of a trial. No one thinking seriously
can be blind to the danger. Life is a constant struggle, and
its result we cannot know beforehand.
Pessimism has the great merit of impressing upon us that we must not
over a progress on the surface forget and lose that which
is after all in the end the kernel and strength of life. In this Tönnies
is right, that if ``essential will'' and all ``Community'' were dissolved,
life would at once become a sham life, an empty form.
The explicit consciousness of what is at stake will be able to
sharpen the sense for the real values and the will to
preserve them. Back we cannot go, and stand still
we cannot either; we must go forward and will go forward, but one
of the most important conditions of real progress is
that we do not lose what has once been attained, and what is
a constant necessity for a healthy life.

\streg{}

% --- p. 158 ---
\opage{158}
\essay[Pagan Seekers of Truth]{``PAGAN SEEKERS OF TRUTH''.}

\begin{center}\emph{F. W. Farrar:} Hedenske Sandhedssøgere. (Seneca --- Epiktet --- Marcus
Aurelius). Translated by Louise Hoffmeyer. With a Preface by Headmaster
Dahlenborg. Copenhagen. 1891.\end{center}

\streg{}

\section*{I}

Through this book Danish readers will be able to make the
acquaintance of some of the most remarkable spirits of the close of antiquity.
It depicts the life of thought and feeling in men
who lived at the time when Christianity was in the process
of spreading in the Roman Empire, but whose whole view of life
was nevertheless wholly untouched by it, --- men who precisely
represent one of the principal forms in which the ancient view of life
appears at the close of antiquity. And since in their manner
of thinking and feeling there shows itself an inwardness, a purity, and
a seriousness which many do not expect to find in authors who
have developed entirely outside any contact with Christianity,
their writings have the significance of an experiment, by whose
help one can test how far the spiritual life could
advance in development under the conditions of antiquity. Whether one
regards Christianity as supernatural revelation or
as the involuntary expression of what has stirred in
the human spirit, it is nonetheless of great interest that it is possible to
make such a comparison.

The authors with whom the book deals belong
to the history of philosophy. But they are not purely theoretical thinkers.
Although they have come out of a definite philosophical school, the
Stoic, and maintain its fundamental thoughts, what is characteristic
of them is the personal character of their thinking. They present
precisely an interesting example of how one and the same
general direction of thought takes shape differently in
% --- p. 159 ---
\opage{159}different personalities. Their interest is a philosophy of life. They
are what Søren Kierkegaard called subjective thinkers.

Their difference from the Greek philosophers of the classical period
rests chiefly on this, that the immediate harmony between the
elements of life, and particularly in human nature, no longer
prevails. They do not have Socrates' assurance that clear
insight is one with the will for what one has known. There has
proved to be a greater distance between ideal and reality
than classical Greekdom assumed. These thinkers are in a
certain sense the expression of a process of dissolution, of the
fact that the classical harmony, which had set itself such
great monuments in art, thought, and political life, had to be
given up. But seen from the other side they mark a
progress, in that their whole conception has an inwardness and
a depth to which Greek thinking did not attain in its
classical period. Between thought and will the life of feeling thrusts itself
forward, a world of inner movements, which for the classical
conception fitted with ease into the frames that thought
could form and the will embrace in its acting, but which now
needed to win its own peculiar expression and demanded
its separate satisfaction. From the contemplation of the
outer existence and from participation in the life of the state, \emph{the
individual} now withdrew in order to immerse himself in an inner world
that was no less rich in events and tasks
than that outer one. Everything now came to depend, for him, on how this inner
world took shape, or rather, how he could shape
it. What did the glory or misfortune, the infinity or narrowness of the outer world
properly concern him, if it did not
reach into this \emph{his} inner world? --- Mood and
disposition now became the decisive thing, not theoretical thinking
or outward-going action. From this follows a quite
peculiar coloring over the writings of these authors, a
shade of mood that is unknown in the literature of classical antiquity.

But at the same time as the field of vision is thus narrowed,
inasmuch as the individual withdraws into himself,
it is widened, as it dawns on him that such an inner life
as he himself leads is possible in \emph{every} other
human being. With the isolation of the individual there goes hand in hand, at the
% --- p. 160 ---
\opage{160}end of antiquity, the development of a general feeling of human unity
and fellowship. The outer differences
between the slave and the free man, between the barbarian and the Hellene
become vanishing and meaningless in comparison with the kinship
in the inner life. Barriers are broken through which in
classical antiquity were assumed to be set by nature itself. It is a fact of
extraordinary historical and ethical importance
that already before the coming of Christianity we had
\emph{universal love of humanity} (caritas generis humani)\footnote{On the remarkable passage in \emph{Cicero} (de finibus V, 23, 65), where this
expression occurs for the first time, I shall make
some historical remarks at a later point. The word love of humanity or
friendliness toward humanity (philantropia) occurs, to be sure, early among the Greeks,
but receives a deeper meaning only among the Stoics, and it is more
indefinite than the expression found in Cicero: love of
the human race.} as
a clearly conscious feeling.

In another respect too the mind is widened. It
embraces not only everything human with its sympathy, but
also the whole of existence, the whole world (Kosmos), with an inward
feeling of admiration and kinship; it feels itself
flowed through by the same power that stirs in the great whole,
conceives the law of conscience as a single, individual form
of the great law of the world, and bows with willing
resignation to the limits that the order of the world once for all sets.
Stoic ethics has from its first origin (about 300
years B.C.) this \emph{religious} character, which already in one of
the first representatives of Stoicism, Cleanthes, found expression in
an inspired hymn.

Elements and germs of the development which Stoic
philosophy thus expresses were, to be sure, already to be found
in Socrates, in his fundamental conception of the ethical.
But the self-knowledge and self-care for which Socrates had
spoken appears among the Stoics, especially the
later Stoics, with whom the present book deals,
considerably deepened. This shows itself not least in this, that whereas
Socrates had indeed impressed the necessity of feeling his
ignorance and of drawing a sharp boundary between what
% --- p. 161 ---
\opage{161}one understood and what one did not understand, with which there also followed
the boundary between what one could do and what one could
not do, --- there is now impressed in a far more penetrating
manner the feeling of one's own imperfection as a condition for
progress. To take a single example: among the signs
that one is making progress, Epictetus mentions (in his ``Handbook''
48, 2) that one sees in oneself an enemy and a traitor. This
is said with an emphasis that is far sharper and indicates
far greater experience than that which found expression in
Socratic irony.

There is in all this so much that could recall
the Christian ideas that it is no wonder that it
has been thought necessary to assume Christian influence on the
later Stoics. It has been supposed, for example, that Seneca should have
stood in connection with the Apostle Paul. The
correspondence that goes under the names of these two men is, however, a
botch, and the whole supposition improbable and unnecessary.
In the title the translator has given the present work,
and which is happier than the title of the original, ``seekers after God,''
since the religious is not the only thing, nor even
the main thing, in these men, there is as it were an echo of
amazement that these men, although they stood outside
Christianity, nevertheless have reached as far as they have, or rather
that they were on this way at all: heathens
--- and \emph{yet} seeking after truth! This contrast does not in
reality exist. The zeal to find the truth
has been quite as great outside Christianity as inside
it. The older antiquity too bears witness to this. The men
here treated only continue in their own way the work of earlier
thinkers, drawing their peculiar practical
consequences from it. --- The title chosen has, however, the advantage that
it will awaken attention to a problem that will gain
ever-increasing importance: the independence of the striving after truth
of the limits of religious confessions.

The English book, which the translator with great
love and care has rendered in Danish, gives good help
and easy access to making the acquaintance of some of the
noblest spirits of antiquity. Although the author is a theologian, he
% --- p. 162 ---
\opage{162}takes up toward the old thinkers an attitude of great freedom from prejudice and sympathy.
He is mild and lenient in his judgment of
a man like Seneca, whose character was the object of
attacks already in antiquity, and whose life certainly did not
always accord with his teaching. He is full
of admiration for Epictetus' self-control, frugality, and
cheerful resignation, and for Marcus Aurelius' inward and
spirited self-contemplations. Somewhat naive and often
disturbing is his juxtaposition of (often unhappily chosen)
passages of Scripture with the sentences of the old philosophers, and his critical
and comparative concluding section is not very profound.
But the strength of the book lies in the interesting and detailed
biographies of the three men and in the extracts from their writings,
which are so abundant that one learns through them to know their
main thoughts. The book must therefore be recommended to anyone who
is interested in the development of the human life of the spirit.

\section*{II.}

On a single point in the present exposition I shall
dwell a moment, because it has in itself great interest,
and because the author here both finds himself in a historical
error and is at the same time led by a misunderstood Christian
interest. --- He speaks of Seneca, Epictetus, and Marcus Aurelius
as quite singular phenomena of their time, as
drops in ``the foaming sea of universal corruption.'' But
he has absolutely no material that can entitle him to
assume a particularly great moral corruption at the close of antiquity.
The justification cannot be sought in the declamations of Seneca and others, any more than
in the utterances of the Apostle Paul
at the beginning of the Epistle to the Romans. The age was in a certain sense
a period of dissolution. But though it did not present many
great figures, it nevertheless presented many pure and noble
characters, who under unfavorable and often painful conditions
remained true to their ideal of life, which was not low. And from the very
recognition that fell to the lot of the men admired also by the author
a counter-proof can be drawn against his
assumption. Thus says a scholar who, with regard to a
question like this, is among the most expert, \emph{Ludwig}
% --- p. 163 ---
\opage{163}\emph{Friedländer} (Darstellungen aus der Sittengeschichte Roms in
der Zeit von August bis zum Ausgang der Antonini. 5. Aufl.
III. p. 676): ``An age that by its own strength raised itself to
higher and purer moral views than the whole earlier
antiquity, --- that not only produced a Musonius [a Stoic contemporary with
Seneca], Epictetus, and Marcus Aurelius, but
in which these proclaimers of a mild, genuinely human
moral doctrine also found the most general admiration and
their teachings general diffusion, cannot have been
a time of the deepest moral decay, as it
has often been called. If there is no standard at all for
morality in a period, even though it be ever so well known, then this holds most of all
for these centuries [the first two after Christ], from which there
exist only isolated reports, partly confined to particular fields,
partly colored or one-sided.'' --- And the same
scholar's investigations prove that, as it stands
with the prejudice of a particularly great moral decay at
the close of antiquity, so it stands with the assumption
of a general decay of the religious life at this time.
Many testimonies prove an inward holding fast to the
old faith and the old religious customs among the people as a
whole. One can no more infer from certain literarily educated circles with regard to this period than one can for the 18th
century infer from the circles influenced by the French Encyclopedists
the religious condition of the whole people. There
showed itself, indeed, precisely at the time here in question, a
strong religious movement, which led partly to a revival of
the old cult, so that, e.g., the oracles began again
to speak after having long been mute, partly to the diffusion
of oriental religious cults in the western regions of the Roman Empire.
Christianity itself and its diffusion is, seen historically,
only a single, though the most important, side of this whole
phenomenon. The hard struggle that Christianity had to endure
also shows that the old faith was not gone. Even after
Christianity had become the state religion, several
centuries of cruel and fanatical persecution were still necessary
to displace the old religion --- from the surface. In
% --- p. 164 ---
\opage{164}manifold disguises the old ideas continued
to work and modified in many ways the Christian
teachings, in a more philosophical form among the learned, in a more
mythological form among the unlearned.

A misunderstanding, I believe, asserts itself here in
the assumption of a great moral and religious decay, in so far as one
means thereby to glorify Christianity. Since
according to the orthodox conception it is not merely to stand in
relation to a particular time, since it is also not the fruit of
particular historical conditions, --- and since according to the same
orthodox conception human sinfulness is at all times
equally bottomless and helpless, it would be precisely for
those who hold this conception questionable to assume a
special decay in the period when Christianity appeared.

\section*{III.}

The three men whose lives and thoughts are presented in the
present book belonged, as already remarked, to the Stoic
school. It was especially the practical view of life which this philosophy
grounded that won it so many adherents
in the Roman world. It asserted the individual's
personality over against the outer world and at the same
time impressed manly resignation toward the course of fate. Both
things were closely connected; for when the highest worth of life lies
in the inner, everything outer becomes in the end indifferent. There was
something here that appealed to the Roman national character, and we
find accordingly a whole series of eminent statesmen from the
time of the Republic among the adherents of the Stoic school. The
Stoic philosopher Panaetius was the friend of Scipio Africanus the Younger
and of Laelius, and Caesar's noblest opponent, Cato the
Younger, belonged decidedly to this direction. In the first times of the Empire
a whole host of men of Stoic education formed an
opposition to the insane tyranny on the throne, an
opposition that often led to death. The manliness
with which so many of these men went to meet death, and which
Tacitus has depicted in a series of famous scenes, was an
expression of the spirit of Stoic philosophy. But Stoicism also had
another influence than that of teaching how one can
% --- p. 165 ---
\opage{165}defy an unfavorable fate and die in a noble way. It soon
worked in no small degree upon the spirit of government and legislation.
A great number of humane measures from the earlier
imperial period are, no doubt rightly, derived from the influence of this school.
And of the three men with whom the present book deals,
two exercised a definite influence on the
use of the power of the state.

\emph{Lucius Annæus Seneca} (died A.D. 65) was Emperor Nero's
teacher and had --- together with the praetorian prefect Burrhus ---
at the beginning of his pupil's reign a great and good
influence on him. Some years later he was overthrown and
had to kill himself at the Emperor's command. Later there
came a reaction against the philosophers; they were repeatedly
(under Vespasian and under Domitian) driven out of Rome. But
better times came again, and the next century saw
even a philosopher on the imperial throne: \emph{Marcus Aurelius Antoninus}
(born A.D. 121, died 180). --- That it was not only in the
higher circles that Stoicism gained entry is seen from the fact that the
third of the great Stoic authors of the imperial period, \emph{Epictetus}
(who lived in the latter half of the first century), was
a slave, who won spiritual freedom by living himself into
this doctrine, of which he afterwards became the proclaimer.

No less great than the difference between the three men's
outer circumstances was the difference between their spirit and their language,
within the common Stoic frame. Seneca is the
most prolific author of the three. He has often more of the rhetorician than
of the philosopher, and a certain sentimentality cannot be acquitted
of him. But a rich experience in the spiritual domain, a real
solicitude as a pastor of souls for his friends, and a serious absorption in
ethical thinking make many of his writings interesting and
instructive, as they are also of great importance
for the knowledge of the Stoic doctrine, since we have no older Stoic writings.
Even if he has perhaps shown weakness of character in his
public conduct, and even if his luxurious surroundings, in the opinion of
some, did not quite agree with the strict doctrine
he proclaimed, there nevertheless runs through his life a
serious striving after the true, and the dignity that perhaps
sometimes failed him in the struggle for existence on the
% --- p. 166 ---
\opage{166}slippery floors of the court, he found in any case when it
came to being able to die.

In great contrast to the rich courtier Seneca, whose
luxurious surroundings had an analogy in his flowering and
rhetorical style, stands the sickly and poor Epictetus,
who restricted both his literary and physical
needs to the least possible. He himself wrote nothing,
but one of his disciples recorded his lectures and made
in turn from this work (of which one half is preserved) an extract,
``Epictetus' Handbook,'' the shortest and simplest presentation
we have of the Stoic doctrine. It is the ethics of the individual that
is here taught. Through endurance and abstinence
man wins himself. The surest way to spiritual freedom
is that of restricting one's needs. Be not bitter if
the dishes at the table of life pass you by, --- but show yourself even
so superior that you let a dish pass by, even if it
is offered to you! And yet it is no dark asceticism that is here preached.
The glory of life and existence Epictetus praises with
inward enthusiasm; for precisely to him who makes himself spiritually free in
the great school of self-control it becomes clear that there is nothing
really evil in the world. And likewise the
most inward love of others is taught, even toward enemies and toward those who
maltreat us.

The third author, Marcus Aurelius, combines Epictetus'
clarity and simplicity with Seneca's brilliant intellect and
rhetorical power, just as he was able to preserve
self-control and gentleness of mind on the imperial throne and to
hold fast enthusiasm for life despite the opposition and the
disappointments he experienced, and despite all the small and the petty
that his honesty and sharp eye did not overlook. His
self-contemplations are in a certain respect the most remarkable book we
have from antiquity. There is something modern in a man
thus quietly for himself giving an account of his thoughts and
moods about life. And there is a depth in the mood,
a peculiar tone over his utterances, that we do not find
in any other writing of antiquity. Stoicism appears at the same time
with greater mildness in him than yet in Seneca and Epictetus,
who both have something of the preacher's severity. Marcus
% --- p. 167 ---
\opage{167}Aurelius had only himself to deal with, and when in the
palace or in the camp he withdrew to his closet,
it was --- despite the sharpness with which he looked at things
in their material nakedness --- not in order to give vent to his melancholy
or his indignation, but in order to strengthen himself in the feeling of
human fellowship and its tasks. His Stoic
cosmopolitanism does not hinder him from energetically feeling himself a
Roman: ``as Antoninus I have Rome, as a human being the whole
world for my native city; only that which benefits these two places
can be good for me.''

\section*{IV.}

It may perhaps be of interest to investigate
how the line of thought and view of life that we meet in these
men, and which in many features diverge so greatly from the
classical Greek, has developed. Only brief indications
can be given here, since, as already remarked, the whole literature of the older
Stoic philosophy has perished, so that we are
restricted to using second-hand sources.

Stoic philosophy was founded by \emph{Zeno} of Cyprus,
who a little after the year 300 B.C. taught in Athens in a colonnade
(stoa), after which the school took its name. It was a time that
was not a little different from that in which Socrates and Plato
were active in Athens. Greece had lost its freedom and
had come under Macedonian rule, after the
individual Greek states had been weakened and the sense for
public and national interests had been lost. At the same time
Alexander's conquests had opened a larger horizon and
paved the way for a new and more comprehensive view of
human relations. The young Macedonian commander had
broken down the barriers between Hellenes and barbarians, not only
by uniting them under one scepter, but also by
founding Greek cities round about in the Orient and bringing about
interaction between Greek and oriental customs and
ideas. On the Greek side people now learned in practice that
there was not the difference in principle between Hellenes and
barbarians which even Alexander's great teacher Aristotle had
maintained. Characteristic is the story (found in a
% --- p. 168 ---
\opage{168}work which, perhaps wrongly, is attributed to Plutarch) that
Aristotle is said to have advised Alexander to treat the Hellenes
as Hellenes, the barbarians as barbarians, but that the young
conqueror, whose ambitious political plans coincided with
what humanity demanded, did not obey his teacher in this. Through
living together and the common fate a common
feeling was now developed; the idea of a universal humanity, which
embraced all nations despite their differences, now
came into being. The conquered nations found comfort in the thought that they
had not succumbed in order to be absorbed into a single other
people, but in order to stand as members of the whole of humanity\footnote{On one single point, which goes far back historically,
classical Greek ethics too led beyond the local and national
limitation of the feeling of humanity, namely in
the relation of hospitality and in general in the relation to strangers, who from
ancient times were regarded as standing under the special protection of the gods.
The duties toward the helpless stranger stood as some of the
highest duties. To show the wandering traveler the way, to give others
fire and water when they lack it, and not to advise them anything but
what one holds beneficial oneself, were duties that were impressed by
a special, venerable tradition. Here, then, the barrier did not hold
that otherwise divided human beings into different ethical worlds. The
stranger and the suppliant stood as a human being, neither as
Hellene nor as barbarian. There was here a germ that could be developed further,
and also was developed.}.
--- The same circle of experiences came to be
made in still larger circles and under more lasting and firmer
conditions, when the whole civilized world was absorbed into the Roman Empire.

Stoic philosophy, which emphasizes the independence of the inner, of reason,
of the will, from the outer, was particularly called to
develop the idea of a universal humanity prepared by way of practice.
National differences had to
stand as something outer and accidental, when one found man's
proper essence not in outer forms and conditions of life, but
in what moves in mind and thought. There shows itself here an
interesting reciprocal relation between historical experience and
philosophical thinking. The discovery of the common human nature
was conditioned and facilitated by the actual breaking down of the
national barriers, and philosophical thinking expressed,
% --- p. 169 ---
\opage{169}so far, only the result of historical experience. But on
the other hand this philosophical thinking had already been
founded earlier. Zeno, the founder of Stoicism, was influenced
by the great movement proceeding from Socrates a century before. According to a story Zeno
is said to have become interested in philosophy by reading Xenophon's work on Socrates.
It was the Socratic thought of self-knowledge as the foundation
of all knowledge and conduct of life that was developed further among the Stoics in
consequences that had not been drawn before. And the
historical experience made after Alexander's conquests
first received its higher significance in that the idea of a
universal humanity was developed. The Stoics deepened this
thought through the whole world-view according to which
existence constitutes a unity, in that the world-soul lives and moves
in the individual beings. The whole world, they taught, is one great
state, common to gods and human beings, animated by one and the same
infinite power and subject to one and the same great law, which
appears to each individual being as its own inner law.
They sought thereby to connect the independence of the individual
personality with devotion to the universal: in that the
individual follows his nature, he follows the universal nature,
in which he is only a single member.

Among \emph{the first} Stoics the personal self-assertion
and independence are emphasized more than devotion to the universal,
and, what is connected with this, within the individual
personality they emphasize reason and the will as the
most important. Their ethics thereby receives a strict and negative character.
The Sage, their ideal, is to be raised not only above all
outer vicissitudes, but also above the movements in his own
inner life, above pleasure and pain, love and anger, above all
affects, which were regarded as something evil, since they abolished
independence, the absolute resting in oneself, and led
to dependence on something that lay outside the will.
The chief virtue was justice, which was conceived as the opposite of
love. --- Later Stoics (\emph{the middle Stoic school}
in the 2nd and 1st centuries B.C.) altered this conception and
deepened the feeling of humanity, indeed first made
the consciousness of humanity into a feeling of humanity. Unfortunately
% --- p. 170 ---
\opage{170}we know very little of this development, and only at second
or third hand from later authors. The Stoic thinker
whose influence has here been especially important is
\emph{Panaetius}, already mentioned as the friend of Africanus the Younger.
On the basis of the natural impulse to social life, which the
predecessors had already assumed, and which Aristotle had expressed
in the well-known proposition about man as a being
fitted for society (zoon politikon), Panaetius developed the ethical
significance of love of humanity as a force that holds together.
Between justice and love he placed
equity, which does not keep to formal right, but
takes all individual considerations into account. And whereas
the older Stoics properly assumed true
friendship only between the ideal ``Sages,'' in whom reason had
reached full development, Panaetius found in
love of humanity a living bond between all human beings, even if they
have not reached the absolute ideal. Of special interest is
it that he has not merely kept to psychological descriptions
and ethical postulates, but has given the fundamental features of a
developmental history of sympathy, which recalls the modern
theory of evolution. In his work ``on the highest good and evil'' (de finibus)
Cicero says (V, 23, 65): ``Of all that is good, nothing is so
admirable and comprehensive as the mutual
connection of human beings, as their community in what is useful to them, and
as the love of the human race, a love which
has its source in procreation, in that the offspring is loved by
the parents and the whole house is held together by the relation
between spouses and between parents and children, but which
then gradually spreads outward, first through nearer
and more distant relations of kinship, then through bonds of friendship,
through relations of neighborhood and through association in civic and
public relations, until it finally embraces the whole human race.''
There is indicated in these words a successive development of the
feeling of humanity from narrower to wider circles, conditioned
by the way in which human beings, by inner impulse and by
the force of circumstances, come to live, work, and suffer together,
a development that agrees entirely with the way in
which modern sociology conceives the matter. The first
% --- p. 171 ---
\opage{171}part of this process of development Aristotle had already depicted
in a clear and interesting way, and he stated the fundamental law
that makes itself felt in the whole of this course of development: that
the involuntary practice in an activity or a relation of life
produces a corresponding feeling of satisfaction therein.
I have therefore (in my ``Etik'') called this important law ``the
Aristotelian principle.'' The last part of the process of development
indicated in the passage in Cicero, the transition from the
national belonging-together to the unity of the whole human race, is due to
conditions and ideas that first appeared after the time of Aristotle and
Alexander, first in the eastern part of the world then known,
later, with the extension of the Roman Empire, in still larger
circles. And it is probable that the exposition in Cicero
rests on thoughts that first saw the light in the Stoic
school, probably in Panaetius\footnote{The Greek philosopher whom Cicero follows in the section from which the
cited passage is taken is \emph{Antiochus} of Ascalon, whom he himself had
heard in Athens. This man belonged to the Academic school, which was
founded by Plato, but he sought, often in a very external manner,
to combine the doctrines of different schools and took over much
both from the Aristotelian and from the Stoic school. The question might then
arise, which of these two sources we are most nearly to
assume as the origin of the cited line of thought. Madvig seems (in his
edition of de finibus, note on III, 19, 62) most nearly to assume a
Stoic origin. Later Bernays has sought to show that it is from
the Aristotelian school that the idea of universal
love of humanity derives. But the most recent investigations (O. Apelt: Beiträge
zur Geschichte der griechischen Philosophie 1891 p. 354, --- A. Schmekel:
Die Philosophie der mittleren Stoa. 1892 p. 398 cf. p. 220) point
to this, that Cicero's authority Antiochus is here most nearly
dependent on Panaetius and on the form he gave to the oldest Stoic
doctrine.}. In this philosopher, and in still
higher degree in his successor \emph{Posidonius}, the
life of feeling and its significance were, on the whole, emphasized far more than among
the first Stoics. People sought to mitigate the sharp
opposition that these had established between reason and the active
will on the one side, feeling and affect on the other
side. This was naturally connected with their seeking to
abolish the absolute opposition that the founders of the school had
% --- p. 172 ---
\opage{172}established between the ideal ``Sages'' and all other human beings.
There has on the whole been something in this movement within
the Stoic school in the last centuries B.C. which is
akin to the influence that Rousseau's assertion of the
power and right of feeling against the other sides of the life of the mind exercised in
the last century both on the development of psychology and ethics
and on literature as a whole.

The change and development that took place in the
middle Stoic school, and by which an increasing
weight was laid on the universal, the generally human as opposed
to the merely individual and particular, on feeling as
opposed to cognition and on love as opposed to
justice, belongs to the most important of what takes place
in history in the last centuries B.C. Both on account
of the successive way in which this development has taken place
and on account of its taking place in the minds
and thoughts of human beings, it does not push itself forward into the foreground of history,
where the noisy political events take place.
But for that reason its significance may be quite as great and lasting.

\emph{The last Stoics}, who are represented by the three Romans
Seneca, Epictetus, and Marcus Aurelius, now do not stand for us
as an inexplicable historical phenomenon. It is easy to
understand how they, as the circumstances of the time developed
after the fall of the Roman Republic, came to
lay a still greater weight on what the
middle Stoic school (Panaetius and Posidonius) had already
emphasized. They could not rightly be called Stoics
if they did not energetically maintain the inner personality's
independence of the course of fate and of outer goods and
evils. In that respect there is no lack in them of
sharp and harsh utterances. But peculiar to them is
the great inwardness, the deep mood, with which they utter
the Stoic thoughts. A single example will show the character
of this spiritual direction, which comes to light especially in a
beautiful and often moving way in that one of these men who
stands highest both as a personality and as an author, Marcus
Aurelius. He says in one place in his diary that human beings are
fitted for common activity, for cooperation, like the body's
% --- p. 173 ---
\opage{173}limbs, and continues: ``This thought will appear more clearly to
you when you frequently say to yourself; I am a \emph{member} (melos) of the
community of spiritual beings, than when you (by a small
change in the letters) merely call yourself a \emph{part} (meros). If
you are merely a part, not a member, \emph{then you do not yet love
human beings from the heart;} doing good to others does not yet delight you
so that it wholly captivates you, but you do it merely
from duty, not yet with the feeling that thereby you do yourself
good.'' (Marcus Aurelius' Notes VII, 13). The same
inwardness appears, as already mentioned, in
the demand to love one's enemies and in the religious conception
of the life of the world. It is the inner richness and fullness that
raises these men above human beings' assaults and lets them
overlook the dark sides of existence. The power they felt in
themselves to overcome the outer, the changing, and the painful did
not stand for them in conflict with the power that stirred in
the life of the world, but was a branch of it; therefore harmony became
their last thought, and they remained Greeks, despite
the sharp opposition between the inner and the outer, which
was so essential a part of their doctrine. Not with defiance,
but with cheerful readiness the imperial philosopher awaits
what fate will bring him, and his cheerfulness breaks
at times into a hymn, as in the famous passage where
it says: ``Everything is harmonious for me that is harmonious
for you, O World. Nothing comes too early or too late
for me, if it comes in due season for you. All that your
seasons bring is fruit for me, O Nature. From you comes
everything, in you is everything, to you everything returns. When the poet
says: dear city of Cecrops\footnote{Cecrops was, according to the legend, the founder of Athens.}, should I not say: dear
city of Zeus!'' (Notes IV, 23). --- The inward
optimism did not rest on any contemplation of outer nature.
The older Stoics had still had an interest in the study of nature
and in historical science, and they had only by fragile
and sophistical grounds (which to some extent recur in the theological
reasonings of more recent times) been able to maintain their
optimistic conception against experience. From all such 
% --- p. 174 ---
\opage{174}outward-going interest the latest Stoics have withdrawn;
therefore they escape great and provoking problems. They
do not solve the task of how the inwardness, cheerfulness,
and power of the inner life can be maintained together with the experience of
suffering and struggle in the world. But this imperfection
is offset by the firm conviction with which they have
pointed to the inner life as our nearest and most essential
world and maintained it as a domain of its own. ``Look
inward,'' says Marcus Aurelius (VII, 59); ``in there is the
fountain of the good, which can always gush forth, if only one digs!''

Men like Epictetus and Marcus Aurelius stand as beautiful
testimonies to the capacity that strong and clear spirits have
to remain true to themselves and preserve their freedom amid
currents of the time that led many to throw themselves into the arms of
mysticism and superstition. The old and
new religions coming from the Orient lured many to themselves. Many felt weary
--- some of dissipations, others of art and literature, others
again of politics. And when one is weary, one bows willingly the
head which one perhaps formerly carried somewhat more proudly than
the inner strength really warranted. % print: til (misprint)
In Epictetus and in Marcus
Aurelius we see that there were men both in the lowest and in the
highest circles who could not and would not give up
themselves, and who also had no need to. Their manly
view of life has not only in its content much that, seen from
the standpoint of our time, is still of great interest and of lasting
importance; but it stands, especially by its relation to the
currents of the time, as a remarkable and instructive phenomenon.

\section*{V.}

It would be a difficult task to point out wherein
properly the life of feeling in men like those here discussed was
different from what Christianity brought with it. Of course
the task becomes easy when one simply makes a new kind of feeling
for every rubric in dogmatics. In that way it will
be easy to show what richness came with Christianity
into the life of feeling; this was in this way, to be sure,
enriched with several new kinds. Far more difficult
is the task for him who is clear that, however important and
% --- p. 175 ---
\opage{175}decisive it may be that new ideas emerge
about the aim and significance of life, the life of feeling nevertheless has
its own laws, its natural fundamental relations, which cannot
be changed or extended indefinitely. It is therefore far
easier to point out what new ideas Christianity
introduced than what change in the life of feeling it brought about.
And the task becomes all the more difficult when one
notes what inwardness and what scope
the life of feeling could develop to in the ``natural''
human being. Yet I am convinced that it was an important, in its
kind unique change, which in its continued existence shows itself to
be independent of the dogmatic ideas that originally
called it forth. I shall try to set forth how
I believe that this change can be described, without my daring
to maintain that I thereby state the only way in which
this question can be answered. A comparison between
Stoicism in its last form of development and Christianity
will lead to what I regard as the decisive point.

The chief fault of Stoicism was the excessive regulation, the
strong emphasis on the necessity of voluntary intervention
in the course of the life of the mind. At the bottom of this lay a distrust of
the involuntary movements in the mind. The Stoic declared
all emotions to be evil. The Sage is raised above the
surge of moods and acts from clear and strict
consciousness. To follow the mood, to let oneself be led by
emotion led, in the opinion of the Stoics, inevitably to
a fall. And what mattered, after all, was this, to stand as
one's own master, oneself to sit at the helm on one's voyage over the restless
sea of life. As so often, here too the fault is connected
with the very greatness and excellence, with that which makes the
Stoic characters unique and gives them the stamp of sublimity.
But the point is precisely this, that the Stoics believed they could
reach this ideal position only by constantly inhibiting the
involuntary urge to devotion. A main principle of Stoic
ethics was and remained that one is to inhibit
the emotional effects of ideas. Things awaken ideas
in us; that cannot be avoided, but we can avoid letting
these ideas gain power over us by constantly saying: it
% --- p. 176 ---
\opage{176}is only ideas (or as the Stoics often said: dogmas)!
If we prevent the ideas from carrying us away, then things have
no power over us. This relation to ideas is in
reality the only thing that stands in our power, and therefore
the only thing decisive.

Among the last Stoics this characteristic
appears sharpest in Epictetus. In Seneca and especially in Marcus
Aurelius it is to a high degree mitigated, without however being
given up. Marcus Aurelius' Notes show how the
involuntary moods rule the mind; there is an inwardness
in devotion which is lacking in the other authors. The
task of depicting the difference between the last Stoicism and
the life of feeling of Christianity would therefore become considerably
more difficult if we were to keep merely to him. In order to bring out
the difference clearly I therefore keep to Epictetus\footnote{It must, however, be well noted that Epictetus' utterances in the ``Handbook''
appear with a sharpness and strictness that is for the most part due to
the concise form. The fuller exposition, of which
the ``Handbook'' is an extract, shows us Epictetus' doctrine in a milder,
more human form. This observation has been made by \emph{Martha:} Les
moralistes sous l’empire romain. 2. ed. Paris 1866. p. 162 f.}.

Since everything is indifferent that does not stand in our own
power, and since the most important of all, our tasks and duties,
belong to what stands in our own power, the Stoic is
not a strongly interested spectator of or participant
in what goes on in the world around him. He knows, indeed,
that it is always possible for him to take what comes in the
right way; the future therefore awakens neither his suspense
nor his curiosity; \emph{what} will happen, he does indeed
not know, but \emph{of what kind} it will be, he knows, namely, that
if it is something that does not stand in his power, it is neither
good nor evil, and if it is something that is subject to his power,
then what is essential depends on his own will, and
what it ought to do, that he knows beforehand. (Epictetus'
Handbook, ch. 32). What Epictetus says about how
one is to behave at a play can be applied to
the Stoic's relation to the great play of life. ``It is,'' says
he, ``for the most part not necessary to go to the theater. But
% --- p. 177 ---
\opage{177}if it is ever opportune to go there, then you must
not be occupied with anything other than yourself, that is to say,
you are merely to wish that what really happens happen, and only
that he win who really wins; for you will then meet no
hindrance. But keep wholly away from shouting, jeering, and
strong agitation, and when you have gone away from there, you are not to
talk much about what has taken place, except in so far as
it can serve for your own strengthening; for from such conduct
it will appear that you had felt wonder at what you saw.''
(Handbook, ch. 33, § 10). Also as concerns real life,
the Stoic wants only that what happens happen, and
in the way it happens, and on this rests his inner welfare
(Handbook, ch. 8). It is grounded on his own
conception. And the same view finds application to
others: of them too he knows that a real evil
cannot befall them. ``When you see one weeping and in
sorrow, then take care that the idea does not carry you away, so
that you think he is struck by misfortune; but make at once a
distinction in your quiet mind, and have the thought ready to hand that
it is not what has happened to him that torments him (another would
perhaps not be tormented by it), but his conception of it that
torments him. In words you are not to hesitate to show him
sympathy, and if need be, sigh with him; \emph{but take care
that you do not also come to sigh in your inner self!}''
(Handbook, ch. 16).

What is here warned against is the full devotion to
and captivation by what really happens. In order that the mind
shall not be set in such strong movement that it cannot
find its equilibrium again, one keeps it away from the living
relation to reality and makes this an indifferent
play. The way the Stoic goes in order to preserve his freedom is
to restrict his spiritual needs, his
need of feeling; so much the surer is he of getting it satisfied. He
makes the fraction in life's sum grow by
diminishing the denominator, not by enlarging the numerator. Here
the negative and passive in Stoic ethics appears
clearly. A certain fear of life makes itself felt. The more
% --- p. 178 ---
\opage{178}this element came to rule, the more the limitation of Stoicism had to show itself.

Psychologically viewed, Christianity in contrast to this is
characterized by boldly setting the life of feeling in the
strongest movement, letting it swing out into the most extreme
forms. For the highest pleasure and the deepest pain
Christianity has room. Its dogmas express the
conception that only he knows life who has lived through
the extremes. The need which human feeling
has to get complete discharge, to give itself vent
in reckless devotion, is satisfied in the Christian cult.
Jesus himself was sorrowful unto death, and he foretells his
disciples that they shall grieve and be sorrowful. And on the other
side a glory and a joy is announced which no sense
can grasp. Here is no anxiety before feeling, no
fear of sighing in one's inmost, any more than of letting
oneself be carried away by wonder. Therefore Christianity did not lead
to separating oneself from existence, but called on people to
live and suffer in its course. All thoughts were directed toward
one great goal, which was soon to reach its realization, and
on whose relation to it the individual's highest woe and weal
depended. The individual was thus really drawn in
among the greatest world events, into the great historical
connection, whose fundamental features revelation had unveiled.
As an indifferent spectator, who had enough in himself, the
individual could not stand here. In the expectation of the great
and near conclusion of the life of the world he was swept into so
violent a surge of hope and fear as the Stoic
Sage, to whom existence could bring nothing, and from whom
it could take nothing, did not know and in any case shuddered
back from. Thereby it won a power over minds which
Stoicism was excluded from, apart from a small circle of
energetic and inward characters. Christianity taught how to
win oneself by forgetting oneself, whereas it was
the Stoic's misfortune that he was always so busy
regulating the involuntary course of the life of the mind that he could not
attain true self-forgetfulness.

But the conditions in the domain of the spiritual
% --- p. 179 ---
\opage{179}life are so many-sided that we have scarcely brought forward and emphasized
a decisive opposition between Stoicism and Christianity
before we are forced to become aware of the likeness that
exists between them. \emph{To Stoicism's sharp dualistic difference
between the inner and the outer, a difference that finally
barred living and active participation in life, there corresponds
Christianity's still more pronounced dualism between the
beyond and ``this world''}. That kingdom, which was expected
shortly, was to come by supernatural way and abolish
the natural life and all its relations. From this there followed for
consistent Christian ethics a negative and passive relation to
natural life and historical development. Only by
reinterpretation has a compromise later been reached, just as there
also were not lacking analogous compromises within the
Stoic school. If Christianity is to mark a lasting
progress in relation to Stoicism, it must depend on whether
the living, passionate expectation, the full devotion
to the oscillations of the life of feeling and the confidence that one
thereby need not lose oneself, but can precisely learn to
know life thereby, --- whether all this can be preserved and
tied to the course of development of real human life
through history, even if the dogmatic ideas that
originally called forth these mental states fall away.
And as to the possibility of this the development of the spiritual life
in modern times seems to bear witness. Even where it stands remote from Christianity
in its dogmatic form, it has nevertheless taken up from
it psychological and ethical elements of the greatest importance.

The well-known saying that the history of the world is
the judgment of the world confirms itself especially in the
fate of the great spiritual currents. They appear from the first with definite
peculiarities and imperfections, which are due to the conditions under
which they came into being. In order to pass over into acquiring lasting
historical significance, they must first pass through a purgatory,
in which what conflicts with what the struggle with
\emph{real} existence demands is separated out. Both Christian and
Stoic ethics have needed such a purgatory, and there
are valuable elements in both which are able to pass through
this test.

\streg{}

% --- p. 180 ---
\opage{180}
\essay[Philosophy as Art]{PHILOSOPHY AS ART.}

\section*{I.}

When one studies law, we call him a jurist, --- when
he studies medicine, a physician, and so on. Quite
in the same way one cannot call a person who studies
philosophy a philosopher. One would do so only if the word philosophy
historically designated a purely theoretical and scientific
endeavor. But the history of philosophy shows that the thinking
treatment of certain problems is indeed a principal determination
of the concept of philosophy, but that there attaches to it also
another, more practical meaning: a \emph{life} in and for thoughts,
or a striving to carry thought through into life, --- in
larger or smaller circles of life, in any case in one's
own life. If this side of the concept of philosophy is taken, no one
without presumption could call himself a philosopher. One
could, with Kant, who returns repeatedly to this point,
distinguish between philosophy as knowledge and philosophy as wisdom.
Kant is careful to specify that there is a
sense of the word in which no one dares to call himself a philosopher.
In the practical sense of the word philosophy lays claim
to other spiritual faculties than the purely intellectual.

Empirical and critical philosophy distinguish more sharply
between the two sides of the concept of philosophy than the
speculative, which in the romantic manner lets the different domains
melt together and wants to comprehend them all in one idea,
without acknowledging the law of the division of labor.
The carrying through of the empirical and critical method has led to
a distinction between subjects that earlier did not come to lead
% --- p. 181 ---
\opage{181}their own separate existence. Philosophical thinking is in our
day no unified, single stream, but has divided into many
streams: psychology, logic, ethics, etc. It is a division
that has proved fruitful; but the immediate enthusiasm
that could easily be awakened when one believed oneself to
comprehend in a single idea the whole content of thought and existence, all
knowledge and all values, like a pearl that was worth
buying with all one's goods, --- such a concentrated
movement of thought the philosophical research of our day cannot offer.
It has given up finding the philosophers' stone or the key that
opens everywhere.

Where the division of labor gains admission, it brings about a
crisis. So it happens in the organism, where the necessity
of concentration in the struggle for existence sets a
limit to differentiation, --- in society, in that the social
problem is essentially connected with the division of labor, ---
and so also in science: for when all thinking
distributes itself over the separate, mutually different
domains, how then is spiritual concentration possible,
the comprehensive conception and the unified activity which life
demands, if it is not to be conducted blindly? --- In my
view this problem presents itself not only to the humane
or naturalistic view of life, but also to the religious.
The religious problem in our day is closely connected
with the fact that religious faith has now become a specialty \emph{alongside
of} science, art, and practical striving, whereas it
originally stood as the epitome of everything that could bring
spiritual and personal nourishment\footnote{Cf. my Etik p. 305---307. (2nd ed. p. 419 --421).}.

The disquieting effects of dispersion and
specialization seem to be in the process of calling forth a reaction.
Whoever follows with attention what is stirring
in the age will be able to sense that a reaction in a mystical
direction is in the air. One may call this movement
mystical: for mystical is not merely obscurity, but
concentration of the intellectual life, absorption into a unity. One needs
condensed nourishment for thought; one is weary of analysis and
% --- p. 182 ---
\opage{182}specialization and wants intuition, a gathering of knowledge into a
guiding idea that can fill the personality and not merely
satisfy a single side of it. One demands that science
shall become art, or that art shall be esteemed the highest
science.

There are many features that the close of the nineteenth
century has in common with the close of the eighteenth century.
Then too the need for concentration and mysticism stirred after
the great deal of analysis, and the consequence was reaction all along the line.
The caricature of this need expressed itself then, just as now, in
the tendency toward spiritism, occultism, and other
secretive movements. The rise of such spiritual currents
has in and of itself something obscure about it. In the year 1790
Kant was asked by a friend about the cause of the prevailing
fanaticism. Kant answered that it was as difficult a
question for the psychologists to answer as it was for the physicians
to state the causes of the prevailing influenza. He
thinks that one should not take the matter too seriously and not
use heroic cures against it, and seems to conceive of a
spiritual remedy for it in analogy with cold water, which
the physicians used for influenza\footnote{Kant's letter is found in \emph{Borovski}: Darstellung des Lebens und
Characters Immanuel Kants. Kønigsberg 1804, p. 227.}. Characteristic of the
limited experience that people in Kant's time still had of
spiritual currents, their nature and laws, is the fact that Kant
did not think of whether there might not be, in the whole character
of the spiritual life as it had formed itself in the latter half of
the century, something that naturally had to
bring about a reaction. In the course of the century that
has since passed we have had the opportunity to see and study the
movements in the spiritual domain, and are familiar with the
thought that a strong current is always in one way or another
motivated by the results or shortcomings of the preceding period.
We have seen so many waves succeed one another that
we think we can infer from the preceding ones roughly what
the following ones will be like. And we can then arm ourselves in
advance to receive what is to come. It
% --- p. 183 ---
% print p. 183: personsonligt (for personligt) -- misprint; the sense is translated
\opage{183}is by no means certain, after all, that the cold-water cure that Kant recommends is
sufficient.

In various authors of the most recent years the
need for intuition, for thoughts in which one can live
personally and concentratedly, for a conclusive view of life, finds vent
in a peculiar, partly very subjective and baroque manner.
Among German authors Lagarde, Nietzsche, and the anonymous
author of ``Rembrandt als Erzieher'' are to be mentioned here. In
France one has in Renan a spirited representative of this
tendency. In a recently published, talented Swedish work by Hans
Larsson (``Intuition'') the same tendency shows itself. Whatever
may be objected against these authors, each of them expresses
in his own way a spiritual necessity that does not
let itself be pushed back. We cannot live on abstractions
and specialties, and the question cannot be put aside
what significance thinking and research then have for our
personal conduct of life.

Although I believe, and have also occasionally already
touched upon, that the problem arises just as well on the positively religious
standpoint in our day as on the philosophical or humane,
I shall nevertheless here keep especially to this latter standpoint.
The question is how philosophy can become art, and in
what sense it can become so; \emph{that} it must become so is
given. And as I began by saying, it lies historically in the
concept of philosophy not merely to be knowledge, but also art
or wisdom. The question is how this can happen under our
conditions.

Most of the authors I mentioned before seem not
to have seen the whole difficulty of the problem. They cut the
knot instead of untying it; they leap over into mysticism
and intuition, do violence to thinking, forget, break off,
or conclude it in an arbitrary manner, indeed perhaps even look
down with romantic contempt on the men of sober thinking and
of special investigations. What interests me in
these authors is therefore more the need that comes
forth in them than the manner in which they themselves seek to have this
need satisfied.

% --- p. 184 ---
\opage{184}I shall try to show in brief outline for what reason
philosophy must become art, and in what way this
art is to be practiced.

\section*{II.}

As our knowledge is now once and for all situated, every
conclusive world-view can only be a hypothesis.
Every comprehensive conception, every interpretation of existence
that we attempt, must establish its validity by being compared
with new, more far-reaching experiences. But since such
experiences present themselves as long as life lasts in us and outside us, the
confirmation by experience can never be concluded in the way that is possible
with more limited tasks. It is always on a
fragmentary basis that we build an interpretation of existence; we are
placed roughly as the philologist who is to decipher a fragmentary
manuscript, or as the historian who from scattered and
detached reports is to construct an event or an
action and its motives. Philosophical research ends in
the art of conjecture, just like textual criticism and
historical research. The characterization of existence that every
world-view aims at cannot then get beyond the
hypothetical. The limit results, the ultimate consequences that a thinker
draws, will always retain a problematic character.

Now not only does art belong to the working out of such an
ultimate hypothesis, but a peculiar
art belongs also to working with it in practice. There is required an
elasticity that constantly goes back from the completed construction
to analysis, and is able to make this swing
back and forth without losing energy. The art is to
unite the enthusiasm of faith with the indefatigability of discussion.
The force of the work of thought must not be weakened by the consciousness
that the result is always problematic. When the given
reality is in conflict with the ideas that we with all our diligence
have formed for ourselves, is it then the fault of the ideas, or is it the
fault of the limited experience? Here the art is this, always
to be ready for a new test. It is easier to draw a
% --- p. 185 ---
\opage{185}line through the confirmation by experience and once and for all to sink into
prophesying and romantic intuition; but that leads to the death
of thought --- the great danger with which all spiritual reaction threatens,
and which it is above all a matter of avoiding. But it
is avoided only by practice in the art of living on ideas.
That is more than mere knowledge; for it makes demands on the energy of the will
and the liveliness of feeling. It is something that cannot
be done once and for all. --- Of that which one needs only
to know in order to be able to practice it, an old German
saying says: ``Das ist keine Kunst, das ist nur eine
Wissenschaft.'' The transition from philosophy as knowledge to philosophy
as art demands a mustering of personal force and
virtuosity which pure knowledge does not require. The point is to
preserve the connection between our thinking and our
practical striving and doings in the manifoldness of life, to make the
transition in such a way that no inner discord comes into our
being.

Briefly I shall point out that an entirely similar
task is set for the believer. Even if the content of his faith is for
him given once and for all and stands unalterably fast,
as proclaimed by an absolute authority, life nevertheless
changes constantly, and the task then becomes that of preserving
the connection between his faith and his practical participation
in the course of life. As the results of thinking there, so
the postulates of faith here are to be held fast without changing character
under the changing conditions. It would be an interesting
task to investigate how this task has been solved from
the side of the Church at the different times. The Church itself is more
content with the way in which it has solved this
task than its critics from highly different
standpoints (as different as Feuerbach's and Kierkegaard's) are.
But here I wished merely to indicate that also on the standpoint
of faith an artistic supplement is necessary, although
it naturally acquires a different character from the philosophical
art, which does not have authority as its presupposition.

A life in ideas need not presuppose that the validity
of these ideas stands dogmatically fast. It is a beautiful
% --- p. 186 ---
\opage{186}saying of Lagarde: ``The true idealists can wait''\footnote{Deutsche Schriften. Gøttingen 1886 p. 481.}. It
is only impatience and laziness that cause dogmatism.
It still goes with many people with ideas as it once
went with everyone with the stars: they were afraid that they
would fall down if they were not securely embedded in fixed
spheres.

There is required not only spiritual freedom and elasticity
in order to live in ideas without making them into immovable dogmas;
but there belongs to it also a mind that is not offended
that the great so often shows itself to research to spring
from the small and inconspicuous. What in the idea, in
the total view stands so free and high has slowly worked its way
forward and is due to contributions from many sides, as the great river
arises from the many small brooks. Analysis, the special
investigation leads us as it were behind the scenes, shows us the
machinery, and there easily arises a blasé attitude that exclaims:
``Is that all? Is that what lies behind it?'' Instead
of rejoicing in that secret of existence which
joins the great and the small, the infinite and the
finite together in an inseparable manner and in constant interplay,
one finds in analysis proof that the value of life is illusion,
or, if one wants to hold fast to the value, one turns away from
analysis and special research with scorn. Also in this
respect artistic virtuosity is needed in order to hold fast the
different points of view that may seem to conflict with one another,
and in order not to lose faith in the gospel that history
proclaims again and again, that even the greatest powers of life
once had to be swaddled and laid in a manger, and that the helping
and redeeming forces often come from a quarter from which
one expected them least.

The purely scientific work, pure thinking lays claim
only to a single side of the personality. But with
the task of working with the ideas in the changing
manifoldness of life, in the confusion of the small conditions without lapsing
either into dogmatism or into a blasé attitude, claim is laid
on the whole personality. Many a one who with enthusiasm
% --- p. 187 ---
\opage{187}threw himself into the scrutiny of the great problems in order to
penetrate them with thought has experienced that there came a
point at which the situation changed, and where he was thrown back
into himself from the more impersonal search for a purely objective solution of
the riddles and had enough to do with
working through in his personal life the thoughts won provisionally,
with cultivating in himself the personal art that belongs
to living in ideas. An enthusiastic thinker from
the time of the Renaissance, Giordano Bruno, has expressed this transition of
thinking into personal life by means of the myth of Actaeon,
the hunter, who after having caught sight of the
divine beauty became the prey of his own hounds: the hounds
are the thoughts that are directed toward the great ideas; when the eye is
opened to the ideal, the thoughts turn inward by reason of the
strong movement in the mind, no longer seek outward.
The problem becomes subjective, personal, no longer purely
objective. And Bruno emphasizes that so it goes where
it is not special and finite objects that thought is directed
toward, but the Divine, the Universal\footnote{Bruno: De gl’ heroici furori. Op. ital. ed. Lagarde. Gøttingen
1888. II. p. 651 f., 723 f.}. That research,
especially where it is directed toward the highest tasks, is transformed into
an eternal striving is a thought in which Bruno stands as
the predecessor of Lessing, Kant, and Goethe. Could the goal be reached
once and for all, there would be no room here for art.

Science itself finds its account in thus
giving room for an art which nevertheless preserves the connection
with it. For this transition to personal art brings with it
that we get a new domain of experience. Our horizon is enriched.
We become ourselves ``instances,'' ourselves objects of experiments.
Experiences are gained about the relation between the thoughts and the
personal forces, about the nutritive value and fruitfulness of the former
and the flexibility and endurance of the latter. They are human
experiences in the true sense of the word. And if we are able in some
way to shape and express them, they will be able to
benefit others. On the basis of
such experiences there will be able to form a comparative philosophy of
% --- p. 188 ---
\opage{188}life\footnote{Cf. my book: Søren Kierkegaard som Filosof. Copenhagen 1892. p.
52 f., 82 ff.}, which will become an analogy to comparative
science of religion. Hitherto one has had a tolerable
respect only for views of life that were common to nations, to
great groups of men. The more dogmatism withdraws
(which the heralded mysticism will, it is to be hoped, not
hinder for too long), the more will individual experiences and
the personal art of living come into their own. I see in this
a point of light for the future.

The diversity of personalities will inevitably have great
influence when philosophy is to pass over into art. What
I maintain is only that art from one essential side precisely
consists in preserving the connection between research and
personality, or rather between that part of the personality's
nature which pure thinking lays claim to and that
which is at work when the thoughts are to be held fast and carried through in
the reality of life. It cannot be otherwise than that the ways
at this point will part because of the personal
differences, and only through a comparative
philosophy of life can there again be a possibility of mutual understanding.

In the art of living everyone is first and foremost thrown back upon
himself. He has a task which he must solve in his own way, and
cannot solve in any other way, if he does not wish to end
in affectation or hypocrisy. --- When Buddha was born, there were
10,000 women who wished to be his nurse. And as Buddha did not
wish to grieve the 9,999, he multiplied himself
in such a way that each woman got her wonder-child to bring up\footnote{I take the image from \emph{Renan}: Feuilles détachées. 4. ed. Paris,
1892, p. 162 f., but I use it somewhat differently from him.}.
So we all have, each for himself, a divine child to rear. And
however many children there are, genuine they all are. It is a matter
for everyone to develop \emph{his} powers, to bring his personal
dispositions and impulses together with the thoughts that the generation
at the given time commands, or with as many of them as he
can appropriate, and through this to cultivate his philosophy of life, to practice
the human art of living in his own way.

% --- p. 189 ---
\opage{189}
\section*{III.}

Hitherto we have considered the philosophical art only in its
relation to ideas, to the ultimate hypotheses of thought. Still
more clearly the \emph{necessity} of this art shows itself where it is
a matter of practical conduct of life. It is a matter here not merely of
contemplation and theory, but of action. We no longer have
to do with theoretical philosophy, with the attempt to
develop a world-view, but with philosophy as ethics.
Also here it shows itself that there is an artistic conclusion
and carrying-through which theory cannot directly determine.
Philosophical ethics can give only the general principles:
the practical application in the single determinate case,
under the entirely concrete circumstances, cannot be directly derived
from them. And what here too will show itself to be
the great art is that of uniting opposites in oneself. In the
personal relation to the results of knowledge it was a matter of
uniting the confidence of faith with constantly continued discussion, the eye
for the great with the sense for the small conditions. In
practical conduct of life it is a matter of the union of
opposites such as memory and hope, absorption in the moment and striving
toward comprehensive aims, assertion of one's own individuality
and devotion to something that lies beyond it,
involuntary unfolding of life and voluntary intervention. Many more
opposed pairs would here on closer investigation
be mentioned, which the personal art is to unite into
unity in a way that no general formula can teach,
since the conditions, both in degree and in quality, are different for
each individual, and since the solution moreover
cannot be reached once and for all, as a proposition stands
fast when one has but once got it proved\footnote{Cf. my book: Søren Kierkegaard som Filosof. p. 44---48.}. --- Only
a couple of these opposed pairs shall I dwell on a little here,
in order to illustrate them by examples.

A danger that threatens when conduct of life is no longer to be
determined by tradition and authority, but by free
knowledge, is that the involuntary and immediate is cowed
% --- p. 190 ---
\opage{190}and crowded out. When philosophy as art is spoken of,
one will perhaps even straightway conceive the matter as if it were now a question
of regulating all manifestations of life according to conscious principles.
When the right of thought to intervene decisively in life
is acknowledged, there will easily arise mistrust or contempt of
the natural, habitual, traditional, and involuntary. One
will perhaps even find in emancipation from this in and of itself an advantage.
Arbitrariness is praised merely because it gives
something new, breaks with the habitual. Conscious reflecting,
reasoning, and regulating therefore easily come to play
too great a role, and this is at the expense of freshness and
immediate strength. The reaction after a rationalistic
period of emancipation finds perhaps its best support in the neglected
need to give oneself up to the involuntary instead of
letting everything be determined by conscious ideas.
The first effects of Christianity in the European world rested
for no small part upon its meeting this need halfway.
The Stoic doctrine, which was the most serious philosophy of life of the time,
suffered precisely from the one-sidedness that consists in distrust of
the involuntary movements of the mind. The Stoic wished above all to
preserve his rational equanimity, did not wish to be tossed about by
moods, not to be forced into passion. He sacrificed fullness and
depth in order to preserve self-mastery. Not to be astonished, not
to sigh in one's inmost being, not to rejoice exceedingly, that was his
maxim. Christianity, on the other hand, gave by its fundamental thoughts
and its cult occasion to live through the strongest
emotions, to come to know the most opposite feelings
and thereby as it were to measure out what life contains of light
and of darkness\footnote{See more fully on this the essay ``Pagan Seekers of Truth''.}. It was the justified reaction of the involuntary
against the voluntary. The Stoic doctrine of the governing
power of thought was nevertheless not wrong. But the relation is more
complex than the Stoic saw. Thought is the rudder, but wind
and wave must drive the ship forward. Here too the art is to
move through the opposites, so that both
come into their own; on the interplay of the voluntary and
the involuntary rests the clarity, health, and depth of life.
% --- p. 191 ---
\opage{191}In more recent times Rousseau's appearance also marks a
reaction of the involuntary against the conscious regulating which
the first period of modern rationalism had set in system.
He discovered a new world, that of the heart, of immediate
feeling, of the involuntary instincts, in opposition to
the culture of the understanding and external, mechanical
education. That he spoke up for nature against culture meant that
he wished once more to bring the immediate, unconscious, original, and
simple to honor over against the consciously formed and articulated.
It was like a new spring that burst forth in the desert; a new
appreciation of life was introduced, and the wave he set in
motion in the spiritual life of mankind was continued through
the romantic movement, whose forerunner he was, into
our own day. In a characteristic manner the
Rousseauan philosophy of life appears in his demand that education shall
be \emph{negative}, that is to say, consist essentially in the removal
of hindrances and excrescences, so that life is allowed to unfold
involuntarily. He maintains that only by this cautious
and indirect relation can one at all come to know the nature of the child
and avoid stunting it by intervening according to
preconceived ideas. Voluntary intervention presupposes, after all,
always such preconceived ideas; and the art, the difficult
art, is precisely this, to make the necessary use of the ideas
that are grounded on the conscious experiences of the past, without violating
or stunting that which has unconsciously, in concealment, sprung forth
and begins to unfold. Therein consists the art of all education,
also of self-education. It finds application to the
divine child that we each have to rear.

Another relation, in which the general formula that
can be given in theoretical ethics likewise does not suffice, so that
the personal art must work out the union of opposites,
is that between assertion of one's own personality
and devotion to others. From many sides there now sounds
the slogan of the ruthless assertion of individuality as
the one thing needful. Humanity, says Lagarde even, must yield
to individuality. But with no less strength the socialist idea
stirs in the consciousness of our time. For the
individual to find the harmonious, that is, the relation between the two opposites harmonious according to his
% --- p. 192 ---
\opage{192}nature\footnote{On the individual proportionality of all ethical demands,
see my essay ``Om Forholdsloven i Etiken.'' (Etiske
Undersøgelser. Copenhagen 1891, p. 53. 56.)} is a great art. To proclaim in strong words the one or
the other member of the opposition as the only thing that
matters is no art (except in the sense that it can
happen with the art of presentation; but that is not the kind of
art that is being spoken of here). Just as little as it is art,
is it science. It is agitation. --- In Ibsen's ``The Master Builder''
there are vivid contributions to illuminating the problem.
I shall be careful not to give any interpretation of
what the play is about. It is, after all, not certain that it
is about anything. I make use only of some traits of the
characters of the principal figures. There appears in Solness a
tendency toward absolute self-assertion, toward a ``bird-of-prey morality'' that
makes his own striving the only law. (With the expression
bird-of-prey morality it must never be forgotten, however, that birds of prey have
their nests, and that one raven does not peck out the eyes
of another.) The self-assertion that first becomes possible
through devotion and sacrifice, and the widening and
deepening of all the powers of the mind that is won only through
absorption in ideals that embrace more than one's own
individuality, --- these Solness does not know. But it nevertheless shows itself that
when he cows and deceives others to reach his goal, then
it is not a sign of power, but precisely of impotence. As
soon as the courage and force of youth (personified
in Hilde --- if it is not a sin to call this beautiful
figure a personification) appear in him again,
he at once does justice --- now, however, too late. And as
long as the master builder's plans embrace the interests of human life,
they succeed for him; when, dissatisfied with that, he wants
to go up into the lonely regions to build castles in the air, he
crashes to the ground. Whether it is Ibsen's meaning that that which hinders
Solness in ruthless self-assertion hinders him \emph{rightly},
is impossible to say, and it does not concern us either.
But with the striking power that is worthy of so great a psychologist
% --- p. 193 ---
\opage{193}he has shown that self-assertion is most easily carried through
against those whom one does \emph{not} know. For those whom we \emph{know} and come
into an intimate relation with, they belong to the nest of birds of prey, like
ourselves, and thereby the provisional limit of bird-of-prey morality is marked
--- not by a ``moral'' law, but by a law of nature.

Mrs. Solness is constructed as a counterpart. She goes
wholly into memory and clinging to the earlier, into
consideration and duties. For sheer duties she neglects her duty.
She will nurse her children because it is her duty, although
her milk in her sickly condition is harmful to
the children. The kindness she shows her surroundings she
herself takes the heart out of by expressly practicing it
as duty. And in the most essential relation in which she
stands, that to her husband, she goes wrong for sheer duties.
Perhaps Solness is right that her life is broken because
that which was her real calling, the rearing of children, does not
fall to her lot. But the many duties do not make it any better in any
case. Better it would be if \emph{she} went into
the great duty of self-assertion, developed fully into
a personality whose sweep and whose warmth could be communicated
to those nearest her. Instead of giving her small
parcelled-out services and considerations --- at bottom only an expression of
a whining self-assertion --- she could then give the highest,
a vigorously striving personality. As it has now turned out,
she lives at his side as a spiritual corpse --- and it
therefore threatens that she will be thrown out of the nest.

I take from the great poetry only what I can use in this
connection. It is the poet's advantage that he
can illuminate the personal problems by giving us a
presentation that includes all individual nuances and all special
circumstances. (Ibsen's greatness shows itself especially in respect
of the individual nuances, not so much in the dramatic
net he wants to draw together over his persons; for it
often has holes). The philosophizing view, which must
keep to general principles and typical examples, stands in a
great debt of gratitude to the great poets, whose
knowledge of human nature takes shape in individual figures. But
on the other hand I believe that one exaggerates the significance
% --- p. 194 ---
\opage{194}of such presentations when one holds that the poet
can straightaway treat or even solve problems. From that
he is already cut off precisely by the purely individual form
in which he gives his presentations. The question
whether what is presented can find application to other
cases remains unanswered; it was perhaps an exception
with respect to characters or situations, or opposing
portrayals can be set up. The problem is not
exhausted by the single poetic presentation. The debate
begins only properly after this, in that the question becomes that of the
relation of what is presented to other possibilities which life can
offer. To keep to the individual presentation and
regard the matter as exhausted by it would be a dogmatism
which is the opposite of that which can often be laid to the charge of the scientific
treatment of questions of life. Life goes as little
into a single constellation as into an abstract principle,
but must be illuminated from both sides. Art can as little
make science superfluous as science can do without
art. What it is a matter of holding fast, --- perhaps
especially in the times that are now coming (if those times come
that are heralded) --- is that even where art takes the place of science
and supplements it, the bond between them cannot be severed. There is a reciprocal relation between the two principal forms
of spiritual activity like that between inhaling and exhaling;
the one process continues the other and rests on the other,
despite the opposition that is between them. Their union
is necessary in order to make life free and healthy. This union
is difficult to attain; but Socrates, the grand master of philosophy
as art, has indeed also said: ``The beautiful is difficult.''

\streg{}

% --- p. 195 ---
\opage{195}
\essay[The Relation between Faith and Knowledge in Its Historical Development]{THE RELATION BETWEEN FAITH AND KNOWLEDGE IN ITS HISTORICAL DEVELOPMENT.}

\begin{center}An Introductory Lecture.\end{center}

\streg{}

The problem of faith and knowledge arises of necessity at a
certain stage of cultural development. We encounter it in the Indian,
the Greek, and the Muhammadan worlds as soon as spiritual, and especially
intellectual, development has reached a certain point. Here, however, we
shall confine ourselves to tracing the main forms in which it has
appeared within the European spiritual development shaped by
Christianity.

\section*{I.}

No instruction on this question can be drawn from the source writings
of Christianity. On the standpoint from which they were composed, no
knowledge distinct from faith exists. The gaze is fixed on one single
object, one single goal, which lays claim to all striving and all
spiritual force. The endeavor to know the truth is here immediately one
with the urge to win the salvation of the soul, and does not go its own
way by its own method. There is no distinction between theory and
practice, since theory and theoretical interest simply do not exist.
When the Deity itself has revealed itself visibly to all --- when faith
makes it possible to attain understanding of the Mystery that contains
all the treasures of wisdom and knowledge,
% --- p. 196 ---
\opage{196}--- every motive for seeking
knowledge by other paths is lacking.

It was natural that faith should from the outset have to be all in all,
leaving no room for other spiritual needs. The first generation still
stood so close to the mighty spiritual movement that led to the founding
of the Christian community that it lived and breathed entirely in the
great memories and visions. Added to this was the fact that this first
generation lived in the belief in the imminent return of Christ. Then
everything was to be concluded, the last judgment rendered, and all human
and earthly arrangements give way to a higher order of things. Who, with
such a catastrophe in view, had the time and composure of mind for
philosophical thinking or scientific inquiry? Science could no more than
any other human endeavor acquire any significance. There was to be no
historical development at all; the goal had been reached, and the end
stood at the door. One does not understand the New Testament and its
relation to cultural endeavor in its various forms unless one holds fast
to this. What happened here with scientific striving was the same as with
the striving for personal freedom and with the entering into marriage
(see 1~Corinthians, ch.~7). How could all such things acquire
significance --- or even become the object of a drive --- when all
thoughts and energies had to be strained to be prepared for the great,
the absolute decision!

\section*{II.}

Time passed; the conclusion did not come. The Church consoled itself with
the thought that with God a thousand years are as one day, and one day as a
thousand years (2~Peter, ch.~3), and soon began to settle down in the world,
to adapt itself to circumstances, and to take up a position toward the
various sides of the endeavors of cultural history. Toward science its
position became this: that it no longer found in faith all requisite
knowledge, but sought to show that the teachings of faith set the crown upon
all human knowledge. The Church appropriated the
% --- p. 197 ---
\opage{197}science the Greeks had developed: Greek astronomy, medicine, and philosophy, and placed its
theology as the spire upon the edifice of science, --- much as it converted the
Greek temples into churches and consecrated, for example, the Parthenon,
the temple of Athena, to ``holy Wisdom.''

Already in the so-called Gospel of John we find, in the opinion of some,
Platonic philosophy and Christian faith joined together. This endeavor is
carried on by the Church Fathers and reaches its completion in the classical
period of theology among the great teachers of the medieval Church. In
\emph{Thomas Aquinas} (born 1227, died 1274), the teacher still honored and
acknowledged in our own day by the Catholic Church, the relation between
faith and science stands as follows. Science investigates all \emph{subordinate} and
\emph{particular} causes, while the thought of the believer concerns the \emph{first}, the \emph{absolute}
cause of all things. Only theology, therefore, teaches us rightly to
understand things; science really teaches us only \emph{what}
things are, not \emph{whence}
they derive. But there is no conflict between faith, which is the principle of
theology, and reason, which is the principle of science. Reason stems from God
just as much as faith does, and faith does not conflict with reason, although
it goes beyond reason.

In poetic form we find this characteristic medieval standpoint in \emph{Dante} (born
1265, died 1321), who was a disciple of Thomas. In his ``Comedia'' the
poet relates how, in order that his soul might be purified, it was granted him
to pass through the world beyond: Hell, Purgatory, and Paradise. It is his
transfigured beloved, Beatrice --- the symbol of theology grounded in heavenly
revelation --- who bestows this favor upon him and sends Virgil, the symbol of
natural science, to lead him through the underworld. But Heaven Virgil dares
not tread; he himself states the reason, with great meekness:

\begin{verse}
for the Ruler to whom the city of Heaven belongs \\
does not wish that I should tread its threshold, \\
because I was a rebel against His laws.
\end{verse}
% ti Herskeren, hvem Himlens Stad tilhører,
% vil ej, at jeg dens Tærskel maa betræde,
% fordi jeg mod hans Love var Oprører.

Therefore, when the wanderers draw near to Heaven, Beatrice reveals herself,
and Virgil disappears.

% --- p. 198 ---
\opage{198}So pious and so humble had natural reason become in the Middle Ages. But its
powers grew, and the close of the Middle Ages is marked by violent discord
between Beatrice and Virgil, faith and knowledge. Here we meet the problem
posed for the first time.

\section*{III.}

The proud edifice that Thomas Aquinas had completed, and that Dante had sung,
soon began to totter on its foundation. Continued reflection discovered that
it had been built of elements that could not be brought into harmony. Already
\emph{William of Occam} (\dag\ 1347) finds a contradiction between natural reason and
the divine omnipotence that is exalted above every law. He teaches that
something may be true in theology which is not true in philosophy: if the
Bible and the Church teach something that is self-contradictory, it must
nevertheless be taken for truth! Theology rests on supernatural authority and
is therefore no science.

Have we then two kinds of knowledge, two kinds of logic? This question, which
lies so near at hand, is not brought forward. Occam himself was a believing
Catholic; he pointed out the discord, but did not pursue its consequences. The
Catholic Church held fast to its holy Thomas and has permitted no other view
of faith and knowledge to assert itself within its precincts. In our own
century no fewer than three Catholic philosophers have been condemned, because
they were either more or less rationalistic than St.~Thomas. For the Catholic
Church there no longer exists any problem here. --- Let us see, then, how the
matter stands within Protestantism.

\section*{IV.}

In modern times life spreads out in a great many different directions and has
a far more variegated and many-sided character than in antiquity and the
Middle Ages. And the Church stands as only one among the many domains of the
movement of life; it
% --- p. 199 ---
\opage{199}is no longer all in all, as in the first community, nor
the crown of the whole, as in the Middle Ages. History is not \emph{first and
foremost} a history of the Church, but \emph{among other things} also a history of
the Church. It is predominantly passively that the Church receives the
repercussions of the new directions of life and seeks to come to terms with them; but they
no longer proceed from it.

The relation of ecclesiastical Protestantism to science is much the same as
Catholicism's. Luther had indeed appealed --- not only to Holy Scripture, but
also to ``clear reasons''; but otherwise he applied to reason Paul's words,
that the woman shall keep silence in the congregation. Reason, then, was to
have Virgil's role, and might not enter the Holy of Holies. The Protestant
Church would gladly have gone on performing ``the divine comedy.'' But Virgil
has proved no longer tractable.

Protestant orthodoxy became a new scholasticism, but more anxious and less
grand than the medieval. --- It has found no Dante. The great Protestant poets
pass it by or go against it. Poetry's relation to the Church forms a parallel
to that of science: the ancient Church has no poetry, Catholicism an
ecclesiastical poetry, Protestantism a poetry standing outside the Church, if
not an anti-ecclesiastical one.

For Protestantism \emph{the eighteenth century} is what the fourteenth was for
Catholicism. But the crisis it brought was far more thoroughgoing, and reached
far deeper, than the Renaissance movement. It is \emph{the birth-time of the
independent humane view of life}. Forerunners there were, of course, many ---
but only now were they understood and acknowledged. Only now did Spinoza, for
example, who for a century and a half had lain ``like a dead dog,'' find a
circle able to think his thoughts and make them their own. Only now did the
theological view of life acquire an opponent of equal standing. I am thinking
here not of any definite, finished system, but of a general direction of
thought and life, which in its broad features was common to such men as
\emph{Lessing} and \emph{Kant}, \emph{Schiller} and \emph{Goethe}. --- As one sees, this new view of life
found at once more than one Dante.

What is distinctive about this standpoint, whose
% --- p. 200 ---
\opage{200}leading thoughts were only
transiently cast into shadow by the Romantic reaction, is --- \emph{seen from the theoretical side}
---, that what it sets over against theological faith is not a
pretended insight into the impossibility of the theological doctrines, but an
assertion of their unprovability. Kant's great achievement was, seen from this
side, to give a scientific doctrine of the limits of science. Not humbled and
cowed, like a rebel who has submitted, not checked from without, but with
independent and clear self-knowledge, human reason here says: so far can I
reach, and no farther! --- The limits of knowledge follow from the nature of
knowledge. By an investigation of the ultimate presuppositions and laws of our
knowledge, the domain within which it can move is staked out. This domain is
beyond surveying, yet it is definitely bounded over against that which by its
nature cannot become an object of our knowledge, constituted as this knowledge
once and for all is.

For Kant and his age the significance of knowledge lay not in its giving us
``the absolute truth,'' but in this, that by means of it we can advance
continually in truth. Eternal seeking, eternal striving became from now on the
watchword:

\begin{verse}
Wer immer strebend sich bemüht, \\
den können wir erlösen.
\end{verse}

Not as though nothing could be found or attained: but every find, every goal
attained, stands as the starting point for new striving. It was a great law of
life and of the world that was here given expression --- the thought of
development in one of its most important forms, a thought to which our own age
has given a realistic underpinning and a closer critical determination, but
which it is that generation's honor to have uttered for the first time with
full and clear consciousness.

But still more important was another principle, which was laid down as the
main point in the humane view of life: \emph{the principle of the free conscience}.
Kant's second great achievement was that he maintained the possibility and the
necessity of a purely humane ethics, and established that conscience can as
little as theoretical inquiry have its limits staked out for it
% --- p. 201 ---
\opage{201}from without.
The measure of all that we are to acknowledge as good, noble, and exalted must
be found within ourselves and be determined by human nature.

When those who stand at this standpoint cannot accept ecclesiastical faith, it
is because it conflicts not only with their understanding, but with their
conscience. One is not merely to assent to a formulation of the truth given
once and for all, but also to bow before doctrines that cannot be brought into
agreement with conscience's measure of good and evil.

Only here does the problem become sharpened, and its real sting come to light.
What from the one side is made the highest duty, in a way the only duty, a
duty whose non-fulfillment brings damnation --- that is rejected from the
other side precisely as conflicting with the clear and definite demand of
conscience! ---

Bishop Martensen saw clearly, from his own standpoint, that the most important
turning point in the controversy over faith and knowledge came with the
founding of ethical humanism at the close of the last century. In his
``Kristelige Etik'' he deals at length with the representatives of this
tendency and their relation to the faith of the Church. The final explanation
of their refusal to accept it he finds, rightly, not in their understanding,
but in their will --- only that this will is, according to his assertion,
defiance against the Holy: ``Against God's holiness they have a secret
antipathy!''

How strange, or rather how appalling, human life must look from the orthodox
standpoint! The most honest seeking and striving, the strictest
conscientiousness, joined with the most unreserved acknowledgment of the
imperfection of what has been attained, must from this standpoint be construed
as antipathy toward holiness. --- If one has not understood why the problem of
faith and knowledge so easily sets the passions in motion, one understands it
here.

% --- p. 202 ---
\opage{202}
\section*{V.}

The new humane standpoint seemed to have a fairly good prospect of prevailing.
The most eminent minds of the age paid homage to it, the leading statesmen
worked in the direction it marked out, within the Church itself orthodoxy was
decidedly weakened, and confessional differences receded before an enlightened
and liberal-minded way of thinking that was common to all. Such were
conditions at the beginning of the century, not only in the Protestant, but
also in the Catholic Church.

But how differently were conditions to develop in the course of the century,
and how changed were they already after the first two decades had run out!

A single trait shows the change plainly. In the year 1799 Schleiermacher
published his ``Speeches on Religion to the Cultured among Its Despisers.'' He was by
no means contending there for any dogmatic faith, but asserted a pantheistic
religiosity and glorified Spinoza. For such a view he wished to win
recognition among the cultured. But when in the year 1821 he saw the third
edition of his book through the press, he declared in the preface that if one
\emph{now} looked about among the cultured, one would rather find it needful to write
speeches to bigots and slaves of the letter, to ignorant and uncharitable
condemners, to the superstitious and the hyper-dogmatic!

Four years earlier (1817) the Holstein pastor Claus Harms had, in commemoration
of Luther's action, published 95 new theses, in which he protests, among other
things, against ``the popes of our time'': reason in the theoretical domain,
and conscience in the practical domain. --- How was it possible that such a
performance could win adherence and mark the beginning of a new period?

What, on the whole, was the reason that the breakthrough of humanism at the
close of the eighteenth century did not come to inaugurate a development
carried on in a straight line down to the present, but was succeeded by its
sharp and conscious opposite?

Since a closer psychological-historical explanation of the mighty reversal
that came to fill so large a part of
% --- p. 203 ---
\opage{203}our century, and made it ``the age of
oppositions,'' may serve to throw further light on the nature of our problem,
we shall enter into it somewhat more closely, but in such a way as to keep
chiefly to the causes that can be found in the very character and
distinctiveness of the humanism of the eighteenth century. Many different
causes were of course at work together. Thus the general political reaction
had its influence on religious development as well; the resurrection of the
principle of authority in the state had also to lead to its being favored in
the religious domain, and the political reactionaries often made use of
religion in the most repellent way. But had there not been deep-lying inner
causes, this outward influence could not have accomplished so much. ---

1. Every reaction is due in large part to an \emph{effect of contrast}.
Just as the eye wearied by light finds rest in darkness, or at least in a
fainter gleam, so on the whole the human spirit cannot endure a strong and
protracted straining of a single side of its nature, but throws itself
involuntarily over to the opposite side in order to find a new position of
rest. With respect to the problem we are speaking of here, it is above all the
\emph{opposition between the active and the passive, the effecting and the receptive
side of our nature} that becomes significant. Not as if in our spiritual life
there could be pointed out states or phenomena in which we are either
absolutely effecting or absolutely receptive: in every form of spiritual life
and in every spiritual state there can, on the contrary, be pointed out both an
activity and a passivity. But these two moments can stand in very different
relations to one another; now the one, now the other has the preponderance.
Thus inquiring and brooding thought, the testing and forward-striving
conscience, and the forming and creating imagination have a preponderantly
active character, and over against them in this respect stand the need for
consolation and self-surrender, and feelings such as piety and resignation.
Where it is a matter of bringing forth new forms of life, the active side of
our nature acquires the preponderance and lays claim to the whole of our
energy.
% --- p. 204 ---
\opage{204}Such a task presented itself precisely to the generation of the latter part
of the eighteenth century. It was the age in which unlimited advance was the
watchword, and inertia was regarded as original sin. But there had to come a
time when the other sides of human nature demanded their satisfaction, and
when men sought rest after the restless labor. And rest was not to be found in
the new forms. They still required all too much active cooperation, made a
constant demand for self-activity. Rest, on the other hand, might perhaps be
found in the old forms, on which in one's onward rush one had turned one's
back, but of which one now discovered how wonderfully they were adapted to the
needs of the life of feeling: for centuries or millennia generation after
generation had deposited its experiences in them, had arranged them so that
uninterrupted labor was not needed to use them. Our activity seeks the new;
with our effecting nature we live in the present; but the passive and
receptive side in us is determined and nourished preponderantly by the
after-effects of the preceding age; here spiritual life draws near to the
unconscious and takes nourishment from it. An example of this from the life of
the individual we have in the power of childhood memories: they work without
our conscious cooperation, and they rise up with overwhelming force
particularly when the mind is weary of active intellectual labor.

It is characteristic that Schleiermacher, precisely at the transition to the
new century, set up his famous definition of religion as the feeling of
unconditional dependence. It was like a program for the direction the
intellectual life was now to take. Humanism was succeeded by confessionalism, free inquiry
and the free conscience by faith in authority.

Are we then consigned to a labor of Sisyphus? Does the stone always roll back
once we have got it up the mountain? --- No; but what lies in the foregoing is
that only that truth can be victorious which is able to penetrate and fill the
whole of life. And that does not happen at a single stroke; it takes time, and
only through many oscillations does it advance. It seems to be a psychological
law that \emph{knowledge develops more rapidly than the life of feeling}. Only slowly
are the results of active
% --- p. 205 ---
\opage{205}and conscious labor deposited in the involuntary life
of feeling. Before that can happen there naturally comes a period in which the
new thoughts grow pale and prove unsatisfying; doubt and unrest take hold of
many, and there seems to be no other way out than to throw oneself back into
the old forms of the life of feeling, if one is not to perish. But this
disharmony will be only transitory: if the new thoughts are valid, if they
really express some of the laws to which our existence is subject, then they
will always work their way forward anew, and the time of unstable equilibrium
will be past. In the long run the life of feeling will develop in greater or
lesser harmony with the life of thought: it was in that way, after all, that
those very old forms and dogmas came into being. Even the periods of reaction
themselves prove to bear the stamp of what they combat; into the old
wineskins new wine will to no small extent be poured, and only one who keeps
to externals will suppose that everything has become again as it was before.

We stand here before a spiritual natural interconnection to which we must bow,
and which we cannot defy. Whoever has come to see it rightly will neither
share the cheaply bought confidence of many in the nearness of victory, nor
lose heart at an apparent setback. He knows that just as the apparently sudden
revolutions signify only that what has long been prepared in secret now steps
forth into the light of day, so the sudden reactions indicate that there were
drives and needs which were only imperfectly satisfied by the seemingly so
brilliant advance, and which therefore must seek satisfaction in the old,
until they have accommodated themselves to the new, or the new to them.

2. The humanism of the eighteenth century not only gave the knowing
and effecting part of our nature a decided preponderance over the feeling and
receptive part, but, in close connection with this, it had a decidedly
\emph{intellectually aristocratic} character. The whole movement was, and by its
nature had to be, confined to a small circle; the people proper were very
little affected by it. At no time, perhaps, have the well-placed classes in
society shown so great a zeal for taking the lead in the
% --- p. 206 ---
\opage{206}work of enlightenment
and progress. But the distance between those who led and those who needed
leading was extraordinarily great. And the development had led to a break with
the \emph{common} faith, the common forms for the highest thoughts. In the Middle
Ages the most deeply brooding schoolman and the simplest serf could meet in
the same confession. Now this unity was shattered, and there had to arise ---
especially since at the same time feeling for the national and the popular took
a strong upswing --- a need to be one with one's people in faith as well. The
reaction against the eighteenth century acquired, seen from this side, an
\emph{intellectually democratic} character. It was, after all --- to take a single
example --- the conviction that ``the faith of the unlearned cannot possibly be
bound to the testimony of the learned'' that led Grundtvig to his
``discovery.''

Like the discord between the different sides or forces in the spiritual life,
the discord within the life of the people will also be a constantly operative
motive for returning to the old faith. Common forms and expressions for the
highest thoughts are created slowly and presuppose that the thoughts have
passed into the flesh and blood of all. Since now the individual differences
in religious need and religious sense must be assumed to rise and increase the
more religious freedom becomes a reality, such a community will become ever
harder to bring about.

A testimony to how great the need for community can be is offered by the
philosophy of religion of \emph{Schelling} and \emph{Hegel}. These men and their followers
were convinced that in their philosophical ideas they had developed the same
content that the dogmas of the Church expressed. The difference was to concern
only the form: what theology has as an idea, in the form of imagination
and feeling, philosophy was to shape into abstract concept, without the
essential being thereby lost. In the dogma of God's becoming man, for example,
there was supposed to be contained the philosophical thought that even the
highest must develop through struggle and oppositions, and that a principle
shows its validity only when it is able to penetrate actual life with all its
discord and all its conflicting relations.

In and for itself it is entirely legitimate to regard the dogmas
% --- p. 207 ---
\opage{207}as symbols.
For the non-believer they cannot be anything else, if they are to be anything
at all. And many who consider themselves very believing, and are so considered
by others, stand in reality at this standpoint. But it will always lead to
dubious consequences when one uses the symbolic interpretation to bring about
a false semblance of community. A real unity cannot be produced by that road.
On such a false semblance rested the last attempt at a reconciliation of faith
and knowledge that has been made, Hegel's philosophy of religion, from which
many expected a ``new era.'' Here Virgil was not to subordinate himself humbly
to Beatrice; they were to walk arm in arm. But this apparent agreement gave
occasion only for the opposition between faith and knowledge to be expressed
in a still far sharper and more pointed form than hitherto, as was done by
\emph{Strauss} and \emph{Ludwig Feuerbach} in Germany and by \emph{Søren
Kierkegaard} among us.
These men were agreed that the content of theological faith was the
incomprehensible; the difference between them consisted in this, that Strauss
and Feuerbach said: It is incomprehensible, and therefore it cannot be
believed! while Kierkegaard said: Precisely because it is incomprehensible, it
can and shall be believed! --- Strauss's words: ``Of false attempts at union
enough has been done; only a sharpening of the oppositions can lead further!''
are characteristic of the course of development in the nineteenth century. A
famous scholar of church history, \emph{F.~C.~Baur}, has rightly found the
distinctive mark of our century in this, that all oppositions --- in science,
religion, and politics --- are worked out to the highest possible degree. That
hangs together naturally with the fact that reaction and revolution in the
various domains have so frequently stood over against one another; thereby
there has arisen a clear consciousness of the nature of the oppositions and an
endeavor to work each of them out separately in its full consequence. Both the
necessity of making a choice between irreconcilable oppositions, and the
character of the paths one has to choose between, dawn more and more upon the
general consciousness. ---

If the mere need for community has thus led to sympathy with the old faith ---
whether this sympathy
% --- p. 208 ---
\opage{208}led to personal adherence to its letter or merely to a
symbolic interpretation of its statements --- then this whole question of
faith and knowledge had to present itself in a particularly serious way to
those who saw its connection with the social question. Even granting that the
scientific criticism of the dogmas leads to their dissolution, the thought will
still arise that these dogmas, however contrary to reason they may be, are for
many their only spiritual support and consolation --- that which keeps up
their hope under life's heavy struggle --- the only thing at their disposal
that can lift their gaze beyond the narrow circle in which their life moves.
It is from this point of view that one of the noblest philosophical thinkers of
our time, \emph{Albert Lange}, the author of ``Geschichte des Materialismus,'' has
expressed himself against Strauss and Feuerbach, with whom he agrees as
regards the purely theoretical side of the matter. It is on the whole
characteristic how great a role the question of religion's usefulness and
indispensability plays in the religious debate. What is at issue in the first
line is not: truth or not truth? --- but: consolation or not consolation? ---
In the debates of earlier times this side of the matter did not come forward
as something on its own, still less as the main thing. Men did not search for
grounds and motives for holding to dogmatic faith in spite of the
difficulties; rather, it was those who gave up dogmatic faith upon whom it fell
to give reasons for their step.

It is surely of very great importance that spiritual community should not
disappear, and that our race should not lose any force that can hold it up in
the struggle for existence. It is with full right that these considerations
become decisive for the manner in which the religious struggle is conducted;
but the struggle itself they can damp only for a time, since they point out no
way to the solution of the problem.

3. Finally, the humane view of life, as the eighteenth century was
able to present it, had a one-sidedly \emph{abstract} and \emph{idealistic} character. It
was stated according to its general principle and with the most important
presuppositions and demands it carries with it, but it lacked sufficient
carrying-through and a real foundation. Practically every
% --- p. 209 ---
\opage{209}advance the
nineteenth century has made has contributed its share toward remedying these
defects to some degree. The magnificent upswing of the natural sciences has
given a firmer basis for the conviction of the law-ordered connection of
nature. In the concept of development a guiding idea has been indicated around
which all our knowledge in the domain of nature and of history can gather.
Historical criticism and comparative religious science have widened and
corrected the material necessary for the treatment of problems in the
philosophy of religion. A more deeply penetrating psychological understanding
of the various sides of human mental life in their relation to one another,
and a renewed discussion of epistemological questions, give the problem of
faith and knowledge a somewhat different character from what it could formerly
have. --- All this is of course only small progress in relation to the
magnitude of the problem. The problem will for long ages accompany the human
race on its way, and the different ways in which it is conceived and treated
will be able to furnish a measure of the degree and direction of spiritual
development at the various times. Here we have only to show how the problem
must present itself in our own day.

\section*{VI.}

In the discussion of faith and knowledge that took place in our literature in
the sixties, great weight was laid on the distinction between faith and
theology, and I myself, as a disciple of Rasmus Nielsen, belonged to the number
of those who attached no slight hopes to it. It plainly does not help us, and
in any case it is not tenable. As soon as faith has a definite content, a
definite object, there must also develop a doctrine about this content, which
seeks to exhibit its inner coherence and its significance. Otherwise neither
the believer himself nor anyone else knows what it really is that he believes
in.

Nowadays, moreover, theology has given up supporting its assertions with
scientific proofs. It has given up all proud dreams of getting Virgil to walk
in leading-strings, and
% --- p. 210 ---
\opage{210}is content when it is left unassailed from the side of
scientific criticism. It declares its dogmas to be \emph{unprovable}, perhaps even
concedes that they present difficulties for the ``natural'' understanding, and
takes essentially the standpoint of maintaining dogmatic faith as a spiritual
necessity of life --- a necessity, be it noted, which it is a duty to have. It
is supposed to be a sign of hardening and of defiance against the Holy when one
does not feel that need which makes reason bow before the incomprehensible.
Therein lies, as I have already had occasion to touch upon, the real sting of
the problem.

On the other hand, no scientific thinker will assert that the absolute
\emph{impossibility} of the content of the dogmas can be strictly proved. In
scientific inquiry we build upon the principle of natural causes. It has helped
us through many a time, and hardly any other scientific principle for the
understanding of existence will ever be able to be set up; --- but it cannot be
proved that \emph{everything} in existence is subject to this principle.

Indeed, we must go a step further still. Simultaneously with the wonderful
advances that have been made in our knowledge of nature and of history, both as
regards particulars and as regards the great connection of things, we have
learned in a far deeper sense than before that \emph{the whole of existence is one
great riddle}. For the philosopher the least phenomenon is in this respect just
as instructive as the greatest and most conspicuous. The simplest causal
relation gives us more than enough to think about --- really encloses the whole
riddle of the world within itself. Everything we understand and explain, we
understand and explain by exhibiting a causal connection: if we only knew what
causal connection really is --- what it is that holds everything in nature
together and makes it that wonderful whole which reveals itself to us more and
more --- then all riddles would be solved! But to such knowledge every prospect
is barred. Our sensations, ideas, and concepts (and thus also the
concept of cause) are quite good aids for finding our way about in existence,
but we have no right to suppose that they can depict for us the
% --- p. 211 ---
\opage{211}innermost essence of existence. We can at most form speculative hypotheses about how
existence might be constituted in itself; but such hypotheses cannot, in the
nature of the case, receive any confirmation from experience.

From this it by no means follows that one bows before every theological oracle.
It follows only that the discussion of faith and knowledge must be conducted in
a somewhat different way from that in which it has hitherto been conducted. Let
us briefly touch on the main questions that will here have to be discussed.

1. When theological faith acknowledges that its assertions are
unprovable, or perhaps even irreconcilable with some of the most important
presuppositions and results of science, it apparently places itself at an
inaccessible and unassailable standpoint. There seems no longer to be any
common ground on which the battle can be fought.

But when theological faith nevertheless --- in order to express its content ---
makes use of ideas and concepts that natural reason and scientific
thinking also employ, then it must be asked whether these concepts are taken in
exactly the same sense. It has proved to be bound up with insuperable
difficulties, even for the boldest speculation, to carry through propositions
and judgments about what lies beyond every possible experience. What, now, can
it really be that enables faith to come off better here than scientific
thinking? \emph{Why do the great difficulties in the dogmas fall away when one passes
over from knowledge to faith?} A creation out of nothing, for example --- a
first cause --- an absolutely good being who is nevertheless the cause of
everything, and therefore also of evil --- such concepts, against which human
thought has collided again and again, until at last it has given up the attempt
to reconcile them with its other concepts --- how can they appear in a wholly
different light because they are declared to be objects of faith and not of
knowledge?

If it is answered (as, for example, by Bishop Martensen in his
Dogmatik) that the rebirth of human nature that takes place in the
believer embraces knowledge as well, so that
% --- p. 212 ---
\opage{212}\emph{that} becomes intelligible which was formerly incomprehensible
or even self-contradictory, then of course the unreborn
cannot have a say about how things look after so radical a change. But since
the believing and reborn do not, after all, entirely give up their connection
with the natural world, and, insofar as they continue to belong to it, must
also continue to use their natural reason, they ought nevertheless to be able to
compare the two kinds of knowledge, the natural and the reborn, and tell us
wherein the difference really consists. Here there is a chapter in dogmatics
which the theologians, in all the many years during which theology has
flourished, have neglected to write, namely: a new or higher logic and theory of
knowledge, which could show us by the help of what new principles one slips
past, in faith, the reefs on which natural thinking constantly runs aground.
Until this chapter is written, we must have the right to examine theological
concepts in the same way as we examine all others, and to pass the same judgment
on inconsistencies and unwarranted applications of concepts as we pass
everywhere else. Such an examination will lead to the result \emph{that the dogma
stands in reality just as powerless before the great riddles as science does}.

2. But --- it will then be said from the theological side (especially
now, when philosophical interest in our theological world seems to have been
laid in the grave with Bishop Martensen) --- it is after all a misunderstanding
to suppose that faith is to satisfy a theoretical need. For faith, thought and
reflection mean nothing, and all the counter-arguments in the world fall
powerless to the ground for the simple reason that no one comes to faith except
the one who seeks consolation for an anguished conscience. It is
\emph{feeling}, particularly the anguished conscience, and no other force, that is at work here.

This is not an agreeable objection, if what is meant is that the speaker himself
cannot do without the consolation that lies in faith; for in that case the
discussion must cease. Who would rob anyone of his consolation? Even the most
zealous chemist would not use for analysis the drink of water
% --- p. 213 ---
\opage{213}with which a human
being is to slake his burning thirst. Every discussion presupposes that the
participants each have so much calm and strength of mind that they can look the
truth in the face, so that one need not, as with the sick, conceal anything.

What consoles the one, moreover, does not console the other, or even arouses his
offense and indignation. As soon as an appeal is made to feeling, the appeal is
more or less to something purely individual. Such an appeal can never be an
argument; for \emph{feeling in and for itself teaches us nothing whatever except our
own state}. The most vivid joy and reassurance that I feel at a piece of news
guarantees nothing about that news's reliability. From my feeling I can
therefore infer nothing with certainty about anything that lies beyond my own
state. Our wishes do not determine what is truth; but it might well be that
H.~C.~Ørsted was right that the full truth always brings its consolation with
it.

The feeling to which appeal is especially made as a support for theological
faith is \emph{the consciousness of sin}. Here now arises the great question whether
the consciousness of sin, as theology conceives it, really is, ethically
regarded, a sound and warranted feeling --- whether there may not be contained
in it elements that are not really ethical, but hang together with barbaric
ideas (of an angry deity who must be propitiated), and which must be
separated out if the consciousness of sin is to become a real ethical feeling of
remorse.

But it is precisely the measure of the ethical that is so different at the two
standpoints! And theology, which formerly referred us from the dogmatic over to
the ethical, now maintains conversely that without dogmatic faith no ethics is
possible. The appeal to the consciousness of sin therefore leads us only in a
circle: for the consciousness of sin is a feeling conditioned by definite
dogmatic ideas, and it is then no wonder if it bears witness in favor
of the dogma. --- We get out of this circle only if theology can show us --- what
has hitherto not succeeded --- that only dogmatic faith can be the
% --- p. 214 ---
\opage{214}foundation of
ethics; indeed, that the theological assumptions are on the whole suited to be
the foundation of an ethics. If, on the other hand, a humane ethics is possible,
and if an unconscious humane ethics should prove always to underlie theological
ethics in its various forms (as one might infer from the fact that when human
beings improve, their gods improve too), then the sharpest sting would be taken
away, and what then remained of the problem might well continue to set feeling
in strong and deep motion, but would no longer split our race into two opposed
camps that clash in irreconcilable strife over the ultimate measure of good and
evil.

3. There is an assertion that has long been waiting impatiently to be
heard from the theological side. It runs roughly thus: ``Faith is built upon
\emph{historical facts;} its content is first and foremost history, and neither some
excogitated system nor mere postulates of feeling. And historical facts one must
take as they are; no argument of reason can shake them.''

If only the word history were here taken in the same sense in which we take it in
historical science! But it goes with the theological use of this concept as with
the use of other concepts (cause, ethics, etc.): despite the identity of the
word, the real meaning of the concepts is very different. Historical is not
everything that stands in old books or is contained in old traditions. A
historical fact may often require penetrating critical inquiry before it can lay
claim to acknowledgment. The great question then becomes whether theology will
submit to the ordinary laws of historical criticism, or whether it lays claim to
certain privileges for its ``facts.'' That it must demand such privileges follows
already from the fact that its source writings contain narratives of miracles
that require faith in order to be acknowledged. A historical inquiry that applies
no other presuppositions than those from which all science must
% --- p. 215 ---
\opage{215}set out, leads ---
as experience has always shown --- to a wholly different conception of the events
at the founding of the Church from the one that theological faith lays down. It
goes, then, with the appeal to history as with the appeal to feeling and to
ethics: it leads in a circle, for by history is meant history \emph{theologically}
conceived and evaluated. ---

\streg{}

% --- p. 216 ---
\opage{216}
\begin{center}

\essay[The Religious Problem]{THE RELIGIOUS PROBLEM.}

Address at the University's Reformation Festival, November 16, 1897.\end{center}

\streg{}

The festival we celebrate is dedicated to the memory of one of the great
breakthroughs in history, in which what had been prepared by the
inner work of thought and conscience, through powerful
personalities, claimed the right to rule in the outer life as well.
The Reformation stands as a testimony that in the
spiritual domain, as in the life of nature, a
development can take place that leads to new forms of life. Against this
it does not militate that the Reformers essentially conceived their work
as a renewal of original Christianity. They stand
here as the men of the Renaissance in general stood. In every domain
people wanted a rebirth of antiquity. In art and in
literature, in thinking and in the study of nature as well as in
religion they appealed to the models of antiquity. The
new art was inspired by the antique; the new philosophy went to
school with Plato, Aristotle, and Lucretius; Copernicus sought
to win recognition for his new astronomical theories
by appealing to similar ideas among the ancient
Pythagoreans, and the atomic theory, which was to play so great a
role in modern physics and chemistry, was presented at first as
a reconstruction of ancient atomism. Machiavelli has in
his ``Discorsi,'' which were written about the same time as the
outbreak of the Reformation, given utterance to the tendency of the age, when he says
that religious and political institutions must from time to time
be led back to their origin, judged and inspirited by
comparison with it.

% --- p. 217 ---
\opage{217}In the absolute sense of the word, a rebirth of the thought or faith
of an earlier age is not possible. As history shows, rebirth
is in reality a new life. That immersion in
antiquity in the domain of art, thinking, and the study of nature
did not bring the antique forth again, but served as
preparation for new forms and directions, is denied by
no one, and it is no different in the domain of
religious life. It could not be otherwise,
unless life were to go backward. The need that led
to a new immersion in antiquity was in reality a new
need, a need that had arisen because new conditions
and necessities had developed in the spiritual
domain.

As far as it was from being the case that the remarkable age
whose memory we celebrate reproduced the antique and the
primitive, just as far is it from being the case that it was once for all to
have found and uttered the truth. ``Truth, thou recognized
fatherland!'' it says in our Reformation cantata. But
truth is great and comprehensive, and no generation can find
it for all time. The very truth that the Reformers
recognized must again and again be tested as to its validity. The
religious problem stands today, and in far sharper
and more conscious form than before, in the first rank among the
problems of the spiritual life. It is now posed in an essentially
different way than at the time of the Reformation. New elements
have been added in the domain of intellectual life, elements
that the Reformers did not know or did not heed, and which give
the problem itself a new significance and set the conditions for its
treatment essentially differently. It would not be in the spirit of
the Reformation if we were to close our eyes to this. But
it will be in the spirit of the Reformation when we undertake
an investigation of the elements that --- as far as we can
see so far --- are of essential influence on religious life and
religious faith.

There is scarcely any of the sciences that are
represented here at the University that could not
give a contribution to such an investigation. I shall from my
standpoint continue what two of my colleagues have begun
% --- p. 218 ---
\opage{218}in addresses at this festival. One of my theological colleagues
has spoken of the significance of historical criticism in
the establishing of the religious tradition, and my astronomical
colleague has spoken of the quiet, but constantly
growing influence of empirical science on the content and
significance of religious ideas.

There are, as far as I can see, four elements that undoubtedly
determine the character that a human view of life, whether
it takes the form of religious faith or not, assumes
at a particular given time. They may be briefly designated as
tradition, observation, thinking, and will.

Religion persists from generation to generation as a venerable
tradition, into which the individual involuntarily lives himself.
It gives form and content to the need to understand
life, its significance and its tasks. A religious
problem arises, as concerns this point, when it is asked
whence the tradition derives, how far one can rely
on its authenticity, and who in cases of doubt is to decide
such questions. In our days what is at issue here is the
relation of historical criticism to ecclesiastical tradition. I
need not remind you that this subject has precisely in recent
years come up here at home too. A still more
burning and deeply penetrating question is whether the
conduct of life that now calls itself by the name of Christianity
is really the same as that which was originally attached to
this name. From this side the religious problem was raised
here at home by our people's greatest religious thinker in the form of
the question: Does the Christianity of the New Testament
really exist?

The observation and knowledge that the individual can win
will necessarily gain influence on his faith. It will
determine the way in which he conceives the
tradition, --- what he really appropriates and what he lets
lie, --- what he takes in the literal sense and what
he conceives as mere symbol. Herein consists what my
astronomical colleague has called the religious significance of
observation.

When I set thinking as a separate element in
% --- p. 219 ---
\opage{219}this connection, I mean by thinking the
spiritual work by which we become clearly conscious of the content of our
ideas and the mutual relation between our
various ideas. Such a spiritual
work goes on uninterruptedly and involuntarily, with greater or
lesser energy. Here a whole series of questions arises.
Can religious assumptions be proved by the way of reason? The
Catholic Church has by its highest authority established that
it is not only possible to prove the truths of so-called natural
religion (especially the existence of God), but also to
furnish a proof of the necessity of revelation and of the infallible Church.
In the Protestant world, where critical
philosophy (without one's being able to ascribe to the Protestant
Church any credit for this) has on the whole displaced
Scholasticism, the question poses itself differently. The question is
not how far religious assumptions can be proved, but whether
they can be refuted. When they cannot be refuted, one considers
oneself entitled to hold fast to them. Against this, however, objection
arises from another side. It is a rule that is
recognized everywhere in science that one must not use
other assumptions and suppositions than those that follow with strict
necessity from critically tested experiences; it
is not enough that a proposition does not conflict with what can be
proved. The struggle stands between these two standpoints: from
the one is maintained the right to hold fast to what cannot
be scientifically refuted; on the other is maintained the obligation
not to assume anything other than what follows with strict
necessity from real observation. From the latter
standpoint it is maintained that the law of parsimony, which
holds for scientific assumptions, ought also to make itself felt and
gradually will make itself felt in the domain of faith and
of view of life, so that one will approach a point
where empirical science becomes the sole foundation for the
thoughts that man has need of in the struggle of life. That
such an approach is taking place there is no doubt.
The question is how far it will be able to go. A spirited
American philosopher, William James, has said that
the carrying out of the law of parsimony in the domain of faith
% --- p. 220 ---
\opage{220}would lead to barbarism. Of this, however, we can know nothing, since the
time when there could be talk of such a thing lies remote
yet, if it ever comes; should it come, it
may bring with it spiritual conditions of life that we now
cannot form any idea of. But for the time being we stand
in the midst of development. Experience brings us constantly new things,
not least in the domains where the questions of life belong,
the biological, the psychological, and the social
domain. All faith concerns, after all, the continuance of life, particularly of the spiritual
life, the preservation and development of the highest
values we have learned to know. And one must constantly fight
for life, not only for the outer, material life, but for
maintaining that life which in our highest and best moments stands
before us as the only life that is worth living. And here
we stand --- despite all that empirical science has brought ---
constantly before the enigmatic and the unsurveyable. New
experiences displace old, new points of view prove
necessary, and the way in which we have arranged things
for ourselves is constantly put to new tests. With suspense must
precisely he who attributes to experience a decisive significance
for his faith follow the continued development of the life around
us and in us:

\begin{verse}
Neu ist immer die Erfahrung, \\
Immer ist dem Herzen bang!
\end{verse}

This consideration leads us over to the last, the most
decisive element: the personal need and longing, which is
unsatisfied by everything given, because it aims at something
that experience only imperfectly lets come into
possession. Through the concept of value religious faith is connected
with the will. Our faith is determined by what for us stands
as the highest value, which gives all else value, and by the
relation we think we find between the world of reality and
this highest value. All value rests on an aim that our
will with greater or lesser consciousness holds fast. As
long as there is a relation of tension between our highest value
and empirical reality, so long there is room for a
religious problem. It is, to be sure, first and foremost our
% --- p. 221 ---
\opage{221}own task to fight for the highest valuable --- for
its victory in us and outside us; but the religious problem
poses itself with the question: is it we alone who fight
for it, or are there powers and laws in existence that work
in harmony with it and in the direction of it? --- It is the
Reformation's greatest merit that it has maintained this as the
fundamental question. The Reformation let the old dogmatics
stand, but laid the whole weight on the inner assent of
personality. In his Catechismus major Luther asks: Qvid est
Deus? (What is God?) And his answer is: that to which you have
full faith and trust. It is the heart's trust alone that makes
something a God or an idol (sola cordis fiducia deum
pariter atqve idolum facit et constituit); and the right God
you will have when your faith and trust is right and genuine. Luther
became the man he became because he knew from his own experience
the anguish that is felt in the struggle for life's highest values,
not only as to whether they would conquer outside us, but first and
foremost as to whether they would conquer in ourselves. There, where the will's
limit lies --- inwardly and outwardly --- there the will asserts
itself still in the form of faith. What Scholasticism and the Roman
Church system could not let come to unfolding, or what
they satisfied too hastily and on too easy terms, ---
the individual's inner need and demand upon himself and upon the world,
that the Reformers placed in the first rank. It is their
great, inalienable merit.

The personal need to assert the values of life is
the factor that finally becomes decisive for the influence
that tradition, observation, and thinking gain. It makes its
choice among the multiplicity of traditions and observations,
emphasizes some and places others in the background, and it
inhibits or promotes the consistency of thinking. It is
the most restless element in the religious. It is toward it
that attention turns more and more, because there has
arisen a vivid sense that it is here that the battle
is finally to be fought. The old gods, from whom the race has turned
at various stages of its development, have never
been refuted. Dynasties of gods die out when there is no longer
need of them. Therefore the treatment of the
% --- p. 222 ---
\opage{222}religious question turns more and more from the domain of speculation to that of
psychology.

But man's tendency to dogmatize readily follows along here too.
The psychology of religion is still very little
cultivated. Only slowly has the conviction become
general that Luther --- at least in a single place --- has so clearly
uttered. And after it had begun to take hold, people thought
they could assume that this decisive element, the
personal need and longing, the innermost will, must be the
same in all and at all times. Psychological observation
does not confirm this assumption. There is a
relation of reciprocal action between this personal factor and the three other
factors: tradition, observation, and thinking determine
the nature and direction of the need, just as the need in turn
gains influence on their development. So it is, once for all,
in the spiritual domain; it is a web in which a
single thread cannot be drawn out without its trembling throughout
the whole. The elements of the spiritual life enter into so intimate
a union that it is impossible to regard any single one of them
as independent and constant. It was therefore an unjustified
hope when people thought, by building on the personal
need as a uniform and constant factor, to be able to erect
a new dogmatic edifice. The problem is far more intricate
than that. We stand here before some of the most delicate and most
composite psychological questions.

One may say that the philosophy of our time is moving
toward becoming an empirical philosophy of personality. By
working in the psychological, epistemological, and ethical
domain, --- the domains that in the last generation
have claimed almost all philosophical interest, it strives
to gain insight into the conditions and forms of the spiritual life,
its interest and its need. The realistic or positivistic
tendency in philosophy is often accused of
denying the right of personality. I believe, on the contrary, that no
tendency recognizes this right with more seriousness and
sincerity than Positivism. Whereas romantic theology and
speculation very quickly got their concept of personality ready
and at once began to build their often so fragile structures
% --- p. 223 ---
\opage{223}of thought upon it, Positivism sees in the concept of personality a
problem that must be illuminated and treated from very many sides, ---
by all the many ways that modern psychology uses in order to
gain understanding of the human life of the mind. And whereas
the romantic doctrine of personality harbored contempt or fear
(now more the one, now more the other) of the lawful
connection of nature and of the determinate conditions that bear
everything in life, the highest as the lowest, it is the cheerful
faith that lies at the ground of modern psychological research
that the value of the spiritual life is not diminished by its being tied
to determinate conditions and laws. In that romantic prejudice
there stirs a secret unbelief. In the personal life there
is an inexhaustible domain for research, and nowhere is
dogmatic closure more questionable. Fortunately the problems
will constantly arise anew.

How far it will ever become possible to gather all
the elements we have here discussed into one great and comprehensive
world-view, which unites the venerableness of tradition
with the consistency of thought, the richness and lawfulness
that observations offer, with the heart's innermost need, it is
a question that is impossible to answer. The only
philosophy the Catholic Church recognizes, Scholasticism, as it
was formed in the 13th century by Thomas Aquinas, and as it
was used by Dante as the frame and element of thought in his
great poem, --- it could for its time yield something of the kind.
Perhaps it has been possible only at this particular point
in history. Perhaps some day many threads that to us
seem to run in opposite directions will be gathered into a wonderful whole,
of whose nature we now cannot have any inkling. Whichever
of these two possibilities we hold to, it will for the time being
not be able to change anything in our way of working. We shall seek to
keep our horizon open, but at the same time take up our tasks
where there is really work to be done. There is in the
most recent time an idealistic tendency to be traced in the philosophical
world. It is justified, in so far as it combats the
ignoring of decisive problems to which the reaction against the
romantic theology and speculation has often led. Its
positive justification it will be able to demonstrate only when it
% --- p. 224 ---
\opage{224}can point out new means for treating the old problems.
But it will contribute to keeping the horizon open, so that the
new doubts are as little able to close it to us as the
old dogmas have been able to. Thought can do without its
solutions, but it cannot do without its problems.

For him whose position in life entails that he not only
works with the problems in the quiet closet, but also
is to lead the young on their way in their striving to find
their footing in the world of thought, for him this consideration
will have its special significance. He will seek in all his conduct to
combine the sense for the striving personality with
the conviction that personality is to be purified and ennobled through
the work of winning the truth and bowing to its demands.
He will wish to be able to do his work in this spirit, and
he will wish that this spirit may rule in the institution of higher learning
it is his honor to belong to.

\streg{}

% --- p. 225 ---
\opage{225}
\begin{center}

\essay[The Struggle between Old and New Ideas]{THE STRUGGLE BETWEEN OLD AND NEW IDEAS.}

\emph{A RETROSPECT}\footnote{Lecture, delivered in the Norwegian Students' Society in Kristiania on April 14,
1895.}.\end{center}

\streg{}

Homer calls man a being that looks forward and
backward. The two things are closely connected, for in order to be able to
look forward, which is the most important thing, one must be able to look
back, and the thoughts and moods with which we face
the future and prepare ourselves for it will necessarily --- whether
we are conscious of it or not --- be determined by
the way in which we conceive the past. And toward the
close of a century there is a special occasion for a
retrospect. To be sure, the division into centuries is a
purely external and civil arrangement, and the events neither of the inner
nor of the outer world fit into the frame
that this division gives. I shall here also try to show
that, spiritually speaking, our century begins before one
began in the almanac to write 18… Perhaps one will some
day say the same of the next century; we who
live in the midst of the time cannot know it, since we are still not able to
distinguish what is the beginning of something new from what
belongs to the old; we must naturally hope that there are
such beginnings in our time. There will nevertheless be
nothing to prevent our using the external occasion
that the civil close of the century gives, to look back on
the course of spiritual development since the century last changed
% --- p. 226 ---
\opage{226}century. A more essential occasion I have for my own
part had for such a retrospect in the
conclusion of a work on the history of philosophy in the last
century. It is on results of this work that I base myself for
the most part in the views that I shall here
advance. But before I, supported by history, go on
to give an account of how we are now situated with regard
to some of the most important problems in the domain of the spiritual
life, I must first send ahead some remarks on the
standpoint and the standard of measure for the development of the spiritual life.

\section*{I.}

If one at all acknowledges that a spiritual development
takes place, not merely in the individual, but also
in the race, so that no more in the inner than in the outer
world does everything remain as it was, then there naturally arises the
question whether by this development a progress is made
or not, --- whether something is lost or gained by the changes
that take place. In the concept of development there lies not necessarily
that something more valuable comes into being than there was before, nor
yet that the values won are preserved. A
development --- in the sense of a lawful series of
changes --- is conceivable which leads to a decline not merely in value,
but also in energy.

In the domain of external nature, as is well known, in
the last half-century the law of the conservation of energy
has been set up and carried through. Physicists are convinced that in
the material world no energy arises and no energy
perishes, but that in every change and
development energy is only converted into new forms. But not even here is
the conservation of energy one with the conservation of value, for one has
indeed discussed the possibility that in our world-system
such a distribution of heat might set in that all motion
and all life would cease, so that that which from our standpoint
(and another standpoint we simply do not have) gives the world
value --- that it can be the home of life and spirit
--- would thereby have fallen away.

% --- p. 227 ---
\opage{227}How are we now situated in the spiritual domain?
I shall not here enter into the much-disputed
question of the relation of the spiritual life to the law of the
conservation of physical energy. However one may take up a position on
this question, there will, as soon as one admits that
a development takes place in the spiritual domain, be
a natural need to ask whether energy and value
are preserved during this development or not. But since
the science of the spiritual world is not, as
the science of the material world is, an exact science,
we cannot expect to arrive at any certain result.
Already with regard to demonstrating the conservation of
the energy that makes itself known in the functions of the spiritual life,
we are very unfortunately situated. We cannot demonstrate that
there is constantly the same possibility of spiritual life present in
the world. And even if we limit the consideration to a single
period, we lack means to convince ourselves that in
the course of development no energy is added or perhaps
new energy arises. The question whether in the course of one or
several centuries a weakening or a conservation
or an increase of energy has occurred, we cannot answer with certainty.
Even if we now confidently assume that the energy
that stands at the disposal of the spiritual life is not diminished,
the question is nevertheless not thereby answered whether this
energy finds application in as valuable a way as
formerly. Morbid development can lay claim to as
much energy as healthy development, a process of dissolution can
consume no less energy than a process of formation;
the magnitudes may be the same, but the signs different.
The mere conservation of spiritual energy will for the time being
be of value, since the value of life also rests on its strength;
but it is not the only standard of measure. Besides the
strength of the spiritual life we ask also about the richness and fullness
of its content. The greater and richer the nuance that it
presents, the more differences and oppositions it is able
to embrace, the greater value we attribute to it. But on the
other hand: only if the richness of content and
diversity of the intellectual life is combined with the capacity to act with concentrated
% --- p. 228 ---
\opage{228}strength is life healthy. Very easily an inverse
relation may arise between strength and fullness. In our valuation we do not
restrict the consideration to the single individual taken by himself alone.
We regard the intellectual life as healthy only when the single
personality stands in close connection --- as giving and
receiving --- with society, the race. We value the
intellectual life of a period by the strength and independence that the single
personality can possess, without the spiritual fellowship
with other personalities being thereby broken. It is an unhappy
time, not only when strength or fullness is lacking in the single
individual, but also when fellowship is lacking,
understanding and harmonious interaction between the individuals
among themselves. When the thoughts and feelings that bear life become
wholly different in the different individuals, indifference
and contempt take the place of sympathy; there is lacking the joy
of giving and the joy of receiving. Precisely the most
eminent individuals, whose intellectual life has significance for large
circles, suffer under this; \emph{they} feel the greatest want, because they
have the possibility of giving more and receiving more than others.
One is spared many problems and many sorrows when
one narrows one's horizon.

Recent psychology has in various ways demonstrated
the possibility that the energy that the spiritual life has possessed
at earlier stages can continue to exist under new forms
at later stages in the development. To this belongs first and foremost
the phenomenon of the displacement of motives demonstrated by Spinoza and the English psychologists.
What at first has interest only and
awakens striving because it is a means to attaining an end,
can become the object of independent interest and gradually
itself come to stand as an end. Thus a thought
that we at first bow to out of obedience to an authority, a thought
whose value for us therefore at first rests solely on our fear or
our trust in the authority, later by its own power,
its own grounding, can stand firm for us, even if
the relation of authority falls away. In society there can develop
institutions that are to fulfill definite ends, but that later
are maintained after the original ends have fallen
away, because it turns out that they can serve new ends. Such
% --- p. 229 ---
\opage{229}displacements of motive go on uninterruptedly in the single
individual and in history, and without them one does not find
connection in development. Another form in which continuity
through all changes can be preserved is what, with
the use of an obvious analogy, has been called mental
chemistry. Mental phenomena can enter into such a
combination that they no longer assert themselves each
separately in their peculiarity, but as it were fuse
into a new phenomenon, which is their common product, just as
water arises by the combination of oxygen and hydrogen. Joy and
sorrow can attach themselves to the same event, and the thought of
this event can then in the end be connected with a
fundamental mood, melancholy, which lets neither the joy by itself
nor the sorrow by itself step forth, but expresses them both
at once. Here an equalization has taken place, whereby
the earlier mental experiences are preserved under other forms.
Displacement and equalization are the two most important ways in
which spiritual energy will be able to be preserved. But thereby
it is not said that the value is the same in the new forms
as in the old. Not only the healthy, but also the diseased
consciousness presents displacements and equalizations.
Also in the doctrine of mental diseases one speaks of psychical
equivalents, thus especially where the tendency to epileptic fits
finds vent in another way, in the form of sudden rage or
surprising change of mood. The one and same sum of
energy that is present can find outlet in many ways.
But even if it can be demonstrated that it is preserved, it is not
said that it is preserved without change of value. Indeed,
displacements and equalizations go on constantly, without conservation
either of the strength or of the value.

One point is here of special importance. At a certain stage of
spiritual development, that of the race as well as of the single individual,
the sole dominion of instinct and authority is past. The horizon
widens; there shows itself a greater multitude of possibilities and
directions both for faith and for action. Hitherto the world had
been limited, and only a single
way of conceiving life and a single way of living it had presented itself.
Now the doors are thrown open to the infinite world with its many
% --- p. 230 ---
\opage{230}points of view and directions of life. Can now strength, security, and
inwardness be preserved, when the content becomes so much the
richer? Can the spiritual fellowship be preserved, when so
many directions in faith and life gain their more or less
free play? --- That was a question which first the century of the Renaissance and
later that of the Enlightenment and the Revolution posed.
It is a fundamental problem that under various forms extends
through the thinking of modern times from the time of the
Renaissance down to our own day. The Renaissance itself did not pose the
problem; it was so absorbed in the new and great things that it
discovered that it did not feel the sting of the problem. The age of the Enlightenment
was likewise so absorbed in its criticism and its reason,
in thrusting the old aside and feeling the joy in the
use of critical and testing reason, that it did not think
whether something had not been lost by the break with instinct and authority
that the new forms of spiritual life would not
be able to replace. Even if the continuity of development could be
demonstrated in spite of all Renaissances and revolutions (as
geology thinks it can demonstrate the continuity of
the earth's development in spite of all revolutions of the earth), and even if it could be
demonstrated that the same energy that expressed itself in the concentrated form
of instinct and faith in authority now expresses itself in a more
dispersed, distributed, and nuanced form, it is nevertheless not said
that these different forms of expression are of equal worth. The
force distributed over many points and appearing in many nuanced forms
has perhaps at no point the
possibility of appearing concentrated, the security and might of
concentration, which it had under the old order of things.
Does not life grow feebler in the course of development, not
because strength vanishes, but because it is distributed and
nuanced, --- just as physicists believe that our world-system will
some day congeal, because the energy is distributed, so that the great rhythms
in the world process are no longer possible?

That is the problem. And when we now take up our
view of the time, I find it peculiar to our century
that this problem has been felt and posed. And \emph{our}
century's beginning must be set where the problem first
appears. But then we must leave the limits of the almanac; for
% --- p. 231 ---
\opage{231}the problem was first posed by \emph{Rousseau, Lessing}, and \emph{Kant}. They
were the first who at once felt the whole justification of
modern criticism and of the widening of the horizon, and at the same time
saw it as a great question how the revolution in
the spiritual life that stood at the door was to be able to take place
in such a way that the health and inwardness of life did not suffer thereby.
They belonged to the age of the Enlightenment, at the same time as they saw its
limits and perceived the necessity of preserving the innermost
kernel of the old faith and the old order, even if
the shells had to be cast away. With these three men the nineteenth
century begins. They posed this its real
problem.

\section*{II.}

The history of philosophical thinking in the past
century presents features that prove that the problem has constantly
been felt by a whole series of eminent spirits, so
that it may rightly be called the problem of our time.

Philosophy can --- especially from a historical standpoint ---
be regarded in two different ways. --- Philosophical ideas
can in the first place be regarded as \emph{symptoms} of what
is stirring in the time. What in most people only moves
as obscure feeling and involuntary need, that takes shape
in the philosopher as an idea, just as in the poet it takes
shape as an image. The impulse to shape in clear thoughts what
moves in the mind is peculiar to philosophical
natures. We therefore often see in the history of philosophy ideas come
forth that the thinker concerned is himself not able to give
full grounding of; they stick their heads for the first time
into existence, emerge from the chaos of experiences and the
obscure world of involuntary tendencies into the daylight of thinking, and
their first form is therefore often imperfect and without firm
grounding. Only later can it be understood how they
hang together with the whole conditions of the time and of development. ---
On the other hand philosophy is an attempt at \emph{scientific
treatment} of unavoidable, consciously posed problems. From
this standpoint the valuation of a philosophy often turns out quite
% --- p. 232 ---
\opage{232}differently from that from the first; but every
significant philosophical endeavor will present elements of both
kinds. And there is moreover evidently an inner connection
between philosophy as symptom and philosophy as
science. Common to both sides is that philosophy is the
human spirit's striving for conscious clarity concerning
its nature and its conditions of life; philosophy is \emph{scientific
self-knowledge}.

In this connection I shall, however, especially make use of the
first standpoint. I find it a remarkable sign of the times
that the problem raised by Rousseau, Lessing, and Kant extends
under various forms through the thinking of the whole century.
Again and again the thought comes forth: We have
left the secure and firm ground of instinct and authority;
criticism and analysis have done their work, a work that could
not be kept away and cannot again be undone;
but on criticism and analysis we cannot live, for life constantly
demands positive content and firm ground under the feet; there
must then be reached a new stage, which unites the positive
strength of the former with the critical many-sidedness and
mobility of the latter! If this is not possible, then life is in dissolution.
We can indeed for a while longer feed on
the after-effects of the instinctive and believing childhood time;
but this capital will in the end be used up, if
new positive forms do not develop. To use an image
that Renan has used in his ``Souvenirs d'enfance'': According to
a legend from Brittany a church bell was once sunk in the sea
when a boat capsized. In calm weather ---
it is told --- one can still hear its tolling from
the bottom of the sea. Thus we still hear the tones from the time
that was united in common and powerful faith and in naive trust in
life; but will not the sound some day die away? will we not
move so far away from its place that its tones cannot
be heard?

It was \emph{Rousseau} who first in earnest raised the problem,
not only in his well-known paradoxes, but in all his writings,
in his inmost line of thought. He was shaken through by it, became
a victim of it. He became homeless, because he neither shared
% --- p. 233 ---
\opage{233}the Conservatives' joy in what exists nor
the Encyclopedists' joy in the new enlightenment. A similar
position was taken by \emph{Lessing} toward the Orthodox and Rationalists
in Germany, and by \emph{Kant} toward Dogmatism and Skepticism.
Both of them are --- especially Kant --- influenced by Rousseau.
None of these men wanted to go back to the old, as little as
they were satisfied with the new. They awaited --- each in his own
way --- the third kingdom. Kant, the greatest thinker of the
three, has most clearly formulated his relation to the two tendencies
from which distance was to be taken. He rejected
like the Skeptics the dogmatic systems, but he did
it not merely through negative criticism. He sought to understand
how the systems could come into being, to understand the spiritual
need and capacity that found vent in them, and that could continue
to live, even if it found vent in another way. In this
going back from the finished systems to their first motives
and in this investigation of how far these motives can
continue to be effective under other conditions and in other ways,
I find Kant's greatest achievement. Thereby at once
continuity with the past was saved, and the prospect of the future
kept open. Kant has thereby set the whole
program of the science of spirit. We come to know the human spirit through
its works. But the \emph{finished works}, the fully executed
forms of faith and systems of thought easily become a dead treasure without
connection with our own life and therefore useless for
it. Therefore we must penetrate to the \emph{forces that have
produced the works,} in order to see whether they still stand at our
disposal, and to see on what conditions they must now
work.

In the century that has passed since Rousseau, Lessing,
and Kant, two principal tendencies can be pointed out in philosophical
thinking. The one I call \emph{Romanticism,} the other
\emph{Positivism}. The one seeks the solution of problems by way of
speculation and contemplation, the other through
establishing facts and reflections on
facts. They stand sharply opposed to one another, since they take their
point of departure at the opposite limits of knowledge.
The one wishes to descend from above, the other to ascend from below. But
% --- p. 234 ---
% print p. 234: paavirket at Kant (at for af) -- misprint; the sense is translated
\opage{234}it would be strange indeed if they did not meet at all. And
a closer investigation also shows that both
tendencies, each in its own way, follow Kant's program for the
science of spirit. Characteristic of both tendencies is a deep
understanding of the Middle Ages and in general of the spiritual life
of earlier times, a dissatisfaction, perhaps even an aversion toward
mere negative criticism, and a view toward a new time, when
positive fullness and critical reflection, common fundamental outlook and
individual striving for freedom shall no longer stand hostilely
over against one another. Throughout, they also conceive
the present period as a time of transition, a
time of upheaval, as a time of unfinished beginnings. Yet
--- strangely enough --- the Romanticists have a greater inclination
to acknowledge the given order of things than
the Positivists.

In \emph{Fichte's Grundzüge des gegenwärtigen Zeitalters} (1806)
the doctrine of the three stages finds perhaps its clearest and most
energetic expression. After the period of instinct and authority
there comes a period that is designated as the period of empty freedom,
whose hallmark is ``Auf- und Ausklärung.''
Glad of enlightenment and criticism, one turns with pride
from the darkness of the past. Negation and egoism rule. But
the thinker looks forward to a new time, when a more thorough
science and a higher art will lead to a new view of life
that lets all spiritual needs, devotion as well
as reflection, come into their own. At about the same time
\emph{Saint-Simon} in France designated the eighteenth century as
a critical, revolutionary period, which was now to be succeeded by an
organic period. His disciples generalized the
opposition between organic and critical periods and found it
recurring at various points of historical
development. An organic period is characterized by firmness and
closedness in faith and thought, and by spiritual fellowship. In
the critical period, on the other hand, doubt and antagonism rule
at the expense of faith and association. Whereas Fichte is
influenced by Kant, and through him by Rousseau, Saint-Simon
seems to have come to his conception rather through
reflections on the
% --- p. 235 ---
\opage{235}course of development of industry and science. On the further development of his ideas among his
disciples Lessing's essay on the education of mankind had
no small influence.

It would be too lengthy to bring forward all the
different forms in which the same great thought
appears among the romantic and positive philosophers. --- \emph{Hegel's}
triads and his whole dialectical method have properly
arisen by schematizing the law of the three stages. At
the basis there lie --- as is seen from notes of Hegel's
youth --- studies in the history of culture and religion. As
happy times Hegel praises those periods in which there is faith
in the kernel of existence, and where a common understanding of life
unites all. Subjective criticism then does not isolate the individual
from the whole and does not make him at odds with himself. But
it was Hegel's conviction that the substantial in the
old faith and in the old order of things could be preserved
through the purgatory of reflection, and in his own philosophy
he thought he had demonstrated it. \emph{Schleiermacher} was more critical,
but yet thought he could demonstrate the same
experiences of feeling from which what is essential in the old faith had
sprung. Neither of the two thinkers would admit that they
were remodeling the old forms of spiritual life according to their needs. Therein
consists precisely what is romantic in them, that they directly
think to live the life of the past over again. --- Within
Positivism there is first and foremost to be mentioned \emph{Auguste Comte's}
formulation of the law of the three stages. The spiritual development
proceeds accordingly through three stages, the theological, the
metaphysical, and the positive. The metaphysical stage is
purely negative and critical; it is a period of dissolution, corresponding
to what Fichte called the period of empty freedom, to what
Saint-Simon called the critical period, and to what Hegel
called the unhappy times, the periods of critical subjectivity.
According to Comte there is a spiritual kinship between the
positive and the theological period; they each have their own
comprehensive conception, and they connect the individual with the race in an
inward, organic way. In \emph{Stuart Mill} and \emph{Herbert Spencer}
we find similar ideas, likewise in \emph{Schelling} and in \emph{Eugen
Dühring}. \emph{Carlyle} is filled with the thought that he lives in
% --- p. 236 ---
\opage{236}one of the most unhappy times in history, in the period of criticism and
mechanism. But he is no less convinced
that ``the old eternal powers live continually,'' even if
every romantic-speculative attempt to rehabilitate the
old forms is, in the manner of the ostrich, to stick one's head in a
bush so as not to see things as they are. He is even
convinced that the bird Phoenix does not die on the pyre, in order then
to come into being anew, but that it is constantly alive, even if
we do not hear the beat of its wings. He lived on the law of the three stages,
in the faith in the third kingdom.

The Romanticists are more inclined to believe in an
identity of the old and the new thoughts, the Positivists more
to conceive the relation between the old and the new as
equivalence. Therefore the Romanticists are more conservative than
the Positivists. But common is the conviction of the
continuity of development. Whereas the age of the Enlightenment and the Revolution
looked down on the Middle Ages as sheer darkness and barbarism,
both Positivists and Romanticists acknowledge its great
significance as an organic period. And whereas the reaction in the
nineteenth century looked back at the age of the Enlightenment and the Revolution
as sheer intellectual pride, savagery, and negation, so
both Romanticists and Positivists acknowledge the necessity of such
critical periods.

\section*{III.}

But --- one might say --- this agreement does not mean
much, when the two tendencies disagree about how the
third stage, the third kingdom, is to be. Here a
decisive opposition emerges. And besides: Is the doctrine of the three
stages, or of organic and critical periods in rhythmic
alternation, acknowledged by all? Are there not those who on this
point agree neither with the Romanticists nor with the Positivists,
but who either think that no third kingdom at all \emph{needs}
to come, or that --- even if it is needed --- no
such thing \emph{can} come?

I shall seek to answer these objections, and take
the last first.

% --- p. 237 ---
\opage{237}Although a comparative consideration of the two
principal tendencies in the last century, which may be designated as
Romanticism and Positivism, finds greater agreement
than one would expect beforehand, with regard to so
essential a problem as the one we here consider, there is
nevertheless another feature which cannot but force itself upon
whoever considers the course of intellectual development in our
century. It is the way in which the oppositions
are sharpened. Not content with the recognition and sympathy
that both Romanticism and Positivism show toward
the old faith and the old order of things, faith in authority
unfolds itself in the course of the century with ever greater
consistency and in ever more extreme forms.
Intermediate standpoints are crushed. Within the Catholic Church
Gallicanism drops away, the infallibility of the Pope is proclaimed, and all
philosophizing other than that which follows in the footsteps of Thomas Aquinas
is rejected. And as Catholic theology goes
back to the thirteenth century, Protestant
theology goes back to the seventeenth century. From standpoints such as
these one feels no need of any third kingdom. The one thing
needful is assumed to have entered into history once and for
all, and all the unrest and distress of the time is derived from men's not
keeping to it. With pity one looks down on
the philosophers who torment themselves with finding the conditions for the
continued existence of the spiritual life, like those who stand on the shore look
on the struggle of the shipwrecked with the waves. On the other
side there are those who think that the work of dissolution, which
the Renaissance and the Revolution have begun, cannot be stopped. They
perhaps even rejoice in the dissolution and regard
the products of dissolution as the highest product of development. Criticism and
negation have burned out in them the capacity and the need for
positive life, for devotion and for spiritual fellowship. They
feel no need of a third kingdom, and even if they feel it,
these ``tired men'' regard it as impossible that it can
come; far rather than strive forward they throw themselves back
into the arms of the old faith.

But over against this undeniable fact, the sharpening of
oppositions in the spiritual domain, we must refer to
% --- p. 238 ---
\opage{238}a psychological law which was already mentioned above:
the law of the displacement of motives. Even if old forms and
traditions are still maintained among us, a closer
consideration will show that essential changes have taken place
in the subjective relation in which individuals stand to them.
It is in any case a great question
--- a new problem, or rather a particular form of the
problem we here treat --- whether the Christianity of our time is
the same as primitive Christianity, and how the two
relate to one another. Thinkers who come from such different
sides as Schopenhauer, Feuerbach, and
Kierkegaard deny it. Only blindness to what is essential
in a practical view of life can deny that here
in any case is a problem. And more I will not
for the time being maintain. --- As regards the Skeptics, it is a matter
of showing that they confuse the old system with the
forces that worked to produce it. Criticism
hits the finished building, not the hands and the energy
that brought it into being. Only experience can decide whether
these same forces are not still living, even if the forms and
conditions for their application are now essentially changed.
The Skeptic himself becomes dogmatic when he denies the new
possibilities. There are forces of which one learns only that
one has them when one uses them. And as the Dogmatist
overlooks the displacement, the metamorphosis, that has
taken place especially in the last centuries, so the
Skeptic overlooks the displacement that is constantly going on among us,
and which makes it that the critical periods too, when they come
again, assume a different character from the previous time. There is
an essential and characteristic difference between Ludvig
Feuerbach's relation to religion and Voltaire's. Such
displacement and fusion of motives is the great household resource
of history in the domain of religious and social development.

As to the nearer nature of the third kingdom,
the new organic period, it is no wonder that
the divergences in the conceptions are great. A construction of
the faith of the future or the social order of the future is impossible.
% --- p. 239 ---
\opage{239}The general features that we can indicate, --- the union of
concentrated strength and nuanced fullness, of mobile criticism
and firm faith, of individual freedom and social organization, ---
are not sufficient for such a construction. And
besides, the constant displacements will prevent one from being able
beforehand to form an idea of the forms of the more
advanced life.

In my view it is also a misunderstanding of
the significance of philosophy when one assumes that its task
should be to give such a construction. It is
the business of philosophy to be a scientific self-knowledge, to seek
clarity about how we are situated in spiritual respects,
how far the forms of our life correspond to the stage we have reached,
and what forces have produced these forms.
The philosopher must be content when he can contribute to the
solution of this task. He sympathizes with the unrest of the
Skeptics, in so far as it is really a seeking, and with the
conservatism of the Dogmatists, in so far as it really
preserves something valuable.

For my own part I have in any case only
wished to establish a problem, not to give a definitive solution.
But with the posing of the problem there is also given the spirit in
which, and the method with which, we must work. I shall
emphasize the points that I consider
the most essential in this respect.

In the first place there lies in the recognition of the problem
a principle one might call the principle of equivalence. What
has filled an essential place within the spiritual
life cannot fall away without replacement. Mere
criticism cannot be the last word, if life is not to
come to a stop. It is the task of the science of spirit to investigate
whether new forms of spiritual energy stand ready to
take their place, if its activity diminishes. It is a matter of
finding out what gives the spiritual life strength and value
at any given time. This task is --- as already
touched upon --- far more difficult than the corresponding task
in the material domain, because we lack an exact
standard of measure, both for the energy and for the value. But it
% --- p. 240 ---
\opage{240}is an unavoidable task, both for the theoretical understanding of
the spiritual life and for its practical continuation. We
understand a personality only when we see what it draws
nourishment from, and how it uses the nourishment that it
takes up, just as we understand what goes on in an
organism only when we follow the metabolic processes in it,
from the taking up of elements from the surrounding world to
the manifold use by the functions of the elements taken up.
We are still only spelling out, in our psychology, a doctrine of
personality, of the conditions for spiritual
centers to form, and of the conditions for their health. Here
the romantic tendency has far too often contented itself with
determining the concept of personality in abstracto and with
setting it up as an object of contemplation, as Plato
set up his ideas. As I understand and favor
the positivistic tendency in philosophy, it aims at
finding the conditions and laws to which personal life,
like everything else in the world, is subject. Through
experiment and self-observation, by the use of physiological and
historical material it will find the laws that
hold together life within the consciousness of the single individual, as well
as the laws according to which the individual's conscious life is
bound up with the life of the whole race. It is so far from being the case that
realistic research will lead to denying the
idea of personality, that it will precisely strengthen the applicability of this idea
by leading to insight into the definite laws, conditions, and
limits for it. The romantic fear or contempt of
the mechanical connection, which once and for all is the bearer
and frame of everything that is to be able to exist and to assert
itself in the world of reality, is one of the most important
hindrances to spiritual development. It is itself a
secret unbelief, in that it rests on the presupposition
that it lowers the value of the intellectual life to be subject to a
lawbound connection. A confident faith in
the idea of personality will bring both trust that it will be able to
assert itself under new forms, and zeal to investigate the
order of things that can lead personal life over from
% --- p. 241 ---
\opage{241}old to new forms. The ideal for all psychological and ethical
research is hereby given. Toward this ideal the
Positivist fixes his gaze, when he wishes to avoid being confused by the
manifoldness of the single investigations. Like every ideal it has
something irrational, that is to say: irrational in the mathematical sense,
about it: by our actual endeavors we can only
reach an approximation, determination of more and more decimals.
But it is the same limitation to which all our knowledge
is subject.

Under this endeavor to provide
the foundation of a philosophy of personality, which can teach us to distinguish
between the true and the false positions and between the
true and false negations, to discover the ``organic'' and
the ``critical'' in their justified forms and distinguish them from
surrogates and falsifications, --- we must be clear about
one thing: that there is so great a difference between human
individualities that they need highly different nourishment,
and that the values cannot become the same for all. The
real nutritive value belongs only to that which suits the particular
individual personality. One will not be able to bring all under
a common external expression of the truth, when this
is to be personal truth for each single one. The mark
of personal truth is, as a Danish thinker has
said, that it is produced anew by the individual himself: ``No
manifestation of life has truth, unless creative
self-activity lies in it.'' No excuse is to be made for the
individual differences which bring it about that everyone shapes
the truth in his own way. As life shows its energy
in the single personality by the fullness of thoughts and
motives that it is able to concentrate on its guiding
aim, so it is a sign that life in the race is
energetic, when diversified, each distinctly marked
personalities appear, who take life and the world each
according to the bent of his nature. The ``critical'' periods set in
in history, when a preceding ``organic'' connection
has become too narrow, too habitual, --- when the horizon needs
to be widened and new experiences and new points of view must be
% --- p. 242 ---
\opage{242}won. Then the future of the race can be saved only through the
individual variations. Every single variation, every single
nuance of the general human theme means possibility
for new experiences, new outlooks and discoveries, which can
benefit the continued life. There appeared in England
somewhat after the middle of the century two books, by two of
England's greatest investigators in this century, which each from
his standpoint pleaded the cause of individual variations. It
was \emph{Stuart Mill's} book on Liberty and \emph{Darwin's} book on the Origin
of Species. Stuart Mill pleaded the cause of individual
freedom, because it was a matter of there being as many
independent starting points for development as possible, and
Darwin demonstrated the individual variations in animals and
plants as a condition for the development of new forms.
When in the most recent time we so often see individuality
asserted in a forced manner, when e.g. in
literature peculiar characters and peculiar situations are presented with
predilection, it is a sign of a need and a
sense that is in the process of awakening, and which may bring
with it much that is not healthy, but which yet only by
being allowed to come under the discipline of life and experience itself
can find its limits. Not by moralizing from without,
but by trust in the capacity of personal life to find its
own rules, to work its way out of its crises to new
organic life, --- and by holding fast the great ideals and
conditions of spiritual health, will it succeed in saving the
true values through the purgatory of criticism and the apparent
dissolution.

Our time presents in the spiritual domain a picture
of motley confusion. There are those who hold on to the
old at any price, --- there are those who with a light heart
give it up, without feeling need of anything new, ---
there are those who grieve that the old times are past,
without being able to hope for anything from the new, --- and there are
finally those who have the program for the new complete in all
its rubrics and impatiently demand that the old shall give
place to the new glories. The philosopher's task is,
through quiet research, to find means to test the strength
% --- p. 243 ---
\opage{243}and the value both of that which goes and of that which comes,
and then perhaps also, through the clarity he may be able to
bring, and the faith in the unsurveyable possibilities of life that he
may be able to strengthen, to give his contribution to the work of the time.
He himself seeks, like Dr. Faust in his brooding, back
to the thought that the world of the spirit is not closed: neither
old dogmas nor new doubts will be able to shut it
against us!

\streg{}
\end{document}
